\documentclass[11pt]{article}
\usepackage[T1]{fontenc}
\usepackage[utf8]{inputenc}
\usepackage{lmodern,amsmath,amssymb,amsthm,mathrsfs,mathtools,booktabs,longtable,array}
\usepackage[margin=1in]{geometry}
\usepackage[colorlinks=true,urlcolor=blue,linkcolor=blue]{hyperref}
\newtheorem{theorem}{Theorem}[section]
\newtheorem{proposition}[theorem]{Proposition}
\newtheorem{lemma}[theorem]{Lemma}
\newtheorem{corollary}[theorem]{Corollary}
\theoremstyle{remark}\newtheorem{remark}[theorem]{Remark}
\DeclareMathOperator{\CH}{CH}
\DeclareMathOperator{\Alg}{Alg}
\DeclareMathOperator{\Pic}{Pic}
\DeclareMathOperator{\Alb}{Alb}
\DeclareMathOperator{\NS}{NS}
\DeclareMathOperator{\End}{End}
\DeclareMathOperator{\Hom}{Hom}
\DeclareMathOperator{\Nm}{Nm}
\newcommand{\Q}{\mathbf Q}\newcommand{\Z}{\mathbf Z}\newcommand{\F}{\mathbf F}
\newcommand{\tun}[1]{#1}
\setlength{\emergencystretch}{2em}
\title{C44\_POS: primitive signs, special axes, and the 1964 programme}
\author{Research continuation of JAC1964 and SIGN\_NEXT}
\date{22 November 2026}
\begin{document}
\maketitle
\begin{abstract}
Working research manuscript. The general fixed-axis sign and the Fermat
Abel--Jacobi injection remain unresolved. New checked results include
an exact axis primitive decomposition, lifts of every embedded smooth
axis (without a cycle-lifting assertion), and an unconditional sign
theorem whenever at most one potential Tate direction remains after
a known positive algebraic span. An extending orthogonal-block criterion
now also proves the sign on actual closed affine-$S_3$ families with
three residual potential Tate directions, including families with a
known positive span of rank six. The unsplit higher-rank gate remains
open. Other
checked findings include closed supersingular axes with their entire
positive kernels controlled, intrinsically nonliftable surface centres
whose primitive nonliftability is still open, lifts of every finite
curve cover as a map, special-cover sign tests, exact integral
arithmetic obstructions, a characteristic-$p$ detector/division theorem,
and geometric extensions of the primitive incidence polarization.
The algebraically trivial Fermat obstruction is reduced to specified
nonzero universal cycles from excess Picard quotients; several entire
classes of singular ample image surfaces are eliminated in rational
Chow. Full Abel--Jacobi injectivity remains open.
The proposed supersingular Fermat test case already satisfies the
required Hodge standard conjecture; its corrected restriction and
weighted Gysin conditions are verified explicitly. A separate
nonsupersingular degree-seven Fermat calculation supplies explicit
new special cycles, all-prime Tate and positivity for the fourfold,
and numerical and 13-adic positivity for its square.
Every use of 032 or Milne's 2026 report remains conditional on those
unrefereed preprints. The separate Milne 2022 transfer assertion is
exposed wherever used. Whole-proof reading events and finite
computations are distinguished in \texttt{CHECK\_EVENTS.tsv}.
\end{abstract}
\tableofcontents
\part{Checked findings and explicitly conditional deductions}
\section{Normalization, scope, and inherited results}
The starting isometry is JAC1964 H9-T7, equation (H9.18):
\[
 A_X\oplus A_B\simeq\CH^m(Y)_{\mathrm{hom},\Q}/\mathcal H_0,
 \qquad Q((z,w),(z',w'))=(-1)^m\int_Xzz'+(-1)^{m-1}\int_Bww'.
\]
Here $A_B=\ker\kappa_*$ on the algebraic middle classes of the axis,
and $h_BA_B=0$. This is an inherited statement, not a new result of
this run. The ambient sign is conditional on 032 Corollary 8.4 and
Milne's report. Weak semidefiniteness and strict positive definiteness
are distinguished throughout; no numerical/homological equivalence
on an arbitrary axis is silently assumed.

Geometric generic parameter points are taken over algebraic closures
of function fields. They are not closed points over
$\overline{\mathbf F}_p$. A monodromy vanishing assertion at the
former does not settle a sign at every latter point.

Every finding retains its stable C44 identifier. A displayed
proposition that assumes an additional unproved input is a checked
conditional implication, not a proof of that input. Prior statements
retain their JAC1964 or SIGN\_NEXT identifiers and are not claimed
as findings of this continuation. The governing source is the attached
\texttt{PROMPT(20261121-231316).md}.

% BEGIN work/axis/axis_contribution.tex
% C44_POS /root/axis_independent; GPT-6 Astra Pro; effort setting not exposed.
% Mathematical deductions for root review; no claim of formal verification.
\section{The exact axis primitive space and the limits of the proposed transfers}

\paragraph{Notation and carried input.}
Put $r=m-1$, so that $\dim X=2r+2$ and $\dim B=2r$.
Write $L_X$ and $L_B$ for multiplication by $h$ and $h_B$.
All cycle groups have rational coefficients, and $\operatorname{Alg}^j$
means the image of the rational cycle-class map for one fixed
$\ell\ne p$. No equality between numerical and homological equivalence
on $B$ is assumed. The carried statement JAC1964 H9-T7, equation
(H9.18), is the isometry
\[
 A_X\oplus A_B\ \simeq\ \mathcal H/\mathcal H_0,
 \qquad
 Q((z,w),(z',w'))=(-1)^m\int_Xzz'+(-1)^r\int_Bww'.
\]
Its proof establishes $L_Bw=0$ for $w\in A_B$, not merely
$L_B^2w=0$. The factorization $B\hookrightarrow Y\hookrightarrow X$
through a smooth pencil member and the injectivity of the latter
Gysin map are essential to that stronger conclusion. We use this
checked input without reproving its height isometry.

\subsection*{C44-A01. The axis gate is the entire remaining primitive sign on the axis}

\textbf{Proposition.}
There is an orthogonal direct sum of rational algebraic classes
\[
 \operatorname{Alg}^{r}(B)\cap P^{2r}(B)
 =\kappa^*\bigl(\operatorname{Alg}^{r}(X)\cap P^{2r}(X)\bigr)
       \mathbin{\perp} A_B,                                  \tag{C44.A01.1}
\]
where $P^{2r}(B)=\ker L_B$ and $P^{2r}(X)=\ker L_X^3$.
On the first summand the signed middle form is
\[
   (-1)^r\int_B\kappa^*a\,\kappa^*a'
       =(-1)^r\int_Xh^2aa'.                                 \tag{C44.A01.2}
\]
For every $0\le j<r$, restriction identifies the primitive algebraic
classes of codimension $j$ on $X$ and $B$, including their signed
Lefschetz pairings.
Consequently, \emph{conditional on the unrefereed 032/Milne input
asserting the Hodge standard conjecture for $X$}, strict positivity
on $A_B$ is equivalent to the full Hodge standard conjecture $I(B)$
for $h_B$. Weak nonnegativity on $A_B$ is equivalent to the version
of $I(B)$ with every positive-definite assertion replaced by
positive-semidefinite. These last two statements must not be
identified unless nondegeneracy on algebraic classes is additionally
known.

\textbf{Proof.}
Hard Lefschetz makes $L_X^2:H^{2r}(X)\longrightarrow H^{2r+4}(X)$
an isomorphism. Its inverse $\beta$ is induced by an algebraic
correspondence because the Lefschetz standard conjecture is known
for abelian varieties, without 032. The self-intersection and
projection formulas give
$\kappa_*\kappa^*=L_X^2$. Hence
$e=\kappa^*\beta\kappa_*$ is a projector and every algebraic middle
class $x$ of $B$ has a unique decomposition
\[
 x=\kappa^*a+w,\qquad a=\beta\kappa_*x,\qquad w\in A_B.       \tag{C44.A01.3}
\]
Both components are rational algebraic classes. The summands are
orthogonal because
$\int_B\kappa^*a\,w=\int_Xa\,\kappa_*w=0$.
H9-T7 gives $L_Bw=0$. The Gysin map
\[
 \kappa_*:H^{2r+2}(B)\longrightarrow H^{2r+6}(X)
\]
is an isomorphism, being dual to weak Lefschetz in degree $2r-2$.
Therefore $L_B\kappa^*a=0$ if and only if
$\kappa_*L_B\kappa^*a=L_X^3a=0$. This proves (C44.A01.1).
The projection formula proves (C44.A01.2).

For $j<r$, weak Lefschetz gives an isomorphism
$\kappa^*:H^{2j}(X)\longrightarrow H^{2j}(B)$.
Its inverse on this degree is the algebraic operator
\[
 (L_X^{2r+2-2j})^{-1}\kappa_*L_B^{2r-2j};                   \tag{C44.A01.4}
\]
indeed its composition with $\kappa^*$ is the identity.
Thus this is an isomorphism on rational algebraic classes too.
For $1\le j<r$, the map $\kappa_*$ in the degree of
$L_B^{2r-2j+1}\kappa^*a$ is an isomorphism by dual weak
Lefschetz in degree $2j-2$. Thus the primitive condition on $B$
is equivalent to
\[
 L_X^{2r-2j+3}a=0,
\]
which is precisely the primitive condition on $X$. For $j=0$
both primitive spaces are $H^0$, so the same assertion holds.
Finally
\[
 (-1)^j\int_B h_B^{2r-2j}\kappa^*a\,\kappa^*a'
 =(-1)^j\int_Xh^{2r+2-2j}aa'.
\]
The conditional positivity assertions now follow from the
orthogonal direct sum and the signed Hodge standard conjecture on
$X$. No positivity assertion on the new summand has been proved.
\hfill$\square$

\paragraph{Programme location and scope.}
This is a sharpened identification of JAC1964 QH9-axis, and concerns
cote 44, Th.~5, scan 004, with the primitive margin on scan 008.
In cote 16 terminology it isolates $I(B)$; the algebraic operator
$\beta$ uses $C_G(X)=B(X)$, not $I(B)$. It does not close
QH9-axis and does not upgrade its requested semidefinite sign to
strict positivity.

\subsection*{C44-A02. Every graph transfer to an abelian target annihilates the axis classes}

\textbf{Proposition.}
Let $D$ be an abelian variety and $f:B\longrightarrow D$ any
morphism. Then $f_*w=0$ in cohomology for every $w\in A_B$.
The same holds for a rational map, interpreted through its unique
extension. Consequently $A_B$ is orthogonal to every cohomology
class pulled back by such a morphism in the complementary degree.
This is unconditional. It does not rule out correspondences on
$B\times D$ which are not assembled from these graph maps.

\textbf{Proof.}
Choose a geometric base point $b$ on $B$. Let
$a_B:B\longrightarrow\operatorname{Alb}(B)$ be the pointed Albanese
map and let $u:\operatorname{Alb}(B)\longrightarrow X$ be induced by
$x\mapsto\kappa(x)-\kappa(b)$.
Weak Lefschetz in degree one identifies $H^1(X)$ with $H^1(B)$.
The identification of degree-one cohomology with the dual of the
Albanese Tate module therefore makes $u$ an isogeny.
After subtracting the constant $f(b)$, the morphism $f$ factors as
$v a_B$. Since $u$ is invertible in the isogeny category, there are
an integer $N>0$ and a homomorphism $t:X\longrightarrow D$ such that
$t u=[N]_D v$. Thus
\[
 [N]_D\circ f=t\circ\kappa+\text{a constant}.
\]
Translations on an abelian variety act trivially on cohomology.
Functoriality of proper pushforward gives
\[
 [N]_{D,*}f_*w=t_*\kappa_*w=0.
\]
Multiplication by $N$ induces an invertible map on rational
$\ell$-adic cohomology, so $f_*w=0$. The orthogonality assertion is
projection formula. A rational map from a smooth variety to an
abelian variety extends by Weil's extension theorem; one exact
modern locator is Milne, \emph{Abelian Varieties}, course notes
version 2.00, I, Theorem 3.2, printed p.~17 (PDF p.~23).\hfill$\square$

\paragraph{Programme location and scope.}
This specifies a limit of the proposed curve-test attack on
QH9-axis. It concerns the incidence higher Picard framework of
cote 44, Th.~1--4, scan 002, and the incidence divisor on scans
008--009. It complements JAC-P19 (no abelian rational domination)
and the already proved annihilation calculation in H9-T13, rather
than replacing either. An algebraic correspondence shifts
cohomological degree by an even number, so $H^1(C)$ by itself
cannot map into $H^{2r}(B)$ through an algebraic correspondence.
The elliptic tensor construction $w\boxtimes P_E$ resolves that
parity issue but its ambient pushforward is zero, exactly as
recorded in H9-T13. The missing input is a genuinely new
pairing-compatible correspondence, not a graph of another
Albanese map.

\subsection*{C44-A03. The finite-field test is absence of potential Tate roots, not one exceptional slope}

\textbf{Proposition.}
Suppose the axis and the embedding are defined over $\mathbf F_q$.
Let
\[
 V_B=\ker\bigl(\kappa_*:H^{2r}(B_{\overline{\mathbf F}_q},
 \mathbf Q_\ell)\longrightarrow H^{2r+4}(X_{\overline{\mathbf F}_q},
 \mathbf Q_\ell)(2)\bigr)
\]
be the untwisted primitive variable cohomology, with the twists in
the target understood by Gysin compatibility. Write $F$ for
geometric Frobenius and $P_B(t)=\det(t-F\mid V_B)$.
Every algebraic axis class is contained in
\[
 T_B=\bigcup_{a\ge1}\ker(F^a-q^{ra}\mid V_B).                \tag{C44.A03.1}
\]
In particular all its normalized Newton slopes are $r$.
If $P_B$ has no root $q^r\zeta$ for a root of unity $\zeta$,
then $A_B=0$.
For $D=\deg P_B>0$, a finite exact sufficient test is
\[
 \gcd\bigl(P_B(t),\ t^N-q^{rN}\bigr)=1,
 \qquad N=\operatorname{lcm}(1,2,\ldots,2D^2).               \tag{C44.A03.2}
\]
Here $P_B\in\mathbf Q[t]$ is obtained by dividing the middle
Frobenius polynomial of $B$ by that of the restricted ambient
summand. No Tate conjecture and no Frobenius semisimplicity
conjecture is needed for this vanishing criterion.

\textbf{Proof.}
A rational algebraic cycle on $B_{\overline{\mathbf F}_q}$ is
defined over a finite extension $\mathbf F_{q^a}$. Its codimension
$r$ cycle class in untwisted cohomology satisfies $F^a w=q^{ra}w$.
This proves (C44.A03.1). In crystalline realization, the same
cycle is a Frobenius eigenclass of slope $r$; equivalently every
root $q^r\zeta$ has normalized valuation $r$ at every prime over
$p$.
If the polynomial has no such root, every kernel in (C44.A03.1)
is zero, which proves the asserted vanishing.

The splitting (C44.A01.3) is Frobenius-compatible after the
polarization is defined over $\mathbf F_q$; its image is isomorphic
to $H^{2r}(X)$. Thus the usual rational Frobenius polynomials of
$X$ and $B$ give $P_B\in\mathbf Q[t]$ without asserting
rationality of a chosen algebraic basis of $A_B$.
If $q^r\zeta$ is a root and $n$ is the order of $\zeta$, the
cyclotomic polynomial of $\zeta$ divides
$q^{-rD}P_B(q^rt)$, so $\varphi(n)\le D$.
For $n=\prod p^a$,
\[
 \frac{\varphi(n)^2}{n}=\prod_{p^a\parallel n}p^{a-2}(p-1)^2
 \ge\frac12:
\]
each factor for $p\ge3$ is at least one, and the sole possible
factor less than one is $1/2$ for $2^1$. Consequently $n\le2D^2$,
so $n\mid N$. This proves (C44.A03.2), including its equivalence
to absence of roots of this form. For $D=0$ the conclusion is
immediate.\hfill$\square$

\paragraph{Why H9-T5 does not supply this test.}
JAC1964 H9-T5 is a theorem about the \emph{odd} variable local
system $R^{2m-1}$ of the pencil and its variable incidence abelian
factor. Its irreducibility hypothesis turns a nonzero injection
$H^1(D)(-(m-1))$ into an isomorphism; one slope outside
$[m-1,m]$ then contradicts that isomorphism. The even-degree
fixed-axis space in (C44.A03.1) is a different representation.
At a closed point the arithmetic Galois group is procyclic, and
an isolated algebraic class need not extend to a flat algebraic
class on the generic axis family. No family-monodromy
irreducibility statement can therefore be substituted into
(C44.A03.1) without a separate spreading theorem.
An out-of-range slope in some other direct summand is compatible
with a slope-$r$ Tate summand. The correct stronger sufficient
Newton condition here is absence of \emph{every} slope-$r$
subspace, while the cyclotomic criterion above is sharper.
This repairs the proposed use of H9-T5 without producing a
witness for every special axis.

\paragraph{Supersingular ambient varieties do not resolve this distinction.}
If $X$ is supersingular, $H^{2r}(X)$ has only slope $r$ because
$H^1(X)$ has slope $1/2$ and the cohomology algebra is exterior.
That identifies the slopes of the ambient summand in
(C44.A01.3), not the slopes or cycle classes of $V_B$.
A nonliftable algebraic class pulled back from $X$ lies in the
ambient summand and is removed by the projector $1-e$. Thus
ambient nonliftability alone gives no nonliftable element of
$A_B$.

\subsection*{C44-A05. The embedded axis itself always has a characteristic-zero lift}

\textbf{Proposition.}
Let $k=\overline{\mathbf F}_p$, let $X/k$ be an abelian variety,
let $L$ be ample, and let $s_0,s_1\in H^0(X,L)$ have a smooth
codimension-two common zero scheme $B$. There is a mixed-characteristic
discrete valuation ring $R$, with finite residue field contained in $k$,
an abelian scheme $\mathcal X/R$, a relatively ample line bundle
$\mathcal L$, and two sections whose common zero scheme
$\mathcal B/R$ is smooth and projective, such that the geometric
special fibre is the given $(X,L,s_0,s_1,B)$. In particular this
applies to every axis in QH9-axis, including supersingular ambient
abelian varieties. The proposition asserts no lifting of a specified
middle-dimensional algebraic class on $B$.

\textbf{Proof.}
Descend the given data to a finite subfield of $k$. The polarized
lifting theorem of Mumford--Norman--Oort gives an abelian scheme
with polarization lifting $(X,\lambda_L)$ over an integral
mixed-characteristic base. An exact statement is Oort,
\emph{Did Earlier Thoughts Inspire Grothendieck?}, Theorem (8.3).1,
p.~29. By finite-presentation descent we may use an integral base
$S$ of finite type over $\mathbf Z$, with a closed point of finite
residue field giving the prescribed geometric special fibre. The
scheme, group laws, and polarization all descend; passing to a
sufficiently large finite-type model preserves smoothness and
properness. This is the usual limit construction, as in Stacks
Project, Lemma 32.10.1, tag 01ZM, Lemma 32.8.9, tag 0C0C,
and Lemma 32.13.1, tag 081F.

The chosen line bundle also lifts after an \'etale neighborhood of
the special point: Olsson, \emph{Compactifications of moduli of
abelian varieties: an introduction}, Lemma 2.1.8 and its proof,
pp.~299--300. Its proof preserves the chosen special-fibre line
bundle. Alternatively, after choosing a representative of the
lifted polarization, the difference from $L$ lies in the dual
abelian variety and lifts after a smooth base change. Take an
integral component through the chosen point and still dominating
$S$; this is possible by flatness of the \'etale neighborhood.
Write $\pi:\mathcal X\to S$ for the resulting family and
$\mathcal L$ for its ample line bundle. For an ample line bundle
on an abelian scheme,
\[
 R^i\pi_*\mathcal L=0\quad(i>0),\qquad
 E=\pi_*\mathcal L\ \text{is locally free and commutes with base change}.
\]
The exact family statement is Olsson, \S2.1.4, p.~299.

The total space $W$ of the vector bundle $E\oplus E$ carries a
universal pair of sections on $\mathcal X\times_S W$. It is
integral of characteristic zero and of finite type over
$\mathbf Z$. The given pair determines a closed point $b\in W$
after a finite residue-field extension. Let $\mathcal Z$ denote
the universal common zero scheme. On $\mathcal Z$, failure of the
two relative differentials to have rank two is closed. Its image
in $W$ is closed because $\mathcal Z\to W$ is proper. The
complement $U$ contains $b$, since $B$ is smooth of codimension
two. The relative Jacobian criterion makes $\mathcal Z_U\to U$
smooth of relative dimension $\dim X-2$, and it is projective.
Thus $U$ is an integral finite-type mixed-characteristic
parameter scheme carrying a smooth projective lift of the
specified geometric fibre.

For completeness one can pass to an arithmetic trait without
changing that geometric fibre. Let $A=\mathcal O_{U,b}$, with
maximal ideal $\mathfrak m$, and put $d=\dim A$. This is an
excellent local domain, $0\ne p\in\mathfrak m$, and its residue
field is finite. Choose $x_1,\ldots,x_{d-1}$ so that
$(p,x_1,\ldots,x_{d-1})$ is $\mathfrak m$-primary. Set
$I=(x_1,\ldots,x_{d-1})$. Some minimal prime $P$ of $I$ does
not contain $p$: otherwise $p\in\sqrt I$ would give
$\sqrt I=\mathfrak m$, contrary to Krull's height theorem.
The local domain $A/P$ has dimension one, since $p$ is nonzero
and generates an ideal with maximal radical in that quotient.
Its normalization is finite by excellence; localization at a
point above its closed point is a discrete valuation ring $R$
with characteristic-zero fraction field and finite residue field.
Pulling back the family along
$\operatorname{Spec}R\to\operatorname{Spec}A\to U$ proves the
claim. The residue extension embeds in $k$, so the geometric
special fibre is the prescribed one.\hfill$\square$

\paragraph{Programme location and remaining obstruction.}
This completes the geometric lifting substep of QH9-axis(a),
in the schemes-over-$\mathbf Z$ framework. The unresolved issue
is surjectivity, for some such lift, of specialization on the
specified rational algebraic middle classes. Smooth proper base
change is an isomorphism on cohomology; it is not a lifting
theorem for cycles. The sign transfer to classes which do lift
is exactly the already settled H9-T8a specialization argument.
Cote 44, Th.~5, scan 004, and the primitive margin on scan 008
therefore retain the same cycle-level gate.

\section{Unresolved axis steps after the ordered attacks}

\paragraph{Lifting.}
For any specified smooth projective mixed-characteristic lift
of $(B,h_B)$ and any rational algebraic class on the geometric
generic fibre specializing to $w\in A_B$, compatibility of
cycle classes with smooth proper specialization carries
$L_Bw=0$ back to the same primitive relation on the generic
fibre. The specialization is an isomorphism on $\ell$-adic
cohomology, so a nonzero $w$ has nonzero generic class, and the
intersection number is unchanged. H9-T8a then gives
$(-1)^r\int_Bw^2>0$. This is the carried lifting argument,
not a surjectivity theorem for Chow specialization.

No proof that every $w\in A_B$ lifts with $B$ is obtained.
Nor is a specific cohomological class in $A_B$ proved
nonliftable in this contribution. Attempts using nonliftable
ambient Tate classes fail because their projection to $A_B$
is zero. Agugliaro, \emph{Examples for the standard conjecture
of Hodge type}, Doc. Math. 31 (2026), Proposition 5.12,
pp.~876--877, supplies a rank-two nonliftable Tate block for
CM lifts of certain abelian varieties; Remark 3.5, p.~870,
explains the passage to arbitrary abelian lifts. Even after
032 supplies algebraicity, these remain ambient examples.
They are not examples for the axis kernel.

\paragraph{Correspondences.}
C44-A02 and JAC1964 H9-T13 eliminate graph/ambient transfers.
A general correspondence from $B$ to an abelian variety could
still act nontrivially, but a splitting preserving the signed
geometric pairing has not been constructed. Merely splitting
an abstract Tate class from $h(B)$ would not identify that
pairing with a positive form.

\paragraph{Newton slopes.}
C44-A03 gives the correct finite-field vanishing test.
No Frobenius polynomial satisfying that test has been exhibited
for every special axis, and axes with $A_B\ne0$ necessarily
fail it. A slope witness in the odd pencil variation does not
settle the even fixed-axis sign.

\paragraph{Current gate.}
The unresolved assertion remains weak nonnegativity of
$(-1)^r\int_Bw^2$ for every algebraic $w\in\ker\kappa_*$.
Conditional on 032/Milne, C44-A01 identifies it with the weak
Hodge-standard assertion for $B$. No wrong-sign cycle is
constructed. No main question is marked closed by these
route boundaries.

% END work/axis/axis_contribution.tex

% BEGIN work/axis/axis_correspondence_obstruction.tex
% C44_POS /root/axis_independent; GPT-6 Astra Pro; effort not exposed.
\subsection*{C44-A02b. Divisor-generated correspondence tests annihilate the axis kernel}

\textbf{Proposition.}
Let $\kappa:B\hookrightarrow X$ be a smooth Lefschetz axis in an
abelian $(2r+2)$-fold, $r\ge2$, and let $A$ be any abelian variety.
Every rational divisor cohomology class on $B\times A$ is the
restriction of a rational divisor cohomology class on $X\times A$.
Hence every class
in the algebra generated by these divisors is an ambient
restriction. Any algebraic correspondence from $B$ to $A$ whose
cohomology class is such a restriction acts as zero on $A_B$.
In particular this holds for every correspondence class in the
divisor-generated algebra of $B\times A$.

\textbf{Proof.}
Choose $b_0\in B$ and the origin of $A$. For a line bundle $D$
on $B\times A$, remove its restrictions from the two factors.
The remaining normalized line bundle is a family of
algebraically trivial line bundles on $A$, parametrized by $B$,
and therefore is the pullback of the Poincar\'e bundle along
$f\times1_A$, for a morphism $f:B\to\widehat A$ with
$f(b_0)=0$. This is the defining universal property of the dual
abelian variety.

The map $\operatorname{Alb}(B)\to X$ induced by
$b\mapsto\kappa(b)-\kappa(b_0)$ is an isogeny, as used and
proved cohomologically in C44-A02. After multiplying by a
positive integer $N$, the map $[N]f$ therefore extends to a
morphism $X\to\widehat A$. Pullback of the Poincar\'e bundle
gives a line bundle on $X\times A$ restricting to the $N$-th
tensor power of the normalized part of $D$.
The factor from $A$ already extends. The factor from $B$ has
a rational algebraic divisor-class extension to $X$ by the
low-degree algebraic inverse in C44-A01, with $j=1$.
These three extensions prove the assertion for rational
cohomology classes of divisors. Taking products proves the
assertion about the divisor-generated algebra.

For the vanishing, write $\rho=\kappa\times1_A$ and
$[\Gamma]=\rho^*\gamma$ for an ambient cohomology class
$\gamma$ of the appropriate degree. For $w\in A_B$, the
projection formula gives
\[
 \begin{split}
 \Gamma_*w
 &=p_{A,*}\rho_*\bigl(p_B^*w\,\rho^*\gamma\bigr)\\
 &=p_{A,*}\bigl(p_X^*(\kappa_*w)\,\gamma\bigr)=0.
 \end{split}
\]
Here the first $p_A$ on the right is the projection from
$X\times A$. This calculation in fact proves the stated
vanishing for every ambient-restriction correspondence, without
the divisor-generation hypothesis.\hfill$\square$

\paragraph{Scope.}
This is a constraint on proposed curve tests in the manner of
cote~44 Th.~4 (scan~002), when used to attack the axis term of
Th.~5 (scan~004). It extends the graph-map obstruction of
C44-A02. It does not exclude general correspondences. In
particular, an extending injection of the kind in C44-A08 must
use correspondence classes involving the new middle
cohomology of $B$; ambient divisor and Poincar\'e operations
alone cannot supply it. It does not assert that any such new
correspondence is nonalgebraic or nonexistent. The universal
Poincar\'e property used above is Milne, \emph{Abelian Varieties},
course notes v.~2.00 (2008), I \S8, printed pp.~36--37
(definition of the dual abelian variety and universal sheaf);
the other inputs are C44-A01 and
C44-A02, already checked in this run.

% END work/axis/axis_correspondence_obstruction.tex

% BEGIN work/axis/axis_forced_tate_root.tex
% C44_POS /root/axis_independent; GPT-6 Astra Pro; effort not exposed.
\subsection*{C44-A03b. A forced central Frobenius root in the first open dimension}

\textbf{Theorem.}
Let $X/\mathbf F_q$ be an abelian variety of dimension $2r+2$,
$r\ge2$, with an ample line bundle $L$ of class $h$. Let
$B\subset X$ be a smooth complete intersection of two members
of $|L|$, defined over $\mathbf F_q$, and put
\[
 V_B=\ker\bigl(\kappa_*:H^{2r}(B,\mathbf Q_\ell)
       \longrightarrow H^{2r+4}(X,\mathbf Q_\ell)(2)\bigr),
 \qquad d=\int_Xh^{2r+2}.
\]
Then
\[
 \begin{split}
 \dim V_B
 &=(2r+1)d+2\binom{4r+3}{2r-1}-\binom{4r+4}{2r}\\
 &=(2r+1)d-\frac1{r+1}\binom{4r+4}{2r}.
 \end{split}
 \tag{C44.A03b.1}
\]
This dimension is odd exactly when $r=2^s-2$ for an integer
$s\ge2$. In these dimensions the Frobenius on $V_B$ has an
eigenvalue $+q^r$ or $-q^r$. In particular the potential Tate
eigenspace is nonzero after an extension of degree at most two,
and $V_B$ has a Newton slope $r$.

For an abelian sixfold, $r=2$, formula (C44.A03b.1) is
\[
                  \dim V_B=5d-165,
\]
which is always odd. Thus the no-potential-Tate-root criterion
C44-A03 can never prove $A_B=0$ for any axis in the first open
dimension. The same obstruction occurs for ambient dimensions
$6,14,30,62,\ldots$.

\textbf{Proof.}
The tangent sequence of the smooth complete intersection gives
\[
 c(T_B)=(1+h|_B)^{-2},\qquad
 \chi(B)=\int_Bc_{2r}(T_B)=(2r+1)d.
\]
Here the Euler/Chern identity can be read entirely in Weil
cohomology: the self-intersection of the diagonal is both the
alternating sum of the Betti numbers, by K\"unneth and Poincar\'e
duality, and the degree of the top Chern class of its normal
bundle $T_B$, by self-intersection (Fulton, \emph{Intersection
Theory}, \S6.3). No characteristic-zero comparison is needed.

Weak Lefschetz gives $b_i(B)=b_i(X)$ for $i<2r$, and Poincar\'e
duality determines the higher Betti numbers from those lower
ones. Since $b_i(X)=\binom{4r+4}{i}$, we obtain
\[
 b_{2r}(B)
 =\chi(B)-2\sum_{i=0}^{2r-1}(-1)^i\binom{4r+4}{i}
 =(2r+1)d+2\binom{4r+3}{2r-1}.
\]
The restriction of $H^{2r}(X)$ is injective and splits off
orthogonally, by $\kappa_*\kappa^*=L_h^2$ and hard Lefschetz.
Subtracting its dimension proves the first equality of
(C44.A03b.1); the second is a binomial identity.
The exterior-algebra description of abelian cohomology is
Milne, \emph{Abelian Varieties}, course notes v.~2.00 (2008),
I Theorem~12.1 (printed pp.~55--56).

Riemann--Roch on $X$ gives $d=(2r+2)!\chi(L)$, so $d$ is even.
Reducing the first formula modulo two gives
\[
 \dim V_B\equiv\binom{4r+4}{2r}
                 \equiv\binom{2r+2}{r}\pmod2.
\]
The second congruence follows from $(1+x)^2=1+x^2$ over
$\mathbf F_2$. More generally, expanding
$(1+x)^n=\prod_{j:\,n_j=1}(1+x^{2^j})$ in binary shows that
$\binom{2r+2}{r}$ is odd exactly when adding $r$ and $r+2$
has no binary carries, equivalently $r\mathbin{\&}(r+2)=0$.
This requires $r$ to be even. Writing $r=2k$ reduces it to
$k\mathbin{\&}(k+1)=0$, which says that $k$ consists solely
of trailing binary ones: $k=2^{s-1}-1$.
Thus the stated parity criterion holds.

The middle intersection form restricts nondegenerately to $V_B$
by the same orthogonal splitting. Frobenius $F$ is a similitude
with multiplier $q^{2r}$, so $q^{-r}F$ is an orthogonal
transformation of $V_B$. Over an algebraic closure of the
coefficient field its eigenvalues occur in reciprocal pairs,
counted with multiplicity. In odd dimension at least one
eigenvalue is fixed by inversion, hence is $+1$ or $-1$.
This argument does not assume semisimplicity. An eigenvalue
always has a nonzero actual eigenspace, and its square equals
one. Consequently $F^2$ has eigenvalue $q^{2r}$ on a nonzero
vector in $V_B$, as claimed. The same characteristic polynomial
occurs in crystalline cohomology: apply Katz--Messing,
\emph{Some consequences of the Riemann hypothesis for varieties
over finite fields}, Invent.~Math.~23 (1974), Theorem~1(i) and
the preceding observation that crystalline cohomology satisfies
its hypotheses, p.~74, to $B$ and $X$, and take their middle
quotient. The corresponding root has normalized $p$-adic
valuation $r$, so it is also a forced central Newton
slope.\hfill$\square$

\paragraph{What this does and does not settle.}
No algebraicity of the forced eigenvector is asserted. If the
Tate conjecture for codimension-$r$ cycles on $B$ holds over
the relevant finite extension, the algebraic variable projector
$1-\kappa^*(L_h^2)^{-1}\kappa_*$ shows that $A_B\ne0$.
That is a further hypothesis about $B$, not a consequence of
032 on the ambient abelian variety. The theorem is compatible
with geometric-generic vanishing in JAC1964 JAC-P09: the
geometric generic parameter is not a finite-field closed point.

For route~(c) of QH9-axis, the theorem rules out the full
fixed-axis variable-part vanishing certificate in these
dimensions. A slope outside the interval in H9-T5 can coexist
with this forced central root. It remains possible to isolate
a smaller factor by an independently justified projector, but
no such universal factor is supplied here. In cote~44 terms,
this constrains the proposed Newton route to Th.~5, scan~004,
and its primitive axis summand; it does not decide its sign.

% END work/axis/axis_forced_tate_root.tex

% BEGIN work/axis/root_explicit_axes.tex
\section{An explicit nonzero axis subspace in every characteristic}

\subsection*{C44-A04. Axes containing fibres of an abelian product}
\begin{theorem}[C44-A04]\label{thm:product-axis}
Let $r=m-1\ge2$, let $S$ and $T$ be abelian varieties of dimensions
$r$ and $r+2$, and put $X=S\times T$. Choose very ample classes
$s$ and $t$ on $S$ and $T$ and let $h=s+t$ on $X$. Let
$B\subset X$ be a smooth complete intersection of two divisors of
class $h$. Suppose that it contains the distinct fibres
$S_i=S\times\{t_i\}$ for $1\le i\le N$. Put
\[
 d_S=\int_Ss^r,\qquad d_T=\int_Tt^{r+2},\qquad
 a=\frac1{d_T}\sum_{j=0}^r(-1)^j(j+1)s^jt^{r-j},
 \qquad w_i=[S_i]_B-\kappa^*a.
 \tag{C44.A04.1}
\]
Then $w_i\in A_B$ and the exact signed Gram matrix is
\[
 \left((-1)^r\int_B w_iw_j\right)_{i,j}
       =(r+1)d_S\left(I_N-\frac1{d_T}{\bf1}{\bf1}^{t}\right).
 \tag{C44.A04.2}
\]
Smoothness of $B$ forces $N\le d_T$. Hence this matrix is positive
semidefinite in every characteristic, without 032. If $N<d_T$ it is
positive definite, and the $w_i$ are linearly independent rational
algebraic classes. In particular a single contained fibre supplies a
nonzero axis class of the required sign.

Smooth such axes containing a prescribed fibre exist after replacing
$s,t$ by sufficiently divisible positive multiples. The two defining
sections may be chosen to give a genuine Lefschetz pencil. Thus the
construction includes products of supersingular elliptic curves and
gives special axes with $A_B\ne0$ over $\overline{\F}_p$ for every $p$.
The assertion concerns the displayed subspace, not all of $A_B$.
\end{theorem}
\begin{proof}
All computations are in rational Weil cohomology, with their usual
cycle representatives. Since $s^{r+1}=0$ and $t^{r+3}=0$, the finite
identity $(1+z)^2\sum_{j=0}^r(-1)^j(j+1)z^j\equiv1\pmod{z^{r+1}}$
gives
\[
                      h^2a=t^{r+2}/d_T=[S_i]_X.
\]
Thus $\kappa_*w_i=0$ by self-intersection, and $w_i\in A_B$.
The normal bundle of $S_i$ in $X$ is trivial of rank $r+2$.
Because $B$ is smooth along $S_i$, its normal sequence is exact:
\[
 0\longrightarrow N_{S_i/B}\longrightarrow\mathcal O_S^{r+2}
   \longrightarrow L_S\oplus L_S\longrightarrow0.
\]
Here $c_1(L_S)=s$. Consequently
$c(N_{S_i/B})=(1+s)^{-2}$ and
\[
                 \int_B[S_i]^2=(-1)^r(r+1)d_S.
\]
Distinct fibres are disjoint, so their intersection is zero.
Restriction of $a$ to any $S_i$ has degree
$(-1)^r(r+1)d_S/d_T$; the same is true of
$\int_B(\kappa^*a)^2=\int_Xa h^2a$.
Expanding each $w_iw_j$ therefore gives (C44.A04.2).

We prove the bound on the number of contained fibres without using a
positivity conjecture. Write the two defining sections in a basis of
$H^0(S,L_S)$ as
$F_\nu=\sum_\alpha\sigma_\alpha\tau_{\nu\alpha}$,
where $\tau_{\nu\alpha}\in H^0(T,L_T)$ and $\nu=1,2$.
Let $Z\subset T$ be the common zero scheme of all coefficients
$\tau_{\nu\alpha}$. Its geometric points are exactly the fibres
contained in $B$. Suppose $T_zZ$ has positive dimension at a point
$z\in Z$. Differentiating the coefficients shows that the normal
surjection along $S\times\{z\}$ factors through
\[
       \mathcal O_S^{r+1}\longrightarrow L_S\oplus L_S
\]
(or through a trivial bundle of still smaller rank). Its kernel
would have rank at most $r-1$ but total Chern class $(1+s)^{-2}$.
Its $r$th Chern class would therefore be both zero, by rank, and
$(-1)^r(r+1)s^r$, whose degree is nonzero. This is impossible.
Thus $Z$ has zero tangent space at every geometric point; it is a
finite reduced scheme. Choose $r+2$ general linear combinations of
the coefficient sections. Since their common base scheme is finite,
successive dimension cuts yield a zero-dimensional complete
intersection containing $Z$. Its degree is $d_T$, so
$N\le\#Z\le d_T$.
The eigenvalues of the matrix in parentheses in (C44.A04.2) are
$1$ with multiplicity $N-1$ and $1-N/d_T$ with multiplicity one.
This proves all the asserted signs and, for $N<d_T$, independence.

For existence take the prescribed smooth centre $S_0=S\times\{0\}$.
For sufficiently positive $L_S,L_T$, the linear system
$H^0(X,I_{S_0}\otimes(L_S\boxtimes L_T))$ separates second jets
at each point off $S_0$, separates the first jets at each pair of
distinct points off $S_0$, and maps onto the conormal first jets
along $S_0$. These are uniform assertions for all such points:
apply relative Serre vanishing to the universal length-bounded
jet subschemes and the fixed coherent ideal $I_{S_0}$ (or use
high powers of product sections). A common sufficiently high
power satisfies all three conditions.
The latter assertion also follows directly from product sections:
sections of $L_S$ generate its fibres and sections of $I_0L_T$
generate the cotangent space of $T$ at $0$ after taking a high power.
For two general sections, the locus on $S_0$ where their normal
derivatives have rank at most one has expected codimension $r+1$;
the evaluation to the two normal covectors is surjective at each
point, so the incidence dimension estimate, with $\dim S_0=r$,
makes this locus empty. Away from $S_0$, jet separation and
Bertini give a smooth codimension-two intersection. Here are the
dimension counts for the pencil condition. Write $n=\dim X=2r+2$.
At a specified point off $S_0$, singularity imposes $n+1$
independent first-jet equations on a section. Degeneracy of its
quadratic part imposes an additional proper equation, by the
second-jet surjectivity. Letting the point vary gives a bad locus
of codimension at least two in the projective linear system.
Singularities at two specified distinct points impose $2(n+1)$
independent equations; their parameter space has dimension $2n$,
again giving codimension at least two. For singularity along the
imposed centre, the normal covector must vanish; this imposes
$r+2$ conditions at a specified point of $S_0$, and $\dim S_0=r$.
Its incidence image has codimension at least two as well. A general
projective line avoids these bad loci. Smoothness of its axis
also prevents the line from being tangent to the discriminant at
a single ordinary quadratic point: such tangency says that the
other pencil section vanishes at that point, putting a singular
member's zero differential into the axis's defining pair.
Thus the line gives a Lefschetz pencil, with distinct singular
values and no singularity on $S_0$. Since $n$ is even, the
nondegenerate quadratic forms are a nonempty open also in
characteristic two. These are the ordinary Lefschetz local models.
All conditions used are nonempty Zariski-open conditions over the
algebraically closed field. Over $\overline{\F}_p$ the chosen
sections are defined over a finite extension. This proves existence.
\end{proof}

\paragraph{C44-A04a: a concrete sixfold normalization.}
Take $X=E^6$, $S=E^2$, $T=E^4$ and
$h=d\sum_{i=1}^{6}[0]_i$, with $d$ sufficiently large for the
construction. For an axis containing $E^2\times\{0\}$,
\[
 d_S=2d^2,\qquad d_T=24d^4,\qquad
 \int_Bw^2=6d^2-\frac1{4d^2}>0.
\]
The integral cycle $\widetilde w=d_T[S]_B-
\kappa^*(t^2-2st+3s^2)$ has
\[
        \int_B\widetilde w^2=144d^6(24d^4-1)>0.
\]
By H9-T7 the corresponding height classes have these exact heights
in its normalization. One may choose $E$ supersingular in every
characteristic. These are positive special-axis classes, not a
classification of nonliftable classes and not a wrong-sign example.

\paragraph{Archive and source location.}
The height interpretation is JAC1964 H9-T7, (H9.18), for cote 44,
Th.~5, scan 004. The primitive qualification is the margin on scan
008. The normal-bundle and self-intersection calculations use Fulton,
\emph{Intersection Theory}, Section 6.3, self-intersection formula,
and the codimension-two blowup normalization already checked in
H9-T7. The existence argument uses the ordinary jet construction
underlying SGA~7~II, XVII, 1.1--1.2 (pp.~213--214), 2.5
(p.~217), 3.2.1--3.2.4 (p.~219), 3.3 (p.~221), 3.7
(p.~234), and 4.2--4.3 (pp.~236--239); its imposed-centre
incidence argument is provided here, not attributed as a separate
statement to SGA. In cote 16 language the calculation verifies $I$
on a specified algebraic primitive subspace. It does not establish
$I(B)$ on its complement or close QH9-axis.

% END work/axis/root_explicit_axes.tex

% BEGIN work/axis/axis_boundary_relation.tex
% C44_POS /root/axis_independent; GPT-6 Astra Pro; effort not exposed.
\subsection*{C44-A04b. The equality case of the product-fibre Gram bound}

\textbf{Proposition.}
In C44-A04, suppose the smooth axis $B\subset S\times T$ contains
$N=d_T$ distinct fibres $S_i=S\times\{t_i\}$, where
$\dim S=r$, $\dim T=r+2$, $s=c_1(L_S)$,
$t=c_1(L_T)$ and $d_T=\int_Tt^{r+2}$. Put
\[
 a=\frac1{d_T}\sum_{j=0}^r(-1)^j(j+1)s^jt^{r-j},
 \qquad w_i=[S_i]_B-\kappa^*a.
\]
Then
\[
 \sum_{i=1}^{d_T}w_i=0
 \tag{C44.A04b.1}
\]
already in the rational Chow group of $B$, for these displayed
representatives. This is the only linear relation in cohomology.
Thus the signed H9-T7 form is positive definite on the actual
span of all the $w_i$, of dimension $d_T-1$. Together with
C44-A04, the form is positive definite on the span of any
collection of contained product fibres, including the boundary
case $N=d_T$.

\textbf{Proof.}
Let $W\subset H^0(T,L_T)$ be the span of the coefficient sections
of the two equations of $B$ in a basis of $H^0(S,L_S)$.
C44-A04 proves that their common zero scheme $Z\subset T$ is
finite and reduced and that its length is at most $d_T$.
It contains the $d_T$ given points, so these are exactly $Z$.
Choose $g=r+2$ general sections $t_1,\ldots,t_g\in W$.
They form a proper zero-dimensional complete intersection on
$T$, of length $d_T$, containing $Z$ scheme-theoretically.
Hence their common zero scheme equals $Z$.

The Koszul resolution of $I_ZL_T$ has terms
$\bigwedge^j k^g\otimes L_T^{1-j}$ for $1\le j\le g$.
For $j\ge2$, Serre duality and ample-line-bundle vanishing on an
abelian variety give
$H^i(T,L_T^{1-j})=0$ for $i<g$. Consequently its degree-zero
cohomology gives
\[
 H^0(T,I_ZL_T)=\langle t_1,\ldots,t_g\rangle.
 \tag{C44.A04b.2}
\]
The same sections are linearly independent, since otherwise
$g-1$ equations would have a zero-dimensional common zero scheme
on a $g$-fold, contrary to the height theorem. Every coefficient
of the equations of $B$ belongs to (C44.A04b.2). Thus these two
equations can be written as $U(t_1,\ldots,t_g)^t$, with a matrix
of sections on $S$ giving
\[
 U:\mathcal O_S^{g}\longrightarrow L_S^{\oplus2}.
\]
This map is surjective everywhere. Indeed, at any point of the
reduced complete intersection $Z$, the differentials of the
$t_j$ form a basis of the conormal space, and smoothness of $B$
along the corresponding entire fibre $S\times\{t_i\}$ says
precisely that $U$ has rank two everywhere on $S$.

Let $K=\ker U$, a rank-$r$ bundle with
$c(K)=(1+s)^{-2}$. On $B$, the vector of the $t_j$ is a section
of $K\boxtimes L_T$. Its zero scheme is exactly $S\times Z$.
Since it has codimension $r$ in the smooth variety $B$, this is
a regular section. The top Chern class formula therefore gives
\[
 \begin{split}
 \sum_{i=1}^{d_T}[S_i]_B
 &=c_r(K\boxtimes L_T)\\
 &=\kappa^*\sum_{j=0}^r(-1)^j(j+1)s^jt^{r-j}
 =d_T\kappa^*a.
 \end{split}
\]
This proves (C44.A04b.1) in Chow. By C44-A04 the signed Gram
matrix is $(r+1)d_S(I-J/d_T)$, whose nullspace at $N=d_T$ is
exactly the constant vector. Any cohomological relation is in
this nullspace. The exhibited relation is therefore the only
one, and the induced form on the quotient is positive definite.
For fewer than $d_T$ fibres the matrix itself is positive
definite by C44-A04.\hfill$\square$

\paragraph{Inputs and scope.}
The ample vanishing and duality used in (C44.A04b.2) are the
standard abelian-variety index theorem; see Milne,
\emph{Abelian Varieties}, course notes v.~2.00 (2008),
I Theorem~11.1(c), printed p.~54. The regular-section formula is
Fulton, \emph{Intersection Theory}, \S14.1. The apparent isotropic
vector in the boundary Gram matrix therefore cannot supply a
nonliftable nonzero axis class. This sharpens the explicit
positive subspace of cote~44 Th.~5, scan~004, without deciding
QH9-axis on its complement.

% END work/axis/axis_boundary_relation.tex

% BEGIN work/axis/axis_asymptotic_centres.tex
% C44_POS /root/axis_independent; GPT-6 Astra Pro; effort not exposed.
\section{Prescribed smooth centres: an unconditional asymptotic axis sign}

\subsection*{C44-A06. Fixed smooth centres become positive after increasing the polarization}

\textbf{Theorem.}
Let $X$ be an abelian variety of dimension $2r+2$, $r\ge2$, over
an algebraically closed field of arbitrary characteristic. Fix an
ample line bundle $L_0$ with class $h_0$, and pairwise disjoint
smooth closed subvarieties $Z_1,\ldots,Z_s\subset X$ of dimension
$r$. For every sufficiently large integer $N$ there is a Lefschetz
pencil in $|L_0^{\otimes N}|$ whose smooth axis $B_N$ contains every
$Z_i$. Put $\kappa:B_N\hookrightarrow X$ and
\[
 a_i=(L_{h_0}^{2})^{-1}[Z_i]_X,
 \qquad w_i=[Z_i]_{B_N}-N^{-2}\kappa^*a_i.
 \tag{C44.A06.1}
\]
These are rational algebraic classes and $w_i\in A_{B_N}$. Define
\[
 \begin{split}
 P_i(N)&=\sum_{j=0}^r(-1)^j(r-j+1)N^{r-j}
       \int_{Z_i}c_j(N_{Z_i/X})h_0^{r-j},\\
 G_{ij}&=\int_X a_i a_jh_0^2.
 \end{split}
 \tag{C44.A06.2}
\]
The exact signed Gram matrix is
\[
 \left((-1)^r\int_{B_N}w_iw_j\right)_{i,j}
       =\operatorname{diag}(P_i(N))-(-1)^rN^{-2}G.
 \tag{C44.A06.3}
\]
In particular it is positive definite for all sufficiently large
$N$, independently of 032. More explicitly, let
\[
 \begin{split}
 d_i&=(r+1)\int_{Z_i}h_0^r>0,\qquad d_{\min}=\min_i d_i,\\
 S_i&=\sum_{j=1}^r\left|(r-j+1)
             \int_{Z_i}c_j(N_{Z_i/X})h_0^{r-j}\right|,\\
 R&=\max_i\sum_j|G_{ij}|.
 \end{split}
\]
Whenever a smooth axis containing these centres exists, the matrix
is positive definite provided
\[
 N\ge1,\qquad N\ge\max_i\frac{2S_i}{d_i},
 \qquad N^{r+2}>\frac{2R}{d_{\min}}.
 \tag{C44.A06.4}
\]
Thus the $w_i$ are nonzero and linearly independent in that range.

\textbf{Proof.}
The inverse Lefschetz operator in (C44.A06.1) is algebraic because
$B(X)$ is known for abelian varieties, without the Hodge standard
conjecture. Since $[B_N]=N^2h_0^2$, self-intersection gives
$\kappa_*\kappa^*=N^2L_{h_0}^2$, so $\kappa_*w_i=0$.
The normal sequence of the smooth inclusions $Z_i\subset B_N\subset X$
is
\[
 0\longrightarrow N_{Z_i/B_N}\longrightarrow N_{Z_i/X}
 \longrightarrow (L_0^{\otimes N}|_{Z_i})^{\oplus2}
 \longrightarrow0.
\]
Therefore
$c(N_{Z_i/B_N})=c(N_{Z_i/X})(1+Nh_0)^{-2}$.
The self-intersection formula and the coefficient identity
$(1+x)^{-2}=\sum_{k\ge0}(-1)^k(k+1)x^k$ show that
$(-1)^r\int_{B_N}[Z_i]^2=P_i(N)$.
For distinct $i,j$ the intersection is zero because the centres
are disjoint. Projection formula gives
\[
 \int_{B_N}[Z_i]\kappa^*a_j=G_{ij},\qquad
 \int_{B_N}\kappa^*a_i\kappa^*a_j=N^2G_{ij}.
\]
Expanding (C44.A06.1) proves (C44.A06.3), with no asymptotic
approximation.

For $N\ge1$ the lower-degree terms in $P_i$ have total absolute
value at most $S_iN^{r-1}$. The first two inequalities in
(C44.A06.4) imply $P_i(N)\ge\tfrac12d_iN^r$.
For a real vector $x=(x_i)$, symmetry of $G$ and
$2|x_ix_j|\le x_i^2+x_j^2$ give
$|x^tGx|\le R\sum_i x_i^2$. Hence the right side of
(C44.A06.3), evaluated at $x$, is at least
\[
 \left(\tfrac12d_{\min}N^r-RN^{-2}\right)\sum_i x_i^2,
\]
which is positive for $x\ne0$ under the last inequality.

It remains to justify existence with the prescribed centres.
Let $Z=\coprod_iZ_i$ and $V_N=H^0(X,I_ZL_0^{\otimes N})$.
For $N$ large, Serre vanishing for $I_Z^2L_0^{\otimes N}$ and
global generation of $(I_Z/I_Z^2)L_0^{\otimes N}$ make the
normal first-jet map from $V_N$ onto the conormal fibre at each
point of $Z$ surjective. Off $Z$, $V_N$ separates second jets at
one point and simultaneous first jets at any two distinct points.
To see uniformity, first choose $a$ for which $I_ZL_0^{\otimes a}$
is globally generated. At any one or two points off $Z$, a general
section of this sheaf is nonzero at all those points. Multiplying
it by sections of a sufficiently high power of $L_0$, which
separates the indicated finite jets, gives the required maps from
$V_N$. High powers of a very ample bundle separate these jets;
handling the finitely many residue classes of $N$ gives the claim
for every sufficiently large $N$.

Set $n=2r+2$. For a pair of sections, rank at most one for the
two normal covectors at a fixed point of $Z$ has codimension
$r+1$ in the pair parameter space. Since $\dim Z=r$, a general
pair has rank two everywhere on $Z$. At a point off $Z$, the
conditions that both section values vanish and their two first
derivatives have rank at most one have codimension
$2+(n-1)=n+1$. Varying the point therefore still leaves a proper
bad locus. A general pair consequently defines a smooth
codimension-two axis.

The second-jet and two-point first-jet surjectivity also give the
Lefschetz property. A singular section at a fixed point off $Z$
imposes $n+1$ independent linear conditions. A degenerate
quadratic part imposes at least one more; after varying the point,
these sections have codimension at least two in $\mathbf P(V_N)$.
Sections singular at two distinct points also have codimension at
least two, since the fixed-pair conditions have codimension
$2(n+1)$ and the pair of points has dimension $2n$.
A general pencil avoids the closures of both loci. Along $Z$ no
member is singular, because the two normal covectors are
independent. Smoothness of the axis also excludes singularities
of a member elsewhere on the axis. Thus every singular member
has just one ordinary quadratic singularity. In characteristic
two an invertible Hessian is possible here because $n$ is even.
These are the local incidence arguments of SGA~7~II, XVII,
1.1--1.2, 3.2.1--3.2.4, 3.3, 3.7, and 4.2--4.3, with the
imposed-centre conditions verified above. The generic member is
smooth by the same first-jet count. All the required opens in the
pencil parameter space are nonempty, so the construction exists.
Over $\overline{\mathbf F}_p$ a chosen pencil is defined over a
finite extension.\hfill$\square$

\paragraph{Scope and programme location.}
This gives a new unconditional positive subspace of the precise
H9-T7 axis height form, in cote~44 Th.~5, scan~004, with the
primitive qualification of scan~008. It allows prescribed smooth
centres whose own lifting is not known, but it does not prove that
the projected classes are nonliftable. Its threshold depends on
the chosen centres, and it changes the axis and the polarization.
It supplies neither a uniform bound over all cycles in one fixed
$B$ nor positivity on the rest of $A_{B_N}$. It therefore does not
close QH9-axis or $I(B)$ in cote~16 terminology. The geometric
existence argument has the same characteristic-two scope as
C44-A04. All normal and self-intersection operations are those of
Fulton, \emph{Intersection Theory}, \S6.3.

% END work/axis/axis_asymptotic_centres.tex

% BEGIN work/axis/axis_product_curves.tex
% C44_POS /root/axis_independent; GPT-6 Astra Pro; effort not exposed.
\section{A fixed-polarization positive axis class from a product of curves}

\subsection*{C44-A07. Products of smooth curves in abelian threefold factors}

\textbf{Theorem.}
Let $A,A'$ be abelian threefolds over an algebraically closed field of
arbitrary characteristic, with ample line bundles $L,M$ of classes
$s,t$. Put $X=A\times A'$, $h=s+t$, and let $B\subset X$ be a smooth
complete intersection of two sections of $L\boxtimes M$.
Suppose that $B$ contains $Z=C\times D$, where $C\subset A$ and
$D\subset A'$ are smooth connected curves. Write
\[
 a=\int_Cs,\quad b=\int_Dt,\quad
 V=\int_As^3,\quad W=\int_{A'}t^3,\quad
 u=a/V,\quad v=b/W.
\]
Then $0<u,v\le1$, and the rational algebraic class
\[
 w=[Z]_B-\kappa^*(L_h^2)^{-1}[Z]_X\in A_B
 \tag{C44.A07.1}
\]
satisfies the unconditional strict lower bound
\[
 \int_Bw^2\ \ge\ ab(2u+2v-3uv)>0.
 \tag{C44.A07.2}
\]
In particular this class is nonzero. When $B$ is a Lefschetz axis,
this is the positive H9-T7 axis sign ($m=3$) on the line spanned by
$w$. Neither 032 nor a lifting of the embedded curves is used.

\paragraph{1. Smoothness forces two independent conormal directions on each curve.}
Set $K_C=H^0(A,I_CL)$ and $K_D=H^0(A',I_DM)$. By the degree-zero
K\"unneth formula and elementary linear algebra, the sections of
$L\boxtimes M$ vanishing on $Z$ form
\[
 K_C\otimes H^0(A',M)+H^0(A,L)\otimes K_D.
 \tag{C44.A07.3}
\]
Let $e_C$ be the generic rank of the evaluation of conormal derivatives
\[
 K_C\otimes\mathcal O_C\longrightarrow N_{C/A}^{\vee}\otimes L|_C,
\]
and define $e_D$ similarly. The smooth normal sequence for $Z\subset
B\subset X$ gives a surjection
\[
 N_{C/A}\boxplus N_{D/A'}\longrightarrow
 (L|_C\boxtimes M|_D)^{\oplus2}.
 \tag{C44.A07.4}
\]
Its first and second summands factor through the derivatives in
(C44.A07.3). If $e_C=0$, restrict (C44.A07.4) to $C\times\{y\}$.
It would be a surjection from a trivial rank-two bundle to
$(L|_C)^{\oplus2}$, up to a constant one-dimensional factor. Its
determinant cannot be invertible because $2a>0$. Thus $e_C>0$,
and similarly $e_D>0$.

Suppose $e_C=1$. Saturate the derivative image in
$N_{C/A}^{\vee}\otimes L|_C$; the saturation is a line subbundle
$P_C$. It has a nonzero global section, so $\deg P_C\ge0$.
Put $Q_C=L|_C\otimes P_C^{-1}$ and $q=\deg Q_C\le a$.
All the derivatives in the first summand of (C44.A07.4) factor
through the line-bundle quotient $N_{C/A}\twoheadrightarrow Q_C$.
If also $e_D=1$, a rank-two quotient
$Q_C\boxplus Q_D$ would map surjectively to the rank-two target in
(C44.A07.4). Taking determinants and restricting to $C\times\{y\}$
would give $q=2a$, a contradiction. Hence $e_D=2$.

Put $c=\deg N_{C/A}=2g(C)-2$ and
$d=\deg N_{D/A'}=2g(D)-2$. The equality $e_D=2$ supplies a nonzero
section of $\det(N_{D/A'}^{\vee}\otimes M|_D)$, so $d\le2b$.
Now (C44.A07.4) factors through a surjection from the rank-three
bundle $Q_C\boxplus N_{D/A'}$. Its kernel is a line bundle. The
degree of the second Chern class of the formal quotient by
$(L|_C\boxtimes M|_D)^{\oplus2}$ must therefore vanish. It is
\[
 qd-2(ad+bq)+6ab=(q-2a)(d-2b)+2ab>0,
\]
since $q\le a$ and $d\le2b$. This is impossible. Consequently
$e_C=e_D=2$, and taking the two determinants gives
\[
 c\le2a,\qquad d\le2b.
 \tag{C44.A07.5}
\]
All saturation arguments here are on smooth curves, so the saturated
subsheaves and the corresponding quotients are vector bundles.

\paragraph{2. A curve whose conormal derivatives have rank two is bounded by a complete intersection.}
Remove the common fixed divisor $H$ from the linear system $K_C$.
It cannot contain $C$: at the generic point of a curve contained
in a common divisor, every derivative of a section divisible by its
local equation lies in at most one conormal direction. Thus the
residual system in $L(-H)$ still vanishes on $C$ and has no fixed
divisor. Two general residual divisors meet properly and contain
$C$ in their effective intersection cycle. If $s_0=s-[H]$, this
shows that $s_0^2-[C]$ is the class of an effective curve cycle.
The class
\[
 s^2-s_0^2=[H](s+s_0)
\]
is effective as a cohomology class: both factors have effective
divisor representatives, and translating one of them generally
makes their intersection proper without changing its class.
Therefore
\[
 s^2-[C]\text{ and }t^2-[D]
 \quad\text{are effective curve classes.}
 \tag{C44.A07.6}
\]
Intersecting with $s,t$ gives $0<u,v\le1$.

\paragraph{3. The precise Fourier bound for effective curves.}
We record the needed form and its characteristic-independent
justification. For a polarized abelian threefold $(A,s)$ let
\[
 \mathcal F_s(x)=\chi(L)^{-1}\varphi_L^*
       p_{\widehat A,*}(p_A^*x\,\exp(c_1(\mathcal P))),
 \qquad T_s(x)=-\mathcal F_s(x)
\]
on degree-four cohomology; $\mathcal P$ is the normalized Poincar\'e
bundle on $A\times\widehat A$. Changing its sign makes no difference
in this degree. The normalized Fourier formulas are
\[
 \mathcal F_s(s^2)=-2s,\qquad
 \mathcal F_s(s\alpha_0)=\alpha_0\quad(s^2\alpha_0=0).
 \tag{C44.A07.7}
\]
These follow directly from Beauville, \emph{The action of
$\mathrm{SL}_2$ on abelian varieties}, JRMS 25 (2010),
Corollary~5.4, p.~261, with $g=3$, and the normalization in
\S4.1, p.~260. That paper works over an algebraically closed field
of arbitrary characteristic (\S1.1).

For every effective curve cycle $F$ on $A$, the class $T_s([F])$
is a nef rational divisor. Indeed it suffices to normalize an
integral component $\nu:\widetilde F\to A$ and use the restriction
homomorphism $\rho:\widehat A\to J(\widetilde F)$. The Poincar\'e
identity on a curve gives
\[
 -\frac12 p_{\widehat A,*}
       \bigl([F]c_1(\mathcal P)^2\bigr)=\rho^*\theta_{\widetilde F}.
 \tag{C44.A07.8}
\]
Thus $T_s([F])=\chi(L)^{-1}\varphi_L^*\rho^*\theta_{\widetilde F}$,
which is nef. For clarity, (C44.A07.8) uses only the canonical
principal polarization of a Jacobian, not complex Hodge positivity.
In a symplectic basis $e_i,f_i$ of $H^1(\widetilde F)$ with
$\int e_if_j=\delta_{ij}$, the universal Poincar\'e tensor is
$\sum_i(e_i\otimes v_i-f_i\otimes u_i)$ and the theta class is
$\sum_i u_iv_i$. Its square pushed along the curve is
$-2\sum_i u_iv_i$, proving the formula. The identification of these
tensors in every characteristic is the autoduality statement of
Milne, \emph{Jacobian Varieties} (2021 corrected reprint of the
1986 chapter), Theorem~6.6, Lemmas~6.8--6.9 and Summary~6.11,
pp.~21--24. Equivalently, pullback to the square of the curve is
the diagonal minus the two base-point fibres, as in Summary~6.11.

Put $\alpha=L_s^{-1}[C]=us+\alpha_0$, with $s^2\alpha_0=0$,
and $Q_\alpha=\int_A s\alpha_0^2$. The inverse is algebraic by the
known Lefschetz standard conjecture for abelian varieties.
Equations (C44.A07.6)--(C44.A07.8) show that
\[
 T_s([C])=2us-\alpha_0,
 \qquad 0\le T_s([C])\le2s
 \tag{C44.A07.9}
\]
in the nef order. We claim that
\[
 0\le -Q_\alpha\le 2u(1-u)V.
 \tag{C44.A07.10}
\]
Here is an algebraic proof of the bound, including the reality
needed to use (C44.A07.9). For a rational divisor class $E$ put
$f=\varphi_L^{-1}\varphi_E\in\operatorname{End}^0(A)$. This is
self-adjoint for the Rosati involution. The positive definite
Rosati pairing $(x,y)=\operatorname{Tr}(xy^\dagger)$ makes left
multiplication by $f$ self-adjoint: cyclicity of the trace gives
$(fx,y)=(x,fy)$. It follows that the minimal polynomial of $f$
has only real roots. Rosati positivity here is the classical
theorem, for example Conrad, \emph{Abelian varieties: geometry,
parameter spaces, and arithmetic} (2005), Theorem~5.2.4,
pp.~45--46; it is unrelated to 032.

Riemann--Roch and the degree polynomial of an endomorphism give
\[
 \det(T+f\mid V_\ell A)
   =\left(\frac{\int_A(Ts+E)^3}{V}\right)^2.
 \tag{C44.A07.11}
\]
Exact locators are Milne, \emph{Abelian Varieties}, course notes
v.~2.00 (2008), Proposition~10.20 (printed pp.~50--51) and
Theorem~11.1(a),(b) (printed p.~54). In particular the monic cubic
on the right before squaring has three real roots, counted with
multiplicity. Write it as $\prod_{i=1}^3(T+\mu_i)$.
For $E=T_s([C])$, both $E$ and $2s-E$ are nef. The mixed
intersection coefficients of a nef divisor with $s$ are
nonnegative, so their cubics are strictly positive at every
positive real $T$. Applied to $f$ and $2-f$, this shows
$0\le\mu_i\le2$.

Expanding (C44.A07.11) using $E=2us-\alpha_0$ gives
\[
 \sum_i\mu_i=6u,\qquad
 \sum_i\mu_i^2=12u^2-6Q_\alpha/V.
\]
The lower bound $\sum_i\mu_i^2\ge(\sum_i\mu_i)^2/3$ and the
upper bound $\sum_i\mu_i^2\le2\sum_i\mu_i$ prove
(C44.A07.10). The identical argument for
$\beta=L_t^{-1}[D]=vt+\beta_0$ gives
\[
 0\le -Q_\beta\le2v(1-v)W,
 \qquad Q_\beta=\int_{A'}t\beta_0^2.
 \tag{C44.A07.12}
\]

\paragraph{4. Exact self-intersection and ambient subtraction.}
The normal sequence gives
$c(N_{Z/B})=c(N_{Z/X})(1+s+t)^{-2}$. Since
$N_{Z/X}=N_{C/A}\boxplus N_{D/A'}$, its degree-two part gives
\[
 \int_B[Z]^2=6ab-2(ad+bc)+cd
              =(c-2a)(d-2b)+2ab\ \ge2ab.
 \tag{C44.A07.13}
\]
This is the ordinary self-intersection formula (Fulton,
\emph{Intersection Theory}, \S6.3), with (C44.A07.5).

Using $s^2\alpha_0=t^2\beta_0=0$, direct multiplication gives
\[
 \begin{split}
 a_Z:=(L_{s+t}^2)^{-1}[C][D]
 &=\tfrac12\alpha_0\beta_0+v s\alpha_0+u t\beta_0\\
 &\quad+uv\bigl(st-\tfrac12(s^2+t^2)\bigr).
 \end{split}
 \tag{C44.A07.14}
\]
Indeed multiplication by $(s+t)^2$ gives
$(us^2+s\alpha_0)(vt^2+t\beta_0)$ term by term.
Consequently
\[
 G:=\int_Xa_Z[C][D]
       =\tfrac12Q_\alpha Q_\beta+u^2v^2VW.
 \tag{C44.A07.15}
\]
The self-intersection identity $\kappa_*\kappa^*=L_h^2$ gives
$\kappa_*w=0$ and
$\int_Bw^2=\int_B[Z]^2-G$. Equations (C44.A07.10) and
(C44.A07.12), together with $ab=uvVW$, imply
\[
 \frac{G}{ab}\le uv+2(1-u)(1-v).
\]
Subtracting this from (C44.A07.13) proves (C44.A07.2).
Finally $2u+2v-3uv\ge u+v-uv>0$ when $0<u,v\le1$.
\hfill$\square$

\paragraph{Scope and programme location.}
The theorem is valid for supersingular factors and for curves
whose embedded support is not known to lift. It does not claim
that $w$ is nonliftable. It establishes the axis height sign on
one explicit nonzero line at the original polarization, not on
the span of all such lines or on all of $A_B$. In cote~44 it
addresses the remaining axis summand in Th.~5 (scan~004), with
the primitive qualification of the letter on scan~008. In
cote~16 terms it supplies a particular positive primitive
algebraic class toward $I(B)$, without settling that condition
or its product version. It strengthens the calculation of a
candidate class; QH9-axis itself remains open in this work.

% END work/axis/axis_product_curves.tex

% BEGIN work/axis/axis_product_grid.tex
% C44_POS /root; GPT-6 Astra Pro; effort not exposed.
\section{A semipositive span at one fixed product polarization}

\begin{theorem}[C44-A09: a grid of product curves]
Let $A,A'$ be abelian threefolds over an algebraically closed field
of arbitrary characteristic, with ample classes $s=c_1(L)$ and
$t=c_1(M)$. Let $B\subset A\times A'$ be a smooth complete
intersection of two members of $|L\boxtimes M|$.
Let $C_1,\ldots,C_e\subset A$ and $D_1,\ldots,D_f\subset A'$
be two nonempty families of smooth connected curves, pairwise
disjoint within each family. Suppose that $B$ contains every
$Z_{ij}=C_i\times D_j$. Put
\[
 \begin{gathered}
 V=\int_As^3,\qquad W=\int_{A'}t^3,\qquad
 a_i=\int_{C_i}s,\quad b_j=\int_{D_j}t,\\
 u_i=a_i/V,\qquad v_j=b_j/W,\qquad
 U_0=\sum_i u_i,\quad V_0=\sum_jv_j,\\
 w_{ij}=[Z_{ij}]_B-\kappa^*(L_{s+t}^2)^{-1}[Z_{ij}]_{A\times A'}.
 \end{gathered}
 \tag{C44.A09.1}
\]
Then $0<U_0,V_0\le1$ and $w_{ij}\in A_B$.
The intersection form $\int_Bww'$ is positive semidefinite on
the entire rational span of the $w_{ij}$. If $e=1$ or $f=1$, its
Gram matrix in these generators is positive definite, so the
generators are independent. All these statements are unconditional.

More precisely, for any real coefficients $x_{ij}$, define
\[
 R_i=\sum_jv_jx_{ij},\qquad
 S_j=\sum_i u_ix_{ij},\qquad
 z=\sum_{i,j}u_iv_jx_{ij}.
\]
The exact axis form satisfies the explicit lower bound
\[
 \begin{split}
 \int_B\left(\sum_{i,j}x_{ij}w_{ij}\right)^2
 \ge VW\left(2\sum_i u_iR_i^2+2\sum_jv_jS_j^2-3z^2\right)\ge0.
 \end{split}
 \tag{C44.A09.2}
\]
A vector of zero norm must therefore have every weighted row sum
$R_i$ and every weighted column sum $S_j$ equal to zero. No assertion
that every such vector has zero norm, or is a cohomological relation,
is made.
\end{theorem}

\begin{proof}
Write $C=\coprod_iC_i$, $D=\coprod_jD_j$ and
\[
 K_C=H^0(A,I_CL),\qquad K_D=H^0(A',I_DM).
\]
The restriction kernel on $C\times D$ is
$K_C\otimes H^0(M)+H^0(L)\otimes K_D$, by the same tensor-kernel
identity as (C44.A07.3). At every pair $C_i\times D_j$, the
normal surjection associated with the smooth axis factors
through the conormal derivatives of these two global spaces.
The rank-zero and rank-one contradictions in part~1 of C44-A07
apply to these spaces themselves: their proof used only this
factorization, saturation on a smooth curve, and the surjectivity
of the normal map. It did not require that all sections vanishing
on the individual component be available. Hence $K_C$ has
conormal rank two generically on each $C_i$, and $K_D$ has rank
two generically on each $D_j$.

Removing the common fixed divisor of $K_C$, no component $C_i$
can be contained in that divisor, since that would leave at most
one conormal direction. Two general residual divisors meet
properly and contain every $C_i$. Restoring the fixed divisor
by the general-translation intersection argument of C44-A07
proves
\[
          s^2-\sum_i[C_i]\ \text{and}\
          t^2-\sum_j[D_j]\quad\text{are effective curve classes}.
 \tag{C44.A09.3}
\]
In particular $0<U_0,V_0\le1$.
Writing $c_i=2g(C_i)-2$ and $d_j=2g(D_j)-2$, the conormal
rank-two determinants also give $c_i\le2a_i$, $d_j\le2b_j$.
The normal-bundle formula consequently yields
\[
 T_{ij}:=\int_B[Z_{ij}]^2
       =(c_i-2a_i)(d_j-2b_j)+2a_ib_j\ge2a_ib_j.
 \tag{C44.A09.4}
\]
Distinct $Z_{ij}$ are disjoint.

\paragraph{A simultaneous Fourier bound.}
Decompose
\[
 L_s^{-1}[C_i]=u_is+\alpha_i,\quad s^2\alpha_i=0,
 \qquad
 L_t^{-1}[D_j]=v_jt+\beta_j,\quad t^2\beta_j=0.
\]
Let
\[
 Q^C_{ik}=\int_As\alpha_i\alpha_k,\qquad
 Q^D_{jl}=\int_{A'}t\beta_j\beta_l,
 \qquad D_u=\operatorname{diag}(u_i),\quad
 D_v=\operatorname{diag}(v_j).
\]
We claim the inequalities of real symmetric matrices
\[
 0\le -Q^C\le2V(D_u-uu^t),\qquad
 0\le -Q^D\le2W(D_v-vv^t).
 \tag{C44.A09.5}
\]
Here $u=(u_i)$ and $v=(v_j)$ are column vectors. Matrix
inequalities mean nonnegative quadratic forms, not entrywise
inequalities.

Use the normalized Fourier map $T_s$ of C44-A07. Each
$T_s([C_i])=2u_is-\alpha_i$ is nef, and (C44.A09.3) gives
$\sum_iT_s([C_i])\le2s$ in the nef order. Let $E_i$ be the
Rosati-self-adjoint endomorphism corresponding to this divisor.
As proved in C44-A07 by Rosati positivity and Riemann--Roch,
nef divisors correspond here to positive self-adjoint operators
in the faithful left-regular representation of
$\operatorname{End}^0(A)\otimes\mathbf R$. To make the
representation comparison explicit, the Tate representation is
faithful, so every root of the minimal polynomial of a rational
self-adjoint endomorphism occurs among the roots of its degree
polynomial. The Riemann--Roch and nef arguments of C44-A07 make
all these roots nonnegative. This is exactly positivity in the
left-regular Rosati representation, and extends to real
self-adjoint elements by continuity. Thus
\[
                        E_i\ge0,\qquad \sum_iE_i\le2I.
\]
For arbitrary real $x_i$, elementary operator Cauchy--Schwarz gives
\[
             \left(\sum_i x_iE_i\right)^2
                         \le2\sum_i x_i^2E_i.
 \tag{C44.A09.6}
\]
For example, let $R$ be the block-row operator with entries
$E_i^{1/2}$ and let $D$ act by the scalar $x_i$ on its $i$-th
summand. Then $RR^*=\sum_iE_i\le2I$, so $R^*R\le2I$, and
\[
 (RDR^*)^2=RD(R^*R)DR^*\le2RD^2R^*.
\]
No commutativity among the $E_i$ is needed.

Let $\tau=\tfrac12\operatorname{Tr}_{V_\ell A}$ on the rational
endomorphism algebra, extended to the real algebra. It is
positive on positive self-adjoint elements: such an element
has a self-adjoint square root in the real finite-dimensional
algebra, and Rosati trace positivity applies to its square.
Also $\tau(I)=3$ and $\tau(E_i)=6u_i$. For
$E(x)=\sum_i x_iE_i$ and $u(x)=\sum_i x_iu_i$, the degree
polynomial calculation (C44.A07.11) gives
\[
 \tau(E(x)^2)=12u(x)^2-\frac6Vx^tQ^Cx.
 \tag{C44.A09.7}
\]
It is a polynomial identity, so is valid for real as well as
rational $x_i$. Applying $\tau$ to (C44.A09.6) proves the
upper bound in (C44.A09.5). The lower bound follows from
$\tau(E(x)^2)\ge\tau(E(x))^2/\tau(I)$, which is
Cauchy--Schwarz for the positive Rosati trace pairing. The
second matrix bound is identical.

\paragraph{The complete Gram matrix and its lower bound.}
The exact inverse and ambient subtraction in C44-A07, applied
bilinearly to each pair, give
\[
 G_{ij,kl}=\tfrac12 Q^C_{ik}Q^D_{jl}
                              +VWu_iu_kv_jv_l,
 \qquad
 Q_B=\operatorname{diag}(T_{ij})-G.
 \tag{C44.A09.8}
\]
The first identity also follows by multiplying the displayed
inverse (C44.A07.14) for $C_i\times D_j$ against
$[C_k][D_l]$. The remaining terms have too high a degree on one
threefold factor and vanish. The second identity uses disjointness
of distinct centres and the same projection formula as C44-A07.

The tensor product of positive semidefinite forms preserves
inequalities. Applying (C44.A09.5) therefore shows
\[
 \begin{split}
 G&\le2VW(D_u-uu^t)\otimes(D_v-vv^t)
                           +VWuu^t\otimes vv^t,\\
 Q_B&\ge VW\left(2D_u\otimes vv^t+
                  2uu^t\otimes D_v-3uu^t\otimes vv^t\right),
 \end{split}
 \tag{C44.A09.9}
\]
where (C44.A09.4) was used in the second line. Evaluation on
$x=(x_{ij})$ is the first inequality of (C44.A09.2).
Its right-hand quadratic form is nonnegative because it equals
\[
 2\sum_i u_i(R_i-z/U_0)^2+
 2\sum_j v_j(S_j-z/V_0)^2+
       (2/U_0+2/V_0-3)z^2.
 \tag{C44.A09.10}
\]
The last coefficient is at least one. All weights are strictly
positive. Thus equality in this lower form forces $z=0$ and
all $R_i,S_j$ zero. If either family has a single component,
the weighted row and column sums determine every $x_{ij}$;
hence the lower form is positive definite. This proves all
claims, including independence in that case.
\end{proof}

\paragraph{Programme consequence and limitation.}
Unlike an assertion that several individual classes have
positive squares, C44-A09 controls every linear combination in
the displayed span at the original polarization. For a
Lefschetz axis it is exactly the H9-T7 sign required for this
part of the height in cote~44 Th.~5, scan~004, with the
primitive qualification on scan~008. It is an unconditional
special-axis result in every characteristic, including
supersingular ambient factors. It neither asserts that all of
$A_B$ is generated by these product curves nor settles the
universal $I(B)$ of cote~16. The proved criterion requires the
entire product grid, not an arbitrary collection of unrelated
product curves. Its intersection and Fourier inputs are the
fully specified classical inputs of C44-A07; no 032 consequence
is used.

% END work/axis/axis_product_grid.tex

% BEGIN work/axis/axis_grid_strict.tex
% C44_POS /root/axis_independent; GPT-6 Astra Pro; effort not exposed.
\subsection*{C44-A09a. A generating curve makes the entire grid strictly positive}

\textbf{Theorem.}
In the setting and notation of C44-A09, suppose at least one
$C_i$ generates $A$ as an abelian variety, or at least one
$D_j$ generates $A'$. Here a curve generates if no translate of
a proper abelian subvariety contains it. Then the complete
Gram matrix of the classes $w_{ij}$ is positive definite.
In particular all the $w_{ij}$ are independent. This holds in
every characteristic without 032. It applies in particular
whenever either threefold factor is geometrically simple.

\textbf{Proof.}
Suppose $C_{i_0}$ generates $A$. The induced homomorphism
$J(C_{i_0})\to A$ is surjective, by the universal property of
the Jacobian (Milne, \emph{Jacobian Varieties}, Proposition~6.1,
printed p.~19 in the 2021 corrected reprint). Its dual
$\rho:\widehat A\to J(C_{i_0})$ has finite kernel, so is finite
onto its image. The explicit formula (C44.A07.8) shows that
$T_s([C_{i_0}])=\chi(L)^{-1}\varphi_L^*\rho^*\theta_{C_{i_0}}$
is ample as a rational divisor class. Consequently the
corresponding Rosati-self-adjoint operator $E_{i_0}$ of C44-A09
is strictly positive and invertible. For example, the nef
argument there makes its spectrum nonnegative; its ample top
self-intersection is nonzero, so the Riemann--Roch degree
polynomial excludes eigenvalue zero. Tate faithfulness makes
this conclusion valid in the left-regular representation as
well as in the Tate representation.

Keep the notation of C44-A09 and define
\[
 A_C=-Q^C,\quad H_C=2V(D_u-uu^t),\quad M_C=H_C-A_C,
 \qquad
 A_D=-Q^D,\quad H_D=2W(D_v-vv^t),\quad M_D=H_D-A_D.
\]
All six matrices are positive semidefinite by (C44.A09.5).
We first prove that $M_C$ is positive definite on
$u^\perp=\{x:\sum_i u_ix_i=0\}$.

Use the faithful real left-regular representation with its
positive Rosati inner product. For real $x=(x_i)$ let
$R=(E_1^{1/2},\ldots,E_e^{1/2})$, let $D$ be the block-diagonal
operator with scalar blocks $x_i$, and put
$F=RDR^*=\sum_i x_iE_i$. The exact trace calculation in
C44-A09 gives
\[
 x^tM_Cx=\frac V6\tau\bigl(2\sum_i x_i^2E_i-F^2\bigr),
 \qquad \tau=\tfrac12\operatorname{Tr}_{V_\ell A}.
 \tag{C44.A09a.1}
\]
The operator inside the trace is
\[
 S=RD(2I-R^*R)DR^*\ge0.
\]
The functional $\tau$ is faithful on positive elements: if a
nonzero positive element is written as the square of its
self-adjoint square root in the finite-dimensional real
algebra, this follows from positive definiteness of the
Rosati trace pairing. Therefore equality in (C44.A09a.1)
implies $S=0$. It follows that
$(2I-R^*R)^{1/2}DR^*=0$, and hence
\[
                    2DR^*=R^*RDR^*=R^*F.
 \tag{C44.A09a.2}
\]
In the $i$-th component this says
$2x_iE_i^{1/2}=E_i^{1/2}F$.
For $i=i_0$, invertibility gives $F=2x_{i_0}I$.
Substituting into the other components gives
$(x_i-x_{i_0})E_i^{1/2}=0$ for every $i$.
Each $E_i$ is nonzero because $\tau(E_i)=6u_i>0$.
All $x_i$ are therefore equal. Such a vector in $u^\perp$
must vanish, since $\sum_i u_i>0$. This proves the asserted
strict positivity of $M_C$ on $u^\perp$.

Let $Q_B$ be the complete grid Gram matrix and let
\[
 K=VW\bigl(2D_u\otimes vv^t+2uu^t\otimes D_v
                         -3uu^t\otimes vv^t\bigr)
\]
be its lower bound in (C44.A09.9). The same calculation, retaining
the slack in the tensor inequality, gives
\[
 \begin{split}
 Q_B&\ge K+\tfrac12(H_C\otimes H_D-A_C\otimes A_D)\\
    &=K+\tfrac12(M_C\otimes H_D+A_C\otimes M_D)\\
    &\ge K+\tfrac12M_C\otimes H_D.
 \end{split}
 \tag{C44.A09a.3}
\]
By (C44.A09.10), $K$ is positive semidefinite and its kernel
consists precisely of arrays with zero weighted row and
column sums. As a tensor subspace this is
$u^\perp\otimes v^\perp$.
On $v^\perp$ one has
$y^tH_Dy=2W\sum_jv_jy_j^2>0$ for $y\ne0$.
Thus $M_C\otimes H_D$ is positive definite on the kernel
of $K$. Inequality (C44.A09a.3) proves that $Q_B$ is positive
definite everywhere. If a $D_j$ generates $A'$, interchange
the factors. Finally every positive-dimensional curve in a
geometrically simple abelian variety generates it up to
translation, proving the last assertion.\hfill$\square$

\paragraph{Programme scope.}
This proves strict positivity for a full family of special-axis
classes at one fixed polarization, strengthening C44-A09 and
the single-class result C44-A07. It controls the corresponding
entire summand of the H9-T7 height for cote~44 Th.~5, scan~004.
It makes no claim that the grid exhausts $A_B$, and supplies
neither a universal correspondence injection nor an individual
nonliftable class on the axis.

% END work/axis/axis_grid_strict.tex

% BEGIN work/axis/axis_supersingular_nonliftable.tex
% C44_POS /root/axis_independent; GPT-6 Astra Pro; effort not exposed.
\section{An explicit supersingular nonliftable class outside the axis kernel}
\subsection*{C44-A10. Individual class nonliftability}

\textbf{Theorem.}
For $E/\overline{\mathbf F}_3$ given by $y^2=x^3-x$, there is an
explicit integral divisor-product cycle $z\in\operatorname{CH}^2(E^4)$
whose nonzero cohomology class cannot specialize from an algebraic
cycle in any characteristic-zero abelian-scheme lift of $E^4$.
If $a$ is its pullback to $X=E^6$ and $\kappa:B\hookrightarrow X$
is a smooth complete intersection of two ample divisors of class $h$,
then $\kappa^*a$ cannot lift even under an arbitrary abstract smooth
projective lift of $B$. Nevertheless,
\[
 \kappa_*\kappa^*a=h^2a\ne0,\qquad
 (1-\kappa^*(L_h^2)^{-1}\kappa_*)\kappa^*a=0.                 \tag{A10.1}
\]
Thus this is individual class nonliftability, but gives no
nonliftable class in $A_B$.

\paragraph{Explicit divisors and cycle.}
Choose $\eta\in\mathbf F_9$ with $\eta^2=-1$ and define
$\iota(x,y)=(-x,\eta y)$, $\pi(x,y)=(x^3,y^3)$.
Then $\iota^2=-1$, $\pi\iota=-\iota\pi$, and $\pi^2=-3$:
there are four points over $\mathbf F_3$, so the Frobenius
polynomial is $T^2+3$, and Cayley--Hamilton plus Tate-module
faithfulness gives the identity. Thus $[3]=-\pi^2$ is purely
inseparable and $E$ is supersingular.
The degree/characteristic-polynomial statements are in Milne,
\emph{Abelian Varieties}, v.~2.00, I Proposition 10.20,
printed pp.~50--52; the point-count identity is in II, \S1.

Choose a basis of $H^1(E,\mathbf F_5)$ with pullback matrices
\[
 I=\iota^*=\begin{pmatrix}2&0\\0&3\end{pmatrix},\qquad
 P=\pi^*=\begin{pmatrix}0&1\\2&0\end{pmatrix}.                 \tag{A10.2}
\]
The invertible $\pi$ interchanges the two distinct $\iota$
eigenspaces, and rescaling using $\pi^2=2$ gives this form.
Do not assume Weil normalization: write $[0]=\lambda e_0e_1$
with $\lambda\in\mathbf F_5^\times$ and use the same basis on all
four factors. Put $D_i=\operatorname{pr}_i^*[0]$ and, for $i<j$,
\[
 R_{ij,u}=(\operatorname{pr}_i+u\operatorname{pr}_j)^*[0]
                    -D_i-(\deg u)D_j,\quad
 u=1,\iota,\pi,\pi\circ\iota,\quad \deg u=1,1,3,3.            \tag{A10.3}
\]
These are integral divisors defined over $\mathbf F_9$.
For $a=2i-2,b=a+1,c=2j-2,d=c+1$, their normalized modulo-five
cross coordinates in the order $(ac,ad,bc,bd)$ are respectively
\[
 (0,1,4,0),\ (0,3,3,0),\ (1,0,0,3),\ (2,0,0,4).
\]
Indeed they are the cross terms in
\[
 (e_a+M_{11}e_c+M_{21}e_d)(e_b+M_{12}e_c+M_{22}e_d),
\]
where $M=1,I,P,IP$. The last matrix is $IP$ because pullback
reverses $\pi\circ\iota$.
Define integral divisors, for each pair of factors, by
\[
 \begin{array}{ll}
 F_{ac}=3R_{ij,\pi}+4R_{ij,\pi\circ\iota},&
 F_{ad}=3R_{ij,1}+R_{ij,\iota},\\
 F_{bc}=2R_{ij,1}+R_{ij,\iota},&
 F_{bd}=R_{ij,\pi}+2R_{ij,\pi\circ\iota},
 \end{array}
 \qquad F_{2i-2,2i-1}=D_i.                                  \tag{A10.4}
\]
Direct multiplication of the four cross columns shows
$\operatorname{cl}_5(F_{ab})=\lambda e_ae_b$.
Equivalently these are the representatives in $\{0,1,2,3,4\}$
of the inverse of the $28$-by-$28$ divisor coefficient matrix,
ordered by the four $D_i$, lexicographic pairs $i<j$, and the
four displayed endomorphisms.

Let $c_{abcd}$, in increasing lexicographic four-tuple order
on $\{0,\ldots,7\}$, be the following 70 integers read rowwise:
\[
\begin{array}{rrrrrrrrrr}
4&4&1&3&4&0&0&3&0&1\\
3&3&3&1&2&2&3&4&2&4\\
3&2&2&1&3&4&2&4&3&4\\
1&4&2&4&1&3&1&4&3&2\\
0&0&3&2&3&1&1&4&2&3\\
2&1&4&4&1&3&4&1&3&0\\
4&2&2&3&1&3&4&4&2&2
\end{array}
\]
Set $z=\sum_{a<b<c<d}c_{abcd}F_{ab}F_{cd}$.
Its modulo-five class is $\lambda^2w$, where
$w=\sum c_{abcd}e_ae_be_ce_d$. This completely specifies the
actual integral cycle, without requiring any conjectural classes.

\paragraph{Exact certificate for the actual rational tensor.}
For $w\in\bigwedge^4V$ let
$\mathcal D_w:\operatorname{End}(V)\to\bigwedge^4V$ send
$T$ to its derivative exterior action $T\cdot w$.
Its kernel is the infinitesimal stabilizer.
Order columns by $E_{ab}$, $a,b=0,\ldots,7$ lexicographically,
and rows by increasing lexicographic four-tuples $J$.
Its entries are
\[
 (\mathcal D_w)_{J,(a,b)}=
 \begin{cases}
 c_J,&a=b,\ a\in J,\\
 (-1)^{\operatorname{pos}_J(a)+\operatorname{pos}_K(b)}c_K,
 &a\ne b,\ a\in J,\ b\notin J,\ K=(J\setminus\{a\})\cup\{b\},\\
 0,&\text{otherwise}.
 \end{cases}                                                \tag{A10.5}
\]
Positions are zero-based, and both tuples are sorted.
The first 64 rows, with entries represented in $\{0,1,2,3,4\}$,
have exact integer determinant
\[
 26723729698312903479574054159209851065213568
                       \equiv3\pmod5.                       \tag{A10.6}
\]
The supplied \texttt{supersingular\_class\_certificate.py}
constructs the divisor matrix and tensor; its
\texttt{supersingular\_class\_certificate.json} supplies both
the divisor inverse and the full inverse of this minor modulo
five. Both inverse products were checked.
A separate target-row implementation of (A10.5) and exact Bareiss
elimination reproduced (A10.6). The root agent independently
rebuilt the matrix by source-column exterior replacement and
reproduced the same exact determinant. Random search selected
the fixed displayed vector; verification of it is exact.

Milne, loc.\ cit., I Theorem 12.1, printed pp.~55--56, gives
$H^*(E^4,\mathbf Z_5)=\bigwedge^*H^1(E^4,\mathbf Z_5)$,
torsion free. Lift the chosen modulo-five basis arbitrarily to a
$\mathbf Z_5$ basis. The derivative matrix of the actual integral
class of $z$ reduces to $\mathcal D_{\lambda^2w}$.
Its indicated minor is congruent to
$(\lambda^2)^{64}\cdot3=\lambda^{128}\cdot3=3$ modulo five.
Consequently the actual $\mathbf Q_5$ tensor has zero infinitesimal
stabilizer. Its integral coordinates need not equal the selected
residue representatives: higher $5$-adic terms cannot change
this unit-minor conclusion.

\paragraph{Characteristic-zero obstruction.}
If the class lifts algebraically to a characteristic-zero abelian
fourfold, spread the generic fibre and cycle to a finitely generated
field and embed it into $\mathbf C$. Smooth proper base change and
Betti--\'etale comparison preserve $H^1$, exterior products and
$\mathcal D$. Thus its rational Betti tensor also has zero
infinitesimal stabilizer after scalar extension.
But an algebraic degree-four class is $(2,2)$.
The real Hodge complex structure $J$ on $H^1$ acts by $+i$ on
$H^{1,0}$ and $-i$ on $H^{0,1}$, satisfies $J^2=-1$, and
annihilates a $(2,2)$ exterior tensor. It is a nonzero stabilizer,
contradiction. Although $J$ need not be rational, the derivative
matrix of a rational Betti tensor is rational: a nonzero kernel
over $\mathbf R$ implies rank deficiency over $\mathbf Q$ and
hence over $\mathbf Q_5$.
For the Hodge statements see Deligne, \emph{Hodge cycles on abelian
varieties}, notes by Milne, revised 1 October 2018, Introduction
printed p.~2 and \S3, Hodge structures, printed p.~23.
No Hodge conjecture is used.

\paragraph{Pullback and arbitrary abstract axis lifts.}
The support $S(t)$ of an exterior tensor $t$ is its least supporting
subspace, equivalently the span of its contractions by three
covectors. It is intrinsic, equivariant and compatible with field
extension. The certified tensor on an eight-dimensional space has
full support, since a nonzero endomorphism vanishing on a proper
support would annihilate it.
Hence $a$ on $E^6$ has intrinsic support of dimension eight and
zero infinitesimal stabilizer on that support.
Any $(2,2)$ lift would have Hodge-circle-invariant support;
$J$ restricted to it still squares to $-1$ and annihilates the
tensor, contradiction. No preservation of the selected $E^4$
factor by the lift is assumed.

Weak Lefschetz gives $H^1(E^6)\simeq H^1(B)$ and
$H^4(E^6)\hookrightarrow H^4(B)$. Therefore the intrinsic cup map
\[
                     \bigwedge^4H^1(B)\longrightarrow H^4(B)
\]
is injective, and its rank and image transport through any abstract
smooth proper lift. If $\kappa^*a$ lifted to an algebraic class
$\gamma$, its $\mathbf Q_5$ realization would remain in this image.
After spreading to a complex fibre, rational Betti $\gamma$ lies
in the rational cup-image too: membership can be tested after
scalar extension to $\mathbf Q_5$.
Its unique rational preimage is $(2,2)$, because cup product is
an injective morphism of Hodge structures. The preimage retains
its eight-dimensional support with zero infinitesimal stabilizer.
The same Hodge-circle contradiction excludes the abstract lift.
Finally hard Lefschetz gives
$L_h^2:H^4(E^6)\simeq H^8(E^6)$, and the nonzero $a$ satisfies
(A10.1) by self-intersection. \hfill$\square$

\paragraph{Transfer obstructions.}
For an axis of dimension $2r$ in an abelian $(2r+2)$-fold,
every $w\in\ker\kappa_*\subset H^{2r}(B)$ satisfies
\[
                   w\,\kappa^*\alpha=0\quad
                         (\alpha\in H^{>0}(X)).
\]
Its Gysin image is zero by projection formula; Gysin above the
middle is injective by dual weak Lefschetz. Products beyond top
degree vanish automatically. Thus the displayed $H^1$ tensor
cannot be recovered from an axis-kernel class by ambient cup
products.
Likewise, if $i:S\hookrightarrow B$ has codimension $c<r$ and
$\eta=i^*\kappa^*\theta$ with $\theta\in H^{2r-2c}(X)$, then
$i_*\eta=\kappa^*\theta\,[S]_B$ is ambient, since weak Lefschetz
makes $[S]_B\in H^{2c}(B)$ ambient. Its variable projection is zero.

\paragraph{Sources, scope and genuine checks.}
Primary URLs:
\texttt{https://jmilne.org/math/CourseNotes/AV.pdf};
\texttt{https://www.jmilne.org/math/Documents/Deligne82.pdf}.
Smooth proper base change is SGA4, Expos\'e XVI, Th\'eor\`eme 2.2,
together with Betti--\'etale comparison in characteristic zero.
C44-A05 lifts geometry; this theorem shows a single algebraic
middle class need not lift. It lies wholly in the ambient summand,
where the sign is covered by the conditional 032 input.
It gives neither a nonliftable $A_B$ class nor a wrong-sign class;
QH9-axis remains open.

The author checked the whole proof on 22 November 2026
(Europe/Berlin), including the actual integral tensor and all
lifting quantifiers. The root independently read all three
complete proof messages and passed the whole proof at 00:49:30 UTC
that date, also reproducing the exact determinant independently.
Both agents are GPT-6 Astra Pro; reasoning effort is not exposed.
The closed divisor formulas (A10.4) were subsequently checked by
hand and exact modulo-five JavaScript. None of these events is
a claimed verification of QH9-axis.

% END work/axis/axis_supersingular_nonliftable.tex

% BEGIN work/axis/axis_product_ci_matching.tex
% C44_POS /root/axis_independent; GPT-6 Astra Pro; effort not exposed.
\section{A sharp fixed-polarization bound for products of complete-intersection curves}
\subsection*{C44-A11. Two is the maximum for smooth complete intersections}

\textbf{Theorem.}
Let $S,T$ be abelian threefolds, $L,M$ ample line bundles of
classes $s,t$, and $B\subset S\times T$ a smooth complete
intersection of two sections of $L\boxtimes M$. Put
$V=\int_Ss^3$, $W=\int_Tt^3$.
Suppose it contains $Z_i=C_i\times D_i$, $1\le i\le N$,
where the $C_i$ are pairwise disjoint smooth scheme-theoretic
complete intersections of two divisors in $|L|$, and the $D_i$
are pairwise disjoint smooth scheme-theoretic complete
intersections of two divisors in $|M|$. Then $N\le2$.
Set
\[
 a=st-\tfrac12(s^2+t^2),\qquad w_i=[Z_i]_B-\kappa^*a.           \tag{A11.1}
\]
Then $w_i\in A_B$ and
\[
                         (w_iw_j)_B=VW(2\delta_{ij}-1).       \tag{A11.2}
\]
At $N=2$, $w_1+w_2=0$ in rational cohomology, and the actual
span is a positive line. For a Lefschetz pencil in the usual
one-node-per-singular-member sense the stronger bound $N\le1$
holds (C44-A11b below).

\paragraph{1. Section spaces and finite maps.}
Put $A=H^0(S,L)$, $F=H^0(T,M)$ and
$U_i=H^0(\mathcal I_{C_i}L)$, $V_i=H^0(\mathcal I_{D_i}M)$.
Each has dimension two by the twisted Koszul sequence
\[
 0\longrightarrow L^{-1}\longrightarrow\mathcal O_S^2
                     \longrightarrow\mathcal I_{C_i}L\longrightarrow0.
\]
Here $H^0(L^{-1})=H^1(L^{-1})=0$, using Serre duality and
ample-line-bundle cohomology vanishing on abelian varieties;
the same argument applies to $M$.
For $i\ne j$, $U_i\cap U_j=0$. Otherwise the two curves'
combined defining equations would use at most three sections
of $L$. Such sections have a common zero: a nowhere-zero
section of $L^{\oplus q}$, $q\le3$, would imply $s^q=0$,
contradicting $\int_Ss^3>0$. A common zero belongs to
$C_i\cap C_j$, impossible. Similarly $V_i\cap V_j=0$.

Let $f_0,f_1\in A\otimes F$ define $B$. Vanishing on the products
gives $f_\nu\in U_i\otimes F+A\otimes V_i$, exactly the kernel
of the tensor product of the two restriction maps.
If $N\ge2$, set $A_0=U_1\oplus U_2$, $F_0=V_1\oplus V_2$.
Direct-sum complements give
\[
 (U_1\otimes F+A\otimes V_1)\cap(U_2\otimes F+A\otimes V_2)
                         =U_1\otimes V_2+U_2\otimes V_1.     \tag{A11.3}
\]
Thus both equations lie in the four-by-four tensor space
$A_0\otimes F_0$.
The four sections in $A_0$ have common zero scheme
$C_1\cap C_2=\varnothing$, and give
$\phi:S\to\mathbf P(A_0^\vee)\simeq\mathbf P^3$ with
$\phi^*\mathcal O(1)=L$. Ampleness excludes positive-dimensional
fibres, so $\phi$ is finite and surjective, of degree $V$.
Likewise $\psi:T\to\mathbf P^3$ is finite surjective of degree
$W$. Both are flat by the Cohen--Macaulay/regular-local
criterion for finite morphisms between smooth equidimensional
varieties.

\paragraph{2. A rank-three Chern obstruction.}
Every further $U_i$ lies in $A_0$, and every $V_i$ lies in $F_0$.
Suppose $\dim(U_i\cap A_0)\le1$. Restricting $f_\nu$ to
$S\times D_i$ places its terms in
$(U_i\cap A_0)\otimes H^0(D_i,M|_{D_i})$.
The derivative in the $S$ direction along $Z_i$ therefore uses
at most one constant conormal generator. Terms whose coefficients
vanish on $D_i$ have zero derivative in that direction.
Since $N_{C_i/S}\simeq(L|_{C_i})^{\oplus2}$, the full normal
surjection required by smoothness of $B$ factors through
\[
 (L|_{C_i})\boxtimes\mathcal O_{D_i}\ \oplus\
 \mathcal O_{C_i}\boxtimes(M|_{D_i})^{\oplus2}
                 \longrightarrow (L\boxtimes M)|_{Z_i}^{\oplus2}.
                                                               \tag{A11.4}
\]
If the first block is zero, add an unused first summand.
This rank-three surjection has a line-bundle kernel.
For $x=c_1(L|_{C_i})$, $y=c_1(M|_{D_i})$, one has $x^2=y^2=0$,
and the virtual second Chern class of the kernel is
\[
                    \left[\frac{(1+x)(1+y)^2}{(1+x+y)^2}\right]_2
                           =2xy.
\]
Its degree is $2VW>0$, contradicting vanishing of a line bundle's
second Chern class. Thus $U_i\subset A_0$; interchanging factors
gives $V_i\subset F_0$.
Each $C_i$ is consequently the scheme-theoretic inverse image
of a line $\ell_i\subset\mathbf P^3$, and each $D_i$ is the
inverse image of a line $m_i$. Surjectivity shows that each
family of lines is pairwise disjoint.

\paragraph{3. Classification of the smooth bilinear pencil.}
The equations define $B_0\subset\mathbf P^3\times\mathbf P^3$
of bidegrees $(1,1),(1,1)$ with
$B=(\phi\times\psi)^{-1}(B_0)$ scheme-theoretically.
The codimension is two, since its finite pullback has codimension
two. Moreover $B_0$ is smooth: at a singular point the two
gradients are dependent, and remain so after composition with
any differential. Every point has a preimage, which would then
be singular in $B$. Each $\ell_i\times m_i$ lies in $B_0$.

Write its equations as $x^{\mathsf t}M_0y=x^{\mathsf t}M_1y=0$.
Smoothness forces the determinant of this $4$-by-$4$ pencil to
be nonzero with four distinct projective roots, in every
characteristic. For a pencil member of rank at most two, its
left and right kernels have dimension at least two. Choose
nonzero $x,y$ in them with $x^{\mathsf t}M'y=0$ for an
independent member $M'$. This gives a point of $B_0$ with
dependent gradients. For a rank-three member, the two kernel
lines yield a singular point exactly when
$x^{\mathsf t}M'y=0$. By the adjugate formula this is precisely
vanishing of the determinant's derivative at that root.
If the determinant were identically zero, use a rank-three
member and the same derivative argument, or the lower-rank
case. Thus neither an identically zero determinant nor a
multiple root is possible.

After choosing an invertible member and changing coordinates,
\[
 M_0=1,\qquad M_1=\operatorname{diag}(\lambda_1,\ldots,\lambda_4),
                      \quad\lambda_i\ne\lambda_j\ (i\ne j).   \tag{A11.5}
\]
A product $\mathbf P(U)\times\mathbf P(V)$ of lines lies in
$B_0$ exactly when $V=U^\perp$ for $M_0$ and $U$ is invariant
under $M_1^{\mathsf t}=M_1$. Distinct eigenvalues make every
invariant two-plane a span of two eigenvectors. Hence all such
products are
\[
 Z_I^0=\mathbf P(\langle e_i:i\in I\rangle)
          \times\mathbf P(\langle e_j:j\notin I\rangle),
                        \qquad |I|=2.                        \tag{A11.6}
\]
Two distinct products are disjoint exactly when their two-subsets
are complementary. Otherwise their subsets and their complements
each share one index, producing an intersection point.
A pairwise-disjoint collection thus has at most two members:
$N\le2$.

\paragraph{4. Gram matrix and actual boundary relation.}
In the Chow ring of $S\times T$, $[Z_i]=s^2t^2$ and
$(s+t)^2(st-\tfrac12(s^2+t^2))=s^2t^2$ because $s^4=t^4=0$.
Thus $w_i$ lies in the Gysin kernel. The normal sequence gives
\[
 [Z_i]_B^2=\int_{Z_i}
       \left[\frac{(1+x)^2(1+y)^2}{(1+x+y)^2}\right]_2=2VW.
\]
Distinct products are disjoint and the ambient projection term
is $\int_{S\times T}a[Z_i]=VW$, proving (A11.2).

At $N=2$, zero norm alone cannot establish a zero class.
On $B_0$ put $a_0=xy-\tfrac12(x^2+y^2)$ and
$w_I^0=[Z_I^0]_{B_0}-a_0|_{B_0}$.
The same calculation with $V=W=1$ gives
$(w_I^0+w_{I^c}^0)^2=0$ and puts the sum in the Gysin kernel.
It is Lefschetz primitive: Gysin in degree six is injective
by dual weak Lefschetz in degree two.
Lift the diagonal model and all coordinate products to a
mixed-characteristic discrete valuation ring by lifting the
distinct $\lambda_i$ with their differences invertible.
The same matrix criterion gives a smooth projective model.
On a complex fibre the sum is an algebraic primitive $(2,2)$
class of square zero, hence zero by Hodge--Riemann positivity.
If needed, spread the finite data to a finitely generated
characteristic-zero field before embedding into $\mathbf C$.
Smooth proper base change proves
$w_I^0+w_{I^c}^0=0$ in the original rational cohomology.
This lifting is only of the explicitly diagonal bilinear
variety, not of the given abelian varieties or curve collections.

The finite flat map $f_B:B\to B_0$ pulls the two products back
to $Z_1,Z_2$ and $x,y$ back to $s,t$. Thus $w_1+w_2=0$.
Each $w_i^2=VW>0$, so the actual span is a positive line;
for $N=1$ positivity is immediate. \hfill$\square$

\subsection*{C44-A11a. The six products give a positive three-dimensional subspace}

Under the $N=2$ reduction define, for all six two-subsets,
\[
 \widetilde Z_I=f_B^*[Z_I^0],\qquad
 \widetilde w_I=\widetilde Z_I-\kappa^*a=f_B^*w_I^0.
\]
They are algebraic cycle classes even if additional inverse-image
curves are singular. Their ambient classes are $s^2t^2$, so the
$\widetilde w_I$ lie in the Gysin kernel.
For distinct noncomplementary $I,J$, write
$I\cap J=\{i\}$, $I^c\cap J^c=\{j\}$.
The products meet in the unique point $(e_i,e_j)$.
Linearizing the two equations gives
$dy_i+dx_j=0$ and $\lambda_i\,dy_i+\lambda_j\,dx_j=0$,
hence $dy_i=dx_j=0$.
The four remaining tangent coordinates on $B_0$ split into
complementary pairs used by the two surfaces. Their intersection
is transverse of multiplicity one. Together with self-intersection
two and complementary disjointness, subtraction of the ambient
term one gives
\[
 (w_I^0w_J^0)=
 \begin{cases}
 1,&J=I,\\
 -1,&J=I^c,\\
 0,&\text{otherwise}.
 \end{cases}                                                \tag{A11a.1}
\]
Each complementary sum vanishes by the explicit lifting argument.
Thus $w_{12}^0,w_{13}^0,w_{14}^0$ form an orthonormal basis of
the span of the six classes. Projection formula for the finite
degree-$VW$ map yields
\[
             (\widetilde w_{1i}\widetilde w_{1j})_B
                   =VW\delta_{ij},\qquad i,j=2,3,4.          \tag{A11a.2}
\]
This is an unconditional positive three-dimensional subspace for
smooth complete intersections. It is not asserted as an H9
Lefschetz-axis example, because of the next obstruction.

\subsection*{C44-A11b. The Lefschetz hypothesis forces $N\le1$}

Suppose every singular pencil member has exactly one ordinary
quadratic singularity. For $N\ge2$, the reduction above contradicts
that requirement.
First the pencil is unique for its scheme-theoretic axis in the
given linear system. If $H=L\boxtimes M$, twisted Koszul gives
\[
 0\longrightarrow H^{-1}\longrightarrow\mathcal O_{S\times T}^2
                    \longrightarrow\mathcal I_BH\longrightarrow0.
\]
Serre duality and ample cohomology vanishing give
$H^0(H^{-1})=H^1(H^{-1})=0$, so
$H^0(\mathcal I_BH)=\langle f_0,f_1\rangle$.
There is no alternative $H$-pencil with the same axis.

For each $i$, the target member
$\sum_j(\lambda_j-\lambda_i)x_jy_j=0$ has its unique singular
point $p_i=(e_i,e_i)$ outside $B_0$, since $\sum_jx_jy_j=1$ there.
In affine coordinates $x_i=y_i=1$ its equation is exactly
\[
                    q_i=\sum_{j\ne i}(\lambda_j-\lambda_i)x_jy_j,
                                                               \tag{A11b.1}
\]
a nonsingular quadratic form on six variables in every
characteristic.
Every preimage under $g=\phi\times\psi$ is a singular point of
the pulled-back member, with quadratic tangent form $q_i\circ dg$.
If $dg$ is not injective, a nonzero vector in its kernel is both
isotropic and polar-radical for this quadratic form. Its
projective quadric is singular; hence the hypersurface does not
have an ordinary quadratic singularity. This argument includes
characteristic two.

If all preimages instead have ordinary quadratic singularities,
$dg$ is invertible at each. The finite flat map is then \'etale
there, so the fibre consists of $VW$ distinct geometric points.
That one member has at least $VW$ nodes. Riemann--Roch makes
$V,W$ positive multiples of $3!$, so $VW>1$.
Thus either a forbidden nonordinary singularity occurs or more
than one node lies in the same member. These points are outside
the axis; blowing it up does not change the contradiction.
Therefore a genuine Lefschetz axis satisfies $N\le1$.
\hfill$\square$

\paragraph{Sources, scope and genuine checks.}
The cohomology and Riemann--Roch inputs are Milne,
\emph{Abelian Varieties}, v.~2.00, I Theorem 11.1(c),(b),
printed p.~54:
\texttt{https://jmilne.org/math/CourseNotes/AV.pdf}.
Nonzero sections exist here, so the unique nonzero cohomology
degree in part (c) is zero; Serre duality gives the negative
vanishings. Flatness uses Stacks Project, Lemma 10.128.1,
Tag 00R4: \texttt{https://stacks.math.columbia.edu/tag/00R4}.
The block identity, Chern obstruction, pencil classification,
coordinate intersections and singularity-count argument are
proved explicitly above.

Both curve families separately must be pairwise disjoint, and
each curve must be cut out by two divisors of the given
polarization. Arbitrary product-cycle collections are not covered.
The $N=2$ three-dimensional positive subspace is a smooth-CI
result; it cannot supply the claimed configuration in the usual
H9 Lefschetz setup. QH9-axis remains open.

The author checked the entire proof including both addenda on
22 November 2026 (Europe/Berlin). The root independently read
all three complete A11 proof parts and both complete addenda,
recording whole-proof PASS at 00:58:45 UTC that date.
Checks included the rank-three Chern obstruction, all-characteristic
matrix criterion, actual boundary relation, transverse products,
and strict Lefschetz qualification. Both agents are GPT-6 Astra Pro;
reasoning effort is not exposed.

% END work/axis/axis_product_ci_matching.tex

% BEGIN work/axis/relative_form_transfer.tex
% C44_POS /root; GPT-6 Astra Pro; effort not exposed.
\section{A relative correspondence criterion that retains the geometric form}

\begin{theorem}[C44-A08: positive-spectrum transfer along an extending correspondence]
Let $R$ be a discrete valuation ring of mixed characteristic $(0,p)$,
with fraction field admitting an embedding into $\mathbf C$.
Fix a prime $\ell\ne p$ for the special-fibre cohomology.
Let $\mathscr X/R$ be smooth projective of relative dimension $n$,
with a relatively ample class $h$, and let $0\le 2r\le n$.
Let $V$ be a relative algebraic direct summand of
$P^{2r}(\mathscr X)$, self-dual for
\[
 q_X(x,y)=(-1)^r\int_{\mathscr X}xyh^{n-2r}.
 \tag{C44.A08.1}
\]
Assume its projector and the inverse of this restricted pairing
are represented by algebraic correspondences over $R$ at the
level of Weil cohomology. Suppose there are abelian schemes
$\mathscr A_\nu/R$ and algebraic correspondences over $R$ giving
an injection on every geometric fibre
\[
 \alpha:V\longrightarrow
       \bigoplus_\nu H^{2s_\nu}(\mathscr A_\nu)(s_\nu-r).
 \tag{C44.A08.2}
\]
A relative split Chow-motive injection is sufficient for this
hypothesis; no such injection is assumed to exist in general.

Then, \emph{conditional on the unrefereed 032/Milne 2026
preprints}, the actual signed geometric form (C44.A08.1) is
positive definite on the rational algebraic classes of $V$ in
the geometric special fibre. In particular, no cycle-lifting
hypothesis is required for those special algebraic classes.
The assertion does not use Milne 2022, Proposition~4.2.
\end{theorem}

\begin{proof}
Choose relatively ample symmetric classes on the abelian schemes,
after a finite extension of the trait if necessary. The algebraic
Lefschetz operators of abelian varieties can be taken relatively.
Explicitly, on a relatively polarized abelian scheme of positive
dimension $g$, with $\theta=c_1(L)$ and $D=\int\theta^g/g!$, the cycles
\[
 L=\Delta_*\theta,\qquad
 \Lambda=\frac{\delta^*(\theta^{g-1})}{D(g-1)!},
 \qquad \delta(x,y)=y-x,
\]
give the raising and lowering operators fibrewise. Their
$\mathfrak{sl}_2$ identities and the usual rational polynomial
formulas give relative cohomological primitive and Lefschetz
projectors. Zero-dimensional factors are treated as Tate summands.
On every even-degree target factor choose the
orthogonal direct sum of its signed primitive Lefschetz pairings,
with positive rational coefficients on the summands, and include
the displayed Tate shifts. Write $Q_A$ for the resulting direct-sum
pairing. It is represented by an algebraic correspondence over
$R$. With the common real normalization of Tate factors, its
complex realization satisfies
\[
                  Q_A(Cu,\overline u)>0\quad(u\ne0),
 \tag{C44.A08.3}
\]
where $C$ is the Weil operator. This is the classical
Hodge--Riemann theorem on each primitive summand, extended by the
orthogonal Lefschetz decomposition. On algebraic classes in the
special fibre the same $Q_A$ is positive definite, conditional
on 032 Corollary~8.4 and Milne's 2026 report. Indeed its primitive
projections are algebraic and each signed primitive form has
that positivity. This is precisely where the unrefereed input
enters.

Pull back this pairing and put
\[
                  q_1=\alpha^*Q_A,\qquad
                  c=q_X^{-1}q_1\in\operatorname{End}(V).
 \tag{C44.A08.4}
\]
The assumed algebraic inverse and the transpose of $\alpha$
make $c$ an algebraic correspondence over the trait, at the
cohomological level required here. The characteristic polynomial
of its Betti action is a polynomial $P(T)\in\mathbf Q[T]$.
We show that every root of $P$ is a strictly positive real number.

On the complex generic fibre both $q_X$ and $q_1$ are
polarizations in the following explicit sense: their associated
Hermitian forms
\[
 H_X(x,y)=q_X(Cx,\overline y),\qquad
 H_1(x,y)=q_1(Cx,\overline y)
\]
are positive definite. For $q_X$ this is the primitive
Hodge--Riemann relation; for $q_1$ it follows from the Hodge
compatibility and injectivity of $\alpha$ and (C44.A08.3).
The algebraic correspondence $c$ preserves Hodge types and the
real structure, hence commutes with $C$. Symmetry of the two
bilinear forms gives
\[
 H_X(cx,y)=H_X(x,cy),\qquad H_X(cx,x)=H_1(x,x)>0
 \quad(x\ne0).
\]
Consequently $c$ is self-adjoint and positive for $H_X$.
It is diagonalizable with strictly positive real eigenvalues,
as claimed. This argument uses the common Weil operator; the
indefinite geometric bilinear form on all Betti classes is not
mistaken for a positive form.

Smooth proper base change and compatibility of cycle classes
with specialization identify the characteristic polynomial of
$c$ in the two geometric fibres. Thus $P(c)=0$ in the special
fibre, and $P(0)\ne0$. Cayley--Hamilton expresses $c^{-1}$ as a
polynomial in $c$ with rational coefficients. Both $c$ and
$c^{-1}$ therefore preserve special rational algebraic classes.

For complete precision no unproved rational dimension assertion
about the space of algebraic classes is needed. For any such
class $z$, let
\[
 U_z=\operatorname{span}_{\mathbf Q}
       \{z,cz,\ldots,c^{d-1}z\},\qquad d=\deg P.
\]
It is finite-dimensional, consists of algebraic classes, and is
stable under $c$ and $c^{-1}$. On $U_z$, $q_1$ is positive definite:
$\alpha z\ne0$ for $z\ne0$, and $Q_A$ is positive on algebraic
classes by the conditional input. Moreover,
\[
 q_1(cx,y)=q_X(c^2x,y)=q_X(cx,cy)=q_1(x,cy).
\]
Thus $c$ is $q_1$-self-adjoint on $U_z$. Its eigenvalues are among
the positive roots of $P$, so $c^{-1}$ is positive for $q_1$.
Finally
\[
                   q_X(z,z)=q_1(c^{-1}z,z)>0
                   \qquad(z\ne0).
 \tag{C44.A08.5}
\]
Since $z$ was arbitrary, this proves the assertion on all special
algebraic classes. The newly constructed inverse $c^{-1}$ is an
actual rational polynomial in $c$; algebraicity of $q_X^{-1}$ was
an explicit hypothesis. No square root,
abstract comparison isometry, or unproved positivity transfer
is inserted.
\end{proof}

\begin{corollary}[C44-A08a: a sufficient correspondence condition for a fixed axis]
Let $B$ be a smooth axis in an abelian $(2r+2)$-fold as in
QH9-axis. Suppose an embedded mixed-characteristic lift from
C44-A05 carries an extending algebraic injection of its variable
middle cohomology into a sum of Tate-shifted even cohomology of
abelian schemes. Then the H9-T7 axis form on $A_B$ is positive
definite, \emph{conditional on the unrefereed 032/Milne 2026
preprints}.
\end{corollary}

\begin{proof}
The variable projector in middle cohomology is
\[
       p_{\rm var}=1-\kappa^*(L_h^2)^{-1}\kappa_*.
 \tag{C44.A08.6}
\]
It is algebraic relatively, since the ambient abelian Lefschetz
inverse is algebraic relatively. It is self-adjoint for the
middle intersection pairing. Its image is $\ker\kappa_*$,
which is primitive by C44-A01; the restriction of the middle
pairing to this image is nondegenerate. Its inverse as a duality
morphism is represented by the transpose self-adjoint projector
and the middle Poincare duality correspondence. Hence all the
hypotheses of C44-A08 hold, and (C44.A08.1) is exactly
$(-1)^r\int_Bww'$, the second summand of H9-T7.
\end{proof}

\paragraph{Exact remaining gate and Grothendieck's frameworks.}
This is a criterion, not an assertion that every axis admits
such a relative abelian correspondence. C44-A05 lifts the
embedded geometry but supplies no injection (C44.A08.2).
C44-A02 excludes graph maps as a source for a nonzero variable
injection. An arbitrary characteristic-$p$ correspondence that
does not extend to a characteristic-zero fibre also does not
meet the criterion. Thus QH9-axis is not closed. The distinction
is the same geometric-polarization issue exposed by the prior
JAC-P20 transfer audit, but the present hypothesis permits an
explicit proof via a positive spectrum. It gives a sufficient
condition for cote~44 Th.~5, scan~004, with the primitive
qualification of scan~008, and proves the relevant $I$ summand
in the terminology of cote~16 scans~049--050.

The cohomological tools are Weil cohomology, Poincare duality,
Lefschetz decomposition, and smooth proper base change (SGA~4,
XVI, Corollary~2.2), together with algebraic correspondences and
Cayley--Hamilton. Relative Fourier/Lefschetz operators on an
abelian scheme are supplied by the same Poincare-bundle
correspondences as Beauville, \emph{The action of SL(2) on abelian
varieties}, JRMS 25 (2010), Theorem~3.3, pp.~258--259, and
Theorem~4.2, p.~260;
their identities are fibrewise identities in Weil cohomology.
The characteristic-zero sign is the Hodge--Riemann bilinear
relation (Voisin, \emph{Hodge Theory and Complex Algebraic
Geometry I}, Theorem~6.32). The only conjectural positivity
input is the expressly conditional 032/Milne assertion on the
abelian special fibre.

% END work/axis/relative_form_transfer.tex

% BEGIN work/axis/axis_frobenius_selection.tex
\section{Selecting closed fibres with exactly the prescribed algebraic part}
\label{sec:c44-a12}

\begin{proposition}[C44-A12: even orthogonal complement and closed-point selection]
Let \(U/\mathbf F_q\) be smooth, geometrically connected and
quasiprojective, and let \(f:\mathcal B\to U\) be smooth projective
of relative dimension \(2r\). Fix \(\ell\ne\operatorname{char}\mathbf F_q\).
Let \(\mathcal V\subset R^{2r}f_*\mathbf Q_\ell(r)\) be a lisse
orthogonal direct summand for the signed middle pairing
\(Q(x,y)=(-1)^r\int xy\).
Suppose that rational relative cycle classes
\(z_1,\ldots,z_s\) lie in \(\mathcal V\), are linearly independent
and have constant nonsingular rational Gram matrix \(G\).
Write
\[
 \mathcal V=\mathcal K\perp\mathcal W,\qquad
 \mathcal K=\sum_i\mathbf Q_\ell z_i .
\]
Assume that \(d=\operatorname{rank}\mathcal W>0\) is even, that the
geometric monodromy closure on \(\mathcal W\) contains
\(\operatorname{SO}(\mathcal W,Q)\), and that every normalized
closed-point Frobenius characteristic polynomial on \(\mathcal W\)
lies in \(\mathbf Q[T]\).
Then a Zariski-dense set of closed points \(u\) has
\[
 \Alg^r(\mathcal B_{\bar u})\cap\mathcal V_{\bar u}
       =\sum_i\mathbf Q z_{i,\bar u}.                    \tag{A12.1}
\]
In particular, if \(G\) is positive definite, the entire algebraic
part of this summand is positive definite at those closed points.
No Tate conjecture, cycle-lifting assertion or 032 input is used.
\end{proposition}

\begin{proof}
Put \(N=\operatorname{lcm}(1,\ldots,2d^2)\).
The elementary cyclotomic degree bound used in C44-A03 says that a
root of unity which is a root of a rational degree-\(d\) polynomial
has order at most \(2d^2\). Thus a Frobenius polynomial here has a
root-of-unity root if and only if
\(\det(F^N-1)=0\), without a semisimplicity assumption.

On the even orthogonal group's identity component the regular function
\[
                     P(g)=\det(g^N-1)
\]
is not identically zero. After passing to an algebraic closure,
a maximal torus has eigenvalues
\(t_1,t_1^{-1},\ldots,t_{d/2},t_{d/2}^{-1}\), and one may choose
every \(t_i^N\ne1\).
The geometric monodromy image therefore contains some element
\(\gamma\) with \(P(\gamma)\ne0\).

Let \(\Gamma\) be the compact arithmetic monodromy image on
\(\mathcal W\). The subset \(P\ne0\) in \(\Gamma\) is open for
the \(\ell\)-adic topology and invariant under conjugation.
Choose an open normal subgroup \(K\) of \(\Gamma\) such that
\(\gamma K\) lies in this subset. The union of the conjugates of
\(\gamma K\) still lies there. It is the inverse image of a
nonempty conjugacy class in the finite quotient \(\Gamma/K\).

The finite quotient is realized by a finite etale Galois cover of
\(U\). Finite-field Chebotarev supplies a closed point whose
Frobenius image lies in that class. For precision concerning the
constant field, the selected class lies in the geometric subgroup
because \(\gamma\) was chosen geometrically. One may first extend
the constants so that the finite quotient cover has a geometrically
connected component and then apply Chebotarev to that component.
Equivalently, one takes extension degrees in the appropriate
constant-field congruence class in the general theorem.
The resulting Frobenius has \(P(F)\ne0\), hence has no
root-of-unity eigenvalue on \(\mathcal W\).
If one uses a rational point over an extension rather than a closed
point over the original field, its Frobenius is a positive power of
the latter. A root-of-unity eigenvalue of the latter would also
give one for that power, so the conclusion descends to its
underlying closed point.

The same argument applies after replacing \(U\) by any nonempty
open subset: its fundamental group surjects onto that of \(U\),
so the monodromy image does not shrink. This proves the claimed
Zariski density.

Now let \(z\) be an algebraic rational cycle class in
\(\mathcal V_{\bar u}\). Its orthogonal projection to \(\mathcal K\)
has coefficients
\[
  (a_1,\ldots,a_s)^{\mathsf t}
       =G^{-1}\bigl(Q(z,z_1),\ldots,Q(z,z_s)\bigr)^{\mathsf t}
       \in\mathbf Q^s .                                  \tag{A12.2}
\]
The rationality follows from intersection numbers of rational
algebraic cycles; it is not merely a statement about
\(\mathbf Q_\ell\)-linear spans.
Thus \(z-\sum_i a_i z_i\) is still an algebraic class and lies
in \(\mathcal W_{\bar u}\).
Every algebraic cycle over \(\overline{\mathbf F}_q\) is defined
over a finite field extension. After the Tate twist its class
is fixed by a positive power of Frobenius.
There are no such nonzero vectors on \(\mathcal W_{\bar u}\),
because its Frobenius has no root-of-unity eigenvalues.
This proves (A12.1), and the sign assertion follows from \(G\).
\end{proof}

\paragraph{Exact Chebotarev source.}
The finite-quotient existence statement, including the constant-field
coset, is Alexei Entin, \emph{Monodromy of hyperplane sections of curves
and decomposition statistics over finite fields}, Theorem 3 in
\S4, author version arXiv:1805.05454v3, pp.~9--11;
published in \emph{IMRN} 2021, pp.~10409--10441,
doi:10.1093/imrn/rnz120.
Only its nonempty-class conclusion over sufficiently large finite
extensions is used. No numerical error coefficient from the
Chebotarev sketch in Stacks, Tag 03W3, is used.

\begin{corollary}[C44-A12a: application to a family of axes]
Suppose a family of axes has constant abelian ambient variety,
with its variable middle summand \(\mathcal V=\ker\kappa_*\).
If a known positive cycle span has even orthogonal complement
with the monodromy hypothesis of C44-A12, then a Zariski-dense
set of closed axes has exactly that algebraic variable span.
The rational-polynomial hypothesis holds whenever the full
middle and ambient Frobenius polynomials are rational:
the polynomial on the complement is their quotient, followed
by removing the factors supplied by the known cycle classes.
\end{corollary}

\begin{proof}
Weak Lefschetz and hard Lefschetz give the orthogonal splitting
of the middle cohomology into the ambient image and the variable
kernel. The ambient image has the ambient Frobenius polynomial.
The classes \(z_i\) have normalized Frobenius eigenvalue one,
so their factor is \((T-1)^s\).
The quotient polynomial is rational. C44-A12 applies.
\end{proof}

\paragraph{Relation to the main question.}
For cote 44, Th.~5, scan 004 and the primitive margin on scan 008,
this is a sufficient construction principle for complete control
of \(A_B\) at actual finite-field points, rather than merely at
a geometric generic parameter point. It does not assert the
monodromy hypothesis for an arbitrary fixed axis.
The odd-rank obstruction of C44-A03b explains why the known
cycle must first be removed in the abelian-sixfold case.

% END work/axis/axis_frobenius_selection.tex

% BEGIN work/scan_w/closed_rank_one_axis.tex
% C44_POS /root/scan_w; complete draft, independent review requested.
\section{Closed axes with a prescribed nonzero rank-one algebraic kernel}
\label{sec:c44-a13}

\begin{theorem}[C44-A13: closed rank-one axes in a supersingular abelian sixfold]
Let $E/\mathbf Q$ be the elliptic curve $y^2=x^3-x$, with origin $o$,
and put
\[
 A=E^6,\qquad H=\sum_{i=1}^6h_i,\qquad
 h_i=\operatorname{pr}_i^*[o],\qquad
 Y=E^2\times\{o\}^4.
\]
For a sufficiently large integer $t$, there is a one-parameter family
of genuine smooth axes
\[
 B_u=D_0\cap D_u\subset A,\qquad D_0,D_u\in|tH|,
 \qquad Y\subset B_u,
 \tag{A13.1}
\]
defined on a nonempty open $U\subset\mathbf P^1_{\mathbf Q}$.
After excluding finitely many primes (depending on this model and a
fixed coefficient prime $\ell$), there are infinitely many primes
$p\equiv3\pmod4$ for which the following holds.
The reduction $A_p=E_p^6$ is supersingular, and a Zariski-dense set
of closed points $u$ of $U_{\mathbf F_p}$ has
\[
 A_{B_u}:=\Alg^2(B_{\bar u})\cap
 \ker\bigl(\kappa_*:H^4(B_{\bar u})\longrightarrow H^8(A_{\bar u})\bigr)
       =\mathbf Q w_u,
 \tag{A13.2}
\]
where $w_u$ is the explicitly algebraic primitive class
\[
 \begin{split}
 \alpha_0={}&\frac14h_1h_2
 -\frac1{12}\sum_{i\in\{1,2\},\,j\in\{3,4,5,6\}}h_ih_j
 +\frac1{12}\sum_{3\leq i<j\leq6}h_ih_j,\\
 w_u={}&[Y]_{B_u}-t^{-2}\kappa^*\alpha_0.
 \end{split}
 \tag{A13.3}
\]
Its actual geometric intersection form is
\[
                    \int_{B_u}w_u^2
                       =6t^2-\frac1{4t^2}>0.
 \tag{A13.4}
\]
Thus the required sign on the entire nonzero $A_{B_u}$ is
unconditional at these closed finite-field axes. Their full variable
cohomology is nonzero, of rank $3600t^6-165$.
The original pencil generated by $D_0,D_u$ may be required to be a
Lefschetz pencil. No Tate conjecture or unrefereed 032 result is used.
The construction does not assert that the entire variable motive
is of abelian type.
\end{theorem}

\begin{proof}
\emph{The family and its two pencils.}
Let $\pi:Z=\operatorname{Bl}_Y A\to A$ and let $F$ be the exceptional
divisor. For $t\gg0$, the line bundle
$L'=\pi^*\mathcal O_A(tH)\otimes\mathcal O_Z(-F)$ is very ample.
Its sections are exactly the sections of $\mathcal I_Y(tH)$.
Choose a general member $\widetilde D_0$ of $|L'|$.
Its image $D_0\subset A$ is smooth and contains $Y$, and
$\widetilde D_0=\operatorname{Bl}_Y D_0$.
Indeed the normal derivative of a section along $Y$ is a section
of $\mathcal O_Y(tH)^{\oplus4}$. For $t$ sufficiently large these
derivatives are generated by the chosen section space; their common
zero locus has expected codimension four, exceeding $\dim Y=2$.
Away from $Y$, Bertini applies. Likewise two general sections give
a $2\times4$ matrix of normal derivatives, whose rank-one locus
has expected codimension three. It misses $Y$, so their original
intersection $B$ is smooth along $Y$ as well as away from it.
These statements can equivalently be obtained by applying the
usual jet-generation bounds to $\mathcal I_Y(tH)$.

Embed $Z$ by $L'$. The induced embedding of $\widetilde D_0$ admits
Lefschetz pencils in characteristic zero: SGA~7~II, XVII,
Theorem~2.5.2, printed p.~217. Choose a general pencil of these
hyperplane sections of $\widetilde D_0$. It gives a family
\[
 \widetilde B_u=\widetilde D_0\cap\widetilde D_u
                  =\operatorname{Bl}_Y B_u.
 \tag{A13.5}
\]
The restriction map from all sections on $Z$ to the complete
section space on $\widetilde D_0$ need not be surjective; the
induced projective embedding is all that is used.
Remove the finitely many auxiliary-pencil critical values and
any additional values for which the original $B_u$ is not smooth.
The open that remains is nonempty by the preceding general-pair
argument. We may also require that the line spanned by
$D_0,D_u$ give a Lefschetz pencil on $A$. This is an open nonempty
condition on general pairs in this system. It follows alternatively
from the Lefschetz condition for the corresponding pencil on $Z$:
smoothness of $B_u$ along $Y$ makes every member of that pencil
smooth along $Y$, so its critical points are away from the
exceptional divisor and coincide with the original critical points.
For a general fixed $\widetilde D_0$, the open set of Lefschetz
lines through that member is nonempty. Choose the auxiliary
pencil generically so that its general member gives such a line;
then remove the finitely many parameters outside this open.
The auxiliary pencil varying the axis and the original pencil
having that axis are different pencils.

All the required nonempty open conditions are defined over
$\mathbf Q$. Choose the first member and the auxiliary pencil
over $\mathbf Q$ in these opens. This proves the asserted
number-field model of (A13.1).

\emph{The primitive class and its norm.}
Write $h=tH$ and $\kappa:B\hookrightarrow A$.
The algebraic ambient projection in degree four is
$\kappa^*(L_h^2)^{-1}\kappa_*$; it is defined since
$\kappa_*\kappa^*=L_h^2:H^4(A)\to H^8(A)$ is an isomorphism.
For the particular class $[Y]=h_3h_4h_5h_6$, direct multiplication,
using $h_i^2=0$, gives
\[
                  H^2\alpha_0=[Y].
 \tag{A13.6}
\]
For clarity, the coefficients of a four-index monomial are
$12(1/12)=1$ when all indices are in $\{3,4,5,6\}$,
$2(3(-1/12)+3(1/12))=0$ when precisely one is in $\{1,2\}$,
and $2(1/4+4(-1/12)+1/12)=0$ when both are.
Consequently $(L_h^2)^{-1}[Y]=t^{-2}\alpha_0$, proving
$\kappa_*w=0$ and (A13.3).

The normal bundle sequence is
\[
 0\longrightarrow N_{Y/B}\longrightarrow
 \mathcal O_Y^{\oplus4}\longrightarrow
 \mathcal O_Y(h)^{\oplus2}\longrightarrow0.
 \tag{A13.7}
\]
Thus $c_2(N_{Y/B})=3h^2|_Y$ and
$[Y]_B^2=3\int_Yh^2=6t^2$.
Orthogonal projection subtracts the ambient norm:
\[
 \int_Bw^2
 =\int_B[Y]_B^2-\int_Yt^{-2}\alpha_0
 =6t^2-\frac1{4t^2}.
\]
This computation is valid over every good characteristic, without
positivity assumptions. In particular $w$ is nonzero and
nonisotropic.

\emph{The exact vanishing system.}
Let $q:\widetilde B\to B$ be (A13.5), and write $j$ and $p$ for
the inclusion and projection of its exceptional divisor
$\mathbf P(N_{Y/B})$. The blowup formula gives
\[
 \begin{split}
 H^4(Z)&=H^4(A)\oplus H^2(Y)(-1)\oplus\mathbf Q_\ell(-2),\\
 H^4(\widetilde D_0)&\simeq H^4(Z),\\
 H^4(\widetilde B)&=q^*H^4(B)\perp j_*p^*H^2(Y)(-1).
 \end{split}
 \tag{A13.8}
\]
The middle isomorphism is weak Lefschetz for the ample
fivefold $\widetilde D_0\subset Z$. The last decomposition is
orthogonal: the mixed pairing is zero by the projection formula,
and the exceptional pairing is the negative of the middle pairing
on $H^2(Y)$. Pullback $q^*$ preserves the middle pairing.
The additive blowup formula is SGA~7~II, XVIII,
Theorem~2.2 and (2.2.2), printed pp.~260--261; the stated maps
and pairings follow directly from the projective-bundle formula
and $j^*j_*=-\xi$, where $\xi=c_1\mathcal O_{\mathbf P(N)}(1)$.

Restriction from $Z$ sends its $H^2(Y)(-1)$ summand onto the
exceptional summand in (A13.8). The remaining exceptional
generator restricts to $j_*\xi$. The excess-intersection formula
for the rank-two normal bundle says
\[
 q^*[Y]_B=j_*(\xi+p^*c_1(N_{Y/B})).
 \tag{A13.9}
\]
Therefore the image of $H^4(\widetilde D_0)$ in
$H^4(\widetilde B)$ is exactly
\[
 j_*p^*H^2(Y)(-1)
  +q^*\bigl(\kappa^*H^4(A)+\mathbf Q_\ell[Y]_B\bigr).
 \tag{A13.10}
\]
There is no assertion that a restriction of an ambient
correspondence detects variable cohomology. The additional
class $[Y]_B$ enters through the exceptional divisor and has
nonzero primitive component $w$.

Let $\mathcal V=\ker\kappa_*$ in the family and put
$\mathcal W=\mathcal V\cap w^\perp$.
Equations (A13.8)--(A13.10) identify $q^*\mathcal W$, isometrically,
with the vanishing cohomology of the auxiliary Lefschetz pencil
on $\widetilde D_0$. In particular this is an identification of
local systems and pairings, not just a dimension comparison.
All maps are algebraic and all these identities persist in the
relative model at good primes.

\emph{Rank and characteristic-zero monodromy.}
Since $H^6/6!=1$, one has $h^6=720t^6$.
From $c(T_B)=(1+h)^{-2}|_B$, the topological Euler number is
$5h^6=3600t^6$. Weak Lefschetz and Poincare duality give
\[
 b_4(B)=3600t^6+330,\qquad
 \operatorname{rank}\mathcal V=3600t^6-165,\qquad
 d:=\operatorname{rank}\mathcal W=3600t^6-166.
 \tag{A13.11}
\]
In particular $d>0$ is even. To exclude finite monodromy in
characteristic zero, the Koszul resolution gives
$\chi(\mathcal O_B)=(2^6-2)t^6=62t^6$.
The groups $H^i(B,\mathcal O_B)$ for $i<4$ have dimensions
$1,6,15,20$, so
\[
 h^{4,0}(B)=62t^6+10,\qquad
 h^{4,0}(\mathcal W)=62t^6-5>0.
 \tag{A13.12}
\]
Removing $w$ changes only type $(2,2)$, and the ambient
$H^{4,0}(A)$ has dimension $15$.

By SGA~7~II, XIX, Proposition~3.4, printed p.~334, finite
monodromy of an even-dimensional Lefschetz pencil over
$\mathbf C$ is equivalent to its vanishing cohomology being
purely of type $(2,2)$. Thus the monodromy here is infinite.
Deligne, \emph{Weil II}, Theorem~4.4.1 and Lemma~4.4.2s,
printed pp.~227--228, then give the full orthogonal geometric
monodromy group on $\mathcal W$.
Removing the additional finitely many parameter values used
above does not reduce this image, since the fundamental group
of a smaller nonempty open curve surjects onto that of the
larger open curve.

\emph{Preserving that monodromy at good primes.}
Spread this one fixed auxiliary pencil, its smooth proper
families, and its finitely many omitted values over
$\mathbf Z[1/M]$. Increase $M$ so that the total discriminant
and other omitted values form a finite etale divisor in the
relative $\mathbf P^1$, and so that all the stated smoothness
and Lefschetz conditions hold. This is possible because there
are only finitely many points and bad reduction conditions
in this fixed characteristic-zero model.

Fix $\ell$. The local system $\mathcal W$ has a stable
$\mathbf Z_\ell$ lattice: intersect it with the torsion-free
integral cohomology lattice. Enlarge the finite exceptional
set of primes to include $2,\ell$ and the prime divisors of
$|\operatorname{GL}_d(\mathbf F_\ell)|$.
For any remaining $p$, every finite congruence quotient of its
compact integral monodromy image has order prime to $p$,
since the kernel of reduction to $\operatorname{GL}_d(\mathbf
F_\ell)$ is pro-$\ell$. Its image is therefore pro-prime-to-$p$.

Apply the prime-to-$p$ specialization isomorphism for a smooth
proper relative curve minus a relative normal-crossing divisor.
In the present case the curve is $\mathbf P^1$ and the divisor
is finite etale. After strict henselization at $p$, SGA~1,
XIII, Corollary~2.9, construction~2.10 and the proof of
Corollary~2.12 identify the prime-to-$p$ geometric fundamental
groups in characteristic zero and $p$, compatibly with the
local system extending over the trait. Hence their monodromy
images on $\mathcal W$ are conjugate. In particular the
characteristic-$p$ geometric closure is still the full
orthogonal group. This explicit argument avoids inferring
infinite monodromy in characteristic $p$ from a Hodge number.
The source's comparison is on pp.~289--292 of the revised
edition (original pp.~390--394). SGA~7~II, XIX, \S3.6,
printed p.~335, uses the same relative-pencil comparison;
the prime exclusions above give the stated sufficient version
without any additional tame-specialization assumption.

There are infinitely many primes $p\equiv3\pmod4$ outside
this finite set. For such $p$, the character sum for
$y^2=x^3-x$ cancels under $x\mapsto-x$, so its trace is zero
and its Frobenius polynomial is $T^2+p$. Its Newton slopes
are both $1/2$, proving supersingularity of $E_p$ and $A_p$.

\emph{Closed-point selection and the whole algebraic kernel.}
The relative class $w$ has the constant positive rational
Gram matrix (A13.4). The middle Frobenius polynomial of each
smooth $B_u$ and that of the constant ambient $A$ are rational;
their quotient, followed by division by the factor $T-1$
of the normalized class $w$, is the rational Frobenius
polynomial of $\mathcal W$.
Apply C44-A12 and C44-A12a, using the even rank in (A13.11)
and the just-proved orthogonal monodromy. A Zariski-dense set
of closed points has no root-of-unity normalized Frobenius
eigenvalue on $\mathcal W$. Every algebraic class has finite
Frobenius orbit, so its $\mathcal W$ component vanishes.
The rational Gram projection from A12 gives the equality
with $\mathbf Qw$, not merely with $\mathbf Q_\ell w$.
Equations (A13.2) and (A13.4) follow.

Finally, the selected closed parameter lifts over an unramified
mixed-characteristic trait in the fixed open parameter model.
The subvariety $Y$ and every divisor in the polynomial
expression (A13.3) lift along with it. Thus all of $A_{B_u}$
lifts in this example. No nonliftable class or general
fixed-axis sign theorem is being claimed.
\end{proof}

\paragraph{Primary sources actually read.}
The original SGA~7~II scan, including XVII~2.5.2,
XVIII~2.2, XVIII~6.6--6.7, XIX~1.4 and XIX~3.4--3.6, is
\url{https://publications.ias.edu/sites/default/files/Number12.pdf}.
The original \emph{Weil II} proof of 4.4.1 and 4.4.2s is in
\url{https://publications.ias.edu/sites/default/files/Number40.pdf}.
SGA~1 XIII~2.9--2.12 and its specialization proof were read in
the revised primary edition
\url{https://arxiv.org/pdf/math/0206203}.
For the finite-field selection, C44-A12 uses Entin's Theorem~3,
whose whole primary proof was independently read in
\url{https://arxiv.org/pdf/1805.05454v3}, pp.~9--11.
This result supplies a new nonzero closed-axis family for cote~44,
Th.~5, scans~004 and~008, while retaining the general axis question.

% END work/scan_w/closed_rank_one_axis.tex

% BEGIN work/scan_w/closed_odd_rank_axes.tex
% C44_POS /root/scan_w; separate extension preserving the checked A13 file.
\section{An odd number of prescribed primitive axis classes}
\label{sec:c44-a13b}

\begin{theorem}[C44-A13b: closed axes with a prescribed odd-rank kernel]
Let $N\geq1$ be odd. Retain the elliptic curve $E$, the sixfold
$A=E^6$, the product polarization $H$, and the class $\alpha_0$
of C44-A13. There is a number field $K$ and a collection of
distinct points $a_1,\ldots,a_N\in E^4(K)$ such that the
pairwise disjoint abelian surfaces
\[
                     Y_i=E^2\times\{a_i\}\subset A_K
\]
have the following property. For all sufficiently large integers
$t$, a nonempty open $U\subset\mathbf P^1_K$ carries smooth axes
$B_u=D_0\cap D_u$, with $D_0,D_u\in|tH|$ and $Y_i\subset B_u$
for every $i$, each the axis of a genuine Lefschetz pencil.
There is a finite set $S$ of rational primes, depending on these
chosen data and a coefficient prime $\ell$, such that, for every
$p\equiv3\pmod4$ outside $S$ and every place $\mathfrak p$ of $K$
above $p$, the reduction of this family has a Zariski-dense set
of closed points $u$ for which
\[
 \begin{split}
 A_{B_u}
 &=\Alg^2(B_{\bar u})\cap\ker(\kappa_*:H^4(B_{\bar u})
                                      \longrightarrow H^8(A_{\bar u}))\\
 &=\bigoplus_{i=1}^N\mathbf Q w_{i,u},\qquad
 w_{i,u}=[Y_i]_{B_u}-t^{-2}\kappa^*\alpha_0.
 \end{split}
 \tag{A13b.1}
\]
The ambient abelian sixfold is supersingular. On this entire
algebraic variable kernel the actual geometric intersection
matrix is
\[
 G_N=6t^2 I_N-\frac1{4t^2}J_N,
 \tag{A13b.2}
\]
where $J_N$ is the matrix with every entry equal to one. In
particular the matrix is positive definite once $24t^4>N$.
Every class in (A13b.1) has a relative algebraic lift in the
chosen incidence model over a mixed-characteristic trait.
The assertions are unconditional; neither Tate nor 032 is used.
\end{theorem}

\begin{proof}
Choose $N$ distinct torsion points of $E^4(\overline{\mathbf Q})$
and a number field $K$ over which they are all rational. This
is only a convenient way to choose the centres; no torsion
property enters their intersection or monodromy calculations.
Translations act trivially on cohomology, so
\[
       [Y_i]_A=h_3h_4h_5h_6=H^2\alpha_0
                 \quad\hbox{for every }i.
 \tag{A13b.3}
\]
All centres have trivial rank-four normal bundle in $A$.

Let $Y=\coprod_iY_i$, $Z=\operatorname{Bl}_Y A$, and let $F$
be the sum of its exceptional divisors. For $t\gg0$,
$L'=\pi^*\mathcal O_A(tH)\otimes\mathcal O_Z(-F)$ is very
ample and the normal derivatives of its section space are
generated along each $Y_i$. The normal derivative of one
section is a row of four entries on a surface; its zero
locus has expected codimension four. For two sections,
rank at most one has expected codimension three. The finite
disjoint union of centres therefore introduces no zero or
rank-drop point for general sections. Bertini off $Y$ gives
smooth original divisors and smooth original intersections.
Choose a general $\widetilde D_0\in|L'|$ and a general
auxiliary hyperplane pencil on its induced projective embedding.
Exactly as in A13, this gives
\[
       \widetilde B_u=\operatorname{Bl}_Y B_u
                       \subset\widetilde D_0.
 \tag{A13b.4}
\]
There are finitely many parameters to remove to make $B_u$
smooth. For the fixed general $\widetilde D_0$, the Lefschetz
lines through it form a nonempty open set of lines in $|L'|$.
The auxiliary pencil can be chosen so that its general
$\widetilde D_u$ determines a line in this open; remove the
finitely many exceptional parameters. The two normal rows
are independent at every point of every centre, so all
original pencil members are smooth along $Y$. Consequently
their critical points are off the exceptional locus and the
original pencil generated by $D_0,D_u$ is Lefschetz as well.
These are nonempty open conditions over $K$, so the choices
may be made over $K$. This construction uses the induced
embedding of $\widetilde D_0$, and does not require surjectivity
onto its complete section space.

By (A13b.3), every $w_i$ belongs to $\ker\kappa_*$. For each
centre the normal bundle sequence is
\[
 0\longrightarrow N_{Y_i/B}\longrightarrow\mathcal O_{Y_i}^4
       \longrightarrow\mathcal O_{Y_i}(tH)^2\longrightarrow0.
\]
Hence $[Y_i]_B^2=6t^2$. Distinct $Y_i$ are disjoint, so
$[Y_i]_B[Y_j]_B=0$ for $i\ne j$. The restriction of the
divisor polynomial $\alpha_0$ to any translate $Y_i$ has
degree $1/4$, and $H^2\alpha_0=[Y_i]_A$. Orthogonal ambient
projection thus gives
\[
           \int_Bw_iw_j=6t^2\delta_{ij}-\frac1{4t^2}.
\]
The eigenvalues of this matrix are $6t^2$ on the sum-zero
subspace and $6t^2-N/(4t^2)$ on the all-ones line. Taking
$t$ still larger so that $24t^4>N$ proves (A13b.2) and the
linear independence of these $N$ classes in every good fibre.

The vanishing-system identification in A13 extends as an
actual orthogonal identification, as follows. The blowup
formula and weak Lefschetz give
\[
 H^4(\widetilde D_0)\simeq H^4(Z)
   =H^4(A)\oplus\bigoplus_{i=1}^N
                  \bigl(H^2(Y_i)(-1)\oplus\mathbf Q_\ell(-2)\bigr).
 \tag{A13b.5}
\]
For $q:\widetilde B\to B$, let $j_i,p_i,\xi_i$ denote the
maps and tautological class on its $i$th exceptional divisor.
Then
\[
 H^4(\widetilde B)=q^*H^4(B)\perp
              \bigoplus_{i=1}^N j_{i*}p_i^*H^2(Y_i)(-1).
 \tag{A13b.6}
\]
Each exceptional summand has pairing minus the pairing on
$H^2(Y_i)$; different exceptional summands are orthogonal
because their supports are disjoint. The excess formula is
$q^*[Y_i]_B=j_{i*}(\xi_i+p_i^*c_1N_{Y_i/B})$. Thus the image
of (A13b.5) in (A13b.6) is exactly the exceptional sum plus
\[
             q^*\bigl(\kappa^*H^4(A)+\sum_i\mathbf Q_\ell[Y_i]_B\bigr).
\]
The ambient pairing is nondegenerate by hard Lefschetz, and
the pairing on the additional primitive span is $G_N$.
Writing $\mathcal V=\ker\kappa_*$, the vanishing cohomology
of the auxiliary pencil is therefore isometric to
\[
 q^*\mathcal W_N,
 \qquad
 \mathcal W_N=\mathcal V\cap\bigcap_{i=1}^Nw_i^\perp.
 \tag{A13b.7}
\]
All maps are relative algebraic maps. This identifies local
systems and their forms, not just their dimensions.

The topology and Koszul calculation for a smooth intersection
of two members of $|tH|$ are unchanged by the prescribed
centres. Since only type $(2,2)$ classes have been removed,
\[
 d_N:=\operatorname{rank}\mathcal W_N=3600t^6-165-N,
 \qquad h^{4,0}(\mathcal W_N)=62t^6-5.
 \tag{A13b.8}
\]
For our sufficiently large $t$, the rank is positive and even,
because $N$ is odd, and the displayed Hodge number is positive.
SGA~7~II, XIX, Proposition~3.4 excludes finite characteristic-zero
vanishing monodromy; Deligne, \emph{Weil II}, 4.4.1 and 4.4.2s
then give full orthogonal geometric monodromy on $\mathcal W_N$.
These are precisely the primary monodromy results used in A13,
applied to the odd-dimensional $\widetilde D_0$ in (A13b.4).

Spread this fixed data over $\mathcal O_K[1/M]$. Enlarge the
finite set of excluded rational primes to include those below
any bad reduction place of the elliptic curve, centres, blowups,
chosen sections or smooth families, any collision of the centre
sections, and any failure of the omitted parameter divisor to
be finite etale. Include also $2$, the fixed coefficient prime
$\ell$, and all prime divisors of $|\operatorname{GL}_{d_N}
(\mathbf F_\ell)|$. There are finitely many such primes.
A stable integral lattice in $\mathcal W_N$ has compact image
whose finite congruence quotients are all prime to every
remaining residue characteristic $p$. The prime-to-$p$
specialization isomorphism of SGA~1, XIII, 2.9--2.12, on
the relative $\mathbf P^1$ minus its finite etale boundary
identifies the geometric images in characteristic zero and
at every remaining place $\mathfrak p$. This is the same
integral-image argument as in A13 and gives full orthogonal
monodromy at every such place. It does not infer
characteristic-$p$ monodromy from the Hodge number alone.

Infinitely many rational primes $p\equiv3\pmod4$ remain.
At every place of $K$ above one of them, the reduction of $E$
is supersingular: its model over $\mathbf F_p$ has Frobenius
polynomial $T^2+p$, and finite residue-field extension preserves
supersingularity. Each $w_i$ is a rational relative algebraic
class fixed by normalized Frobenius. The normalized Frobenius
polynomial on $\mathcal W_N$ is rational, obtained by dividing
the polynomial of the full middle cohomology by that of the
ambient cohomology and by $(T-1)^N$. C44-A12 and A12a now
apply with the nonsingular rational Gram matrix $G_N$ and
the even rank $d_N$. They give a Zariski-dense set of closed
points with no root-of-unity Frobenius eigenvalue on
$\mathcal W_N$. All algebraic classes have finite Frobenius
orbit, so the whole algebraic kernel equals the known span.
The coefficients in this span are rational because they are
$G_N^{-1}$ times the vector of rational intersection numbers.
This proves the direct-sum equality in (A13b.1).

Finally, a selected closed parameter lifts after an unramified
extension of the completed local trait of the parameter model.
All $Y_i$, the divisors $h_j$, the chosen incidence variety,
and the polynomial expression for $w_i$ lift over this trait.
Thus every rational cycle class in (A13b.1) has an algebraic
lift. The assertion concerns these classes and does not claim
that every individual representing cycle deforms. In
particular these examples still do not produce nonliftable
primitive classes.
\end{proof}

\paragraph{Parity and scope.}
The oddness of $N$ is exactly what makes the complementary
orthogonal rank even. For even $N$, an odd-dimensional
orthogonal transformation always has an eigenvalue $1$ or
$-1$, so the root-of-unity-free open used in A12 is empty.
This construction therefore gives no analogous exact-span
conclusion for even $N$. The result is an existence theorem
for dense sets of closed points in explicit incidence models;
it does not specify one numerical parameter, does not settle
the arbitrary-axis sign, and does not assert that the full
variable motive is of abelian type.

\paragraph{Source locators.}
The proofs of the three monodromy inputs and the blowup maps
are those read and located in C44-A13: SGA~7~II, XVII~2.5.2
(p.~217), XVIII~2.2 (pp.~260--261), XIX~3.4 (p.~334);
\emph{Weil II}, 4.4.1--4.4.2s (pp.~227--229); SGA~1,
XIII~2.9--2.12 (revised pp.~289--292). Closed-point selection
is the separately checked C44-A12/A12a, using Entin, Theorem~3,
Section~4 (author arXiv:1805.05454v3, pp.~9--11).

% END work/scan_w/closed_odd_rank_axes.tex

% BEGIN work/axis/nonliftable_surface_centres.tex
% C44_POS /root/axis_independent; GPT-6 Astra Pro; effort not exposed.
\section{Smooth nonliftable centres and the primitive cancellation}
\subsection*{C44-A14. A smooth surface with an intrinsic lifting obstruction}

\textbf{Theorem.}
Let $k=\overline{\mathbf F}_3$ and $E/k$ be the supersingular elliptic
curve $y^2=x^3-x$. There is a smooth connected surface
$i:S\hookrightarrow A=E^4$, defined over a finite subfield of $k$,
with the following properties:
\[
 h^1(S,\mathcal O_S)=4,\qquad
 i^*:H^1(A,\mathbf Q_5)\xrightarrow{\sim}H^1(S,\mathbf Q_5).
                                                               \tag{A14.1}
\]
The intrinsic top cup form
\[
 f_S:\bigwedge^4 H^1(S,\mathbf Q_5)\longrightarrow
                 H^4(S,\mathbf Q_5)\simeq\mathbf Q_5(-2)
                                                               \tag{A14.2}
\]
has zero infinitesimal stabilizer in
$\operatorname{End}(H^1(S,\mathbf Q_5))$. Consequently $S$ admits
no abstract smooth projective lift to characteristic zero, even
after extension of the residue field.

Embed $S'=S\times\{0\}^2$ in $X=E^6$. For every sufficiently high
power of an ample line bundle on $X$, there are genuine Lefschetz
axes $B\subset X$ containing $S'$. For each such axis, the
uncorrected algebraic class
\[
                   \gamma=[S']_B\in H^4(B,\mathbf Q_5(2))
                                                               \tag{A14.3}
\]
cannot lift to an algebraic class under any abstract smooth
projective characteristic-zero lift of $B$. The primitive class
\[
 w=\gamma-\kappa^*(L_h^2)^{-1}[S']_X\in A_B,
 \qquad h=c_1(L),\quad B\in|L|\cap|L|,                         \tag{A14.4}
\]
is nonzero for all sufficiently high powers, by C44-A06.
The argument below does \emph{not} prove that $w$ is nonliftable.
Its known obstruction vanishes exactly under this projection.

\paragraph{1. The tensor input, including the lifting quantifier.}
Use the explicit integral divisor-product cycle
\[
 z=\sum_{a<b<c<d}c_{abcd}F_{ab}F_{cd}\in\operatorname{CH}^2(A)
\]
of C44-A10. The integers $c_{abcd}$ and the integral divisors
$F_{ab}$ are specified there. Its modulo-five cohomology tensor
is $\lambda^2u_0$ in an eight-dimensional exterior algebra, where
$\lambda\in\mathbf F_5^\times$ and the derivative map
\[
 \mathcal D_{u_0}:\operatorname{End}(\mathbf F_5^8)
                     \longrightarrow\bigwedge^4\mathbf F_5^8,
 \qquad T\longmapsto T\cdot u_0
                                                               \tag{A14.5}
\]
has a $64$-by-$64$ minor with determinant $3$ modulo five.
This is the explicit certificate in
\texttt{supersingular\_class\_certificate.json}; its integer
determinant and two independent reconstructions are recorded in
C44-A10. We use the actual residue class of the integral cycle,
not a claim about unspecified higher $5$-adic coordinates.

Every integral degree-four class whose reduction equals
$\operatorname{cl}_5(z)$ therefore has zero infinitesimal
stabilizer over $\mathbf Q_5$. Indeed, lift the selected
modulo-five basis to a $\mathbf Z_5$ basis. Milne,
\emph{Abelian Varieties}, I Theorem 12.1, printed pp.~55--56,
identifies integral cohomology with the exterior algebra and
gives torsion freeness. The same minor of the actual derivative
matrix is a $5$-adic unit. Changing its higher $5$-adic entries
cannot destroy this conclusion.

\paragraph{2. A split ample bundle with the required second Chern class.}
For divisor classes $x,y$ modulo five,
\[
                  xy=3\big((x+y)^2-x^2-y^2\big).
                                                               \tag{A14.6}
\]
The sum of a line bundle $M$ and its inverse has
$c_1(M\oplus M^{-1})=0$ and $c_2(M\oplus M^{-1})=-c_1(M)^2$.
Thus a direct sum of such pairs realizes
$\operatorname{cl}_5(z)$ as its second Chern class. For complete
explicitness, for each summand $cF_{ab}F_{cd}$ put
$x=c_1(\mathcal O(F_{ab}))$ and $y=c_1(\mathcal O(F_{cd}))$.
Take the following numbers of inverse pairs, using the least
representatives in $\{0,1,2,3,4\}$:
\[
 \begin{array}{c|c}
 \text{line bundle in the inverse pair}&\text{number of pairs}\\ \hline
 \mathcal O(F_{ab}+F_{cd})&-3c\bmod5\\
 \mathcal O(F_{ab})&3c\bmod5\\
 \mathcal O(F_{cd})&3c\bmod5 .
 \end{array}
                                                               \tag{A14.7}
\]
Each pair has zero first Chern class, so different pairs have no
cross term in the second Chern class. Equations (A14.6)--(A14.7)
give the asserted residue class.

Fix an ample line bundle $H$ on $A$. Replace each line summand
$M_j$ by
\[
                        L_j=M_j\otimes H^{\otimes5n_j}
\]
for a sufficiently large positive integer $n_j$. Every $L_j$
can then be chosen very ample. These changes preserve all Chern
classes modulo five. Write
\[
 \mathcal E=\bigoplus_{j=1}^R L_j,\qquad
 c_2(\mathcal E)\equiv\operatorname{cl}_5(z)\pmod5 .             \tag{A14.8}
\]
Here $R\ge2$ is a finite even integer. Each summand is very ample,
so its sections prescribe arbitrary first jets at any one point
of the smooth variety $A$.

\paragraph{3. A smooth determinantal surface, without an unproved
transversality assumption.}
Choose a general homomorphism
\[
                         \psi:\mathcal O_A^{R-1}\longrightarrow
                                      \mathcal E .
                                                               \tag{A14.9}
\]
All its entries vary independently in the corresponding spaces
of global sections. At a fixed point of $A$, evaluation and
first-jet evaluation are surjective. The space of $R$-by-$(R-1)$
matrices of rank at most $R-3$ has codimension
$3\cdot2=6>4$. The incidence dimension estimate consequently
excludes this rank locus for a general $\psi$.

The matrix stratum of rank exactly $R-2$ is smooth of codimension
two. At a point of that stratum the normal space is
$\operatorname{Hom}(\ker\psi,\operatorname{coker}\psi)$, of
dimension two. The map from the four-dimensional tangent space
of $A$ to this normal space can be prescribed arbitrarily by
the first jets of the entries of $\psi$. Failure of its
surjectivity is the rank-at-most-one locus for a $2$-by-$4$
matrix, which has codimension three. Including the two
conditions on the value gives a bad incidence of codimension
five. Varying the point of the fourfold still leaves a proper
bad locus in the parameter space. Together with the preceding
rank estimate, this proves that a general $\psi$ has a smooth
rank-drop locus
\[
                    S=\{x:\operatorname{rank}\psi_x=R-2\}
                                                               \tag{A14.10}
\]
of pure dimension two, if nonempty. These dimension counts use
first jets and hold in characteristic three without generic
smoothness for a possibly inseparable map.

The locus is nonempty. Otherwise $\psi$ is everywhere of rank
$R-1$ and its cokernel is a line bundle, whence
$c_2(\mathcal E)=0$. This contradicts (A14.8), since the certified
modulo-five tensor is nonzero.

For clarity, its cycle class can be computed directly. On
$A\times\mathbf P^{R-2}$ the incidence
\[
                    \{(x,[v]):\psi_x(v)=0\}
                                                               \tag{A14.11}
\]
is the zero locus of a section of
$\operatorname{pr}_A^*\mathcal E\otimes
 \operatorname{pr}_{\mathbf P}^*\mathcal O(1)$.
There is no deeper rank stratum, so (A14.11) is isomorphic to
$S$: the kernel line is unique and varies algebraically.
It is a regular zero locus of codimension $R$. If
$\xi=c_1(\mathcal O(1))$, pushforward of its top Chern class
takes the coefficient of $\xi^{R-2}$ in
$\prod_{j=1}^R(c_1(L_j)+\xi)$. It follows that
\[
                    [S]_A=c_2(\mathcal E).
                                                               \tag{A14.12}
\]
This proves the class formula directly, without relying on a
generality assertion in a degeneracy-locus formula.

\paragraph{4. Connectedness and coherent irregularity.}
The maximal-minor complex gives an exact sequence
\[
 0\longrightarrow
       \big((\det\mathcal E)^{-1}\big)^{\oplus(R-1)}
   \longrightarrow \mathcal E\otimes(\det\mathcal E)^{-1}
   \longrightarrow I_S\longrightarrow0.                      \tag{A14.13}
\]
One can verify exactness locally here. Where an $(R-2)$-minor
is invertible, row and column operations reduce $\psi$ to an
identity block and a single two-entry column $(a,b)$.
The rank-drop ideal is the regular sequence $(a,b)$, and
(A14.13) reduces to its Koszul resolution plus split acyclic
summands. Away from $S$ a maximal minor is invertible and the
complex is split exact. This covers $A$.

Both left terms are direct sums of inverses of ample line
bundles: the summands of
$\mathcal E\otimes(\det\mathcal E)^{-1}$ are
$(\bigotimes_{k\ne j}L_k)^{-1}$. On an abelian fourfold an
anti-ample line bundle has zero cohomology in every degree
strictly less than four. This follows from Serre duality,
$\omega_A\simeq\mathcal O_A$, and ample vanishing; see Milne,
loc.\ cit., I Theorem 11.1(c), printed p.~54.
The long exact sequence of (A14.13) gives
\[
          H^0(A,I_S)=H^1(A,I_S)=H^2(A,I_S)=0.
                                                               \tag{A14.14}
\]
The sequence $0\to I_S\to\mathcal O_A\to\mathcal O_S\to0$
therefore gives
\[
             H^0(S,\mathcal O_S)=k,\qquad
             H^1(S,\mathcal O_S)\simeq H^1(A,\mathcal O_A).
                                                               \tag{A14.15}
\]
Thus $S$ is connected and $h^1(\mathcal O_S)=4$.
As a smooth connected variety over an algebraically closed
field, it is irreducible.

\paragraph{5. The intrinsic cup form and the necessary
\'etale comparison.}
Set $V=H^1(A,\mathbf Q_5)$, of dimension eight, and
$u=[S]_A\in\bigwedge^4V$, suppressing Tate twists after choosing
a nonzero generator. Equations (A14.8) and (A14.12), together
with the unit-minor certificate, imply
\[
                         \ker\mathcal D_u=0.
                                                               \tag{A14.16}
\]
Let $\operatorname{vol}:\bigwedge^8V\to\mathbf Q_5$ be the
trace pairing, and define
$f_u(\alpha)=\operatorname{vol}(u\wedge\alpha)$ for
$\alpha\in\bigwedge^4V$.
This covariant four-form also has zero infinitesimal stabilizer.
Indeed the derivative of the contragredient action is
\[
                  T\cdot f_u=f_{\,T\cdot u-(\operatorname{tr}T)u}.
                                                               \tag{A14.17}
\]
If it is zero, nondegeneracy of the exterior trace pairing gives
$T\cdot u=(\operatorname{tr}T)u$. Scalar multiplication by the
identity acts on $\bigwedge^4V$ by four. Hence
$T-(\operatorname{tr}T)I/4$ annihilates $u$ and is zero by
(A14.16). Taking traces in dimension eight gives
$\operatorname{tr}T=2\operatorname{tr}T$, and therefore $T=0$.

In particular the contraction radical of $f_u$ is zero.
Otherwise a nonzero endomorphism whose image is a nonzero line
in this radical would annihilate $f_u$.
Projection formula identifies $f_u$ with
\[
     (v_1,\ldots,v_4)\longmapsto
                 \int_S i^*v_1\,i^*v_2\,i^*v_3\,i^*v_4 .
                                                               \tag{A14.18}
\]
Every element of $\ker i^*:H^1(A,\mathbf Q_5)\to H^1(S,\mathbf Q_5)$
would thus belong to this radical. Consequently $i^*$ is injective.

Coherent irregularity alone would not justify an \'{e}tale
isomorphism. The additional required bound is
\[
 \dim H^1(S,\mathbf Q_5)
       =2\dim\operatorname{Pic}^0_{S/k,\mathrm{red}}
       \le2\dim T_0\operatorname{Pic}_{S/k}
       =2h^1(S,\mathcal O_S)=8.                              \tag{A14.19}
\]
For the equality, the Kummer sequence identifies the rational
$5$-adic Tate module of the Picard variety with
$H^1(S,\mathbf Q_5(1))$; the N\'eron--Severi torsion disappears
after tensoring with $\mathbf Q_5$. This is explicitly stated
in SGA~1, Expos\'e XI, Corollary 6.6 and the paragraph following
it, revised p.~235, original pp.~305--306.
The tangent-space identity is recalled in XI, Remarks 6.11,
revised p.~237, original p.~308. Dimension of a group scheme
at its identity is at most the dimension of its tangent space.
The injection and (A14.19) prove (A14.1), and (A14.18) then
proves that the intrinsic form (A14.2) has zero infinitesimal
stabilizer.

\paragraph{6. Abstract nonliftability of the surface.}
Suppose $S$ were the special fibre of a smooth projective
mixed-characteristic lift. Smooth proper base change preserves
the rational cohomology and cup tensor. Spread the generic fibre
to a finitely generated characteristic-zero field and embed that
field into $\mathbf C$. Betti--\'etale comparison identifies
the complex fibre's rational top cup form after scalar extension
with (A14.2). Its infinitesimal stabilizer therefore remains
zero. This is a rank condition on a rational derivative matrix,
so it can be checked over $\mathbf Q_5$ or over $\mathbf R$.

On the real $H^1$ of a smooth projective complex surface, the
Hodge circle has a nonzero infinitesimal generator $J$ with
$J^2=-I$. The top cohomology $H^4$ has type $(2,2)$, so the
circle acts trivially on it. Compatibility of cup product with
Hodge structures implies that $J$ annihilates the top fourfold
cup form. This contradicts its zero infinitesimal stabilizer.
No embedded lift, preservation of the selected abelian
fourfold, or Hodge conjecture is assumed.

All choices involved finitely many equations over $k$.
The nonempty parameter opens used above have $k$-points; each
chosen point and all defining data lie in a finite extension of
$\mathbf F_3$. The surface therefore descends to a finite
field. This proves the first part of the theorem.

\paragraph{7. The class on an axis and the exact cancellation.}
Embed $S'=S\times\{0\}^2$ in $X=E^6$. The imposed-centre
construction in C44-A06 gives, for every sufficiently high
power of an ample polarization, a genuine Lefschetz pencil
whose smooth axis $B$ contains $S'$. Weak Lefschetz identifies
$H^1(B,\mathbf Q_5)$ with
\[
       H^1(X,\mathbf Q_5)=V\oplus U,\qquad\dim V=8,\quad\dim U=4.
\]
For the class $\gamma=[S']_B$ consider the intrinsic form
\[
 f_\gamma(v_1,\ldots,v_4)=\int_B\gamma v_1v_2v_3v_4.
                                                               \tag{A14.20}
\]
It is the pullback of $f_S$ from $V$, because restriction to
$S'$ kills the last two elliptic factors. Its radical is
exactly $U$, and its induced form on the quotient by the radical
has zero infinitesimal stabilizer.

If $\gamma$ lifted to an algebraic class in an arbitrary
abstract smooth projective lift of $B$, the generic rational
Betti class would have type $(2,2)$. Cup product and top trace
would make the transported form (A14.20) Hodge-circle invariant.
Its radical is intrinsic and hence stable under that circle.
The quotient of $H^1$ by the radical inherits a weight-one
Hodge structure. Its real infinitesimal generator still
satisfies $J^2=-I$ and is nonzero, while annihilating the induced
four-form. This contradicts the certified zero stabilizer.
Thus the \emph{cohomology class} $\gamma$ is nonliftable;
the claim is stronger than failure to lift its chosen support.

Hard Lefschetz on $X$ and algebraicity of its inverse give
(A14.4), a rational algebraic element of the exact kernel
$A_B=\ker\kappa_*$. C44-A06 makes it nonzero and positive for
all sufficiently high powers. However, for every
$w\in\ker\kappa_*$ and every four $v_i\in H^1(X)$, projection
formula gives
\[
       \int_B w\,\kappa^*v_1\kappa^*v_2\kappa^*v_3\kappa^*v_4
          =\int_X\kappa_*w\,v_1v_2v_3v_4=0.                  \tag{A14.21}
\]
Every $H^1(B)$ class comes from $X$. Consequently the intrinsic
four-form of the projected class $w$ is identically zero.
The entire obstruction (A14.20) is carried by the subtracted
ambient class.

The same cancellation occurs for differences of translates.
Take two distinct points of the last $E^2$ and the disjoint
surfaces $S'_1,S'_2$. Their ambient cycle classes agree because
translation acts trivially on cohomology. A sufficiently ample
axis contains both, and
$[S'_1]_B-[S'_2]_B\in A_B$ is nonzero by the positive Gram
matrix of C44-A06. Both individual supports and both
uncorrected classes are nonliftable, but their four-forms agree
and cancel. Nonliftability of the supports does not exclude a
different characteristic-zero cycle representing the projected
or difference cohomology class. No such exclusion is established
here. \hfill$\square$

\paragraph{Sources and scope.}
The explicit tensor and arithmetic certificate are C44-A10,
with files \texttt{axis\_supersingular\_nonliftable.tex} and
\texttt{supersingular\_class\_certificate.json}.
Primary sources used above are:
\[
 \begin{gathered}
 \texttt{https://jmilne.org/math/CourseNotes/AV.pdf},\\
 \texttt{https://arxiv.org/pdf/math/0206203},\\
 \texttt{https://www.jmilne.org/math/Documents/Deligne82.pdf}.
 \end{gathered}
\]
The Milne locators are I Theorems 11.1(c) and 12.1,
printed pp.~54--56. The SGA~1 locators are XI 6.6, revised
p.~235/original pp.~305--306, and XI 6.11, revised p.~237/original
p.~308. The Hodge-circle description and the necessary Hodge
type of algebraic classes are Deligne's revised notes,
Introduction p.~2 and \S3 p.~23. Smooth proper base change is
SGA~4 XVI Theorem 2.2. The rank-stratum, Koszul and incidence
top-Chern arguments have been given in full above.

This construction certifies a smooth abstractly nonliftable
centre and an individually nonliftable class on a liftable
axis. It does not produce a certified nonliftable class in
$A_B$, and it does not settle the universal QH9-axis sign.

% END work/axis/nonliftable_surface_centres.tex

% BEGIN work/scan_w/axis_one_tate_direction.tex
% C44_POS /root/scan_w. Root's determinant idea; complete proof and all-ell repair.
\section{One remaining potential Tate direction determines the axis sign}
\label{sec:c44-a15}

\begin{theorem}[C44-A15: a positive span and one potential Tate direction]
Let $r\geq1$ and let $X/\mathbf F_q$ be an abelian variety of dimension $2r+2$,
let $L$ be an ample line bundle with $h=c_1(L)$, and let
$\kappa:B\hookrightarrow X$ be a smooth complete intersection
of two members of $|L|$. In particular this applies to
the smooth axes in QH9-axis. For each prime $\ell\ne p$, set
\[
 V_\ell=\ker\bigl(\kappa_*:H^{2r}(B_{\overline{\mathbf F}_q},
             \mathbf Q_\ell(r))\longrightarrow
        H^{2r+4}(X_{\overline{\mathbf F}_q},\mathbf Q_\ell(r+2))\bigr),
 \qquad q_B(x,y)=(-1)^r\int_Bxy.
 \tag{A15.1}
\]
Fix an initial coefficient prime $\ell_0\ne p$ for the definition
of $\Alg^r(B)$ as the image of rational cycles, and put
$A_B=\Alg^r(B)\cap V_{\ell_0}$.
Suppose a finite-dimensional rational subspace $K\subset A_B$
has a positive definite Gram matrix $G$ for $q_B$; the case
$K=0$, with $\det G=1$, is allowed.

There is a common algebraic projection of cycle representatives
onto $V_\ell$ for every $\ell\ne p$. Use it on a rational basis
of $K$, and extend the finite ground field so that the resulting
cycles are all defined over it. Their span $K_\ell\subset V_\ell$
has the same Gram matrix $G$ in every realization. Write
\[
       W_\ell=K_\ell^\perp\cap V_\ell,
 \qquad
       \Phi_W(T)=\det(T-F\mid W_\ell)\in\mathbf Q[T],
 \tag{A15.2}
\]
where $F$ is normalized Frobenius on the Tate-twisted cohomology.
The polynomial is independent of $\ell$. Let $\tau$ be the total
\emph{algebraic multiplicity} of its roots that are roots of
unity. If
\[
                              \tau\leq1,
 \tag{A15.3}
\]
then $q_B$ is positive definite on all of $A_B$.
Moreover $A_B=K$ or $A_B=K\perp\mathbf Qz$ for a class $z$ of
strictly positive square. When $\tau=0$, only $A_B=K$ occurs.
The given cycles spanning $K$ need not lift to characteristic zero.
No Tate conjecture, Frobenius semisimplicity, standard conjecture
$D(B)$, or unrefereed 032 result is assumed.
\end{theorem}

\begin{proof}
All determinants below are ordinary Gram determinants modulo
squares, with no factor $(-1)^{d(d-1)/2}$.

\emph{One algebraic projector for all coefficient primes.}
The self-intersection formula and hard Lefschetz give
\[
 \kappa_*\kappa^*=L_h^2:
       H^{2r}(X)\xrightarrow{\sim}H^{2r+4}(X)(2).
 \tag{A15.4}
\]
There is a rational algebraic correspondence $\Gamma$ on $X$
which realizes $(L_h^2)^{-1}$ simultaneously for every $\ell\ne p$.
Indeed replace $h$ by the symmetric rational divisor
$(h+[-1]^*h)/2$, which has the same cohomology class, and apply
K\"unnemann's hard Lefschetz isomorphism for Chow motives with
$g=2r+2$ and $i=2r$. Its inverse, composed with the indicated
Chow--K\"unneth projectors, is one correspondence $\Gamma$.
The source is \emph{Inventiones Mathematicae} 113 (1993),
Theorem~5.2, printed p.~96, with the primitive decomposition
in Theorem~5.1 on the same page. Only the absolute theorem over
the finite ground field is needed here. No application of its
relative formulation to a mixed-characteristic base is being made.

Consequently the correspondence
\[
                 P=\Delta_B-\kappa^*\Gamma\kappa_*
 \tag{A15.5}
\]
acts on middle cohomology as the orthogonal projector onto
$V_\ell$ for every $\ell\ne p$. These statements concern its
cohomological action; no idempotence claim on all Chow groups
is needed. The ambient summand is nondegenerate, since its
pairing is $(-1)^r\int_X ab h^2$, nondegenerate by hard
Lefschetz. Thus $V_\ell$ is itself nondegenerate.

Let $C_1,\ldots,C_k$ be cycle representatives for a rational
basis of $K$. Replace them by $P_*C_1,\ldots,P_*C_k$ and retain
the names. The initial $\ell_0$-classes are unchanged. Their
intersection numbers, hence their rational Gram matrix $G$,
are unchanged as well. Now their classes lie in $V_\ell$ in
\emph{every} realization, and $\det G\ne0$ makes them independent
in every realization. This explicit replacement is important:
vanishing of $\kappa_*C_i$ in one realization alone would not
justify vanishing in all realizations without another argument.
After a finite ground-field extension all these actual rational
cycles and correspondences are defined over that field, so
normalized Frobenius fixes $K_\ell$ pointwise.

The characteristic polynomial of $V_\ell$ is the quotient of
the normalized middle characteristic polynomial of $B$ by that
of $H^{2r}(X)$; the latter is the ambient summand. These
polynomials are rational, independent of $\ell$, and have all
complex roots on the unit circle, by Deligne, \emph{Weil I},
Theorem~1.6, printed p.~276. Dividing further by $(T-1)^k$
gives $\Phi_W$. Passing to a finite field extension raises
eigenvalues to a common positive power and does not change
their total root-of-unity multiplicity.

\emph{A positive rational determinant for the full variable space.}
Use the embedded polarized characteristic-zero lift of $B\subset X$
from C44-A05. Choose a complex embedding of a field of definition
of its generic fibre. Let $V_{\rm B}$ be the rational Betti
kernel of its Gysin map and let $q_{\rm B}$ be the signed
intersection form in (A15.1). Smooth proper base change and
Betti--etale comparison identify
\[
           (V_\ell,q_B)\simeq
              (V_{\rm B},q_{\rm B})\otimes_{\mathbf Q}\mathbf Q_\ell
                \quad\hbox{for every }\ell\ne p.
 \tag{A15.6}
\]
One lift is used for all these comparisons. The kernel and its
pairing are geometric, so this step does not require the cycles
$C_i$ or the special-fibre correspondence $\Gamma$ to lift.

The Betti kernel is primitive for $h|_B$. Indeed
$\kappa_*(h|_B\,v)=h\kappa_*v=0$, and Gysin on
$H^{2r+2}(B)$ is injective, being dual to the weak Lefschetz
isomorphism on $H^{2r-2}$. Hence $h|_B\,v=0$.
The Hodge--Riemann relations make $q_{\rm B}$ positive on
type $(r,r)$ and give opposite signs on successive pairs of
off-diagonal Hodge types. Every negative real block comes
from a conjugate pair of distinct Hodge types and has even
dimension. The negative index is therefore even. It follows
that, in any rational Betti basis,
\[
                       D:=\det(q_{\rm B})>0.
 \tag{A15.7}
\]
In particular $D$ represents the determinant squareclass of
$q_B|_{V_\ell}$ simultaneously in every $\mathbf Q_\ell$.
This only uses the sign parity on the full variable space;
it does not declare $q_{\rm B}$ positive definite there.

\emph{Removing the known span by rational intersection numbers.}
For any cycle representing $x\in A_B$, first replace it by
its image under $P$, then subtract its orthogonal projection
onto the $C_i$. The coefficients of that projection are
\[
     G^{-1}\bigl(q_B(P_*x,C_1),\ldots,q_B(P_*x,C_k)\bigr)^{\rm t},
 \tag{A15.8}
\]
where $x$ denotes the chosen cycle in this formula. They are
rational intersection numbers. The resulting rational cycle
$Z$ has its class in $W_\ell$ for every $\ell\ne p$ and
represents exactly the residual component of $x$ at $\ell_0$.
Thus the same representative can be used in every realization.

Every algebraic cycle over $\overline{\mathbf F}_q$ is defined
over a finite extension. Its normalized Frobenius orbit is
therefore finite. Its class in $W_\ell$ is annihilated by
$F^e-1$ for some $e>0$, so it belongs to the sum of the
root-of-unity eigenspaces. If $\tau=0$, every residual class
vanishes, giving $A_B=K$. Assume henceforth that $\tau=1$ and
that some residual class $z=\operatorname{cl}_{\ell_0}(Z)$ is
nonzero; otherwise the same conclusion holds.

\emph{The simple potential Tate line and its complement.}
Because $\Phi_W$ is rational and has just one root of unity,
counted with multiplicity, that root is $1$ or $-1$ and is
simple. Replace $F$ by $F^2$ and call the resulting isometry
$S$. In every realization its eigenvalue $1$ is simple and
it has no other root-of-unity eigenvalue. Its one-dimensional
$1$-eigenspace $T_\ell$ is nondegenerate: the other generalized
eigenspaces are orthogonal to it, so degeneracy on this line
would contradict nondegeneracy of $W_\ell$. We have
\[
              W_\ell=T_\ell\perp U_\ell,
 \tag{A15.9}
\]
with $S|_{U_\ell}$ having no eigenvalues $1$ or $-1$.
In particular $\dim U_\ell$ is even, since eigenvalues of an
orthogonal transformation occur in reciprocal pairs away
from $1$ and $-1$.

At $\ell_0$, the class $z$ spans $T_{\ell_0}$. Therefore
\[
                           a:=q_B(Z,Z)\in\mathbf Q^*
 \tag{A15.10}
\]
is nonzero. Since it is an actual rational intersection number,
the class of this same cycle $Z$ is nonzero in every realization.
By (A15.8) it belongs to $W_\ell$ in each realization, and by
its finite Frobenius orbit it spans $T_\ell$. Thus $a$ is the
Gram determinant of $T_\ell$ in the basis supplied by $Z$
for every $\ell\ne p$. This proves the required simultaneous
nonvanishing and kernel membership without $D(B)$.

\emph{The non-Tate determinant from a skew-adjoint operator.}
Let
\[
  \psi(T)=\det(T-S\mid U_\ell)\in\mathbf Q[T],\qquad
  A=S-S^{-1}\quad\hbox{on }U_\ell.
\]
The polynomial $\psi$ is independent of $\ell$: the
characteristic polynomial after squaring is determined by
$\Phi_W$, and one divides out its simple factor $T-1$.
The operator $A$ is invertible and skew-adjoint for $q_B$.
Consequently $b(u,v)=q_B(Au,v)$ is a nondegenerate alternating
form. Its matrix determinant is a Pfaffian square, giving
\[
                   \det(q_B|_{U_\ell})\equiv\det A
                      \pmod{(\mathbf Q_\ell^*)^2}.
 \tag{A15.11}
\]
Indeed $\det b=\det A\,\det(q_B|_{U_\ell})$ is a square,
and an element and its inverse have the same squareclass.
Reciprocal eigenvalue pairing gives $\det(S|_{U_\ell})=1$.
As the dimension is even, it follows that
\[
       \delta:=\det A=\psi(1)\psi(-1)\in\mathbf Q_{>0}.
 \tag{A15.12}
\]
To verify the strict sign, all roots of $\psi$ have complex
absolute value one and avoid $1,-1$. Pairing each root with
its complex conjugate gives positive factors
$|1-\lambda|^2$ and $|1+\lambda|^2$ in these two evaluations.
If $U_\ell=0$, both the determinant and the empty product
are $1$. The calculation uses characteristic polynomials
only, and requires no semisimplicity on $U_\ell$.

\emph{The missing-$p$ squareclass argument fixes the real sign.}
Combining (A15.6), the orthogonal decomposition by $K_\ell$,
(A15.9), and (A15.11), for every prime $\ell\ne p$ we obtain
\[
            a\,\det G\,\delta/D\in(\mathbf Q_\ell^*)^2.
 \tag{A15.13}
\]
The map
\[
 \mathbf Q^*/(\mathbf Q^*)^2\longrightarrow
        \prod_{\ell\ne p}\mathbf Q_\ell^*/(\mathbf Q_\ell^*)^2
 \tag{A15.14}
\]
is injective. One proof is that a nontrivial quadratic extension
of $\mathbf Q$ has an inert prime outside any fixed finite
set. The exact squareclass implication is also explicitly
proved and cited in Ancona--Marmora, \emph{Commentarii
Mathematici Helvetici} 100 (2025), proof of Corollary~4.4,
printed p.~629, referring to Ireland--Rosen, second edition,
\S5.2, Theorem~3. Applying it to (A15.13) gives
\[
                     a=\frac{D}{\det G\,\delta}\,c^2
                       \quad\hbox{for some }c\in\mathbf Q^*.
 \tag{A15.15}
\]
All three rational factors in the denominator and numerator
on the right are positive. Hence $a>0$.

Every nonzero algebraic residual direction has just been
shown positive. If one such class $z$ exists, every other
residual class is its rational multiple: the multiple is
$q_B(z',z)/q_B(z,z)\in\mathbf Q$, and equality follows in
the one-dimensional nondegenerate $T_{\ell_0}$. Thus
$A_B=K\perp\mathbf Qz$; otherwise $A_B=K$. Positivity
on the whole space follows in either case.
\end{proof}

\begin{corollary}[C44-A15a: the simple rank-one gate]
For any fixed axis in the theorem, if the normalized variable
Frobenius polynomial has at most one root of unity counted
with algebraic multiplicity, its entire axis form has the
required positive sign unconditionally. No algebraic
generator of the possible Tate line needs to be exhibited.
If that line contains a nonzero algebraic class, the argument
proves its positive square; it does not assert its algebraicity.
\end{corollary}

\begin{proof}
Take $K=0$ in C44-A15. The convention $\det G=1$ gives
exactly the asserted conclusion.
\end{proof}

\paragraph{A sharper remaining gate in the abelian sixfold case.}
For $\dim X=6$ and $\dim B=4$, Gauss--Bonnet, weak Lefschetz
and Poincare duality give
\[
        \dim V=5\int_Xh^6-165=3600\chi(X,\mathcal O_X(h))-165,
 \tag{A15.16}
\]
which is odd. All non-root-of-unity eigenvalues of its
orthogonal Frobenius occur in reciprocal pairs. Thus the
total root-of-unity multiplicity is odd. C44-A15a settles
the unavoidable multiplicity-one case: the fixed-axis sign
question not covered by this test now requires multiplicity
at least three. More generally, after a known positive span
of dimension $k$ has been removed, the residual multiplicity
has parity $\dim V-k$. An even-dimensional known span in a
sixfold axis can therefore be completed by this theorem
when its residual multiplicity is one. This sharpens a
cohomological sufficient test; it does not prove that every
axis satisfies the multiplicity bound.

\paragraph{Programme locator and conditional height consequence.}
This result concerns cote~44, Theorem~5, scan~004, its primitive
margin on scan~008, and the exact H9-T7 form on QH9-axis.
It sharpens route~(c) from exclusion of every potential Tate
direction to a sufficient sign test allowing one such residual
direction after a known positive span. The axis sign proved
here is unconditional. By the decomposition in C44-A01, the
corresponding full-height and $I(B)$ consequences still use
the abelian positivity input conditional on the unrefereed
032/Milne2026 results. Neither those conditional consequences
nor this sufficient test settle the arbitrary-axis case.

\paragraph{Primary source checks and scope.}
K\"unnemann's actual scanned Theorems~5.1--5.2 on printed
p.~96 and their surrounding construction were read at
\url{https://gdz.sub.uni-goettingen.de/download/pdf/PPN356556735_0113/LOG_0011.pdf}.
The base-scheme convention on p.~86 was also checked: the
quoted relative formulation is over a field, and only its
absolute special-fibre case is used above. Deligne's actual
Theorem~1.6, printed p.~276, was read at
\url{https://www.numdam.org/item/PMIHES_1974__43__273_0.pdf}.
The complete proof of Ancona--Marmora Corollary~4.4, printed
p.~629, and the signature discussion in Proposition~4.5,
p.~630, were read at
\url{https://ems.press/content/serial-article-files/51595?nt=1}.
The theorem uses only the elementary squareclass consequence
of that paper, not its CM or supersingularity hypotheses.
The geometric lift is C44-A05; no lifting of an individual
test cycle or of the known positive span is used. A bound on
the dimension of algebraic classes alone would not replace
(A15.3), which concerns the whole potential Tate part of the
Frobenius characteristic polynomial.

% END work/scan_w/axis_one_tate_direction.tex

% BEGIN work/axis/constrained_moment_monodromy.tex
% C44_POS /root/axis_independent; GPT-6 Astra Pro; effort not exposed.
\section{Constrained axes: fourth moments without lifting the centres}
\label{sec:c44-a18}

\begin{theorem}[C44-A18: orthogonal monodromy with prescribed centres]
Let $k$ be a finite field of odd characteristic, let $X/k$ be an
abelian variety of dimension $2r+2$, $r\geq2$, and let $H$ be a
very ample line bundle with $h=c_1(H)$. Fix pairwise disjoint
smooth geometrically connected closed subvarieties
$Z_1,\ldots,Z_s\subset X$, all of dimension $r$, and put
$Z=\coprod_i Z_i$; the empty collection is allowed.
For all sufficiently large integers $t$, after a finite extension
of $k$ there is a smooth member
$D_0\in|I_Z\otimes H^{\otimes t}|$ such that the following holds.

There is a nonempty open subset $\Omega$ of the affine space
\[
             E_t=H^0(X,I_Z\otimes H^{\otimes t})
\]
whose points $a$ give smooth complete intersections
$B_a=D_0\cap(a=0)$, and for which the original pencil spanned
by $D_0$ and $(a=0)$ is a genuine Lefschetz pencil on $X$.
For $\ell\ne\operatorname{char}k$, let
\[
 \begin{split}
 \mathcal V_a&=\ker\bigl(\kappa_{a*}:H^{2r}(B_{\bar a},
          \mathbf Q_\ell(r))\longrightarrow
                   H^{2r+4}(X_{\bar k},\mathbf Q_\ell(r+2))\bigr),\\
 b_i&=(L_h^2)^{-1}[Z_i]_X,\qquad
 w_i=[Z_i]_{B_a}-t^{-2}\kappa_a^*b_i,\\
 \mathcal K&=\sum_i\mathbf Q_\ell w_i,
 \qquad \mathcal W=\mathcal K^\perp\cap\mathcal V,
 \qquad Q(x,y)=(-1)^r\int_{B_a}xy.
 \end{split}                                                    \tag{A18.1}
\]
The classes $w_i$ are rational algebraic relative classes, their
constant signed Gram matrix is positive definite, and
$\mathcal V=\mathcal K\perp\mathcal W$ is a splitting of lisse
sheaves. The geometric monodromy closure on $\mathcal W$ is the
full orthogonal group, provided $\operatorname{rank}\mathcal W>8$;
this rank condition also holds for all sufficiently large $t$.
The varieties $Z_i$ and their cycle classes are not assumed to
lift to characteristic zero. No 032 input is used.
\end{theorem}

\begin{proof}
We distinguish the integer $s$ counting the centres, a section
$a\in E_t$, and the cardinality $Q_0$ of a finite extension field.

\emph{The family and the nondegenerate known summand.}
C44-A06 supplies a nonempty open set of pairs of sections with
smooth common zero locus containing $Z$, whose original pencils
are Lefschetz. Its imposed-centre jet proof gives this for every
sufficiently large $t$. Projection of this open set onto the first
section space contains a nonempty open. Choose $D_0$ in this open,
after extending the finite field, and take the corresponding
nonempty open $\Omega\subset E_t$ of second sections. Thus $D_0$
is smooth, and the assertion about the original pencil is part of
the chosen open condition.

The algebraic inverse Lefschetz correspondence on the abelian
variety supplies the classes $b_i$ simultaneously for all
coefficient primes. Self-intersection and hard Lefschetz identify
the ambient image in the middle cohomology of $B_a$ as a
nondegenerate constant summand. Its orthogonal complement is
$\mathcal V$. The cycle $Z_i\times\Omega$ inside the universal
axis gives a relative cycle class; the ambient correction in
(A18.1) is also relative. Hence every $w_i$ is a section of
$\mathcal V$, fixed by normalized arithmetic Frobenius.

The exact Gram formula of C44-A06 is constant in this family and
positive definite for large $t$. In particular the $w_i$ are
independent and $\mathcal K$ is nondegenerate. It follows that
$\mathcal W$ is a lisse nondegenerate orthogonal direct summand.
Nothing here assumes that the stated classes exhaust the fixed
part. That will follow from the monodromy calculation.

In the calculation below we use the untwisted sheaves, denoted by
$\mathcal V^{\rm un}$ and $\mathcal W^{\rm un}$, of weight $2r$.
The $s$ known lines in $\mathcal K^{\rm un}$ have eigenvalue
$Q_0^r$ over a field of cardinality $Q_0$.

\emph{Four-point evaluation outside the centres.}
Choose an integer $c$ such that $I_Z\otimes H^{\otimes c}$ is
globally generated. For $t\geq c+3$, evaluation from $E_t$ is
surjective at any set of at most four distinct geometric points
of $X\setminus Z$. Here is a direct argument, also over a finite
extension $k'/k$ with $Q_0=\#k'>4$. At each such point the
evaluation of $H^0(I_ZH^c)\otimes k'$ is nonzero. The union of at
most four proper linear hyperplanes does not cover this vector
space, so there is a section $e$ nonzero at all these points.
For each selected point $x_i$, very ampleness gives a hyperplane
section through each $x_j$, $j\ne i$, but not through $x_i$.
Their product, padded to degree $t-c$ by a hyperplane section
nonzero at all the points, vanishes at the other points and not
at $x_i$. Multiplying these products by $e$ proves simultaneous
surjectivity. This also proves the assertion geometrically, and
surjectivity descends to $k'$.

Set
\[
                     U=D_0\setminus Z.
\]
It is smooth and geometrically connected of dimension $2r+1$.
The same four-point evaluation statement holds for points of $U$,
because it was proved using sections on $X$. No completeness
claim about the induced linear system on $D_0$ is required.

\emph{A trace function with the centre contribution removed.}
For every $a\in E_t(k')$, including singular or zero sections,
define the rational number
\[
 T(a,k')=\#\bigl(U\cap(a=0)\bigr)(k')-\frac{\#U(k')}{Q_0}.
                                                               \tag{A18.2}
\]
For $a\in\Omega(k')$ we claim the uniform estimate
\[
 T(a,k')=\operatorname{Tr}
       (F_{k',a}\mid\mathcal W^{\rm un}_{\bar a})
                   +O(Q_0^{r-1/2}).                            \tag{A18.3}
\]
The implied constant depends on the fixed geometry, not on $k'$
or $a$ in the smooth open.

Indeed $B_a$ is a smooth ample divisor of the smooth
$(2r+1)$-fold $D_0$. The trace formula, weak Lefschetz in degrees
below $2r$, and its Poincare-dual Gysin statement in degrees above
$2r$, give
\[
 \#B_a(k')-\frac{\#D_0(k')}{Q_0}
       =\operatorname{Tr}(F_{k',a}\mid\mathcal V^{\rm un}_{\bar a})
                       +O(Q_0^{r-1/2}).                       \tag{A18.4}
\]
For clarity, the terms in degrees greater than $2r$ on $B_a$
cancel the corresponding terms of degrees greater than $2r+2$
on $D_0$, divided by $Q_0$. In the middle, the restriction image
of $H^{2r}(D_0)$ has the same Frobenius polynomial as
$H^{2r+2}(D_0)(1)$, by hard Lefschetz; its removal leaves the
variable summand. The remaining terms have weights at most
$2r-1$. Moreover $H^{2r}(X)\simeq H^{2r}(D_0)$ by weak
Lefschetz, so this variable summand is exactly $\mathcal V$ in
(A18.1). The ranks of the low-degree groups are constant and
Deligne's weight bound supplies the stated uniform error.
This is the argument of Katz's Key Lemma 2.2.6, with $D_0$ as
the ambient variety.

Since the $Z_i$ are smooth geometrically connected projective
$r$-folds, the trace formula and the Weil bounds give
\[
                      \#Z(k')=sQ_0^r+O(Q_0^{r-1/2}).           \tag{A18.5}
\]
Subtracting $(1-Q_0^{-1})\#Z(k')$ from (A18.4) gives (A18.2).
On the cohomology side, removing $\mathcal K^{\rm un}$ subtracts
exactly $sQ_0^r$. Since $Q_0^{r-1}=O(Q_0^{r-1/2})$, equations
(A18.4)--(A18.5) prove (A18.3).

\emph{The exact fourth moment.}
Write $M=\#U(k')$, $v=Q_0^{-1}$ and let $a$ be uniformly random
in the finite vector space $E_t(k')$. For $x\in U(k')$, put
$I_x(a)=1$ if $a(x)=0$, and $I_x(a)=0$ otherwise. Four-point
evaluation implies that these variables are independent for
every set of at most four distinct points, and each is a
Bernoulli variable of parameter $v$. Thus
$T(a,k')=\sum_x(I_x(a)-v)$, and elementary expansion gives the
exact identity
\[
 \begin{split}
 \frac{1}{\#E_t(k')}\sum_{a\in E_t(k')} T(a,k')^4
   &=3M(M-1)\sigma^4+M\mu_4,\\
 \sigma^2&=v(1-v),\\
 \mu_4&=v(1-v)^4+(1-v)v^4.
 \end{split}                                                   \tag{A18.6}
\]
Only coincidence patterns $(2,2)$ and $(4)$ survive: a singleton
has expectation zero. As $U$ is a dense open in the smooth
geometrically connected $(2r+1)$-fold $D_0$,
$M=Q_0^{2r+1}+O(Q_0^{2r+1/2})$. Consequently
\[
 Q_0^{-4r}\frac{1}{\#E_t(k')}
       \sum_{a\in E_t(k')}T(a,k')^4
                   =3+O(Q_0^{-1/2}).                         \tag{A18.7}
\]
For example, the second term in (A18.6), after this normalization,
is $O(Q_0^{-2r})$.

Let $G$ be the geometric monodromy closure of $\mathcal W$ and
$W$ its fibre after extension of the coefficient field to an
algebraic closure. The moment is
\[
        M_4(G,W)=\dim\bigl(W^{\otimes2}\otimes
                                  (W^\vee)^{\otimes2}\bigr)^G.
\]
Purity makes $G$ reductive. Katz's Lemma 2.1.4 identifies this
integer with the limsup of the normalized fourth trace moments
over the smooth parameter open. Equation (A18.3) changes those
moments by $o(1)$: both the trace and $T$ on this open are
$O(Q_0^r)$, and their difference is $O(Q_0^{r-1/2})$.
Since fourth powers are nonnegative, adding all the omitted
sections can only increase the sum. Therefore (A18.7) gives
\[
                         M_4(G,W)\leq3.                       \tag{A18.8}
\]
This is precisely the use of Katz's Second Variant Lemma 2.1.6.
No estimate on traces of singular members is used or needed.
Because $G\subset O(W,Q)$ and $\dim W>8$, the three standard
orthogonal contraction invariants are independent. Hence
$M_4(G,W)\geq3$, and equality holds.

\emph{A nonzero reflection in the residual summand.}
For large $t$ the imposed-centre jet argument also produces a
section whose restriction to $D_0$ has exactly one ordinary
quadratic singularity at a chosen point $x\in U$, and is smooth
everywhere else, including along $Z$.
Here one prescribes the value and first derivatives to vanish at
$x$, and chooses a nondegenerate quadratic term on $T_xD_0$;
odd characteristic is used here. The system of sections with
this fixed second jet has arbitrary first jets at a second
point outside $Z\cup\{x\}$ for large $t$. A further singular
point therefore imposes $2r+2$ conditions while varying in
dimension $2r+1$, so a general section has no such point.
Along $Z_i$, its normal first derivative is a section of the
rank-$(r+1)$ bundle
$N^*_{Z_i/D_0}\otimes H^t|_{Z_i}$. This bundle is generated by
the constrained sections even after the prescribed jet at $x$;
its zero condition has codimension $r+1>r=\dim Z_i$, so a
general section has no singular point there either.
The required uniform generation follows as in C44-A06 by
twisting the ideal and using bounded jets; equivalently,
hyperplanes vanishing to order three at $x$ and nonzero at a
second point preserve the fixed jet while generating the
indicated first jets. Thus these dimension counts apply to
actual linear systems.

Choose a transverse one-parameter deformation of this nodal
member, by a section nonzero at $x$. Its total space is regular
at the node, and proper over the local trait; its punctured
generic point can be chosen in $\Omega$, since $\Omega$ is
dense open. SGA 7 II, Expose XV, Theorem 3.4(i),(iii), printed
pp.~195--196, gives a global vanishing cycle $\delta$ with
\[
                       \int\delta^2=2(-1)^r.                 \tag{A18.9}
\]
In particular $\delta\ne0$. The same theorem says that the
quadratic inertia character is surjective when the total space
is regular. An inertia element acts by the Picard--Lefschetz
reflection in $\delta$.

Every ambient class and every $w_i$ is fixed by this local
monodromy. The reflection formula and $\delta\ne0$ then imply
that $\delta$ is orthogonal to both the ambient image and
$\mathcal K$. Thus $\delta\in W$, and this is a nontrivial
reflection on $W$, not merely on the entire middle cohomology.
Katz's Theorem 1.4.6 now applies to (A18.8) and proves
$G=O(W,Q)$.

Finally the rank of the variable cohomology tends to infinity
with $t$: C44-A03b gives
\[
 \operatorname{rank}\mathcal V
  =(2r+1)t^{2r+2}\int_Xh^{2r+2}
       +2\binom{4r+3}{2r-1}-\binom{4r+4}{2r}.
                                                               \tag{A18.10}
\]
Subtracting the fixed integer $s$ proves the remaining rank
assertion. Restriction to any further nonempty parameter open
does not change geometric monodromy, since its fundamental
group surjects onto that of the original smooth open.
\end{proof}

\begin{corollary}[C44-A18a: closed axes through arbitrary smooth centres]
With the hypotheses and sufficiently large $t$ of C44-A18,
there is a Zariski-dense set of closed parameters $a\in\Omega$
for which the signed axis form is positive definite on the
\emph{whole} rational algebraic variable space $A_{B_a}$.
Put $d=\operatorname{rank}\mathcal W$.
If $d$ is even, then
\[
                         A_{B_a}=\sum_i\mathbf Q w_i.
                                                               \tag{A18.11}
\]
If $d$ is odd, the right side is either all of $A_{B_a}$ or has
codimension one, and the additional line has positive square.
No lifting of the centres, Tate conjecture or 032 input is used.
\end{corollary}

\begin{proof}
The normalized Frobenius polynomial of $\mathcal W$ is rational:
it is the middle polynomial of $B_a$ divided by the ambient
polynomial and by $(T-1)^s$. In even rank C44-A12 selects a dense
set of closed fibres with no root of unity on $\mathcal W$,
which gives (A18.11). In odd rank C44-A16 selects a dense set
with total root-of-unity multiplicity one. The positive Gram
matrix of C44-A06 and the unconditional determinant criterion
C44-A15 then give precisely the stated alternative and sign.
These statements apply to the original finite-field axes, not
only to geometric generic parameter points.
\end{proof}

\paragraph{Application to the nonliftable centres and the remaining gate.}
Take the intrinsically nonliftable surface $S$ of C44-A14,
embedded as $S\times\{0\}^2\subset E^6$ in characteristic three.
There is just one prescribed centre. Formula (A18.10) makes the
residual rank even, so C44-A18a yields actual closed genuine
Lefschetz axes with $A_B=\mathbf Q w_S$ and $Q(w_S,w_S)>0$.
For every such axis the uncorrected class $[S]_B$ remains
nonliftable under any smooth projective characteristic-zero
lift, by C44-A14. This does \emph{not} certify that $w_S$ itself
is nonliftable: the ambient correction cancels the known
fourfold cup tensor exactly. Thus the construction removes the
monodromy/lift-of-centre obstruction, but the individual
primitive-class lifting question remains open.

\paragraph{Primary source locators and Grothendieck's programme.}
The moment machinery used above is N.~Katz,
\emph{Larsen's alternative, moments, and the monodromy of
Lefschetz pencils}, in \emph{Contributions to Automorphic Forms,
Geometry, and Number Theory} (2004), pp.~521--560:
Theorem 1.4.6 and its proof on p.~535; Lemmas 2.1.4--2.1.6
on pp.~540--541; Key Lemma 2.2.6 with complete proof on
pp.~544--545; the four-point calculation in 2.2.7--2.2.13 on
pp.~545--550. Author PDF:
\url{https://web.math.princeton.edu/~nmk/LAMMLF.pdf}.
The exact global norm and nontrivial inertia assertions are
SGA 7 II, Expose XV, Theorem 3.4(i),(iii), pp.~195--196;
the original-pencil existence uses the imposed-centre extension
proved in C44-A06, with SGA 7 II, Expose XVII, 2.4 and
3.7--4.3 as its ordinary-pencil input.

This settles the signed axis summand of JAC1964 H9-T7 for the
closed axes selected in C44-A18a. In cote 44, Th.~5, scan 004,
it gives the remaining axis sign for these actual special
pencils, on the primitive domain specified by the margin on
scan 008. Any assembly of the full height with the separate
abelian primitive term still uses the unrefereed 032/Milne 2026
input for that separate term. The result does not prove the
sign for every fixed axis and supplies no wrong-sign class.
In the SGA and Weil-cohomology framework, its new point is a
full orthogonal local system over the constrained parameter
space, constructed without a characteristic-zero model of its
prescribed centres.

% END work/axis/constrained_moment_monodromy.tex

% BEGIN work/axis/axis_realization_kernel.tex
% C44_POS /root/axis_independent. Author contribution; independent review pending.
\section{A common cycle-class kernel and a legitimate crystalline lifting test}
\label{sec:c44-a18b}

\begin{proposition}[C44-A18b: realization independence after one potential Tate direction]
Let $X/\mathbf F_q$ be an abelian variety of dimension $2r+2$,
and let $\kappa:B\hookrightarrow X$ be a smooth intersection of two
divisors of the same ample class $h$. Let $\Gamma$ be an absolute
Chow correspondence realizing the inverse of
$L_h^2:H^{2r}(X)\longrightarrow H^{2r+4}(X)(2)$, and put
\[
 P=\Delta_B-\kappa^*\Gamma\kappa_*,\qquad
 \mathscr C=P_*\mathrm{CH}^r(B_{\overline{\mathbf F}_q})_{\mathbf Q},
 \qquad q_B(z,z')=(-1)^r\deg(zz').
 \tag{A18b.1}
\]
The same correspondence is used in every realization. Choose actual
cycles $C_1,\ldots,C_s\in\mathscr C$ with nonsingular rational Gram
matrix $G=(q_B(C_i,C_j))$; $s=0$ is allowed. After extending the finite
ground field, assume these cycles and the correspondences are defined
over it. Let $K_\ell$ be their span and
\[
 V_\ell=\ker\kappa_*\subset H^{2r}(B,\mathbf Q_\ell(r)),\qquad
 W_\ell=V_\ell\cap K_\ell^\perp.
\]
Suppose the total algebraic multiplicity of root-of-unity roots of
the normalized Frobenius polynomial $\Phi_W$ is at most one.
Then the kernels, on the rational group $\mathscr C$, of every
$\ell$-adic cycle-class map ($\ell\ne p$) and of the rational
crystalline cycle-class map coincide.

Here a crystalline class over $\overline{\mathbf F}_q$ is tested after
any finite field of definition and extension of its Witt fraction
field; faithful scalar extension makes the vanishing condition
independent of that choice. The assertion is about projected cycles,
and does not assert realization independence for all cycles on $B$.
No positivity assumption on $G$, no Tate conjecture and no result
conditional on 032 are needed.
\end{proposition}

\begin{proof}
K\"unnemann's absolute Chow-motivic hard Lefschetz theorem
(\emph{Inventiones Mathematicae} 113 (1993), Theorem~5.2,
printed p.~96; Theorem~5.1 gives the decomposition used there)
supplies one $\Gamma$. If necessary replace the divisor by its
symmetric rational symmetrization, which has the same cohomology
class. Because this is a motivic inverse, it acts as the required
inverse in rational crystalline cohomology as well as in every
$\ell$-adic realization. We use only the absolute theorem over a
field, not a relative theorem over a mixed-characteristic base.
The self-intersection formula $\kappa_*\kappa^*=L_h^2$ then shows
that $P$ is the orthogonal cohomological projector onto $V$ in
each realization. Its action on Chow groups need not be idempotent.
Both the ambient summand and its orthogonal complement $V$ are
nondegenerate, by hard Lefschetz and Poincar\'e duality.

Intersection numbers are rational and agree with cup-product
traces in each realization. Thus the $C_i$ have the same matrix
$G$, are independent in each realization, and give nondegenerate
subspaces $K$. The characteristic polynomial on $V$ is the quotient
of the normalized middle polynomial of $B$ by that on
$H^{2r}(X)$. Katz--Messing, \emph{Some consequences of the Riemann
hypothesis for varieties over finite fields}, \emph{Inventiones
Mathematicae} 23 (1974), Theorem~1(i), printed p.~74, identifies
these characteristic polynomials in the reasonable cohomology
theories, including rational crystalline cohomology. Dividing by
$(T-1)^s$ gives the same rational polynomial $\Phi_W$ in all of them.
For crystalline cohomology over $\mathbf F_{q'}$, normalized
Frobenius means $q'^{-r}$ times the $q'$-power crystalline Frobenius,
which is linear over $W(\mathbf F_{q'})[1/p]$ and is an isometry.
Passing to a finite extension raises all normalized eigenvalues
to a common power and preserves their total root-of-unity
multiplicity.

We record the needed linear algebra without a semisimplicity
assumption. Let $F$ be an isometry of a nondegenerate symmetric
space over a characteristic-zero field, with rational
characteristic polynomial having at most one root of unity,
counted with algebraic multiplicity. If such a root exists it is
$1$ or $-1$ and is simple. Its eigenline $T$ is nondegenerate:
the generalized eigenspaces for every other eigenvalue are
orthogonal to $T$, since an isometry pairs only reciprocal
generalized eigenvalues. Were $T$ degenerate, its nonzero vector
would be orthogonal to the whole space. Moreover every vector
annihilated by some $F^e-1$ belongs to $T$; if there is no
root-of-unity root, every such vector is zero. In particular a
finite-orbit vector of square zero is zero.

Let $Z\in\mathscr C$, and first suppose its class vanishes in one
of the realizations under consideration. The cycle-theoretic
equalities
\[
                  q_B(Z,C_i)=0\quad(1\leq i\leq s),
                  \qquad q_B(Z,Z)=0
 \tag{A18b.2}
\]
follow in $\mathbf Q$. In every other realization, $Z$ lies in
$V$ because it is an actual projected cycle, and it lies in
$K^\perp$ by (A18b.2). Choose a finite field of definition for $Z$
as well. Its normalized Frobenius orbit is finite, and its square
is zero by (A18b.2). The preceding linear algebra on $W$ forces
its class to vanish. This proves all kernel equalities at once.
\end{proof}

\begin{corollary}[C44-A18c: a necessary filtered test for a class, not just a representative]
Keep the assumptions of C44-A18b. Let $R$ be a complete discrete
valuation ring of mixed characteristic with finite residue field,
and suppose that, after an allowed finite residue-field extension,
there is a smooth projective lift
\[
          \mathcal B\lhook\joinrel\longrightarrow\mathcal X
                       \longrightarrow\operatorname{Spec}R
 \tag{A18c.1}
\]
of the polarized embedded complete intersection $B\subset X$,
where $\mathcal X$ is an abelian scheme and $\mathcal B$ is cut
out by two relative divisors of its lifted line bundle. Write
$K_R=\operatorname{Frac}(R)$ and $K_0=W(k)[1/p]$. For
$Z\in\mathscr C$, suppose that its class in one $\ell$-adic
realization is the specialization of the class of a rational
algebraic cycle on $\mathcal B_{\overline K_R}$. Then, under
crystalline--de Rham comparison for this lift,
\[
 \operatorname{cl}_{\mathrm{cris}}(Z)\otimes1
       \in \operatorname{Fil}^r
                H^{2r}_{\mathrm{dR}}(\mathcal B_{K_R}/K_R).
 \tag{A18c.2}
\]
Finite extensions of $R$ may be made to define the generic cycle;
membership in the displayed filtration descends by faithful
field extension. The conclusion is only a necessary condition
for this specified embedded lift. It gives neither a converse
nor a lift-independent failure of (A18c.2).
\end{corollary}

\begin{proof}
After a finite extension let $\beta$ be a generic-fibre cycle
whose class specializes to $\operatorname{cl}_\ell(Z)$.
On the generic abelian variety choose an algebraic inverse
$\Gamma_\eta$ of $L_h^2$ and form the analogous algebraic
projector $P_\eta$. Its $\ell$-adic action is the orthogonal
projector onto the generic Gysin kernel. Smooth proper base
change, together with compatibility of Gysin and cup products,
identifies that projector with the special cohomological
projector $P$. Consequently $\beta'=P_{\eta*}\beta$ has the
same specialization of its $\ell$-adic class as $\beta$.

Let $\gamma=\operatorname{sp}(\beta')$ be the specialization
as an actual rational Chow cycle. The specialized correspondence
$\operatorname{sp}(\Gamma_\eta)$ acts as the inverse of $L_h^2$
also in crystalline cohomology: crystalline--de Rham comparison
intertwines its action with the generic de Rham inverse.
That inverse is unique. Thus $\operatorname{cl}_{\mathrm{cris}}
(\gamma)$ belongs to $V_{\mathrm{cris}}$, and applying the
original special-fibre $P$ leaves it unchanged. No lifting of the
original special-fibre correspondence $\Gamma$ has been assumed.

The actual cycle $P_*\gamma-Z$ belongs to $\mathscr C$ and has
zero $\ell$-adic class. C44-A18b gives
\[
 \operatorname{cl}_{\mathrm{cris}}(Z)
   =\operatorname{cl}_{\mathrm{cris}}(P_*\gamma)
   =\operatorname{cl}_{\mathrm{cris}}(\gamma).
 \tag{A18c.3}
\]
The right side corresponds to the de Rham class of $\beta'$,
which belongs to $\operatorname{Fil}^r$ because $\beta'$ is an
algebraic codimension-$r$ cycle. This proves (A18c.2).

For the comparison invoked here, see Berthelot--Ogus,
\emph{$F$-isocrystals and de Rham cohomology. I},
\emph{Inventiones Mathematicae} 72 (1983), Theorem~2.4,
printed pp.~169--170, Corollary~2.5 and base-change
Proposition~2.7, printed p.~171. Corollary~3.7 and
Remark~3.7.1, printed pp.~178--179, give compatibility with
the Chern classes of vector bundles over the same possibly
ramified trait. This also gives the cycle compatibility used
above: take the horizontal closure of each generic integral
cycle, resolve its structure sheaf by vector bundles on the
regular projective total space, and use the degree-$r$ leading
term of its Chern character. Flatness of the horizontal closure
over the trait identifies its fibre cycle with Chow
specialization. The identical argument on the product gives
compatibility for specialized correspondences.
Compatibility with cycle
classes and specialized correspondences is the standard cycle
compatibility; Emerton, \emph{A $p$-adic variational Hodge
conjecture and modular forms with complex multiplication},
\S2, (2.1) and the paragraph immediately before Conjecture~2.2,
printed pp.~4--5, explicitly states the cycle compatibility,
and Remark~2.5, printed pp.~6--7, the correspondence compatibility.
For complete ramified traits, the same comparison is stated in
Lieblich--Olsson, \emph{Fourier--Mukai partners of K3 surfaces in
positive characteristic}, \emph{Annales scientifiques de
l'\'Ecole normale sup\'erieure} 48 (2015), \S2.7, printed p.~1007,
with cup-product and trace compatibility. The cycle compatibility
comes from the Chern-class construction; it is not a variational
Hodge conjecture or an assertion of its converse.
\end{proof}

\paragraph{Exact scope for the programme.}
For the closed axes selected in C44-A18a, these results identify
the crystalline class of a projected algebraic $\ell$-adic class
independently of its chosen projected representative. In particular,
for the axes containing the nonliftable smooth surface of C44-A14,
it is legitimate to test the crystalline filtration of
$w_S=P_*[S]$ along any specified embedded lift. What has not been
proved is that this class fails the filtration test for every
possible lift. Nonliftability of $S$ itself, or of the uncorrected
class $[S]$, supplies no such assertion: the ambient cup tensor
cancels under $P$, as proved in C44-A14a. These results repair a
possible class-versus-representative gap in route (a) of
JAC1964 QH9-axis; they do not close that row or prove a new
universal sign assertion. In Grothendieck's language they concern
the Weil realizations of the primitive axis correspondence in
cote~44, Th.~5, scan~004 and the primitive warning on scans~008--009,
without assuming condition $D(B)$ for all of $B$.

% END work/axis/axis_realization_kernel.tex

% BEGIN work/scan_w/equivariant_axis_blocks.tex
% C44_POS /root/scan_w. No universal lifting of a special Tate projector asserted.
\section{Several potential Tate directions separated by geometric symmetries}
\label{sec:c44-a19}

\begin{theorem}[C44-A19: the orthogonal block criterion]
Let $B\subset X/\mathbf F_q$ be the smooth intersection of two
members of one ample line bundle as in C44-A15, with
$\dim X=2r+2$. Suppose an embedded polarized lift
$\mathcal B\subset\mathcal X$ over a mixed-characteristic trait
has a characteristic-zero generic fibre admitting a complex
embedding. On its relative variable middle cohomology suppose
there are rational algebraic correspondences $e_1,\ldots,e_b$
whose realizations are pairwise orthogonal self-adjoint
idempotents summing to the identity. In particular they
preserve this variable cohomology on both fibres and in every
coefficient realization. Put
\[
 V_{i,\ell}=e_iV_\ell,\qquad
          q=(-1)^r\int_B(-)(-).
 \tag{A19.1}
\]
Assume that the normalized Frobenius polynomial on every
$V_{i,\ell}$ is rational and independent of $\ell\ne p$.
Let $K\subset A_B$ be a known positive rational algebraic
subspace stable under the $e_i$; $K=0$ is allowed. Use the
common cycle projection described below to realize
$K_i=e_iK$ in each $V_{i,\ell}$, and let $\tau_i$ be the total
root-of-unity multiplicity on $V_{i,\ell}\cap K_{i,\ell}^\perp$.
If
\[
                            \tau_i\leq1\quad(1\leq i\leq b),
 \tag{A19.2}
\]
then $q$ is positive definite on the whole rational algebraic
axis space $A_B$. Its rational dimension satisfies
\[
     \dim K\leq\dim A_B\leq\dim K+
                              \#\{i:\tau_i=1\}.
 \tag{A19.3}
\]
Every additional nonzero algebraic direction has positive
square; the existence or liftability of those directions is
not asserted. No Tate conjecture, Frobenius semisimplicity,
standard conjecture $D(B)$, or 032 input is used.
\end{theorem}

\begin{proof}
All Gram determinants use the ordinary convention of A15.
The relative algebraic correspondences give rational Hodge
endomorphisms on the Betti variable space of the generic
fibre. Their images $V_{i,\mathrm B}$ are orthogonal
nondegenerate rational Hodge substructures. They are primitive
for the inherited ample class: the entire variable space is
primitive by the dual weak-Lefschetz argument of A15.
On each image the Hodge--Riemann relations make the signed
form positive on type $(r,r)$, while every negative real
summand arises from a conjugate pair of off-diagonal Hodge
types. Its negative index is therefore even. Consequently
\[
              D_i:=\det(q|V_{i,\mathrm B})>0.
 \tag{A19.4}
\]
For a zero-dimensional image take $D_i=1$. Smooth proper
base change and Betti--etale comparison for these same
relative correspondences identify this determinant
squareclass with that on $V_{i,\ell}$ for every $\ell\ne p$.
The special algebraic classes themselves need not lift.

Here is the promised correction of cycle representatives.
Use the common special-fibre algebraic projector
\[
                  P=1-\kappa^*(L_h^2)^{-1}\kappa_*
 \tag{A19.5}
\]
from A15, constructed using K\"unnemann's absolute Chow
Lefschetz inverse. Apply $e_iP$ to chosen representatives
of $K$. At the initial coefficient prime these images span
$K_i$. Choose a rational basis from these images. Its
positive rational Gram matrix $G_i$ makes the chosen
cycles independent in every realization; their classes
are in $V_{i,\ell}$ in every realization. This choice
does not infer realization-independent vanishing from
vanishing at one prime. After a finite extension all
these cycles and correspondences are defined, and their
normalized Frobenius action on $K_{i,\ell}$ is trivial.
The residual polynomial is the given block polynomial
divided by $(T-1)^{\dim K_i}$, hence remains rational and
independent of $\ell$. Field extension preserves $\tau_i$.

For an arbitrary class in $A_B$, first apply $P$ to its
cycle representative, then apply each $e_i$, and subtract
its orthogonal projection to the chosen $K_i$ using
$G_i^{-1}$ and rational intersection numbers. This gives
actual rational cycles in the residual block for every
$\ell\ne p$, whose initial classes sum to the given
class modulo $K$. Distinct blocks remain orthogonal.

The determinant proof of A15 now applies separately to
each block, as can also be seen directly. An algebraic
residual class has finite Frobenius orbit. If $\tau_i=0$
it is zero. If $\tau_i=1$, the unique cyclotomic root is
simple and is $1$ or $-1$. For $S=F^2$ its eigenspace
$T_{i,\ell}$ at $1$ is a nondegenerate line; write its
orthogonal complement as $U_{i,\ell}$. A nonzero residual
cycle $z$ at the initial prime has
$a=q(z,z)\in\mathbf Q^*$, so its realization is nonzero
and spans $T_{i,\ell}$ at every prime.

On $U_{i,\ell}$, the operator $S-S^{-1}$ is invertible
and skew-adjoint. The determinant of the resulting
alternating pairing is a square. If $\psi_i$ is the
common rational polynomial of $S$ on $U_{i,\ell}$, this
gives, exactly as in A15,
\[
 \det(q|U_{i,\ell})\equiv
      \delta_i:=\psi_i(1)\psi_i(-1)>0
                    \pmod{(\mathbf Q_\ell^*)^2}.
 \tag{A19.6}
\]
Positivity of $\delta_i$ follows by pairing conjugate
unit-circle eigenvalues, none equal to $\pm1$.
Consequently
\[
                 a\det(G_i)\delta_i/D_i
                  \in(\mathbf Q_\ell^*)^2
                         \quad\hbox{for every }\ell\ne p.
 \tag{A19.7}
\]
The missing-$p$ rational squareclass injection used in
A15 makes this a rational square. The other displayed
factors are positive, so $a>0$. Pairing any further
residual cycle with $z$ gives its rational coefficient
on the same line. Each block therefore adds at most
one positive rational direction. Orthogonality proves
both the sign and (A19.3).
\end{proof}

\begin{corollary}[C44-A19a: a tame elementary two-group criterion]
Let $R$ be a mixed-characteristic discrete valuation ring
with finite residue field $k$ of odd characteristic, and
let $(\mathcal X,\mathcal L)/R$ be a polarized abelian
scheme of relative dimension $2r+2$. Suppose the constant
group $G=(\mathbf Z/2\mathbf Z)^s$ acts on $\mathcal X$
and $\mathcal L$ is given an actual $G$-linearization.
Let $B\subset X_k$ be smooth and be cut out by two
simultaneous eigenvectors
\[
       s_0\in H^0(X_k,L)_{\chi_0},\qquad
       s_1\in H^0(X_k,L)_{\chi_1}.
 \tag{A19.8}
\]
They may have the same character. Assume a complex
embedding of the generic field as in A19.

Then the actual embedded $G$-axis lifts. Its rational
orthogonal character projectors are
\[
        e_\chi=\frac1{|G|}\sum_{g\in G}\chi(g)g^*,
             \qquad\chi:G\longrightarrow\{1,-1\}.
 \tag{A19.9}
\]
Their normalized Frobenius polynomials on variable
cohomology are rational and independent of $\ell\ne p$.
For any known positive $G$-stable rational algebraic
subspace $K$, condition (A19.2) on these character
blocks implies positivity on all of $A_B$.
In particular distinct character blocks can each
contribute one unknown potential Tate direction.
No extension of a special-fibre Frobenius projector
is part of this hypothesis or conclusion.
\end{corollary}

\begin{proof}
The ample special line bundle $L$ has a nonzero section.
Its polarization has finite kernel, so the index theorem
for line bundles on abelian varieties gives
$H^j(X_k,L)=0$ for $j>0$. Upper semicontinuity gives the
same vanishing on the generic fibre. Coherent base change
therefore makes $H^0(\mathcal X,\mathcal L)$ a finite
free $R$-module with surjective reduction onto
$H^0(X_k,L)$. These are Milne's Theorems 13.3(c) and
4.2(c),(e), author notes pp.~25 and 6--7, respectively.

Because $2$ is invertible in $R$, the operators (A19.9)
also act on this module. Choose arbitrary lifts of
$s_0,s_1$ and apply $e_{\chi_0},e_{\chi_1}$, respectively.
The resulting sections are actual eigensection lifts
with their prescribed reductions. Their common zero
locus $\mathcal B$ is $G$-stable. Along its special fibre
the two relative differentials have full rank, since
$B$ is smooth of codimension two. Thus $\mathcal B/R$
is smooth along the special fibre. The nonsmooth locus
is closed and proper over $R$; if nonempty its image
would contain the closed point. It is therefore empty.
This proves smoothness of the entire embedded lift.

The graph correspondences of the $G$-action are relative
algebraic correspondences. They commute with restriction,
Gysin and multiplication by $h$, and so preserve the
variable middle cohomology. An automorphism's adjoint
for cup product is its inverse. Here every element is
its own inverse, so (A19.9) are self-adjoint,
pairwise orthogonal rational idempotents summing to one.
No restriction $\ell\nmid|G|$ is needed: the coefficient
fields have characteristic zero even when $\ell=2$.

We prove degreewise rationality of the equivariant
Frobenius traces, rather than infer it from the total
characteristic polynomial. All group elements are
defined over $k$, hence commute with Frobenius. For
each $g\in G$ and each positive integer $n$, base change
$B$ to $k_n=\mathbf F_{q^n}$ and descend its base change
to $k_{2n}$ by the cocycle taking the nontrivial Galois
element to $g$. Effective projective Galois descent
gives a smooth projective twist $B^{(g,n)}/k_n$.
Its Frobenius on cohomology is $g^*F^n$ (the inverse
convention gives the same operator since $g=g^{-1}$).
Deligne's Weil I, Theorem 1.6, gives rational,
coefficient-independent characteristic polynomials
degree by degree for this twist. In particular
\[
       \operatorname{Tr}(g^*F^n\mid H^j(B_{\bar k}))
 \tag{A19.10}
\]
is rational and independent of $\ell$. The same argument
applies to $X$. Twisting is legitimate as a projective
scheme: finite Galois descent is effective here, and
the given linearized ample line bundle descends with
the cocycle. See Stacks, Section 39.25, Lemma 39.25.3
in the version read for this proof.

Weak Lefschetz and hard Lefschetz identify the ambient
middle image with $H^{2r}(X)$ in a $G$- and
Frobenius-equivariant way. Subtract its traces in
(A19.10) and divide by the Tate-normalizing factor.
Taking the rational average (A19.9) now gives all
power traces of Frobenius on each $e_\chi V_\ell$.
Newton identities recover a rational characteristic
polynomial independent of $\ell$. Its degree is also
common, by the equivariant smooth proper comparison
with the generic Betti fibre. The hypotheses of A19
are now established, including the two requirements
that are not automatic for an arbitrary special
Frobenius eigenspace: an actual extending self-dual
algebraic block and a common rational polynomial.
\end{proof}

\begin{corollary}[C44-A19b: finite-group matrix idempotents]
In A19a replace the elementary two-group by any finite
constant group $G$ with $p\nmid|G|$. Assume the given
$G$-action and linearized ample bundle extend over $R$.
Let $U_R$ be a free rank-two $R$-representation of $G$,
and suppose a $G$-equivariant map
$U_k\longrightarrow H^0(X_k,L)$ has a smooth
codimension-two common zero locus $B$.
Then this actual embedded axis and its $G$-action lift.
Every set of pairwise orthogonal self-adjoint rational
idempotents in $\mathbf Q[G]$ summing to one satisfies
the geometric and rational-polynomial hypotheses of A19.
Self-adjointness here is for the group-algebra involution
$g^*=g^{-1}$. The idempotents need not be central.

For $G=S_3$ and a transposition $\sigma$, one such set is
\[
 \begin{aligned}
 e_1&=\tfrac16\sum_{g\in S_3}g,&
 e_{\mathrm{sgn}}&=\tfrac16\sum_{g\in S_3}\mathrm{sgn}(g)g,\\
 e_{\mathrm{std}}&=1-e_1-e_{\mathrm{sgn}},&
 e_+&=\tfrac12e_{\mathrm{std}}(1+\sigma),\qquad
 e_- =\tfrac12e_{\mathrm{std}}(1-\sigma).
 \end{aligned}
 \tag{A19.11}
\]
The four orthogonal blocks are $e_1,e_{\mathrm{sgn}},e_+,e_-$.
The last two have the same Frobenius polynomial, that of
the standard multiplicity module. If the trivial block
has no potential Tate direction and the sign and standard
multiplicity modules each have at most one, then the
axis sign is positive even though its total potential
Tate multiplicity can be three.
\end{corollary}

\begin{proof}
The coherent base-change argument of A19a is unchanged.
Lift the given map $U_k\to H^0(X_k,L)$ arbitrarily to
an $R$-linear map. Average it by
\[
                 \varphi\longmapsto
       \frac1{|G|}\sum_{g\in G}g\varphi g^{-1}.
 \tag{A19.12}
\]
Since $|G|$ is a unit in $R$, its reduction is the
original equivariant map. The same Jacobian and properness
argument shows that its common zero locus is a smooth
embedded lift. This does not assume that an arbitrary
residue-field representation lifts: $U_R$ is given.

Every group-algebra element defines a relative graph
correspondence, and the adjoint involution is $g\mapsto g^{-1}$.
Thus the stated algebraic identities give the required
orthogonal blocks in every realization. For rationality,
the twisting argument (A19.10) works for every finite-order
$g$: use the cyclic extension of degree $\operatorname{ord}(g)$
over $k_n$, choosing the cocycle or its inverse according
to the Frobenius convention. Hence all
$\operatorname{Tr}(g^*F^n\mid H^j)$ are rational and
independent of $\ell$. For $e=\sum a_gg$ with $a_g\in\mathbf Q$,
\[
       \operatorname{Tr}(F^n\mid eV)
                    =\operatorname{Tr}(eF^n\mid V)
                    =\sum_g a_g\operatorname{Tr}(gF^n\mid V),
 \tag{A19.13}
\]
because $e$ commutes with Frobenius and is idempotent.
Newton identities again give a common rational polynomial.
Centrality of $e$ is unnecessary.

The formulas (A19.11) are self-adjoint orthogonal
idempotents by the elementary character decomposition
of $\mathbf Q[S_3]\simeq\mathbf Q\oplus\mathbf Q\oplus M_2(\mathbf Q)$.
On the standard representation
$\{(x_1,x_2,x_3)\in\mathbf Q^3:\sum x_i=0\}$ with
its usual positive dot product, the two $\sigma=(12)$
eigenlines are spanned by $(1,1,-2)$ and $(1,-1,0)$;
their squared lengths are $6$ and $2$. They tensor with
the same multiplicity module, on which Frobenius acts.
The geometric pairing on the isotypic part is this
dot product tensor a rational symmetric pairing on
that module. Thus the two block polynomials agree,
and their relative scalar factors are positive.
Applying A19 to the four blocks proves the assertion.
\end{proof}

\paragraph{An explicit source of the geometric hypotheses.}
For odd residue characteristic take the good-reduction
model of $E:y^2=x^3-x$, put $\mathcal X=\mathcal E^6$,
and use $\mathcal L=\boxtimes_{i=1}^6\mathcal O_{\mathcal E}(t[0])$.
Coordinate sign changes give a constant elementary
two-group and its natural linearization on this line
bundle. Every smooth complete intersection of two
simultaneous eigensections in its special fibre
therefore has the lift and character correspondences
required in A19a. For instance, for even $t=2d$,
the invariant subsystem is the pullback of
$|\mathcal O(d,\ldots,d)|$ on $(\mathbf P^1)^6$ under
the coordinate double covers; a general pair transverse
to the finitely many branch strata gives a smooth
complete intersection for sufficiently large $d$.
This describes actual geometry, not a postulated
relative projector. It does not by itself verify
(A19.2) for a specified finite-field member.

\paragraph{Scope of the additional test.}
A19 can settle a residual block of total potential
Tate multiplicity two or more when its directions
are separated by these extending rational orthogonal
projectors and every block satisfies (A19.2).
It does not settle two potential directions in one
character block. No explicit fixed-axis member of
total multiplicity at least three satisfying all of
these conditions is certified here. In particular
one must not declare such a member solely from the
existence of sign symmetries. A pencil spanned by
invariant sections can have simultaneous orbits of
nodes; the displayed smooth complete intersections
are not being asserted to be axes of genuine
Lefschetz pencils without a separate check.

For cote 44, Theorem 5, scan 004 and the primitive
qualification on scan 008, the criterion controls
the same H9-T7 axis form. The positivity conclusion
is unconditional under its stated geometric and
Frobenius tests. As in A15, assembly with the separate
abelian term in the full height remains conditional
on the unrefereed 032/Milne2026 input. The arithmetic
rank-two obstruction of A17 is not removed for an
axis lacking an extending orthogonal decomposition.

% END work/scan_w/equivariant_axis_blocks.tex

% BEGIN work/scan_w/affine_s3_axis_geometry.tex
% C44_POS /root/scan_w. Geometry and parity input for C44-A21.
\section{An affine symmetric-group family of genuine Lefschetz axes}
\label{sec:c44-a20}

\begin{theorem}[C44-A20: the affine $S_3$ axis family]
Let $k$ be a finite field of characteristic $p>3$, let $E/k$
be an elliptic curve, and, after a finite extension of $k$,
choose a nonzero point $T\in E[3](k)$. On
\[
       X=E^2\times E^2\times E^2,
       \qquad(a,b;u,v)\in X,
\]
put
\[
 \begin{aligned}
 R(P,Q)&=(Q,-P-Q),& S(P,Q)&=(Q,P),\\
 r(a,b;u,v)&=(a+T,b;Ru,Rv),&
 s(a,b;u,v)&=(-a,-b;Su,Sv).
 \end{aligned}                                                   \tag{A20.1}
\]
These generate a faithful action of $G=S_3$, with
$r^3=s^2=1$ and $srs=r^{-1}$. Its linear tangent
representation is $2\varepsilon\oplus2\rho$, where
$\varepsilon$ is the sign representation and $\rho$
the rational two-dimensional standard representation.
The order-three elements act freely. The loci fixed by
the three involutions are pairwise disjoint, each a
union of sixteen translates of an abelian surface.

Choose the product very ample bundle
$M=\boxtimes_{i=1}^6\mathcal O_E(3[0])$ on $X$ and put
\[
               H=\left(\bigotimes_{g\in G}g^*M\right)^{\otimes2}.
 \tag{A20.2}
\]
It has a natural $G$-linearization, and every stabilizer
acts trivially on its fibres. Allow either $Z=\varnothing$
or the following prescribed centre. Choose a point
$c\in E^2\times E^2$ with a free orbit for the linear
$S_3$-action on the last four factors, and put
\[
                  Z=\coprod_{g\in G}Z_g,
             \qquad Z_g=E^2\times\{gc\}.
 \tag{A20.3}
\]
These six surfaces are pairwise disjoint and avoid all
fixed loci. All these choices are defined after a finite
extension of $k$.

For every sufficiently large $t$, let $L=H^t$ and
\[
        \mathcal E_t=\operatorname{Hom}_G
                  (\rho,H^0(X,I_ZL)).
 \tag{A20.4}
\]
There is a nonempty open $\Omega\subset\mathcal E_t$
whose parameters give smooth fourfolds $B\subset X$,
the common zero loci of the two component sections.
Their original two-dimensional systems are genuine
Lefschetz pencils: every singular member has exactly
one ordinary quadratic singularity, outside the smooth
axis, and the generic member is smooth.
The group acts faithfully on their parameter line.

Let $h=c_1(L)$, let $F$ be one component of $X^s$, and set
\[
                 d_6=\int_Xh^6,\qquad d_2=\int_Fh^2.
\]
For the middle variable cohomology $V$, the multiplicities
of $\mathbf1,\varepsilon,\rho$ are exactly
\[
 \begin{split}
 m_1&=\frac{5d_6}{6}-8d_2-12,\\
 m_\varepsilon&=\frac{5d_6}{6}+8d_2-39,\\
 m_\rho&=\frac{5d_6}{3}-57.
 \end{split}                                                     \tag{A20.5}
\]
Their parities are respectively even, odd, odd, and
all exceed eight for large $t$.
With the six centres the positive known primitive span
$K$ is the regular rational representation of $G$;
the residual multiplicity ranks have parities odd,
even, odd and still exceed eight for large $t$.

The family also admits, in its parameter-space boundary,
a member with exactly one free orbit of six ordinary
nodes, and a member with exactly one involution-fixed
orbit of three ordinary nodes. All nodes can be chosen
outside $Z$, with no other singularities. Each admits
a transverse one-parameter smoothing within $\mathcal E_t$
with regular total space at every node and geometric
generic parameter in $\Omega$. The stabilizer of a node
in the orbit of three acts by its nontrivial character
on that node's vanishing-cycle line.

Finally the action, the linearized bundle, the prescribed
centres if present, and each of these smooth axes admit
a common mixed-characteristic lift, carrying their
actual $G$-action and their two-dimensional generator
representation. No algebraic class in the remaining
variable cohomology is asserted to lift.
\end{theorem}

\begin{proof}
\emph{The affine action and its fixed loci.}
The matrices $R,S$ satisfy $R^3=S^2=1$ and $SRS=R^{-1}$.
The translations in (A20.1) satisfy the same relation,
since $3T=0$ and $s$ negates the first coordinate.
An $r$-fixed point would have $a+T=a$, which is impossible.
The same is true of $r^2$.
The $s$-fixed locus is
\[
          E[2]^2\times\{(P,P):P\in E\}
                         \times\{(Q,Q):Q\in E\}.
 \tag{A20.6}
\]
It has sixteen components, all abelian surfaces.
The other two involutions are conjugate to $s$ and
have the same description up to an automorphism.
A point fixed by two distinct involutions would be
fixed by an element of order three. Hence their fixed
loci are disjoint, and these are all points with
nontrivial stabilizer. At every involution-fixed point
the tangent eigen-dimensions are $(2,4)$ for $(+1,-1)$.
Translation does not change tangent representations.

The norm tensor in (A20.2) has the standard permutation
linearization. Squaring makes its fibre character at
each order-two stabilizer trivial. There are no other
nontrivial stabilizers. It is very ample: all factors
are very ample, hence globally generated, and their
tensor product is very ample. A point $c$ outside the
finitely many proper fixed loci in $E^4$ has the required
free orbit. The surfaces (A20.3) are distinct parallel
fibres, hence disjoint. No point of $Z_g$ can be fixed
by an involution, because $gc$ is not fixed by it.

\emph{Compatible jets in the parameter space.}
For large $t$, sections of $I_ZH^t$ have arbitrary
bounded jets at any fixed number of points outside $Z$,
and generate the normal first jets along $Z$.
The bounds can be chosen uniformly in the points.
One elementary construction first chooses $c_0$ with
$I_ZH^{c_0}$ globally generated. A section of this
sheaf can be chosen nonzero at any prescribed bounded
set outside $Z$. Multiply it by products of hyperplanes
to interpolate the required values and derivatives;
the degree needed is bounded by the number and orders
of the jets. For normal jets along $Z$, use the
generators of $(I_Z/I_Z^2)H^{c_0}$, and the same
hyperplane interpolation. Increasing $c_0$ gives
these generators by Serre vanishing. Extra hyperplane
factors vanishing to order three at an already fixed
orbit preserve its second jets while providing all
the stated jets at another orbit.

Apply this construction to whole $G$-orbits, of length
at most six, and to two such orbits when needed.
Specify compatible translated jets, then average a
lift in $H^0(I_ZH^t)\otimes\rho^*$ by $|G|^{-1}\sum g$.
Because $p>3$, averaging is exact. Thus evaluation
from (A20.4) surjects onto every compatible jet target.
At a free point there are no restrictions on the two
component jets. At a point fixed by $s$, use an
$s$-eigenbasis of $\rho$ and local tangent coordinates
with two plus and four minus directions. The allowed
jets have the form
\[
 \begin{array}{c|c|c}
       &\text{value}&\text{first derivative}\\ \hline
 f_+&\text{arbitrary}&\text{arbitrary in the two plus directions}\\
 f_-&0&\text{arbitrary in the four minus directions}.
 \end{array}                                                    \tag{A20.7}
\]
Second derivatives have the corresponding even or odd
parity. These are exactly the compatible targets;
they are not extra assumptions on a generic section.

\emph{Smoothness of the axis and its fixed curves.}
At a free point off $Z$, singularity of the common
zero locus imposes two value conditions and rank at
most one on a $2$ by $6$ derivative matrix. The latter
has codimension five, so the total codimension is seven,
greater than the six-dimensional point space.
At an involution-fixed point, the common zero condition
is $f_+=0$. The two derivative rows belong to different
eigenspaces, so rank drops precisely if one row is zero.
The two bad conditions have codimensions $1+2=3$ and
$1+4=5$, both greater than the fixed-locus dimension two.
Along a prescribed centre, both values and tangential
derivatives vanish. The two normal rows are arbitrary
in a four-dimensional conormal fibre; rank at most one
has codimension three, greater than the centre dimension
two. These incidence counts show that the common zero
locus of a general parameter is smooth everywhere.
They also show that on each component $F$ of $X^s$
the fixed locus $B^s$ is the smooth divisor $(f_+|F=0)$
of class $h|F$. Thus $B^s$ is a union of sixteen smooth
ample curves. The same holds for the other involutions.

\emph{Why the original pencil is genuinely Lefschetz.}
The standard representation has faithful projective
$S_3$-action. There are only eight geometric points of
$\mathbf P(\rho)$ with nontrivial stabilizer: the two
order-three eigenlines and the six involution eigenlines.
For each of these finitely many lines $u$, a general
member $D_u=(\varphi(u)=0)$ is smooth.
At a free point this follows from arbitrary first jets:
a singularity imposes seven conditions while the point
varies in dimension six. At a point fixed by an
involution, an invariant eigenmember has arbitrary
value and two plus derivatives, hence three conditions
over a surface. An anti-invariant eigenmember vanishes
there identically, but has four arbitrary minus
derivatives, hence four conditions over that surface.
For any other line the two components of $u$ in (A20.7)
are nonzero, and its value and all six derivatives
are arbitrary. Along $Z$ independence of the two
normal rows makes every member smooth. Thus all eight
exceptional members can simultaneously be required
smooth on a nonempty open of $\mathcal E_t$.

Now vary a nonexceptional $u$. A member singular at a
free point imposes seven conditions; a degenerate
quadratic part imposes one more. The pair $(x,u)$ has
dimension seven, so a general parameter has no such
degenerate singularity. A singularity at a point with
involution stabilizer imposes seven conditions by
(A20.7), while $(x,u)$ has dimension three, and is absent.
Two singular points in distinct free $G$-orbits impose
fourteen independent first-jet conditions. Their points
and the parameter $u$ vary in dimension thirteen, so
a general pencil avoids them. If the second point is
$gx$ in the same orbit, singularity there is equivalent
to singularity of $\varphi(g^{-1}u)$ at $x$. Since $u$
has trivial projective stabilizer, these are two
independent components of $\varphi$, so again fourteen
conditions, now over a seven-dimensional parameter
space $(x,u)$. This case is absent as well.

The smooth axis excludes a singular point on the base
locus. The nondegenerate Hessian and the nonzero
parameter derivative at each remaining singularity
make the critical pairs isolated and reduced. They
form a finite scheme, because $X\times\mathbf P^1$ is
proper. The generic pencil member is therefore smooth,
and each singular member has just one ordinary node.
There are in fact singular members: the universal
first jet is a section of
$J^1(L)\boxtimes\mathcal O_{\mathbf P^1}(1)$ of rank
seven on $X\times\mathbf P^1$. Since $\Omega_X^1$ is
trivial, its top Chern number is
\[
            \int_{X\times\mathbf P^1}(h+\eta)^7
                        =7\int_Xh^6>0.
 \tag{A20.8}
\]
The transverse zeros just described have this length.
All required good conditions are open and their bad
incidences have the smaller dimensions verified above.
Their intersection is the nonempty open $\Omega$.

\emph{The two boundary degenerations.}
At a chosen free point off $Z$, prescribe the two
values to vanish, their derivative matrix to have
rank one, and the remaining equation after eliminating
the first to have a nondegenerate quadratic part in
five variables. The compatible orbit jets prescribe
exactly six ordinary nodes. At a chosen point $x$ fixed
by $s$, outside $Z$, prescribe
\[
      f_+(x)=0,\qquad df_+(x)=0,\qquad df_-(x)\ne0,
 \tag{A20.9}
\]
and a nondegenerate quadratic part for $f_+$ on
$\ker df_-$, with two plus and three minus variables.
This is an allowed even quadratic form. It gives one
ordinary node at $x$ and at each of its two translates,
and no additional prescribed singularities.

In either case the compatible second-jet evaluation
is surjective. Keeping these jets fixed, the remaining
point and normal-jet incidences have exactly the same
strict dimension inequalities as above away from the
chosen orbit. A general choice is consequently smooth
away from that orbit, including along $Z$. Choose an
equivariant deformation changing the surviving nodal
equation's constant term; at a fixed node this is the
allowed $f_+$ value. Its total space is regular at
every node. A general such one-parameter direction has
its generic parameter in the dense open $\Omega$.

For completeness the stabilizer character can be
calculated algebraically. Near the fixed node eliminate
the odd equation using its nonzero odd derivative.
The remaining smooth five-dimensional total space has
an invariant Morse function with tangent eigen-dimensions
$(2,3)$. An equivariant etale Morse change of coordinates
reduces it to an invariant nondegenerate quadratic form.
Here is a direct construction in odd characteristic:
choose eigen-coordinates by averaging; write the
function as $x^tA(x)x$ with symmetric $A$ and average
so $A(sx)=sA(x)s$. With $A_0=A(0)$, the matrix
$A_0^{-1}A(x)$ has a unique square root near the identity
on an etale neighbourhood, since the derivative of
$P\mapsto P^2$ there is multiplication by two.
Uniqueness makes this root $A_0$-self-adjoint and
$s$-equivariant. The change $y=P(x)x$ is etale and
puts the function into the form $y^tA_0y$.

Compactify the smooth affine quadratic fibre by one
extra coordinate. Its dimension is four and $s$ acts
on the six quadratic coordinates with determinant
$(-1)^3=-1$. It therefore negates the volume element
in the centre of the even Clifford algebra and
exchanges its two primitive idempotents. These label
the two families of maximal generators of the quadric.
Their cohomology-class difference generates primitive
middle cohomology and is negated by $s$. The affine
quadric calculation identifies this line with the
local vanishing line, so its stabilizer character is
nontrivial. Exact primary locators are SGA 7 II,
XII, 1.4--1.5, pp.~65--66; Proposition 1.12, p.~69;
Theorem 3.3, pp.~75--77; 3.5--3.7, pp.~78--81;
and XV, 2.2.3--2.2.5, pp.~182--185. XV, Theorem 3.4,
pp.~195--196 supplies the global nonzero norm and
surjective quadratic inertia for the regular smoothing.
Thus the two orbit representations of vanishing lines
are respectively $\mathbf Q[G]$ and
$\operatorname{Ind}_{\langle s\rangle}^G\varepsilon|_{\langle s\rangle}$.

\emph{Exact equivariant Euler calculation.}
The tangent bundle of $X$ is trivial and
$N_{B/X}=L|B\oplus L|B$. Consequently
\[
                  \chi(B)=\int_Bc_4(T_B)=5d_6.
 \tag{A20.10}
\]
For a finite-order automorphism of order prime to $p$,
its alternating cohomological trace equals the Euler
characteristic of its fixed locus. This is
Deligne--Lusztig, \emph{Annals of Mathematics} 103
(1976), Theorem 3.2, printed p.~119, in the tame case.
Since $r$ is free, its trace is zero. Each of the
sixteen fixed curves of $s$ is a smooth ample divisor
on an abelian surface, so has Euler characteristic
$-\int_Fh^2=-d_2$. All the components $F$ are translates
and have the same intersection number. Thus the
equivariant Euler character on the three conjugacy
classes $(1,s,r)$ is
\[
                            (5d_6,-16d_2,0).
 \tag{A20.11}
\]
Its multiplicities are
$5d_6/6-8d_2$, $5d_6/6+8d_2$, $5d_6/3$.

Translations act trivially on $H^1(X)$, whose rational
$G$-representation is $4\varepsilon\oplus4\rho$.
The exterior-power generating characters are
\[
 (1+z)^{12},\qquad(1+z)^4(1-z)^8,\qquad(1+z^3)^4.
 \tag{A20.12}
\]
The degree-four coefficients are $(495,-17,0)$;
their irreducible multiplicities are $(74,91,165)$.
The alternating sum of the coefficients in degrees
zero through three is $(-165,-5,-3)$, whose
multiplicities are $(-31,-26,-54)$.
Weak Lefschetz identifies the lower cohomology of
$B$ with that of $X$, and Poincare duality identifies
the upper cohomology with their duals. All the rational
$S_3$ representations are self-dual. Subtracting twice
this lower alternating character and the ambient
$H^4(X)$ character from (A20.11) gives (A20.5).
The sum of dimensions is $5d_6-165$, as in A03b.

Riemann--Roch on the abelian sixfold gives
$d_6=720\chi(L)$. Hence $5d_6/6$ and $5d_6/3$ are
even integers, and $8d_2$ is even. This proves the
claimed parities. As $L=H^t$, the leading terms are
positive multiples of $t^6$, while $d_2$ grows only
as $t^2$. Each multiplicity therefore exceeds eight
for all sufficiently large $t$.

\emph{The optional six positive classes.}
Put $h_0=c_1(H)$, let $A_0=(L_{h_0}^2)^{-1}[Z_1]_X$
be the common rational algebraic ambient correction,
and set
\[
       w_g=[Z_g]_B-t^{-2}\kappa^*A_0,\qquad
       d=\int_{Z_1}h_0^2>0,\qquad
       c_0=\int_XA_0^2h_0^2.
 \tag{A20.13}
\]
All $Z_g$ have the same ambient cohomology class. Their
normal bundles in $X$ are trivial of rank four, so the
normal exact sequence gives
$c_2(N_{Z_g/B})=3t^2h_0^2|Z_g$. The disjoint-centre
Gram calculation of A06 becomes the exact matrix
\[
                      3t^2d\,I_6-t^{-2}c_0J_6.
 \tag{A20.14}
\]
It is positive definite for large $t$; for example
$3t^4d>6|c_0|$ suffices. Thus the six classes are
independent. The common ambient correction is
$G$-invariant, so $G$ permutes them regularly.
Removing $K$ subtracts multiplicities $(1,1,2)$
from (A20.5), proving the residual parity assertion.

\emph{The actual equivariant lift.}
Lift the elliptic curve over a mixed-characteristic
henselian trait. Since $p>3$, its finite etale
three-torsion lifts $T$ after a finite extension;
the integer matrices in (A20.1) then give the same
affine action on the abelian scheme. The specified
product bundle $M$ extends using the zero section of
the elliptic scheme. Apply the relative norm and square
to that actual extension; its reduction is exactly
the prescribed $H$ in (A20.2), including its linearization.
When centres are present, lift
$c$ by smoothness of the abelian scheme; its free
orbit and the six disjoint fibre centres remain so
over the trait. Properness excludes a new collision
whose closure would meet the special fibre.
For large $t$, relative Serre vanishing for the
centre ideal gives reduction-surjectivity of
$H^0(I_ZH^t)$. With no centres this also follows
from the abelian line-bundle index theorem.
Average an arbitrary lift of the map (A20.4) by
$(1/6)\sum_{g\in G}g\varphi g^{-1}$. It is an actual
equivariant lift of the two generators. Smoothness
of the special complete intersection and properness
then give a smooth lift exactly as in A19b.

For a fully explicit ambient model one may use
$E:y^2=x^3-x$, over a number field containing its
three-torsion and then localize or complete at a
prime above any $p>3$. At primes $p\equiv3\pmod4$
the ambient $E^6$ is supersingular. The argument above
works over all $p>3$, after the indicated finite
field and trait extensions, without relying on
supersingularity or on an extension of any unknown
special algebraic cycle.
\end{proof}

\paragraph{Role in the programme.}
A20 supplies the geometric hypotheses, ranks and
two reflection configurations for A21; it does not
assert their monodromy conclusion without that proof.
Together with A19b and A21, the odd sign and standard
multiplicity modules can account for three residual
potential Tate directions. With the optional centres
there is already a nonzero positive regular span of
rank six, but the residual multiplicity is still
three; A15 alone cannot settle that residual space.
These are genuine original Lefschetz pencils in
the same ample system, so their axes apply to
cote 44, Theorem 5, scan 004 and its primitive
qualification in scan 008. No assertion of the
universal fixed-axis sign or of algebraicity of
the three remaining potential directions is made.

% END work/scan_w/affine_s3_axis_geometry.tex

% BEGIN work/axis/s3_multiplicity_monodromy.tex
% C44_POS /root; GPT-6 Astra Pro; effort not exposed.
% A19, A20 and this entire proof have passed independent whole-proof readings.
\section{Three residual Tate directions separated by an actual symmetric-group action}
\label{sec:c44-a21}

The affine action and its smooth-pencil existence assertions in the following
statement are the explicit geometric construction C44-A20. Its checked
equivariant jet argument supplies the geometric input for the moment calculation.

\begin{theorem}[C44-A21: monodromy on the three rational multiplicity spaces]
Let $k$ be a finite field of characteristic $p>3$ and let $G=S_3$ act on
$X=E^6$ by the affine action of C44-A20.  Its action on the tangent space has
representation $2\varepsilon\oplus2\mathrm{Std}$, where $\varepsilon$ is the
sign representation and $\mathrm{Std}$ the rational two-dimensional standard
representation.  An element of order three acts freely; the loci fixed by the
three involutions are disjoint unions of abelian surfaces.  Fix a sufficiently
large power $L$ of the $G$-linearized ample bundle in that construction.

Allow either no prescribed centre, or the single free $G$-orbit of six
abelian-surface fibres $Z_g$ specified there.  Write $Z$ for their union in
the second case.  Let
\[
 \mathcal E=H^0(X,I_Z\otimes L\otimes\mathrm{Std})^G,
 \qquad
 \Omega\subset\mathcal E                                       \tag{A21.1}
\]
be the nonempty smooth open of C44-A20: a parameter $a$ gives a smooth axis
$B_a$ cut out by its two component sections, and their original pencil on
$X$ is a genuine Lefschetz pencil.  In the case without centres, set $K=0$.
With the six centres, let $K$ be their primitive projected cycle span.  Its
signed Gram matrix is positive definite by C44-A06; it is a constant copy
of the regular rational representation of $G$.

Let $V$ be the middle Gysin kernel, with its pairing $q(x,y)=\int_{B_a}xy$
and twist $(2)$, and put $W=K^\perp\cap V$.  Over every coefficient field
$\mathbf Q_\ell$, $\ell\ne p$, decompose
\[
 W=(\mathbf1\otimes M_1)\perp(\varepsilon\otimes M_\varepsilon)
                      \perp(\mathrm{Std}\otimes M_{\mathrm{Std}}).
                                                                  \tag{A21.2}
\]
The multiplicity sheaves are nondegenerate orthogonal lisse sheaves.
For $L$ sufficiently large their ranks exceed eight, and each has full
orthogonal geometric monodromy.  This is an unconditional assertion; it
uses no Tate conjecture, cycle lifting, or 032 input.
\end{theorem}

\begin{proof}
We give the moment argument directly over the finite field.  All schemes,
centres and group elements are defined over $k$, after the finite extension
allowed in C44-A20.  For the calculation, remove the twist $(2)$ and give
the multiplicity sheaves their natural weight four.  We write $Q$ for the
cardinality of a finite extension $k'/k$.

The splitting (A21.2) and its forms are actual lisse constructions.
The finite group algebra is semisimple over every $\mathbf Q_\ell$,
including $\ell=2,3$.  Its three rational irreducible representations are
absolutely irreducible and self-dual with symmetric forms: use norm one
on the one-dimensional representations and the usual dot product on
$\{(x_1,x_2,x_3):\sum x_i=0\}$ for the standard representation.
The $G$-invariant nondegenerate pairing on $W$ is orthogonal between
distinct isotypic summands, and on each summand it is the chosen form
on $\rho$ tensor a symmetric form on $\operatorname{Hom}_G(\rho,W)$.
Schur's lemma and nondegeneracy show that the latter form is itself
nondegenerate.  These operations commute with the family and Frobenius.
The fixed positive Gram projector onto $K$, when present, also commutes
with $G$ and the family.  This justifies all the orthogonal multiplicity
sheaves used below.

\paragraph{The free quotient and four-point evaluation.}
Remove from $X$ the loci fixed by nonidentity elements and also $Z$, if
present, and call the resulting smooth open $U$.  The order-three elements
have no fixed points, so the only other removed sets are surfaces.  The
quotient
\[
                            Y=U/G                               \tag{A21.3}
\]
is smooth and geometrically connected of dimension six, and $U\to Y$ is
a geometrically connected finite etale $G$-torsor.  The equivariant vector
bundle $L\otimes\mathrm{Std}$ descends to a rank-two vector bundle on $Y$.
A point $y\in Y(k')$ has a Frobenius conjugacy class $h_y\subset G$.

For every set of at most four distinct points of $Y(k')$, evaluation of
$\mathcal E\otimes k'$ onto the corresponding four descended vector-bundle
fibres is surjective.  To see this uniformly, their inverse images in $U$
are at most twenty-four distinct geometric points.  For $L$ sufficiently
large, sections of $I_ZL$ evaluate surjectively at every such collection.
One may first twist until $I_Z$ is globally generated and then multiply
by products of hyperplanes separating at most twenty-four points, as in
C44-A18.  The bound is independent of the positions of these points.
Tensor by $\mathrm{Std}$ and average over $G$.  Since $p>3$, taking
$G$-invariants is exact, and evaluation remains surjective onto the
invariants of the orbit-fibre target.  These invariants are precisely the
descended fibres.  The map is defined over $k'$; surjectivity after
algebraic closure implies surjectivity over $k'$.

For a uniformly random $a\in\mathcal E(k')$, let $I_y(a)$ be one if the
section vanishes at $y$, and zero otherwise.  Thus $I_y$ is Bernoulli with
parameter $v=Q^{-2}$, and any collection of at most four distinct such
variables is independent.  This remains true when the fibre of the
$G$-torsor at $y$ is nonsplit: the descended vector-bundle fibre is still
a two-dimensional $k'$-vector space.

\paragraph{The character-weighted trace function.}
Let $\rho$ be any of the three rational irreducible representations
$\mathbf1,\varepsilon,\mathrm{Std}$, with character $\chi_\rho$ and dimension
$d_\rho=1,1,2$, respectively.  Define, for every section, including singular
ones,
\[
 T_\rho(a,k')=\sum_{y\in Y(k')}\chi_\rho(h_y)(I_y(a)-Q^{-2}).
                                                                  \tag{A21.4}
\]
For smooth parameters in $\Omega$ this has the uniform approximation
\[
 T_\rho(a,k')=\operatorname{Tr}(F_{k',a}\mid M_\rho^{\mathrm{un}})
                                  +O(Q^{3/2}).                   \tag{A21.5}
\]
Here and below the constants may depend on the fixed large $L$, but not
on $k'$ or on the parameter in the smooth open.

We verify this approximation with the ambient contribution included.
For a $G$-stable scheme $T$, put $N_g(T)=\#\operatorname{Fix}(gF_{k'};T)$.
On a free $G$-orbit stable under Frobenius, the character projector gives
\[
 \frac1{6}\sum_{g\in G}\chi_\rho(g^{-1})N_g(T)
       =\sum_{y\in(T/G)(k')}\chi_\rho(h_y).                       \tag{A21.6}
\]
Indeed identify an orbit with $G$: Frobenius acts by a right translation,
and the sum over all fixed pairs $(g,x)$ is six times the character of
that translation.  All the characters here are real and invariant under
inversion, so the choice of inverse Frobenius convention gives the same
identity.  Applying (A21.6) to the zero locus in $U$ and to $U$ identifies
(A21.4) with the character average of
$N_g(B_a\cap U)-Q^{-2}N_g(U)$.

Weak Lefschetz for the two ample sections, its dual Gysin assertion, and
hard Lefschetz on $X$ give, for each $g$,
\[
 N_g(B_a)-Q^{-2}N_g(X)
       =\operatorname{Tr}(gF_{k',a}\mid V^{\mathrm{un}})
                                      +O(Q^{3/2}).               \tag{A21.7}
\]
To check the degrees: $H^i(B_a)$ for $i>4$ identifies by Gysin with
$H^{i+4}(X)(2)$, so its contribution cancels against the corresponding
term of $Q^{-2}N_g(X)$.  In degree four, the ambient image of $H^4(X)$
has the same Frobenius and $G$-action as $H^8(X)(2)$ by $L_h^2$; its
removal leaves $V^{\mathrm{un}}$.  All remaining terms have absolute
value $O(Q^{3/2})$.  The finite-order action of $g$ does not affect this
bound, and the smooth family has fixed cohomological ranks.  Equivalently
one may apply the ordinary trace formula to the finite-field twist by
$g$, which is smooth projective.

Each removed involution-fixed surface meets $B_a$ in a smooth curve in
the chosen open.  Its contribution to $N_g(B_a)$ is $O(Q)$, whereas that
to $Q^{-2}N_g(X)$ is $O(1)$.  All such estimates are uniform.  Their
removal therefore fits within the error in (A21.7).

If the six centres are present, their quotient is one abelian surface
$Z_0$, and its torsor is split because all six components are defined
individually over $k$.  Their subtraction from the character average is
\[
 d_\rho(1-Q^{-2})\#Z_0(k')
                         =d_\rho Q^2+O(Q^{3/2}).                 \tag{A21.8}
\]
This is exactly the leading trace of the $d_\rho$ constant weight-four
lines in the $\rho$-multiplicity of the regular representation $K$.
Subtracting them leaves $M_\rho$.  With no centres this step is empty.
Finally, averaging the right side of (A21.7) against $\chi_\rho/6$
extracts the multiplicity trace, with no extra factor $d_\rho$:
character orthogonality gives
$\frac16\sum_g\chi_\rho(g^{-1})\operatorname{Tr}(gF\mid V)
=\operatorname{Tr}(F\mid\operatorname{Hom}_G(\rho,V))$.
This proves (A21.5).

\paragraph{An exact fourth moment for each multiplicity sheaf.}
Put $w_y=\chi_\rho(h_y)$ and $S_j=\sum_yw_y^j$.  The finite-etale
character sheaf on $Y$ and the compact-support trace formula imply
\[
                        S_2=Q^6+O(Q^{11/2}),\qquad S_4=O(Q^6).
                                                                  \tag{A21.9}
\]
Here the leading coefficient of $S_2$ is
$\dim(\rho\otimes\rho)^G=1$; the connected geometric torsor has full
geometric group $G$, and every representation is absolutely irreducible
and orthogonal.  All other compact-support cohomology has weight at most
eleven.  The second bound also follows immediately from boundedness of
the three integer characters and $\#Y(k')=O(Q^6)$.  This is the usual
finite-torsor form of the Weil-II weight estimate, not a supposition of
equidistribution without an error term.

Four-wise independence, applied to (A21.4), gives the exact identity
\[
 \begin{split}
 \frac1{\#\mathcal E(k')}\sum_a T_\rho(a,k')^4
     &=3v^2(1-v)^2(S_2^2-S_4)+\mu_4S_4,\\
 \mu_4&=v(1-v)^4+(1-v)v^4,
 \qquad v=Q^{-2}.
 \end{split}                                                     \tag{A21.10}
\]
The terms with a singleton index vanish; the remaining index patterns
are the two pairs and the quadruple.  The character values may be
negative, but they occur only through their squares or fourth powers in
this formula.  By (A21.9),
\[
 Q^{-8}\frac1{\#\mathcal E(k')}\sum_aT_\rho(a,k')^4
                                  =3+O(Q^{-1/2}).                 \tag{A21.11}
\]
For example, the fourth-diagonal contribution is $O(Q^{-4})$ after
normalization.

The geometric monodromy group of $M_\rho$ is reductive, by purity.
On the smooth open, (A21.5) differs from the true trace by $O(Q^{3/2})$,
while both functions are $O(Q^2)$.  Thus normalized fourth moments differ
by $o(1)$.  Katz's Lemmas 2.1.4--2.1.6, printed pp.~540--541, identify
the limsup on that open with the fourth monodromy moment and bound it
above by (A21.11), since adding the omitted nonnegative fourth powers
can only increase the sum.  Hence
\[
                  M_4(G_\rho,M_\rho)\leq3.                      \tag{A21.12}
\]
The three independent orthogonal contraction tensors give the reverse
inequality when the rank exceeds eight.  Therefore the moment is
exactly three.  No trace bound on singular fibres is required.

\paragraph{Reflections on the actual multiplicity spaces.}
The equivariant jet construction in C44-A20 gives a degeneration with
one free orbit of six ordinary double points, away from every fixed
locus and centre, and no other singularity.  A transverse one-parameter
smoothing has regular total space at all six nodes.  SGA~7~II,
Expose~XV, Theorem~3.4(i),(iii), pp.~195--196, supplies mutually
orthogonal vanishing cycles of square two, with one inertia element
acting as minus one on their span and as one on its orthogonal
complement.  The span is the regular representation of $G$.
Consequently this inertia acts as a reflection on the one-dimensional
vanishing part of $M_1$.

There is also a degeneration with one orbit of three nodes, each fixed
by an involution, and no other singularity.  At one node the even
component of the section has zero value and zero first derivatives
along the fixed surface, while the odd component has nonzero normal
derivative.  Eliminating the latter equation leaves an ordinary
quadratic singularity in five variables, with involution eigenspaces
of dimensions two and three.  The involution acts by its determinant,
namely minus one, on the local vanishing line, as established in the
local quadratic calculation of C44-A20.  Thus the three vanishing
cycles form
\[
                  \operatorname{Ind}_{C_2}^{S_3}\varepsilon
                          =\varepsilon\oplus\mathrm{Std}.       \tag{A21.13}
\]
The same regular smoothing and SGA~7 formula make inertia act as minus
one on this span.  It is a reflection on both $M_\varepsilon$ and
$M_{\mathrm{Std}}$.

These cycles lie in $W$, not merely in the full middle cohomology.
Every ambient class and every prescribed cycle is fixed by inertia.
The nonzero orthogonal node cycles are linearly independent; the
Picard--Lefschetz formula therefore forces each of them to be orthogonal
to the ambient image and to $K$.  The involution character and the
regular-orbit decompositions above hence apply inside $W$ itself.

Each multiplicity monodromy is orthogonal, reductive, of rank greater
than eight, has fourth moment three, and contains a reflection.
Katz's Theorem~1.4.6, printed p.~535, now gives the full orthogonal group
in all three cases.  The rank growth follows from C44-A20's equivariant
Euler calculation.  Passing to the smaller open imposing the original
Lefschetz-pencil condition does not change geometric monodromy.
\end{proof}

\begin{corollary}[C44-A21a: actual closed axes with three residual Tate directions]
For each sufficiently large polarization in C44-A20 there is a
Zariski-dense set of closed parameters in $\Omega$ for which the entire
rational algebraic axis space $A_B$ is positive definite.  Without the
six centres, the total root-of-unity multiplicity of normalized variable
Frobenius is exactly three.  With the six centres, their known positive
span has dimension six and the total root-of-unity multiplicity on its
orthogonal complement is exactly three; the full variable multiplicity
is nine.  In particular the latter examples have $\dim_{\mathbf Q}A_B\geq6$.
All these assertions are unconditional and concern genuine original
Lefschetz pencils over finite fields.
\end{corollary}

\begin{proof}
The three rational multiplicity Frobenius polynomials are independent of
$\ell$.  Indeed the traces of $gF^n$ are rational and independent of
$\ell$ by the finite-field twists used in C44-A19, and character averaging
followed by Newton identities gives the assertion.  The known regular
cycle summand, if present, is removed by its rational Gram projector.

Let $H$ be the joint geometric monodromy closure in the product of the
three orthogonal groups.  C44-A21 shows that each projection is onto;
therefore each projection of $H^0$ is onto the corresponding special
orthogonal group.  For every even-dimensional factor, the condition of
having no root of unity is detected by a nonzero regular polynomial on
that special orthogonal group; for every odd-dimensional factor, the
condition of having a unique simple root of unity is likewise detected
by the polynomial of C44-A16.  One common exponent suffices: if $D$ is
the largest of the three ranks, use
$M=\operatorname{lcm}(1,\ldots,2D^2)$, since a root of unity of degree at
most $D$ has order at most $2D^2$.

Each inverse image of this nonempty open is a nonempty open in the
connected group $H^0$.  Their intersection is nonempty.  The finite
quotient and constant-field Chebotarev argument of C44-A12 and A16 thus
selects a Zariski-dense set of closed parameters meeting all three
conditions simultaneously.  Independence of the three monodromy factors
is not assumed; surjectivity of the separate projections suffices.

The parity computation of C44-A20 is
\[
 \begin{array}{c|ccc}
 &\operatorname{rank}M_1&\operatorname{rank}M_\varepsilon
                                  &\operatorname{rank}M_{\mathrm{Std}}\\\hline
 Z=\varnothing&0&1&1\\
 \#\{Z_g\}=6&1&0&1
 \end{array}\pmod2.                                               \tag{A21.14}
\]
Thus the residual potential Tate multiplicity on $W$ is, respectively,
$0+1+2\cdot1=3$ or $1+0+2\cdot1=3$.

Choose an involution $s\in G$.  The four rational self-adjoint,
mutually orthogonal idempotents
\[
 e_1,\quad e_\varepsilon,\quad
       e_{\mathrm{Std}}\frac{1+s}{2},\quad
       e_{\mathrm{Std}}\frac{1-s}{2}                              \tag{A21.15}
\]
sum to one.  They are actual correspondences extending with the affine
action and the equivariant axis over the mixed-characteristic model of
C44-A20.  Their images on $W$ have Frobenius polynomials respectively
those of $M_1,M_\varepsilon,M_{\mathrm{Std}},M_{\mathrm{Std}}$.
Every such image has at most one residual potential Tate direction.
The known cycles split into the same orthogonal blocks, with positive
Gram matrices.  C44-A19 applies its common-projector determinant proof
to these four extending blocks, so their full algebraic subspaces are
positive.  Their orthogonal sum is all of $A_B$.

The six known classes, when imposed, are independent by their positive
definite Gram matrix, proving the last dimension assertion.  No
algebraicity is asserted for any of the three additional potential Tate
directions, and no Tate conjecture is needed for positivity of those
which are algebraic.
\end{proof}

\paragraph{Scope, sources and the main gate.}
These axes have three residual root-of-unity directions even after the
explicit known positive span is removed.  They therefore fall outside
the multiplicity-one sufficient condition C44-A15; the new reason for
their sign is the four actual extending rational projectors.  The
construction removes a nonempty further class of axes from the original
QH9-axis gate.  It does not prove that every axis has such projectors,
prove a general Tate conjecture, or exhibit a wrong-sign cycle.

The monodromy inputs are Katz, \emph{Larsen's alternative, moments, and
the monodromy of Lefschetz pencils} (2004), Theorem~1.4.6 p.~535 and
Lemmas~2.1.4--2.1.6 pp.~540--541, and SGA~7~II, Expose~XV,
Theorem~3.4 pp.~195--196.  The finite-torsor estimate (A21.9) is the
compact-support trace formula with Deligne's Weil-II weight theorem,
\emph{La conjecture de Weil II}, Theorem~3.3.1 and its compact-support
Corollary~3.3.4, printed p.~206, applied to the finite-monodromy weight-zero character
sheaf.  The torsor averaging and exact fourth-moment identities are
proved above rather than borrowed as an unstated equivariant theorem.

In Grothendieck's language this supplies the signed axis term of
cote~44, Th.~5, scan~004, for these closed special pencils, with the
primitive margin of scan~008 enforced.  The separate ambient abelian
term in H9-T7, and hence any assembled full-height conclusion using it,
remain conditional on the unrefereed 032/Milne 2026 preprints.  The
present axis sign itself is unconditional.  Its setting is the actual
relative correspondence algebra and SGA~7 monodromy, not an abstract
choice of a positive form on a numerical-motive summand.

% END work/axis/s3_multiplicity_monodromy.tex

% BEGIN work/axis/axis_realization_blocks.tex
% C44_POS /root/axis_independent. Extends the common-kernel test, not a liftability theorem.
\section{The crystalline test with several symmetry-separated Tate directions}
\label{sec:c44-a18d}

\begin{corollary}[C44-A18d: a common kernel for finite-group blocks]
Let $B\subset X/\mathbf F_q$ and $P$, $\mathscr C$ be as in
C44-A18b. Let a finite constant group $G$ act on this embedded
complete intersection, preserving the ample class. Suppose
$e_1,\ldots,e_b\in\mathbf Q[G]$ are self-adjoint pairwise
orthogonal idempotents summing to one, where the involution
on the group algebra sends $g$ to $g^{-1}$. Let $K$ be a
$G$-stable rational algebraic subspace of the variable
cohomology at one initial coefficient prime, and suppose
its rational intersection form is nondegenerate and each
$K_i=e_iK$ is nondegenerate. A positive definite $K$, including
$K=0$, satisfies this assumption.

Use actual cycles $P_*e_{i*}P_*C$ to represent each $K_i$,
choose a rational basis in every block, and extend the finite
field to define all these cycles. Suppose that in each
orthogonal residual block
\[
       W_{i,\ell}=e_iV_\ell\cap K_{i,\ell}^{\perp}
\]
the total normalized Frobenius root-of-unity multiplicity
is at most one. Then all rational crystalline and
$\ell$-adic cycle-class kernels coincide on $\mathscr C$.
The necessary filtration test of C44-A18c therefore holds
for a projected class along every specified embedded
polarized lift of $B\subset X$.

The finite-group action and its projectors need not extend
along the lift used for that necessary test. Their role in
the present assertion is on the special fibre. In particular
the assertion applies to the C44-A21 axes, including those
with residual potential Tate multiplicity three. It does
not assert that a projected class on any such axis is
nonliftable.
\end{corollary}

\begin{proof}
The group action commutes with Gysin, restriction and
$L_h^2$, and hence with its unique inverse in every
cohomological realization. Thus it commutes with the
cohomological action of $P$. All the stated idempotent
and adjoint identities hold in rational crystalline
cohomology as well as in every $\ell$-adic realization.
The actual projected basis cycles have the same rational
nondegenerate Gram matrices in all these realizations;
they supply common nondegenerate subspaces $K_i$ and
orthogonal residual spaces.

The Frobenius polynomials of the blocks are common rational
polynomials also in crystalline cohomology. Indeed, for
each finite-order $g\in G$ and each $n>0$, the finite-field
twist by $g$ is smooth projective and its Frobenius acts
as $gF^n$, with the inverse cocycle chosen if the Frobenius
convention requires it. Katz--Messing, Theorem~1(i),
printed p.~74, applies to the twist and identifies its
degreewise characteristic polynomials in crystalline and
$\ell$-adic cohomology. Subtract the ambient traces, then
take the rational linear combination specified by $e_i$.
Since $e_i$ commutes with Frobenius,
\[
  \operatorname{Tr}(F^n\mid e_iV)
       =\operatorname{Tr}(e_iF^n\mid V).
\]
Newton identities give the same rational block polynomial
in all realizations. Removing $(T-1)^{\dim K_i}$ preserves
this assertion. This is the finite-twist argument of
C44-A19b with Katz--Messing also included; no independence
conjecture for cycle classes has been used.

If $Z\in\mathscr C$ has zero class in one realization,
set $Z_i=P_*e_{i*}Z$. Its class in that realization is
zero, so all its pairings with the actual basis cycles
of $K_i$ and its self-pairing vanish in $\mathbf Q$.
In every other realization its class lies in $W_i$,
has finite normalized Frobenius orbit, and has square
zero. The simple potential-Tate-line argument of
C44-A18b forces that class to vanish. Summing over $i$
gives $\operatorname{cl}(Z)=0$ in every realization,
because $\sum_i P e_i$ acts as the identity on $V$.
This proves the common-kernel assertion.

The proof of C44-A18c uses only this common-kernel
assertion on the special projected cycle group, and
an algebraic Lefschetz inverse on the generic abelian
variety. It does not use an extension of the special
finite-group action. The same proof therefore gives
the stated filtration test along any embedded polarized
lift under consideration.
\end{proof}

\paragraph{Scope in Grothendieck's programme.}
For cote~44, Th.~5, scan~004 and the primitive instruction
on scans~008--009, this supplies a representative-independent
crystalline necessary condition on the actual axis cycle
classes in the new three-direction families. The sign of
those families is separately proved by C44-A19--A21.
Failure of the filtration condition for every possible
lift remains unproved, so route (a) of QH9-axis is not
closed by this corollary.

% END work/axis/axis_realization_blocks.tex

% BEGIN work/axis/axis_primitive_semiregularity.tex
% C44_POS /root/axis_independent. Coherent calculation only; arithmetic lifting remains open.
\section{Primitive coherent semiregularity of a prescribed smooth surface}
\label{sec:c44-a22}

\begin{proposition}[C44-A22: surjectivity from the coherent Gysin kernel]
Let $X$ be an abelian sixfold over an algebraically closed field
of arbitrary characteristic, let $H$ be ample, and let
$S\subset X$ be a fixed smooth projective surface. For every
sufficiently large integer $t$, and for every smooth complete
intersection $\kappa:B\hookrightarrow X$ of two members of
$|H^t|$ containing $S$, put $N=N_{S/B}$ and
\[
 V^{3,1}_{\mathrm{coh}}
  =\ker\left(\kappa_*:
          H^1(B,\Omega_B^3)\longrightarrow H^3(X,\Omega_X^5)\right).
 \tag{A22.1}
\]
Here $\kappa_*$ is the coherent Hodge-cohomology Gysin map.
Restriction followed by the natural conormal exterior quotient
induces a surjection
\[
 V^{3,1}_{\mathrm{coh}}
       \longrightarrow H^1(S,N^*\otimes K_S).
 \tag{A22.2}
\]
Consequently its dual is an injection
\[
 H^1(S,N)\lhook\joinrel\longrightarrow
                     (V^{3,1}_{\mathrm{coh}})^*.
 \tag{A22.3}
\]
This is the coherent primitive component of the usual dual
description of the semiregularity map. No degeneration of a
Hodge--de Rham spectral sequence or mixed-characteristic
obstruction theorem is assumed in this statement.
The bound on $t$ depends on $(X,H,S)$ and is uniform in
the smooth complete intersection containing $S$.
\end{proposition}

\begin{proof}
Write $L=H^t$, $I=I_{S/B}$, and $F=\Omega_B^3$.
Choose a nonzero invariant volume form on $X$; its only role
is to trivialize $K_X$ and identify $\Omega_X^5$ with $T_X$.
Adjunction gives $K_B=L^2|B$. The tangent sequence, tensored
with $K_B$, gives
\[
 0\longrightarrow F\longrightarrow (L^2|B)^{\oplus6}
              \longrightarrow (L^3|B)^{\oplus2}
              \longrightarrow0.
 \tag{A22.4}
\]
Its terms are locally free. Tensoring with $I$ is therefore
still exact.

\emph{The required vanishing is uniform in $B$.}
The twisted Koszul resolutions on $X$ are
\[
 \begin{split}
 0&\longrightarrow\mathcal O_X\longrightarrow L^{\oplus2}
                   \longrightarrow L^2
                   \longrightarrow\kappa_*(L^2|B)\longrightarrow0,\\
 0&\longrightarrow L\longrightarrow (L^2)^{\oplus2}
                   \longrightarrow L^3
                   \longrightarrow\kappa_*(L^3|B)\longrightarrow0.
 \end{split}
 \tag{A22.5}
\]
Ample line bundles on an abelian variety have no higher
coherent cohomology, in every characteristic. Thus
\[
 H^j(B,L^3|B)=0\ (j>0),\qquad
 H^j(B,L^2|B)\simeq H^{j+2}(X,\mathcal O_X)\ (j>0).
 \tag{A22.6}
\]
For $t$ sufficiently large, Serre vanishing on the fixed
$S$ and fixed ideal $I_{S/X}$ gives
\[
 \begin{gathered}
 H^j(S,L^a|S)=0\quad(j>0,\ a=2,3),\qquad
 H^1(X,I_{S/X}L^a)=0\quad(a=2,3),\\
 H^2\bigl(S,\Omega_S^1\otimes K_S^{-1}\otimes L^2|S\bigr)=0.
 \end{gathered}
 \tag{A22.7}
\]
In particular $H^0(B,L^a|B)\to H^0(S,L^a|S)$ is
surjective: already the sections from $H^0(X,L^a)$
surject onto the latter space. The restriction exact
sequences and (A22.6)--(A22.7) now give
\[
 \begin{split}
 H^1(B,IL^3)&=H^2(B,IL^3)=0,\\
 H^j(B,IL^2)&\xrightarrow{\sim}H^j(B,L^2|B)
                                  \quad(j=1,2).
 \end{split}
 \tag{A22.8}
\]
Here and below a line bundle next to $I$ means its
restriction to $B$.

\emph{Restriction of the differential forms is surjective.}
Apply cohomology to (A22.4) and to its tensor product
with $I$. By (A22.6) and (A22.8), the commutative map
\[
 H^2(B,IF)\longrightarrow H^2(B,F)
 \tag{A22.9}
\]
identifies with the isomorphism
$H^2(B,IL^2)^{\oplus6}\to H^2(B,L^2|B)^{\oplus6}$.
It is therefore an isomorphism. The restriction sequence
$0\to IF\to F\to F|S\to0$ implies that
\[
                  H^1(B,F)\longrightarrow H^1(S,F|S)
                     \quad\hbox{is surjective}.
 \tag{A22.10}
\]

The rank-two conormal sequence for $S\subset B$, on
taking third exterior powers, has the exact quotient
\[
 0\longrightarrow\det(N^*)\otimes\Omega_S^1
       \longrightarrow F|S\longrightarrow N^*\otimes K_S
       \longrightarrow0.
 \tag{A22.11}
\]
By adjunction its left term is
$L^2|S\otimes K_S^{-1}\otimes\Omega_S^1$.
The last vanishing in (A22.7) proves that
$H^1(S,F|S)\to H^1(S,N^*\otimes K_S)$ is surjective.
Together with (A22.10), this already proves surjectivity
from the full $H^1(B,F)$.

\emph{The ambient quotient can be killed without changing the target.}
Let
\[
  \theta:H^1(B,F)\longrightarrow H^1(B,L^2|B)^{\oplus6}
 \tag{A22.12}
\]
be the map in (A22.4). By (A22.8), the image of
$H^1(B,IF)\to H^1(B,F)$ maps onto the entire target of
$\theta$: in the corresponding ideal-tensored sequence
$H^1(B,IF)\to H^1(B,IL^2)^{\oplus6}$ is surjective,
and the right vertical map is an isomorphism. Every
class from $H^1(B,IF)$ restricts to zero on $S$.

For precision, the coherent Gysin map in (A22.1) is
$\theta$ followed by the double Koszul boundary
\[
 H^1(B,L^2|B)\otimes H^0(X,\Omega_X^5)
        \xrightarrow{\sim} H^3(X,\Omega_X^5).
 \tag{A22.13}
\]
The boundary is an isomorphism by (A22.5)--(A22.6).
This identification follows directly from the local
formula, which also fixes the compatibility needed
here. If $B$ is locally cut out by $f_0,f_1$, the
adjunction injection in (A22.4) sends a lifted
$3$-form $\widetilde\omega$ to
$\widetilde\omega\wedge df_0\wedge df_1$ in
$\Omega_X^5\otimes L^2|B$. The double boundary sends it to
the local-cohomology class represented by
\[
       \frac{\widetilde\omega\wedge df_0\wedge df_1}{f_0f_1}.
 \tag{A22.14}
\]
This is exactly the coherent Gysin construction:
Stacks Project, Section~50.23, Tag~0G82, Remark~50.23.1
and Remark~50.23.7. A simultaneous change of orientation
would multiply both maps by the same sign and would
not change their kernel.

Now take a target class in $H^1(S,N^*\otimes K_S)$,
and lift it to $v\in H^1(B,F)$ using (A22.10)--(A22.11).
Choose $u$ in the image of $H^1(B,IF)$ with
$\theta(u)=\theta(v)$, as is possible by (A22.12).
Then $v-u$ is in the Gysin kernel by (A22.13), and
has the same image on $S$ as $v$. This proves (A22.2).
Serre duality on the smooth surface identifies
$H^1(S,N^*\otimes K_S)^*$ with $H^1(S,N)$ and gives
(A22.3).
\end{proof}

\paragraph{What the proposition does and does not supply.}
The dual-of-restriction description used here is the algebraic
description of semiregularity; compare Dan--Kaur,
\emph{Semi-regular varieties and variational Hodge conjecture},
\emph{Comptes Rendus Math\'ematique} 354 (2016), \S2.1,
printed p.~298. The proof above supplies its own global
restriction argument and does not appeal to a
characteristic-zero variational theorem. The ample-vanishing
input is Milne, \emph{Abelian Varieties}, Chapter~I,
Theorem~11.1(c), printed p.~54 in the author notes.

This calculation applies to the smooth centres in C44-A14
and to the high-degree genuine axes containing them whose
existence is proved in C44-A06 and whose closed fibres are
selected in C44-A18. It shows that the coherent obstruction
map remains injective after its target is restricted to the
dual of the coherent Gysin kernel; discarding the ambient
quotient loses no coherent normal obstruction in this range.

An additional theorem would be needed to conclude that a
rational crystalline primitive class lying in
$\operatorname{Fil}^2$ along an arbitrary possibly ramified
mixed-characteristic lift forces an all-orders lift of the
given surface. Such a theorem has not been proved here.
In particular no passage from a rational filtered class to
an integral deformation of this particular support is being
made. C44-A18b--A18d repair the separate issue of changing
cycle representatives; they do not supply this obstruction
compatibility or a converse to the filtration test.

For cote~44, Th.~5, scan~004 and the primitive warning on
scans~008--009, the proposition isolates a precise remaining
deformation step in route (a) of QH9-axis. It neither proves
that the primitive class of the nonliftable surface is itself
nonliftable nor settles the universal axis sign.

% END work/axis/axis_primitive_semiregularity.tex

% BEGIN work/scan_w/higher_dimensional_s3_axes.tex
% C44_POS /root/scan_w; GPT-6 Astra Pro; reasoning effort not exposed.
% Dimensional extension of the independently checked A19--A21 mechanism.
\section{The symmetric-group construction in every higher even dimension}
\label{sec:c44-a23}

The following is a dimensional generalization of the checked construction
C44-A19--A21.  It is not a separately staged progress-clock event.  We give
the changed geometry, signs, moment powers and representation parities;
the common-projector determinant theorem remains exactly C44-A19.

\begin{theorem}[C44-A23: higher-dimensional symmetry-separated axis signs]
Let $r\geq2$, let $k$ be a finite field of characteristic $p>3$, and let
$E/k$ be an elliptic curve.  After a finite extension choose
$0\ne T\in E[3](k)$ and put
\[
 n=2r+2,\qquad a=2r-2,\qquad
 X=E_{\rm sign}^{\,a}\times E_{\rm std}^{\,2}\times E_{\rm std}^{\,2}.
 \tag{A23.1}
\]
On each standard pair use
$R(P,Q)=(Q,-P-Q)$ and $S(P,Q)=(Q,P)$.
Let the order-three generator $\rho$ translate the first sign coordinate
by $T$, act trivially on the other sign coordinates, and act by $R$ on
both standard pairs.  Let the involution $\sigma$ negate every sign
coordinate and act by $S$ on both standard pairs.  These define an
affine $G=S_3$ action.

Set
\[
 M=\mathop{\boxtimes}_{j=1}^{n}\mathcal O_E(3[0]),\qquad
 H=\left(\bigotimes_{g\in G}g^*M\right)^{\otimes2},\qquad L=H^t.
 \tag{A23.2}
\]
For all sufficiently large $t$, the equivariant standard section pairs
have a nonempty smooth parameter open whose members cut out smooth
axes $B\subset X$ of dimension $2r$, and whose original pencils are
genuine Lefschetz pencils.  One may either impose no centre, or impose
the following six centres.  Let $S_0=E^r$ be the first $r$ sign factors,
and choose a point $c$ in the remaining factors with free orbit under
the induced linear $G$ action.  After another finite extension put
\[
                       Z_g=S_0\times\{gc\}\quad(g\in G).
 \tag{A23.3}
\]
These are six disjoint abelian $r$-folds.  Their primitive projected
classes span a positive definite constant regular representation $K$;
without centres set $K=0$.

Write $V=\ker\kappa_*$ for the variable middle cohomology of the axis,
with its twist $(r)$ and signed pairing
$q_B(x,y)=(-1)^r\int_Bxy$, and put $W=V\cap K^\perp$.
The three rational multiplicity sheaves of $W$ for
$\mathbf1,\varepsilon,\mathrm{Std}$ have ranks greater than eight for
large $t$, and each has full orthogonal geometric monodromy.

Define bits $\alpha_r,\beta_r\in\{0,1\}$ by
\[
 (\alpha_r,\beta_r)=
 \begin{cases}
 (0,0),&r\text{ odd},\\
 \left(\binom{r-1}{r/2-2},\binom{r-1}{r/2-1}\right)\bmod2,
       &r\text{ even},
 \end{cases}
 \tag{A23.4}
\]
where an out-of-range binomial coefficient is zero.
There is a Zariski-dense set of closed parameters for which the entire
rational algebraic axis space $A_B$ is positive definite and the total
root-of-unity multiplicity of normalized Frobenius on $W$ is exactly
\[
 \tau(W)=
 \begin{cases}
 \alpha_r+3\beta_r,&K=0,\\
 2-\alpha_r+\beta_r,&\dim K=6.
 \end{cases}
 \tag{A23.5}
\]
With centres, the total multiplicity on $V$ is $6+\tau(W)$ and
$\dim_{\mathbf Q}A_B\geq6$.  These conclusions are unconditional.
They do not assert algebraicity of any residual potential Tate line.
\end{theorem}

\begin{proof}
We verify every dimensional change to C44-A20--A21.

\paragraph{The action, fixed loci, and actual equivariant lifts.}
The identities $R^3=S^2=1$ and $SRS=R^{-1}$ are unchanged.  Negation
conjugates translation by $T$ to translation by $-T$, so the affine
generators satisfy the same relations.  Both order-three elements act
freely.  The fixed locus of $\sigma$ has
\[
                 2^{2a}=2^{4r-4}
 \tag{A23.6}
\]
components, each an abelian surface: the sign coordinates are two-torsion
and both standard pairs lie on their diagonals.  Fixed loci of distinct
involutions are disjoint, since an intersection would contain a point
fixed by an order-three element.  At a fixed point the tangent
eigenspace dimensions are $2$ and $2r$.

The specified norm-square $H$ is very ample, has a natural
$G$ linearization, and has trivial stabilizer character on every fixed
surface.  The last assertion follows from squaring the order-two fibre
character.  The standard representation has one even and one odd
component at an involution.  Thus, at such a point, the even section
has arbitrary value and two arbitrary tangent derivatives, whereas
the odd section has zero value and $2r$ arbitrary normal derivatives.
The corresponding second jets have exactly the invariant parity.

There are at least $r$ sign factors, because $2r-2\geq r$.
The translation by $T$ is inside $S_0$, so the quotient by these factors
carries the stated linear action.  Its free locus is nonempty because
it contains the two standard factors.  Consequently the six fibres
in (A23.3) are disjoint, form a regular orbit, and avoid all involution
fixed loci.  They may all be defined over a finite extension.

Lift $E$ to a smooth elliptic scheme over a mixed-characteristic trait,
and lift its prime-to-$p$ three-torsion section after an allowed finite
extension.  The integer matrices and affine translation extend.  The
\emph{specified} product bundle $M$ extends by the zero section, and
the same norm and square extend $H$.  The point $c$ lifts by smoothness
after completion; its six distinct fibres remain disjoint.  For large
$t$, relative Serre vanishing and averaging by $1/6$ lift the actual
standard section pair, including the prescribed vanishing if the
centres are imposed.  Properness of the nonsmooth locus then gives a
smooth embedded polarized lift of every chosen smooth special axis.
This is precisely the lift needed for C44-A19; it is not a claim that
an arbitrary polarization on a supersingular product extends.

\paragraph{Smooth axes and genuine original pencils.}
The uniform compatible-jet argument of C44-A20 works in dimension $n$:
first globally generate the fixed ideal by an ample twist, then use
products of separating hyperplanes for a bounded number of points and
jets, and average.  Along the smooth prescribed centres, the same
Serre argument applies to their conormal jets.  The bounds depend on
the fixed data, but are uniform in points and orbit configurations.

At a free point the bad complete-intersection jet has two zero values
and a rank-at-most-one $2$-by-$n$ derivative matrix.  Its codimension
is $2+(n-1)=n+1>n$.  On a fixed surface, the bad even derivative has
codimension $1+2=3>2$, and the bad odd derivative, together with the
zero even value, has codimension $1+2r>2$.  Along each $r$-dimensional
centre, the two rows in its $(r+2)$-dimensional normal space fail to
have rank two in codimension $r+1>r$.  These estimates prove the
nonempty smooth-axis open and ensure that every pencil member is
smooth along the centres.

The faithful projective standard action on the pencil has eight
exceptional parameters: two order-three eigenlines and the six
involution eigenlines.  They can all be chosen to give smooth members.
At a free point a critical section has codimension $n+1>n$.
For an involution eigenmember the critical-point codimensions on its
fixed surface are $3$ for the even member and $2r$ for the odd member;
both exceed two.  A member with neither eigencharacter has all $n$
derivatives available, so has critical-point codimension $n+1$ there.
These arguments cover all eight members, even at fixed loci of the
other involutions.

At a nonexceptional parameter, a degenerate critical point imposes
$n+2$ conditions on the $(n+1)$-dimensional point--parameter space.
Two critical points in distinct free orbits impose $2(n+1)$ conditions
on a space of dimension $2n+1$.  If they lie in one free orbit, the
parameter and its translate are distinct; the two independent section
components then both have zero first jet at one point, again imposing
$2(n+1)$ conditions, now on dimension $n+1$.  A fixed-point critical
pair imposes $n+1$ conditions on dimension three.  Each of these bad
loci is avoided.  Smoothness of the base prevents a critical point on
the base; a nondegenerate Hessian and the nonzero pencil derivative
make the remaining critical scheme reduced and finite.  Properness
gives finiteness without any unaccounted boundary points.

This scheme is nonempty.  The first-jet bundle on $X$ has an extension
of $L$ by $L^{\oplus n}$, so the top Chern number for the universal
pencil on $X\times\mathbf P^1$ is
\[
 \int_{X\times\mathbf P^1}(h+\eta)^{n+1}
                =(n+1)\int_Xh^n>0.
 \tag{A23.7}
\]
Thus every singular member has exactly one ordinary double point,
and the original pencil is a genuine Lefschetz pencil.

\paragraph{The two node orbits and their signed norms.}
At a free orbit one may prescribe a rank-one complete-intersection
jet and a nondegenerate remaining quadratic form in $n-1=2r+1$
variables.  This gives exactly one orbit of six nodes.  At an
involution-fixed point, prescribe zero value and derivative for the
even component and a nonzero odd derivative.  Elimination of one odd
variable leaves an invariant quadratic form with eigenspace dimensions
\[
                              (2,2r-1).
 \tag{A23.8}
\]
A nondegenerate such form gives exactly one orbit of three nodes.
In both cases the same bounded-jet independence at other points and
the preceding strict codimension inequalities exclude every other
singularity.  Varying the remaining constant term gives a transverse
smoothing regular at all the nodes; its general member lies in the
smooth-axis open and may also be required to lie in the genuine-pencil
open.

The equivariant etale Morse and even-quadric calculation in C44-A20
applies in every odd number $2r+1$ of quadratic variables.  The
involution determinant in (A23.8) is $(-1)^{2r-1}=-1$; after
compactification it exchanges the two generator families of the
even-dimensional quadric, hence negates its primitive line.  The
three-node vanishing representation is therefore again
$\operatorname{Ind}_{C_2}^{S_3}\varepsilon=\varepsilon\oplus\mathrm{Std}$.
The free orbit gives the regular representation.  SGA~7 supplies
pairwise orthogonal node cycles with ordinary cup square
$2(-1)^r$, hence \emph{signed} square two for $q_B$.  A common quadratic
tame inertia acts by minus one on their span.  It gives reflections
on the trivial and sign multiplicities from the free orbit and on the
standard multiplicity from the three-node orbit.  Ambient classes and
prescribed centre classes are fixed by inertia, so these node cycles
lie in the residual $W$, by their nonzero norms and independence.

\paragraph{The prescribed positive regular span.}
Write $h_0=c_1(H)$ and
\[
 d=\int_{Z_1}h_0^r>0,\qquad
 A_0=(L_{h_0}^2)^{-1}[Z_1]_X,\qquad
 c_0=\int_XA_0^2h_0^2.
 \tag{A23.9}
\]
The inverse is represented by the absolute abelian Chow correspondence
of C44-A15.  All six parallel fibres have the same ambient class.
Their primitive cycles are
$w_g=[Z_g]_B-t^{-2}\kappa^*A_0$.
Their normal bundle in $X$ is trivial, and the two defining divisors
give
$c(N_{Z_g/B})=(1+th_0|Z_g)^{-2}$.
Thus their signed Gram matrix is
\[
        (r+1)t^r d\,I_6-(-1)^r t^{-2}c_0J_6.
 \tag{A23.10}
\]
It is positive definite, for example when
$(r+1)t^{r+2}d>6|c_0|$.  Hence the $w_g$ are independent and form a
constant regular $G$ representation in every realization.  This also
checks the sign change when $r$ is odd.

\paragraph{Fourth moments in the changed dimension.}
Remove the fixed surfaces and the centres from $X$ and take the free
quotient $Y$.  It is a smooth geometrically connected $n$-fold with
connected $G$-torsor.  Four quotient points still require at most
twenty-four geometric points upstairs.  Equivariant evaluation and
descent therefore give the same four-wise independent Bernoulli
variables of parameter $v=Q^{-2}$ as in C44-A21.

For each character $\chi$ define
$T_\chi=\sum_{y\in Y(\mathbf F_Q)}\chi(h_y)(I_y-Q^{-2})$.
The trace identity in C44-A21, now in degree $2r$, is
\[
 T_\chi=\operatorname{Tr}(F\mid M_\chi^{\rm un})
                                  +O(Q^{r-1/2}).
 \tag{A23.11}
\]
Indeed the terms above degree $2r$ cancel by Gysin against the ambient
terms shifted by four, the ambient middle image cancels by $L_h^2$,
and every remaining term has weight at most $2r-1$.  The removed fixed
surfaces meet $B$ in curves, costing only $O(Q)$.  A centre quotient is
one $r$-fold with split torsor, so its subtraction is
$\dim(\chi)Q^r+O(Q^{r-1/2})$, exactly the regular $K$ contribution.
The character average still has factor $1/6$ and no extra dimension
factor.

The finite-character local system gives
$S_2=Q^n+O(Q^{n-1/2})$ and $S_4=O(Q^n)$.
The exact Bernoulli fourth moment remains
\[
                 3v^2(1-v)^2(S_2^2-S_4)+\mu_4S_4.
 \tag{A23.12}
\]
After division by $Q^{4r}$ it tends to three; the diagonal contribution
is $O(Q^{-2r})$.  The trace error in (A23.11) tends to zero after
normalization by $Q^r$.  The nonnegative fourth powers on the singular
parameter locus give the same upper bound on the smooth open.  Purity,
the three orthogonal contraction tensors, the displayed reflections
and Katz's rank-greater-than-eight theorem yield full orthogonal
monodromy in each multiplicity space.  Their ranks grow like a positive
multiple of $t^n$, as the next calculation also shows.

\paragraph{An exact representation-ring parity calculation.}
The equivariant Euler characters of the smooth axis at the three
conjugacy classes $1,\sigma,\rho$ are
\[
 \left((2r+1)\int_Xh^n,
              -2^{4r-4}\int_Fh^2,0\right),
 \tag{A23.13}
\]
where $F$ is any fixed-surface component.  The first entry follows
from $c(T_B)=(1+h)^{-2}$; the second is the Euler characteristic of
the fixed smooth curves; the third is zero because $\rho$ acts freely.
The primary tame fixed-point trace theorem used in C44-A20 applies
unchanged.  All multiplicities of this Euler representation are even:
$\int_Xh^n=n!\chi(L)$, $n\geq6$, and the factor $2^{4r-4}$ gives
the required divisibility in the trivial and sign multiplicities.
Their leading $t^n$ coefficients are positive.

By weak Lefschetz and duality, removing the ambient cohomology changes
this Euler representation by twice the alternating lower cohomology
and by $H^{2r}(X)$.  Consequently the multiplicity parities on $V$
are precisely those on $H^{2r}(X)$.
In the rational representation ring modulo two, write $U=\mathrm{Std}$.
The relations $U^2=1+\varepsilon+U$ and $\varepsilon U=U$ give
$U^4=1+\varepsilon+U$ modulo two.  Since
$H^1(X)=4(r-1)\varepsilon+4U$ as a $G$ representation, its full exterior
algebra has generating series
\[
 \lambda_z H^1(X)
   =(1+\varepsilon z)^{4(r-1)}(1+Uz+\varepsilon z^2)^4
   \equiv(1+z^4)^{r-1}
             \bigl(1+(1+\varepsilon+U)z^4+z^8\bigr)\pmod2.
 \tag{A23.14}
\]
For odd $r$, the coefficient of $z^{2r}$ is zero.  For $r=2s$ it has
sign and standard coefficients $\binom{r-1}{s-1}$; its trivial
coefficient is
\[
 \binom{r-1}{s}+\binom{r-1}{s-1}+\binom{r-1}{s-2}
       \equiv\binom{r-1}{s-2}\pmod2,
 \tag{A23.15}
\]
because $\binom{2s}{s}$ is even for $s\geq1$.
This proves (A23.4): the three multiplicity parities on $V$ are
$(\alpha_r,\beta_r,\beta_r)$.
Removing the regular $K$ subtracts multiplicities $(1,1,2)$, leaving
$(1-\alpha_r,1-\beta_r,\beta_r)$.

\paragraph{Simultaneous Frobenius selection and the full algebraic sign.}
The rational multiplicity Frobenius polynomials are common across
coefficient primes by the finite-order twists of C44-A19.
Each projection of the connected joint geometric monodromy group is
the appropriate special orthogonal group.  The even-rank no-root
selector of C44-A12 and odd-rank simple-root selector of C44-A16 have
nonempty inverse-image opens in this one connected group.  Their
intersection is nonempty.  The same constant-field finite-quotient
Chebotarev argument as in C44-A21 therefore gives a Zariski-dense set
of closed parameters at which all three multiplicities are minimal
at once.  No independence of the three factors is asserted or needed.

The parity list gives (A23.5), remembering that the standard
multiplicity contributes twice to $W$.  The four actual relative
self-adjoint projectors
$e_1,e_\varepsilon,e_U(1+\sigma)/2,e_U(1-\sigma)/2$
each have residual potential Tate multiplicity at most one.
Their embedded polarized lifts have already been constructed, and
the positive $K$ splits into positive orthogonal pieces.
C44-A19 proves the sign on each complete algebraic block, so on their
sum $A_B$.  The positive Gram matrix proves the six-dimensional
lower bound when centres are imposed.  This proves the theorem.
\end{proof}

\paragraph{Concrete dimensions and scope.}
For $r=2$ this recovers C44-A21.  For every odd $r\geq3$, the six-centre
construction has residual multiplicity exactly two and full potential
Tate multiplicity eight.  For $r=2^j\geq4$, all binomial coefficients
in (A23.4) are odd: modulo two,
$(1+x)^{2^j-1}=1+x+\cdots+x^{2^j-1}$.
Thus the construction without centres has multiplicity four, while
the six-centre construction has residual multiplicity two and full
multiplicity eight.  All these entire algebraic axis forms are
positive, even though the unsplit residual block has more than one
potential Tate direction.

The arguments are valid for every $E$ in characteristic $p>3$.
Taking $E:y^2=x^3-x$ and $p\equiv3\pmod4$ gives the same supersingular
ambient examples as C44-A20, after finite extensions.  No claim is made
that any residual direction contains an algebraic cycle, or that these
known centres or their primitive classes are nonliftable.

The primary inputs and their scope are unchanged: the absolute abelian
Chow Lefschetz inverse of K\"unnemann, Theorem~5.2 p.~96; tame
fixed-point traces of Deligne--Lusztig, Theorem~3.2 p.~119; SGA~7~II,
Expose~XII, Theorem~3.3 and the subsequent affine-quadric comparison,
and Expose~XV, Theorem~3.4 pp.~195--196; Katz's Theorem~1.4.6 and
Lemmas~2.1.4--2.1.6; and the finite-character weight estimate of
Weil~II, Theorem~3.3.1 pp.~204--205 and Corollary~3.3.4 p.~206.
These theorems apply in the indicated arbitrary dimensions; no
additional equivariant monodromy theorem or mixed-characteristic
cycle-lifting converse is being imported.

For cote~44, Th.~5, scan~004 and the primitive instruction on scans
008--009, this extends the same sufficient family for QH9-axis to
higher-dimensional genuine pencils.  The axis sign is unconditional;
the separate ambient abelian term in a full height conclusion remains
subject to the unrefereed 032/Milne~2026 assumptions already recorded
in the programme.

% END work/scan_w/higher_dimensional_s3_axes.tex

% BEGIN work/scan_w/totally_real_axis_blocks.tex
% C44_POS /root/scan_w; GPT-6 Astra Pro; reasoning effort not exposed.
\section{A totally real coefficient field can separate several Tate directions}
\label{sec:c44-a24}

\begin{theorem}[C44-A24: one residual direction over a totally real field]
Let $B\subset X/\mathbf F_q$ be a smooth intersection of two members of
one ample line bundle, with $\dim X=2r+2$, and set
\[
 V_\ell=\ker\kappa_*\subset H^{2r}(B,\mathbf Q_\ell(r)),\qquad
 q_B(x,y)=(-1)^r\int_Bxy.
 \tag{A24.1}
\]
Assume the polarized embedded complete intersection has a smooth
mixed-characteristic lift with a complex generic fibre as in C44-A19.
Let $e$ be an actual relative rational algebraic correspondence whose
actions on the variable cohomologies are self-adjoint idempotents, and
put $U=eV$.  Suppose a totally real number field $F$ acts on $U$ through
actual relative rational algebraic correspondences, with identity $e$,
and every element of $F$ is self-adjoint for $q_B$.  The field and
projector identities are required on both fibres and in the coefficient
realizations; no assertion that they hold as equalities of Chow cycles
is necessary.

Assume the normalized Frobenius on $U$ has a common polynomial
$\Phi_F(T)\in F[T]$: on every component $U_v$ for a finite place
$v\mid\ell\ne p$, its characteristic polynomial is the image of
$\Phi_F$ in $F_v[T]$.  The $F_v$-ranks are the common $F$-rank of the
Betti space.  All special correspondences are defined after an allowed
finite extension of the finite ground field.

Let $K$ be a known positive rational algebraic subspace of the initial
realization of $U$, stable under $F$; $K=0$ is allowed.  Use the common
cycle representatives constructed below to realize it everywhere.
Suppose the residual polynomial on $K^\perp\cap U$ has at most one
root of unity, counted with algebraic multiplicity, over any, equivalently
every, embedding of $F$ into an algebraic closure.  Then $q_B$ is
positive definite on the complete rational algebraic subspace of $U$.
If $d=[F:\mathbf Q]$ and $\tau_F$ is that residual multiplicity, then
\[
 \dim_{\mathbf Q}K\ \leq\
 \dim_{\mathbf Q}\operatorname{Alg}(U)
       \ \leq\ \dim_{\mathbf Q}K+d\tau_F.
 \tag{A24.2}
\]
When the residual algebraic part is nonzero it is an $F$-line and its
$F$-valued form is totally positive.  The cycles on that line need not
lift.  No Tate conjecture, standard conjecture $D(B)$, Frobenius
semisimplicity, CM action, or 032 input is used.

The same assertion applies to several orthogonal relative blocks with
possibly different totally real fields.  If they sum to $V$, positivity
on their complete algebraic subspaces proves positivity on all of $A_B$.
\end{theorem}

\begin{proof}
We keep the field $F$ distinct from the normalized Frobenius operator
$\mathsf F$.  After a finite extension all chosen cycles and
correspondences are defined, so $\mathsf F$ commutes with $F$.

\paragraph{The actual $F$-valued form and its cycle values.}
Choose a rational basis $f_1,\ldots,f_d$ of $F$ and its trace-dual basis
$f_1^\vee,\ldots,f_d^\vee$.  On the Betti space define
\[
             b(x,y)=\sum_{j=1}^d q_B(f_jx,y)f_j^\vee.
 \tag{A24.3}
\]
Then $\operatorname{Tr}_{F/\mathbf Q}(a b(x,y))=q_B(ax,y)$ for
every $a\in F$.  Nondegeneracy of the trace pairing, commutativity of
$F$, and self-adjointness imply that $b$ is symmetric and $F$-bilinear,
and that $q_B=\operatorname{Tr}_{F/\mathbf Q}b$.
It is nondegenerate because $q_B$ is nondegenerate on $U$.
The same construction over $F\otimes\mathbf Q_\ell$ gives its local
forms $b_v$; their trace is the corresponding rational pairing.

For actual rational cycles in the block, (A24.3) lies in $F$ itself:
each coefficient is a rational intersection number of actual cycles.
These elements of $F$ have exactly the same images in every $F_v$.
This gives more information than the rational trace of a norm; no
realization independence of arbitrary vanishing classes is assumed.

To put all representatives in the block, first use the common absolute
special-fibre projector
\[
       P=1-\kappa^*(L_h^2)^{-1}\kappa_*
 \tag{A24.4}
\]
from C44-A15 and then apply $e$.  These are actual rational algebraic
correspondences.  The field operations in (A24.3), and below, mean
applications of chosen rational correspondences realizing the
specified elements of $F$, with $P$ and $e$ inserted if necessary.
Their cohomological identities hold everywhere by hypothesis.  Chow
idempotence is not required.

Choose an $F$-basis $C_1,\ldots,C_s$ of the known space $K$, represented
by these actual projected cycles, and put
\[
                         G=(b(C_i,C_j))\in M_s(F).
 \tag{A24.5}
\]
It is nonsingular.  Positivity of the rational form on $K$ makes every
real conjugate of $G$ positive definite, so $\det G$ is totally positive.
The same matrix in every $F_v$ makes the actual cycles $F_v$-independent
and gives common nondegenerate subspaces $K_v$ of rank $s$.
After defining the finitely many $f_jC_i$ over a finite field, their
normalized Frobenius is the identity.  The common residual polynomial
is therefore $\Phi_F(T)/(T-1)^s$.

For any additional algebraic cycle in $U$, subtract its $b$-orthogonal
projection onto $K$, using $G^{-1}$ and the elements of $F$ in (A24.3).
This operation is a finite rational combination of actual cycle
correspondences.  It gives one actual rational cycle whose class lies
in $K_v^\perp\cap U_v$ in every realization, and represents the intended
residual class initially.  This is the field-valued version of the
common-representative repair in C44-A19.

\paragraph{A totally positive Betti determinant at every embedding.}
For a real embedding $\sigma:F\hookrightarrow\mathbf R$, let
$U_\sigma=U_{\mathrm B}\otimes_{F,\sigma}\mathbf R$.
The relative algebraic correspondences preserve Hodge types.  The
space $U_\sigma$ is therefore a real Hodge substructure of primitive
middle cohomology.  Self-adjointness makes the distinct real-embedding
summands orthogonal, and the restriction of $q_B$ is $\sigma b$.
The Hodge--Riemann relations are positive on type $(r,r)$; all negative
real directions come in conjugate off-diagonal pairs.  Thus the
negative index of $\sigma b$ is even for \emph{each} $\sigma$, not
only for their rational trace form.
In any $F$-basis of the Betti space this proves that
\[
                            D=\det_F b_{\mathrm B}\in F^*
              \quad\hbox{is totally positive}.
 \tag{A24.6}
\]
This is an ordinary $F$-determinant; it is not the determinant of the
rational trace form, and no discriminant-of-$F$ factor is omitted.
For a zero-dimensional block the conclusion is empty.

Smooth proper base change and Betti--etale comparison respect these
same relative correspondences.  They are consequently $F_v$-linear
and preserve $b_v$, which is characterized by its trace identities.
It follows that $D$ gives the determinant squareclass of $b_v$ for
every finite $v\nmid p$.  The known special cycle space $K$ does not
need to extend to the Betti fibre for this assertion.

\paragraph{One potential Tate line over $F$, with common cycle values.}
The root-of-unity condition is unchanged under any embedding of the
coefficient field: an extension of that embedding to the algebraic
closure permutes roots and preserves roots of unity and their
multiplicities.  If the multiplicity is one, its root is $1$ or $-1$.
Indeed the polynomial at a real embedding has real coefficients,
so any nonreal root of unity would have a second conjugate.  The
corresponding equality at $1$ or $-1$ is then an equality in $F$,
and its simple multiplicity is common to every place.

If $\tau_F=0$, every algebraic residual class is zero: its normalized
Frobenius orbit is finite.  Suppose instead $\tau_F=1$ and set
$S=\mathsf F^2$ on each residual local space.  Its eigenline $T_v$ at
one is simple and nondegenerate.  For completeness, an isometry pairs
only reciprocal generalized eigenspaces.  If the simple line at one
were degenerate, it would be orthogonal to the whole space.  Every
finite-orbit vector lies on this line, without any semisimplicity
assumption on the remaining generalized eigenspaces.

Let $z$ be an actual residual cycle nonzero at the initial coefficient
prime.  It is nonzero in at least one $F_v$ component, and the
nondegenerate line there gives
\[
                              a=b(z,z)\in F^*.
 \tag{A24.7}
\]
Because this is one nonzero element of $F$, its image is nonzero at
\emph{every} finite place.  Therefore the same actual cycle is nonzero
and spans $T_v$ at every $v\nmid p$.  This is the needed repair of the
cross-place nonvanishing issue; it is stronger than just knowing that
the rational self-intersection has a given sign or trace.

\paragraph{The non-Tate determinant and the missing finite places.}
Let $N_v=T_v^\perp$ inside the residual space.  On $N_v$ the isometry
$S$ has no eigenvalue $1$ or $-1$.  Its common $F$-polynomial is denoted
$\psi(T)$.  It is obtained from the common residual polynomial by
squaring eigenvalues and removing the simple factor $T-1$.
The dimension of $N_v$ is even, and $\det S=1$ there, because all
eigenvalues occur in reciprocal pairs.

The operator $A=S-S^{-1}$ is invertible and skew-adjoint for $b_v$.
The form $b_v(Ax,y)$ is therefore nondegenerate alternating.  Its
determinant is a square.  Consequently
\[
 \det b_v|N_v\equiv \det A
       =\delta:=\psi(1)\psi(-1)
                              \pmod{(F_v^*)^2}.
 \tag{A24.8}
\]
The equality for $\det A$ uses the even dimension and $\det S=1$;
it includes the empty complement, with $\delta=1$.
The common rational Frobenius polynomial on the underlying block is
the field norm of $\Phi_F$.  Purity therefore puts every root of every
real conjugate of $\psi$ on the unit circle.  None is $1$ or $-1$,
so conjugate reciprocal pairs show that $\delta$ is totally positive.

The orthogonal decomposition into $K_v$, $T_v$ and $N_v$ gives
\[
              \frac{a\det(G)\delta}{D}\in(F_v^*)^2
                             \qquad(v\nmid p).
 \tag{A24.9}
\]
The map
\[
 F^*/(F^*)^2\longrightarrow
             \prod_{v\notin\mathcal S}F_v^*/(F_v^*)^2
 \tag{A24.10}
\]
is injective for any finite set $\mathcal S$ of finite places.  Here
is the exact proof needed.  If $c$ is a nonsquare in $F$, the quadratic
extension $F(\sqrt c)/F$ is nontrivial.  Number-field Chebotarev gives
infinitely many unramified primes with nontrivial Frobenius, and one
can choose such a prime outside $\mathcal S$.  At that prime the
extension is a quadratic field extension, so $c$ is not a local square.
Apply this with $\mathcal S=\{v:v\mid p\}$ to (A24.9).

The ratio in (A24.9) is thus a square in $F^*$.  All its other factors
are totally positive by (A24.5), (A24.6) and (A24.8).  Hence $a$ is
totally positive.  The form on the rational space $Fz$ is
\[
                          q_B(fz,fz)
                            =\operatorname{Tr}_{F/\mathbf Q}(a f^2),
 \tag{A24.11}
\]
which is positive definite.

Finally, if $y$ is any other actual residual cycle, then
$b(y,z)/a\in F$.  At every place both $y$ and $z$ lie on the same
line and $z$ is nonzero, so
$y=(b(y,z)/a)z$ cohomologically in the initial realization and in all
the other realizations.  Thus the complete residual algebraic space
has $F$-dimension at most one.  Adding the known positive space $K$
proves positivity and (A24.2).  Orthogonality proves the stated
multiple-block version.
\end{proof}

\begin{corollary}[C44-A24a: the polynomial hypothesis from a finite symmetry algebra]
Suppose the relative action and embedded equivariant lift come from
a finite constant group $G$ as in C44-A19b.  Let $e\in\mathbf Q[G]$
be self-adjoint, and suppose the corner $e\mathbf Q[G]e$ contains a
totally real field $F$, with unit $e$, fixed pointwise by the adjoint
involution $g\mapsto g^{-1}$.  Then the common $F$-valued Frobenius
polynomial required in C44-A24 exists.  In particular that theorem
applies once its known-positive-space and residual multiplicity tests
are satisfied.
\end{corollary}

\begin{proof}
The relative Betti $F$-module is finite-dimensional and free, and
comparison gives its common rank in every $F_v$ realization.
For $f_j$ in a rational basis of $F$, the element $f_j=e f_j e$ is a
rational linear combination of group elements.  The finite-field
twists in C44-A19b give rational, coefficient-independent numbers
\[
             t_{j,n}=\operatorname{Tr}_{\mathbf Q_\ell}
                       (f_j\mathsf F^n\mid U_\ell)
                         \qquad(n\geq1).
 \tag{A24.12}
\]
In detail the trace on $U$ equals the trace of the same group-algebra
element times Frobenius on $V$; each group-element trace is obtained
degreewise from its smooth projective twist, and the ambient traces
are subtracted.  Centrality of $e$ or of $f_j$ in the whole group
algebra is unnecessary.

Trace transitivity for a free $F\otimes\mathbf Q_\ell$-module gives
\[
      t_n:=\sum_jt_{j,n}f_j^\vee\in F,
       \qquad t_n\mapsto\operatorname{Tr}_{F_v}
                              (\mathsf F^n\mid U_v)
                                  \quad\hbox{for every }v\mid\ell.
 \tag{A24.13}
\]
Indeed pairing both sides with every $f_j$ under the field trace
recovers (A24.12), and the trace pairing is nondegenerate.
Newton identities in characteristic zero now reconstruct one
polynomial $\Phi_F\in F[T]$ with the required images.  This is an
actual reconstruction of the $F$-coefficients, rather than an
inference from rationality of the norm polynomial alone.
\end{proof}

\paragraph{The dihedral algebra giving a nonsplit rational example of the criterion.}
For $G=D_5=\langle\rho,\sigma:\rho^5=\sigma^2=1,
\sigma\rho\sigma=\rho^{-1}\rangle$, put
\[
 e_1=\frac1{10}\sum_g g,\qquad
 e_\varepsilon=\frac1{10}\sum_g\varepsilon(g)g,\qquad
 e_0=1-e_1-e_\varepsilon,\qquad
 e_\pm=\frac12e_0(1\pm\sigma).
 \tag{A24.14}
\]
These are the four self-adjoint orthogonal rational projectors.
The nontrivial rational simple component is
\[
              e_0\mathbf Q[D_5]\simeq M_2(F),
                         \qquad F=\mathbf Q(\sqrt5).
 \tag{A24.15}
\]
One direct verification uses $\alpha^2+\alpha-1=0$ and the matrices
\[
       \rho\longmapsto
          \begin{pmatrix}0&-1\\1&\alpha\end{pmatrix},\qquad
       \sigma\longmapsto
          \begin{pmatrix}0&1\\1&0\end{pmatrix}.
 \tag{A24.16}
\]
They satisfy the dihedral relations.  The two real embeddings of $F$
give the two distinct absolutely irreducible two-dimensional
representations; together with $1,\varepsilon$ their squared degrees
sum to ten.  These explicit split representations verify (A24.15).
The central element $\rho+\rho^{-1}$ acts as $\alpha$ and is
self-adjoint.  Thus
$e_\pm\mathbf Q[D_5]e_\pm\simeq F$ is fixed pointwise by the adjoint
involution.

If a geometric axis supplies this actual relative symmetry algebra,
each of the two $e_\pm$ blocks can have residual $F$-rank one, hence
residual rational potential Tate dimension two.  C44-A24 still proves
its complete algebraic sign.  The available rational group-algebra
projectors cannot split either corner further, since a field has no
nontrivial idempotents.  Consequently this case is not obtained by
applying the rational multiplicity-one test C44-A19 to the supplied
group-algebra projectors.  No nonexistence of other geometric
projectors is asserted, and no smooth-pencil or closed-point example
is inferred from this algebra calculation alone.

\paragraph{Sources and scope.}
The geometric comparison and common-cycle inverse are the checked
inputs of C44-A15 and A19.  The new number-field squareclass argument
uses only the quadratic case of Chebotarev: Milne, \emph{Class Field
Theory}, Chapter~V, Theorem~3.23 and Corollary~3.24, printed p.~164,
or \emph{Algebraic Number Theory}, Theorem~8.31, printed p.~147.
The finite-place deletion is proved explicitly in (A24.10); the
rational-only statement in Ancona--Marmora is not being cited as a
number-field theorem.  The trace-dual reconstruction, every-embedding
determinant sign and dihedral algebra calculation are proved above.

For cote~44, Th.~5, scan~004 and the primitive warning on scan~008,
this supplies an additional sufficient test for the actual QH9-axis
form.  It requires an extending self-adjoint field of correspondences
and a verified common field-valued polynomial.  It does not promote
a special Frobenius eigenspace to a relative factor.  The axis sign
is unconditional under these explicit hypotheses; the separate
ambient abelian contribution to a full height theorem retains the
032/Milne~2026 dependency already stated in the programme.

% END work/scan_w/totally_real_axis_blocks.tex

% BEGIN work/scan_w/low_rank_tate_obstructions.tex
% C44_POS /root/scan_w; next gate after A15. No geometric counterexample claimed.
\section{The rank-two and rank-three potential Tate gate}
\label{sec:c44-a17}

\begin{proposition}[C44-A17: determinant parity for a fully algebraic Tate block]
Retain the fixed axis and the known positive span $K$ of C44-A15.
Let $\mathcal T_{\ell_0}$ be the full generalized
root-of-unity Frobenius subspace of the residual space
$W_{\ell_0}$, of dimension $\tau$. Suppose that this
\emph{entire cohomological subspace} is spanned over
$\mathbf Q_{\ell_0}$ by algebraic cycle classes.
Then the rational algebraic residual space is of dimension
$\tau$, its intersection form is nondegenerate, and its
ordinary Gram determinant is positive. In particular its
negative index is even. Thus:
\[
\begin{array}{c|c}
 \tau&\text{possible signatures }(s_+,s_-)\\ \hline
 1&(1,0)\\
 2&(2,0)\text{ or }(0,2)\\
 3&(3,0)\text{ or }(1,2).
\end{array}
 \tag{A17.1}
\]
Full algebraicity is an explicit hypothesis here; it is
not deduced from the Frobenius polynomial or assumed via
the Tate conjecture. No standard conjecture $D(B)$ is used.
\end{proposition}

\begin{proof}
The generalized root-of-unity part is orthogonal to the
remaining generalized eigenspaces: the inverse of a root
of unity is again a root of unity. Hence its pairing is
nondegenerate. Choose $\tau$ actual cycles spanning this
part and apply the common algebraic projector and rational
orthogonal subtraction of A15. Their Gram matrix $G_T$
is a nonsingular rational matrix. In every coefficient
realization these same cycles have matrix $G_T$, so they
are independent, have finite Frobenius orbit, and fill the
root-of-unity part, whose dimension is independent of
$\ell$. Their common finite field of definition also shows
that Frobenius on this part is semisimple of finite order.
This is a consequence of the spanning hypothesis, not an
assumption of semisimplicity on all middle cohomology.

Every other residual algebraic cycle lies in this same
subspace. Pairing with the selected basis and multiplying
by $G_T^{-1}$ gives rational coefficients, proving equality
of the rational cycle-class space with this rational span.
The nondegenerate Gram matrix also constructs an actual
special-fibre Chow projector onto it: use the selected
cycles and their rational intersection-dual cycles.
Their compositions are the identity matrix of intersection
numbers. Its image is a direct sum of $\tau$ Tate motives
and has these same independent cycle realizations, including
the crystalline realization. This observation requires
neither an extension of the projector over a trait nor
independence of an arbitrary homological equivalence relation.

Choose an even exponent killing all the roots of unity.
On the orthogonal non-Tate complement the operator
$S-S^{-1}$ is invertible and skew-adjoint. The proof of
A15 gives a positive rational number
$\delta=\psi(1)\psi(-1)$ representing its determinant
squareclass in every $\mathbf Q_\ell$, $\ell\ne p$.
The one embedded geometric lift gives the positive
rational Betti determinant $D$ of the full variable space.
Thus
\[
             \det G_T\equiv D/(\det G\,\delta)
                     \pmod{(\mathbf Q_\ell^*)^2}
                    \quad\text{for every }\ell\ne p.
\]
The missing-$p$ squareclass injection of A15 makes this
a global rational squareclass equality. Its right side
is positive. Hence $\det G_T>0$, and the negative index
of $G_T$ is even, giving (A17.1).
\end{proof}

\paragraph{The exact remaining Hilbert-symbol condition.}
For a rational nondegenerate form $q_T$ of rank two or
three with positive ordinary determinant, write
$\epsilon_v(q_T)$ for its usual Hasse invariant.
Its two possibilities in (A17.1) are distinguished by
\[
 q_T>0\quad\Longleftrightarrow\quad
 \epsilon_\infty(q_T)=1
 \quad\Longleftrightarrow\quad
 \epsilon_p(q_T)=\prod_{\ell\ne p}\epsilon_\ell(q_T).
 \tag{A17.2}
\]
The product is finite. These equivalences follow from
$\epsilon_\infty=(-1)^{s_-(s_--1)/2}$ and Hilbert reciprocity.
They isolate a $p$-adic invariant; determinant parity does
not determine it. The necessary local classification and
product formula are Ancona--Marmora, Propositions~4.1--4.2
and Theorem~4.3, printed p.~629.

\begin{proposition}[C44-A17a: sharp arithmetic ambiguity away from one prime]
For every prime $p$ there are positive definite and negative
definite rational binary forms $q_+$ and $q_-$ with the same
positive determinant such that
\[
                   q_+\otimes\mathbf Q_\ell\simeq
                   q_-\otimes\mathbf Q_\ell
                         \quad\text{for all }\ell\ne p.
 \tag{A17.3}
\]
There are likewise rational ternary forms of signatures
$(3,0)$ and $(1,2)$ satisfying (A17.3) and containing a
common positive line. Frobenius may be set to the identity
in all these arithmetic examples. Thus all coefficient
primes away from $p$, positive determinant, and a fully
Tate Frobenius polynomial do not by themselves distinguish
the two signatures in (A17.1).
These are quadratic-space examples, not algebraic varieties
or counterexamples to a standard conjecture.
\end{proposition}

\begin{proof}
If $p=2$, take $d=u=1$. If $p\equiv3\pmod4$, take
$d=p,u=1$. If $p\equiv1\pmod4$, take $d=p$ and choose
a prime $u\equiv3\pmod4$ with $(u/p)=-1$; Dirichlet's
theorem and the Chinese remainder theorem give such $u$.
Set
\[
                   q_+=\langle1,d\rangle,
 \qquad            q_-=\langle-u,-d/u\rangle.
 \tag{A17.4}
\]
Both determinants equal $d>0$, and the stated real signs
are immediate. Their finite Hasse invariants satisfy
\[
       \epsilon_v(q_+)=1,
 \qquad\epsilon_v(q_-)=(-u,-d/u)_v=(-u,-d)_v.
 \tag{A17.5}
\]
The last equality uses $(a,-a)_v=1$.

For $p=2$, the quaternion $(-1,-1)$ is ramified precisely
at $2$ and infinity. For $p\equiv3\pmod4$, the symbol
$(-1,-p)$ is negative at $p$ and infinity and positive
elsewhere: at $p$ use $(-1/p)=-1$, while at $2$ the
odd-unit formula gives $1$ since $-p\equiv1\pmod4$.
For $p\equiv1\pmod4$, the symbol $(-u,-p)$ is negative
at $p$ because $(-u/p)=(u/p)=-1$. At the auxiliary prime
$u$ it is positive, since quadratic reciprocity gives
\[
                     (-p/u)=-(p/u)=-(u/p)=1.
\]
At $2$ it is positive because $-u\equiv1\pmod4$;
at infinity it is negative; at every other finite prime
both coefficients are units and the symbol is positive.
Thus in all cases the only discrepancies in (A17.5)
are at $p$ and infinity.

The equality of dimension, determinant and Hasse invariant
over every $\mathbf Q_\ell$, $\ell\ne p$, proves (A17.3)
by local classification. Add $\langle1\rangle$ to both
forms to obtain the ternary examples and their common
positive line. The identity is an isometry with entirely
Tate characteristic polynomial in either case.

The formulas used above are independently checked in
Voight, \emph{Quaternion Algebras}, version~1.0.5,
\S12.4, especially (12.4.10) and (12.4.14), pp.~189--190.
The exact three choices of quaternion presentation above
also appear in Example~14.2.13, p.~214, after the proof
of Proposition~14.2.7 and Dirichlet Theorem~14.2.9,
pp.~213--214. Thus the local ambiguity is present at
every residue characteristic, not just at $2$.
\end{proof}

\begin{proposition}[C44-A17b: a Tate isocrystal need not be an admissible subobject]
There is, for every prime $p$, a weakly admissible filtered
$\varphi$-module over $\mathbf Q_p$ with a perfect symmetric
filtered $\varphi$-pairing and a nondegenerate rank-two
$\varphi$-fixed summand whose induced filtration is not
weakly admissible. Its Frobenius characteristic polynomial
can be rational with all complex roots of absolute value one.
Thus these properties of the full filtered module do not
make the special Tate projector a filtered projector.
\end{proposition}

\begin{proof}
Let $D$ have basis $t_1,t_2,e,f$, with symmetric pairing
\[
 q(t_i,t_j)=-\delta_{ij},\qquad q(e,f)=1,
 \qquad q(e,e)=q(f,f)=0,
 \qquad q(t_i,e)=q(t_i,f)=0.
\]
Choose the root $\lambda\in\mathbf Q_p$ of
\[
                       T^2-p^{-1}T+1
 \tag{A17.6}
\]
with valuation $-1$. Such a root exists because
$1-4p^2$ is a square in $\mathbf Q_p$ by Hensel's lemma;
choose its square root sufficiently close to $1$ and use
$(1+\sqrt{1-4p^2})/(2p)$. The other root is
$\lambda^{-1}$ of valuation $1$. Define
\[
 \varphi t_i=t_i,\qquad \varphi e=\lambda e,
 \qquad \varphi f=\lambda^{-1}f.
\]
This preserves $q$. The rational Frobenius polynomial is
$(T-1)^2(T^2-p^{-1}T+1)$; its complex roots all have
absolute value one. The last two are not roots of unity,
since their sum is the nonintegral rational number $1/p$.

The line $L=\langle t_1+e+f/2\rangle$ is isotropic.
Put
\[
 \operatorname{Fil}^1D=L,\quad
 \operatorname{Fil}^0D=L^\perp,\quad
 \operatorname{Fil}^{-1}D=D,
 \qquad \operatorname{Fil}^2D=0.
 \tag{A17.7}
\]
This filtration is self-dual for $q$. Its Hodge weights
are $1,0,0,-1$, so $t_H(D)=0=t_N(D)$.

Write $T=\langle t_1,t_2\rangle$, $E=\langle e\rangle$
and $F=\langle f\rangle$. Since the three eigenvalues
$1,\lambda,\lambda^{-1}$ are distinct, every
$\varphi$-stable subspace is $T'\oplus E'\oplus F'$
with $T'\subset T$, $E'$ either $0$ or $E$, and $F'$
either $0$ or $F$. For such a subspace $D'$ the induced
filtration has
\[
 t_H(D')=\mathbf1_{L\subset D'}-
                    \mathbf1_{D'\not\subset L^\perp},
 \qquad
 t_N(D')=-\dim E'+\dim F'.
 \tag{A17.8}
\]
If $E'=E$ and $F'=0$, the first indicator is zero and
the second is one, since $q(L,e)=1/2$; hence
$t_H(D')=-1=t_N(D')$. If $E'=0,F'=F$, one has
$t_H(D')\leq0<1=t_N(D')$. If both are present,
$t_N(D')=0$ and the second indicator is one, so
$t_H(D')\leq0$. If neither is present, the first
indicator is zero and again $t_H(D')\leq0=t_N(D')$.
This checks all subspaces and proves weak admissibility.

On the nondegenerate fixed subspace $T$, however, the
induced filtration is
\[
 \operatorname{Fil}^1T=0,\qquad
 \operatorname{Fil}^0T=\langle t_2\rangle,
 \qquad\operatorname{Fil}^{-1}T=T.
\]
Consequently $t_H(T)=-1$ whereas $t_N(T)=0$.
It is not weakly admissible. Equivalently the orthogonal
projector onto $T$ sends $L$ to $\langle t_1\rangle$
and fails to preserve the filtration. This failure
occurs despite the full module's self-dual filtered
structure and weak admissibility.
\end{proof}

\paragraph{What the rank-two theorem really supplies.}
Ancona, \emph{Standard conjectures for abelian fourfolds},
Theorem~8.1 (author arXiv:1806.03216v3, p.~25), does
\emph{not} assume a given CM action: it proves positivity
for a rank-two motive with quadratic form in mixed
characteristic, full numerical algebraic dimension two
on the special fibre, and a polarizing Betti form.
Proposition~8.3, p.~26, deduces the necessary all-realization
algebraicity and nondegeneracy under those hypotheses.
Proposition~10.8, pp.~32--33, derives a quadratic
$p$-adic field action using the orthogonal group of the
\emph{admissible rank-two object}. The isotropic-line
argument and admissibility are used explicitly in
excluding a split quadratic algebra. No such admissible
rank-two object is supplied merely by a Frobenius
projector on the special fibre; A17b shows the inference
from full weak admissibility is false.

The higher-rank signature theorem of Ancona--Marmora,
Theorem~3.2 and Remark~3.3, printed p.~626, does assume
a relative number-field action of full degree with
adjoint involution. Its remark explicitly places the
removal of those action hypotheses in Conjecture~3.1.
We do not transfer that theorem to an arbitrary
special Tate block. Neither numerical semisimplicity
nor a special-fibre identification with a sum of Tate
motives specifies the geometric form's real signature.

\paragraph{Status of the next axis gate.}
For cote~44, Theorem~5, scans~004 and~008, A17 gives a
determinant constraint when the residual block is
actually algebraic. A17a proves that the resulting
arithmetic constraints still allow both remaining
rank-two and rank-three signatures; A17b identifies
the precise missing filtered-subobject step in the
attempt to apply the rank-two $p$-adic theorem.
The residue-characteristic Hilbert invariant in
(A17.2) is not computed here for an arbitrary axis.
There is no geometric wrong-sign example and no claim
that the residual rank-two or full sixfold
multiplicity-three gate has been closed.

\paragraph{Primary sources actually read.}
Ancona's Theorem~8.1, Proposition~8.3, the complete
quadratic-form reduction in \S9, the comparison
discussion in \S10, and the complete proof of
Proposition~10.8 were read in the downloaded primary
author version
\url{https://arxiv.org/pdf/1806.03216v3}.
Ancona--Marmora's Theorem~3.2/Remark~3.3 and the full
proof of Corollary~4.4 were read at
\url{https://ems.press/content/serial-article-files/51595?nt=1}.
The explicit Hilbert symbols and the quaternion
construction of Voight cited above were read in
\url{https://jvoight.github.io/quat-book-v1.0.5.pdf}.
The filtered-module construction A17b is proved
directly above; no geometric realization theorem
for this example is invoked. In particular it does not
come with a compatible rational Betti factor whose form
polarizes its Hodge structure. It refutes only the
filtered-subobject inference, not the combined geometric
hypotheses of Ancona's theorem.

% END work/scan_w/low_rank_tate_obstructions.tex

% BEGIN work/scan_w/odd_tate_selection_axes.tex
% C44_POS /root/scan_w; odd-rank companion to A12 and application of A15.
\section{Selecting one potential Tate direction and positive spans of arbitrary rank}
\label{sec:c44-a16}

\begin{lemma}[C44-A16: odd orthogonal selection]
Let $U/\mathbf F_q$ be smooth, geometrically connected and
quasiprojective of positive dimension. Let $\mathcal W$ be
a lisse $\mathbf Q_\ell$-sheaf, $\ell\ne p$, of odd rank
$d\geq1$, endowed with a nondegenerate symmetric form
preserved by arithmetic monodromy. Suppose that its geometric
monodromy closure contains $\operatorname{SO}(\mathcal W)$
and that every closed-point Frobenius characteristic
polynomial lies in $\mathbf Q[T]$.
Then a Zariski-dense set of closed points has exactly one
root-of-unity Frobenius eigenvalue counted with algebraic
multiplicity. This eigenvalue is simple and is $1$ or $-1$.
No semisimplicity or algebraicity of this line is asserted
or required.
\end{lemma}

\begin{proof}
Put $M=\operatorname{lcm}(1,\ldots,2d^2)$, which is even.
If a root of unity of order $n$ occurs in a rational degree-$d$
polynomial, then $\varphi(n)\leq d$ and $n\leq2d^2$.
For completeness, the elementary latter bound follows from
\[
 \frac{\varphi(n)^2}{n}
       =\prod_{a^e\parallel n}a^{e-2}(a-1)^2\geq\frac12,
 \tag{A16.1}
\]
where the product is over prime powers: the possible factor
less than one is $a=2,e=1$, whose value is $1/2$.
Consequently all root-of-unity eigenvalues of any Frobenius
in the statement are killed by the exponent $M$.

For an orthogonal matrix $g$ define the conjugacy-invariant
regular function
\[
 f(g)=[T^1]\det\bigl(TI-(g^M-I)\bigr),
 \tag{A16.2}
\]
where $[T^1]$ means the coefficient of $T$. In odd dimension
every orthogonal matrix has an eigenvalue $1$ or $-1$.
Thus $g^M$ has eigenvalue one, and $f(g)\ne0$ means precisely
that its eigenvalue one has algebraic multiplicity one.
This is a statement about a characteristic polynomial, so
Jordan blocks elsewhere cause no difficulty.

The restriction of $f$ to $\operatorname{SO}_d$ is not zero.
Over an algebraic closure a maximal torus has eigenvalues
\[
             1,t_1,t_1^{-1},\ldots,t_{(d-1)/2},t_{(d-1)/2}^{-1}.
\]
Choose each $t_i^M\ne1$. Then the zero eigenvalue of
$g^M-I$ is simple, so $f(g)\ne0$. For $d=1$ the function
on $\operatorname{SO}_1$ is simply $1$.
Zariski density of the geometric monodromy image in its
closure supplies a geometric image element $\gamma$ with
$f(\gamma)\ne0$.

Let $\Gamma$ be the compact arithmetic image. By continuity
and conjugacy invariance there is an open normal subgroup
$J\subset\Gamma$ such that the inverse image of the
conjugacy class of $\gamma$ in $\Gamma/J$ is contained in
$f\ne0$. This finite quotient comes from a finite etale
Galois cover of $U$. Its selected conjugacy class lies in
the geometric subgroup, because $\gamma$ does.

Apply finite-field Chebotarev with the constant-field coset
condition, as in C44-A12. More explicitly, if the finite
quotient cover has constant-field degree $\nu$, choose
sufficiently large extension degrees divisible by $\nu$;
the chosen geometric class is then in the permitted coset.
Entin, Theorem~3 in \S4, author arXiv:1805.05454v3,
pp.~9--11, supplies points over such extensions with
Frobenius in that class. At those points $f(F)\ne0$.
By the cyclotomic bound and (A16.2), their total
root-of-unity multiplicity is exactly one.

If this point is presented as rational over an extension,
its Frobenius is a positive power of the Frobenius of its
underlying closed point. Passing to a positive power
preserves the total multiplicity of root-of-unity
eigenvalues: roots of unity remain such, and a number
whose positive power is a root of unity was already one.
Thus the underlying closed point has the same multiplicity
one. The construction works on every nonempty open subset
of $U$, whose fundamental group surjects onto that of $U$;
the monodromy image is unchanged. The selected closed
points are therefore Zariski dense. Finally, rationality
of their polynomial forces its unique root of unity to
be rational, hence $1$ or $-1$.
\end{proof}

\begin{corollary}[C44-A16a: closed supersingular axes with positive spans of arbitrary rank]
Let $N\geq0$ be any integer. With $E:y^2=x^3-x$, $A=E^6$
and $H=\sum h_i$ as in C44-A13, choose $N$ distinct points
$a_i\in E^4(K)$ over a finite number field and put
$Y_i=E^2\times\{a_i\}$. For $N=0$ take $K=\mathbf Q$
and impose no centre. For every sufficiently large $t$
there is an incidence family on a nonempty open
$U\subset\mathbf P^1_K$ of smooth genuine Lefschetz
axes $B_u=D_0\cap D_u$, $D_0,D_u\in|tH|$, containing all
the $Y_i$, with the following property.

After finitely many rational primes have been excluded,
for every remaining $p\equiv3\pmod4$ and every place of
$K$ above $p$, a Zariski-dense set of closed points of the
reduced parameter curve has the required positive axis
sign on the \emph{whole} $A_{B_u}$, unconditionally. The
ambient sixfold is supersingular. The explicit classes
\[
   w_i=[Y_i]_{B_u}-t^{-2}\kappa^*\alpha_0
\]
span a positive $N$-dimensional rational subspace $K_u$
with Gram matrix $6t^2I_N-J_N/(4t^2)$, provided
$24t^4>N$. At the selected points:
\begin{itemize}
\item if $N$ is odd, $A_{B_u}=K_u$;
\item if $N$ is even, $A_{B_u}=K_u$ or
      $A_{B_u}=K_u\perp\mathbf Qz_u$, with $q_B(z_u,z_u)>0$.
\end{itemize}
The known span $K_u$ has relative algebraic lifts. In the
even-$N$ case, no assertion is made that the possible
additional Tate line is algebraic or that a class on it
lifts. For $N=0$ the conclusion is positivity of a kernel
of dimension zero or one; nonvanishing is not claimed.
\end{corollary}

\begin{proof}
The disjoint-centre construction of A13b applies for any
$N$ up to, and including, its monodromy calculation. To
make the lack of a parity assumption at those steps
explicit, take $Z=\operatorname{Bl}_{\coprod Y_i}A$,
$L'=\pi^*\mathcal O_A(tH)(-\sum F_i)$, choose a general
$\widetilde D_0\in|L'|$, and vary a general auxiliary
hyperplane pencil on its induced embedding. For $N=0$
this just means $Z=A$ and $L'=\mathcal O_A(tH)$.
The normal-row ranks four and the rank-drop codimension
three on each two-dimensional centre give the original
smooth axes, and the nonempty open of Lefschetz lines
through $\widetilde D_0$ imposes the original-pencil
condition as in A13b. These statements do not involve
the parity of $N$.

Every centre has $H^2\alpha_0=[Y_i]_A$, normal self-intersection
$[Y_i]^2_{B_u}=6t^2$, and $\int_{Y_i}\alpha_0=1/4$.
Distinct centres are disjoint. These equalities give
the displayed Gram matrix, positive for $24t^4>N$.
The blockwise blowup formula and excess formula give
\[
 \mathcal W_N=\ker\kappa_*\cap\bigcap_iw_i^\perp
       \simeq \hbox{the auxiliary-pencil vanishing system},
 \qquad
 \operatorname{rank}\mathcal W_N=3600t^6-165-N.
 \tag{A16.3}
\]
The isomorphism is isometric, as shown in A13b; disjointness
makes all exceptional blocks orthogonal. The formula remains
valid with an empty exceptional sum when $N=0$.
For every $N$ and sufficiently large $t$ the rank is positive
and
\[
                  h^{4,0}(\mathcal W_N)=62t^6-5>0.
\]
SGA~7~II, XIX, Proposition~3.4 and \emph{Weil II},
4.4.1--4.4.2s therefore give full orthogonal geometric
monodromy in characteristic zero, independently of parity.

Spread this fixed incidence model and its finite omitted
parameter divisor over $\mathcal O_K[1/M_0]$, retaining
smooth families, disjoint centres, and finite etale
boundary. Exclude the finitely many rational primes below
all bad places or collisions, together with $2$, $\ell$
and the prime divisors of
$|\operatorname{GL}_{3600t^6-165-N}(\mathbf F_\ell)|$.
The stable-integral-image argument of A13 and A13b, using
the prime-to-$p$ specialization isomorphism of SGA~1,
XIII, 2.9--2.12, then identifies the geometric image at
every remaining place with the characteristic-zero image.
It applies equally in odd and even ranks. Thus the full
orthogonal closure persists at every remaining place.
Infinitely many rational primes $p\equiv3\pmod4$ remain,
and $E$ is supersingular at all places above them.

The Frobenius polynomial on $\mathcal W_N$ is rational:
divide the normalized middle polynomial of $B_u$ by
the ambient factor from $A$ and by $(T-1)^N$ supplied
by the known relative cycles. If $N$ is odd, the rank
in (A16.3) is even and C44-A12 supplies dense closed
points with no residual root of unity, giving
$A_{B_u}=K_u$ exactly as in A13b.
If $N$ is even, that rank is odd. C44-A16 supplies
dense closed points with residual cyclotomic
multiplicity one. C44-A15, with the known positive
span $K_u$, proves positivity of the entire algebraic
kernel and the stated dimension alternatives. This
last application does not infer an algebraic cycle
from a Frobenius root.

The selected closed parameter lifts in the fixed open
incidence model after an unramified extension of a
completed local trait. The centres and the divisor
polynomial $\alpha_0$ lift with it, giving relative
algebraic representatives of the known $w_i$. This
proves exactly the asserted liftability of $K_u$.
No construction or lifting of a possible additional
class is part of the argument.
\end{proof}

\paragraph{Programme scope and source.}
This extends the whole-axis sign family for cote~44,
Theorem~5, scan~004 and its primitive margin on scan~008
to prescribed positive spans of every rank. The axis conclusion is
unconditional. As in A15, the corresponding full-height
and $I(B)$ conclusions through C44-A01 retain the separate
unrefereed 032/Milne2026 abelian input. The general fixed-axis
case is not proved. The finite-field selection uses Entin's
complete Theorem~3 and proof in \S4, read at
\url{https://arxiv.org/pdf/1805.05454v3}, pp.~9--11, with
the constant-field congruence condition explicitly retained.
All geometric and monodromy source locators are those
verified in A13 and A13b; the sign of a possible final
algebraic line is supplied by the independently checked A15.

% END work/scan_w/odd_tate_selection_axes.tex

% BEGIN work/axis/axis_ring_obstruction.tex
% C44_POS /root/axis_independent; GPT-6 Astra Pro; effort not exposed.
\section{The cohomology-ring boundary of the tensor lifting method}
\subsection*{C44-A14a. Orthogonal freedom of the variable middle summand}

\textbf{Proposition.}
Let $X$ be an abelian variety of dimension $g=2r+2$, $r\ge1$,
over an algebraically closed field, and let
$\kappa:B\hookrightarrow X$ be a smooth complete intersection of
two ample divisors of the same class $h$. Work in rational
$\ell$-adic cohomology, with $\ell$ different from the
characteristic, and suppress Tate twists. Put
\[
 V=H^1(X),\qquad
 W=\ker\bigl(\kappa_*:H^{2r}(B)\to H^{2r+4}(X)\bigr),\qquad
 q(w,w')=(-1)^r\int_Bww'.
                                                               \tag{A14a.1}
\]
The ambient image is the graded algebra
\[
 R=\operatorname{im}\kappa^*
       \simeq\frac{\bigwedge^\bullet V}{\operatorname{Ann}(h^2)},
 \qquad
 \operatorname{Ann}(h^2)=\{\alpha:\alpha h^2=0\}.              \tag{A14a.2}
\]
Every cohomology group of $B$ is ambient except for the direct
summand $W$ in degree $2r$. If $\omega_B$ is the top class with
$\int_B\omega_B=1$, all remaining multiplication is described by
\[
   w\alpha=0\quad(w\in W,\ \alpha\in R^{>0}),\qquad
   ww'=(-1)^r q(w,w')\omega_B.                               \tag{A14a.3}
\]
Let $\operatorname{Sp}(V,h)$ denote the stabilizer of the
nondegenerate degree-two exterior tensor $h$.
There is a faithful action
\[
 \operatorname{Sp}(V,h)\times\operatorname O(W,q)
          \lhook\joinrel\longrightarrow
  \operatorname{Aut}_{\rm gr}\bigl(H^\bullet(B),h|_B,\int_B\bigr)
                                                               \tag{A14a.4}
\]
under which the first factor acts on $R$ through exterior powers
and fixes $W$ pointwise, while the second factor fixes $R$ pointwise.
In particular, the infinitesimal stabilizer of the cup ring,
polarization and \emph{any collection} of variable classes
contains the nonzero Lie algebra $\mathfrak{sp}(V,h)$.

\textbf{Proof.}
The exterior algebra description of $H^\bullet(X)$ is Milne,
\emph{Abelian Varieties}, I Theorem 12.1, printed pp.~55--56.
Self-intersection gives
\[
                         \kappa_*\kappa^*\alpha=h^2\alpha.
                                                               \tag{A14a.5}
\]
For $j\le2r$ weak Lefschetz makes $\kappa^*:H^j(X)\to H^j(B)$
injective, and it is an isomorphism for $j<2r$.
Hard Lefschetz makes $L_h^2$ injective in all these degrees.
Consequently both kernels in (A14a.2) are zero there.

For $j>2r$, the Gysin map
\[
                    \kappa_*:H^j(B)\longrightarrow H^{j+4}(X)
\]
is an isomorphism: its Poincare dual is the weak Lefschetz
isomorphism in degree $4r-j<2r$. The map $L_h^2$ onto
$H^{j+4}(X)$ is surjective. One way to see the last assertion
is to apply hard Lefschetz in the dual degree; equivalently,
it is the upper-half surjectivity in each Lefschetz string.
Equation (A14a.5) therefore makes $\kappa^*$ surjective in
these degrees and identifies its kernel with
$\ker L_h^2$. This proves (A14a.2) in every degree. The
annihilator is an ideal, as required for the displayed algebra
quotient.

In degree $2r$, $L_h^2:H^{2r}(X)\to H^{2r+4}(X)$ is an
isomorphism. Thus self-intersection gives the orthogonal
decomposition
\[
                        H^{2r}(B)=R^{2r}\perp W.
                                                               \tag{A14a.6}
\]
The form on $R^{2r}$ is
$(-1)^r\int_X\alpha\beta h^2$, which is nondegenerate by
hard Lefschetz, so $W$ is nondegenerate too.

If $w\in W$ and $\alpha$ has positive ambient degree,
projection formula gives $\kappa_*(w\kappa^*\alpha)=0$.
The product is above the middle degree; the Gysin isomorphism
just proved makes it zero. Products past the top degree are
zero as well. The product of two degree-$2r$ variable classes
is determined by their top trace, proving (A14a.3).

An element of $\operatorname{Sp}(V,h)$ acts on the exterior
algebra, preserves $h^2$ and its annihilator, and therefore
acts on $R$. It preserves the top trace because
$h^g\ne0$ generates the one-dimensional top cohomology of $X$:
fixing $h$ fixes that generator. In particular it preserves
the form on $R^{2r}$ and fixes $\omega_B$.
Extend this action by the identity on $W$.
Similarly extend an isometry of $(W,q)$ by the identity on
$R$. Formulas (A14a.3) verify every mixed or variable product;
ambient products are already preserved. The two actions commute.
Their product is faithful because the first factor is
faithful in degree one and the second is faithful on $W$.
This proves (A14a.4).

Finally, any collection of elements of $W$ is fixed by the
entire first factor. Its Lie algebra has dimension
$g(2g+1)>0$. Thus the asserted infinitesimal stabilizer is
nonzero, for every choice of those elements.
\hfill$\square$

\paragraph{Exact consequence for the lifting attack.}
C44-A10 and C44-A14 obstruct lifting by exhibiting an
intrinsic fourfold cup tensor on an eight-dimensional
$H^1$ space whose infinitesimal stabilizer is zero.
The proposition proves that this criterion cannot be
recovered from the graded cup ring of an axis, its
polarization, and variable classes alone. Those data retain
the full ambient symplectic symmetry acting in degree one.
In particular, replacing the A14 centre by several centres,
or replacing its class by a difference with zero ambient
pushforward, cannot retain that same zero-stabilizer
certificate through ordinary cup products.

This is a limitation of that specific tensor method.
It does not assert that a variable class lifts, that a
formal Hodge structure on this algebra is realized
geometrically, or that the primitive intersection form is
positive. Constraints involving the period variation,
the Hodge numbers of a lift, algebraic correspondences
outside the cup ring, or integral and crystalline data
are additional information and are not addressed by
(A14a.4).

\paragraph{Programme locator.}
The variable summand is the exact $A_B$ cohomological
container in JAC1964 H9-T7, hence the remaining term in
cote~44, Theorem~5, scan~004, with the primitive margin
on scan~008. This proposition sharpens the obstruction
to route (a), the universal lifting attack, and to
cup-product versions of route (b). It supplies no
additional instance of the desired sign and does not
close QH9-axis.

% END work/axis/axis_ring_obstruction.tex

% BEGIN work/covers/contribution.tex
\section{Special cover strata and the old axis summand}

\subsection{C44-C01: the old variable classes do not disappear}
\begin{proposition}[C44-C01]\label{prop:old-axis-cover}
Use JAC1964 H9-T18. Write $A$ for the abelian fourfold, $B$ for its
surface axis, $X_0=\operatorname{Bl}_B A$, $M_0=A\times\mathbf P^1$,
and $k_0:X_0\hookrightarrow M_0$ for the incidence embedding.
For a finite separable cover $\pi:C\to\mathbf P^1$ of degree $b$,
unramified over the critical values, put
$X_C=X_0\times_{\mathbf P^1}C$, $M_C=A\times C$, and let $g:X_C\to X_0$
and $k_C:X_C\hookrightarrow M_C$ be the maps. On algebraic middle classes,
\[
 \Alg^2(X_C)\cap\ker k_{C*}
 =g^*(\Alg^2(X_0)\cap\ker k_{0*})\ \mathbin\perp\
   (\Alg^2(X_C)\cap\ker k_{C*}\cap\ker g_*).
 \tag{C44.C01.1}
\]
The first summand is canonically the old axis space
$A_B=\ker(\kappa_*:\Alg^1(B)\to H^6(A)(1))$ and has form
\[
          (w,w')\longmapsto-b\int_Bww'.                 \tag{C44.C01.2}
\]
It is positive definite, by the surface Hodge index theorem, without 032.
Consequently the remaining H9-T18 sign is exactly the sign on the second
summand in (C44.C01.1). Statements that generic covers add no algebraic
middle classes refer to this new part (or the full trace kernel), and do
not say that the entire H9-T18 variable space vanishes.
\end{proposition}
\begin{proof}
The finite flat map $g$ has $g_*g^*=b$. Projection formula makes
$P=b^{-1}g^*g_*$ a self-adjoint algebraic idempotent. The Cartesian
incidence square gives
$k_{C*}g^*=(1_A\times\pi)^*k_{0*}$, and proper functoriality gives
$k_{0*}g_*=(1_A\times\pi)_*k_{C*}$. Thus $P$ preserves
$\ker k_{C*}$, and its image there is exactly the first summand of
(C44.C01.1); its kernel is the displayed second summand. Both are
algebraic subspaces, since $P$ and $1-P$ are algebraic. Self-adjointness
proves orthogonality.

For the old space use the codimension-two blowup decomposition from
H9-T7. A middle class has unique form $q^*z+e_*p^*w$ with
$z\in H^4(A)$ and $w\in H^2(B)$. With
$t=[\mathrm{point}]\in H^2(\mathbf P^1)$, incidence pushforward is
\[
 k_{0*}(q^*z+e_*p^*w)=(hz+\kappa_*w)+zt.
\]
The exceptional divisor is $B\times\mathbf P^1$ as a variety over
$A\times\mathbf P^1$; this proves the second term's formula, while
the first follows from $[X_0]=h+t$. Kunneth independence forces $z=0$
and $\kappa_*w=0$ in the kernel. This is the asserted axis identification.
H9-T7 equation (H9.24) gives its old form $-\int_Bww'$. Pullback under
$g$ multiplies intersections by $b$, proving (C44.C01.2). Its strict
positivity is precisely the surface input used in H9-T8. No sign input
on $X_C$ has been used.
\end{proof}

\paragraph{Archive and predecessor implications.}
This sharpens the wording of the cover gate for cote 44, Th.~5,
scan 004 and JAC1964 QH9-cover-primitivep. It uses H9-T7, H9-T8,
H9-T18 and finite-flat push-pull; it does not re-prove their height
identification. Cote 16, scan 050, language is algebraic Gysin and
restriction with their correct Lefschetz weights. In Chow motives,
$b^{-1}g^*g_*$ is the old-part projector; the unresolved piece is its
complement inside the primitive middle projector.

\subsection{C44-C02: every nonextremal fixed splitting stratum of triple covers}
\begin{theorem}[C44-C02]\label{thm:maroni}
Fix the abelian pencil of H9-T23 and integers
\[
                  1\le e\le f<2e,
                  \qquad g=e+f-2\ge2.
\]
In every characteristic, let $U_{e,f}$ be the parameter space of smooth
connected triple covers of the fixed target $\mathbf P^1$, unramified
above the pencil's critical set, with framed Tschirnhausen bundle
$\mathcal O(e)\oplus\mathcal O(f)$. It is a nonempty irreducible open
subset of the space of twisted binary cubics. At its geometric generic
point, every algebraic class in the new middle trace kernel of the
pullback pencil is zero. The exact norm-height formula and radical are
\[
 h_L(z,w)=\tfrac13h_K(\operatorname{cor}z,\operatorname{cor}w),
 \qquad \operatorname{rad}h_L=
       \operatorname{cor}^{-1}(\operatorname{rad}h_K).       \tag{C44.C02.1}
\]
For abelian-fourfold pencils the full homological height is therefore
semipositive, conditionally on 032/Milne as in H9-T8. The vanishing
assertion and norm identity themselves do not use 032.

For $d=f-e>0$ the unframed stratum has codimension $d-1$ in the full
fixed-target smooth triple-cover stack. Thus this treats in particular
the first unbalanced stratum for even $g\ge6$ (a divisor), and for odd
$g\ge9$ (codimension two). It does not assert anything for every closed
point of these strata or for the extremal splitting $f=2e$.
\end{theorem}
\begin{proof}
For $E=\mathcal O(e)\oplus\mathcal O(f)$, Wood's cubic-algebra
equivalence, in the exact form already used in H9-T24 equations
(H9.136)--(H9.139), gives the affine coefficient space
\[
 P_E=H^0\bigl(\mathcal O(2e-f)\oplus\mathcal O(e)
                \oplus\mathcal O(f)\oplus\mathcal O(2f-e)\bigr).
\]
All degrees are nonnegative, so $\dim P_E=2g+8$. The smooth,
geometrically connected, and specified unramified conditions define an
open $U_E$. Each point has the required quotient bundle by the cubic
equivalence. It remains to show nonemptiness and the monodromy property.

Set
\[
             a=2e-f-1\ge0,\qquad c=f-e\ge0.
\]
Then $2a+3c=g$, $a+c+1=e$, and $a+2c+1=f$. Apply H9-T23 with
these particular parameters (its proof allows arbitrary nonnegative
$a,c$, not only the balanced choices used later in H9-T24).
In characteristic different from three, the Kummer normalization has
the relative eigensheaves
$\mathcal O\oplus\mathcal L^{-1}\oplus\mathcal L^{-2}(D_2)$,
where $\deg\mathcal L=e$ and $\deg D_2=a+1$. Its quotient dual is
$\mathcal O(e)\oplus\mathcal O(2e-a-1)=E$.
All exponent-two points and the line bundle are fixed over the moving
branch parameter, so this is an isomorphism of relative bundles.
In characteristic three, there are $a+1$ simple poles and $c$ double
poles. The exact normalized basis calculation (H9.144) gives quotient
dual degrees $a+c+1=e$ and $a+2c+1=f$, again with constant pole
divisors. After choosing a framing, in each case this constructs an
actual morphism from the H9-T23 parameter curve into $U_E$.

The universal pullback pencil is smooth proper, and
$1-\frac13Q^*Q_*$ defines a lisse middle-cohomology direct summand,
exactly as in H9-T24 equations (H9.146)--(H9.147). Its restriction to
the constructed curve has no nonzero finite-orbit vectors by H9-T23's
middle-convolution or Fourier calculation. A finite-orbit vector for
the fundamental group of $U_E$ would remain such after restriction to
this curve. Hence none exists on $U_E$. Every algebraic class at the
geometric generic point is defined over a finite extension of the
parameter function field and consequently has a finite monodromy
orbit. Its new projection must vanish. The actual-cycle push-pull and
admissibility argument in H9-T24 then gives (C44.C02.1) and the radical.
This uses precisely the old result, with the newly specified parameter
stratum, rather than a generic conclusion specialized to a closed point.

For codimension, $\dim\operatorname{Aut}(E)=d+3$ if $d>0$ and is
$4$ if $d=0$: compute $H^0(\End E)$ from its four line summands.
Thus the $d>0$ stratum has dimension $2g+5-d$, while the balanced
stratum has dimension $2g+4$. Smooth separable covers have finite
automorphism groups, so the quotient-stack dimension calculation is
valid also in characteristics two and three. The difference is $d-1$.
The first nonbalanced value is $d=2$ for even $g$ and $d=3$ for odd
$g$; $f<2e$ gives exactly the genus bounds stated above.
\end{proof}

\paragraph{Sources and scope.}
Wood, \emph{Parametrizing quartic algebras over an arbitrary base},
Theorem 2.1 and equation (1), printed pp.~1073--1074; JAC1964 H9-T23,
(H9.127)--(H9.132), and H9-T24, (H9.136)--(H9.147). These are the
precise predecessor inputs. The new conclusion concerns geometric
generic points of special Maroni strata, over algebraic closures of
their parameter fields. It is not a proof for all special covers over
$\overline{\F}_p$. Archive target: cote 44, scan 004, Th.~5 and its
fibration margin; predecessor row QH9-cover-primitivep is narrowed in
this specified stratum scope, not closed.

% END work/covers/contribution.tex

% BEGIN work/covers/boundary_strata.tex
% C44_POS /root; candidates undergoing independent whole-proof review.
\subsection{C44-C03: characteristic-two conductor strata with a simple pole}
\begin{theorem}[C44-C03]\label{thm:AS-strata}
In characteristic two, fix an Artin--Schreier double-cover pole partition
\[
 d_i=2a_i+1\quad(1\le i\le s),\qquad
 s\ge2,\quad a_i\ge0,\quad a_1=0,
 \qquad g=\sum_i(a_i+1)-1\ge2.
\]
Fix an abelian-fourfold Lefschetz pencil as in H9-T22, and require the
branch locus of the covers to avoid its critical values. At the geometric
generic point of the indicated irreducible conductor stratum, there is
no algebraic class in the new middle trace kernel of the pulled-back
pencil. Consequently
\[
 h_L(z,w)=\tfrac12h_K(\operatorname{cor}z,\operatorname{cor}w),
 \qquad \operatorname{rad}h_L=
 \operatorname{cor}^{-1}(\operatorname{rad}h_K).
 \tag{C44.C03.1}
\]
The identity and vanishing are unconditional. Semipositivity of the
resulting height uses H9-T8 and is conditional on the unrefereed
032/Milne inputs. The assertion includes the geometric generic point
of the hyperelliptic $2$-rank $g-1$ divisor for every $g\ge2$.
It is a statement at geometric generic points, not at every closed
point of the stratum.
\end{theorem}
\begin{proof}
Order the poles and choose a coordinate for which zero and infinity
are smooth values of the fixed pencil. We construct one witness
family with its distinguished simple pole fixed at $t=0$; this
fixed-position family is not asserted to be dense in the stratum.
Put the other
poles at distinct nonzero $b_2,\ldots,b_s$, away from the critical set
$\Sigma$ of the pencil. Choose a rational function $f_0$, regular at zero
and infinity, with exactly these poles and orders $d_2,\ldots,d_s$.
The family
\[
                  w^2-w=f_0(t)+\lambda/t,
                  \qquad \lambda\in\mathbf G_m
 \tag{C44.C03.2}
\]
lies in the specified conductor stratum. Its normalization is smooth
and proper over the parameter. Here is the required local verification
also for a higher pole. If $z$ is a local parameter at a pole of order
$2a+1$, set $v=z^{a+1}w$. The integral local equation becomes
\[
 v^2-z^{a+1}v=z^{2a+2}f(z).
\]
The right side has first term $c z$, with $c$ a unit on the parameter
space. At the unique point $z=v=0$ its derivative in $z$ is $c\ne0$.
The equation is regular there, finite flat of degree two, and agrees
with the extension off the pole, so it is the normalization. Away from
the poles the derivative in $w$ is a unit. This proves the relative
smoothness and justifies smooth proper base change.

Let $V$ be the variable odd local system of the original pencil,
of rank $r>0$, with $q=\#\Sigma$. Its irreducibility and rank-one tame
transvections are the inputs used in H9-T21--T22. Change coordinates
to $x=1/t$, and put
\[
 G=\mathcal L_\psi(f_0(1/x)),\qquad
 F=(\operatorname{inv}^*V)\otimes G.
\]
Their middle extensions to $\mathbf A^1_x$ are irreducible Fourier
sheaves: $G$ has a finite wild pole, and $F$ is irreducible of rank
$r\ge2$. Both are unramified at infinity. The exact character-part
Leray calculation of H9-T22 identifies the new middle cohomology
with
\[
 H^3(A)\otimes\operatorname{NFT}_\psi(G)
       \ \oplus\ \operatorname{NFT}_\psi(F)
 \tag{C44.C03.3}
\]
over the parameter open. This is a calculation of the actual
smooth-proper cohomology, not a substitution of extension by zero at
the old transvections. The outer Leray terms are deck invariant; the
two displayed terms are its nontrivial deck character.

Set $R=\sum_{i=2}^s(d_i+1)$. The Fourier ranks are $R$ and $q+rR$,
by the finite conductor formula, Katz, \emph{Exponential Sums and
Differential Equations}, Lemma 7.3.9(2). Theorem 7.3.8(3) gives
irreducibility of the transforms. At parameter zero, Corollary
7.4.3.1 and Theorem 7.4.4(2)--(3), equivalently Theorem 7.5.4(5)--(7),
give their respective local representations
\[
 J(1,2)\oplus J(1,1)^{R-2},\qquad
 J(1,2)^{r}\oplus J(1,1)^{q+rR-2r}.
 \tag{C44.C03.4}
\]
Each is infinite because it contains a nontrivial unipotent block.
An irreducible representation with infinite image has no nonzero
finite-orbit vector: the span of the orbit would be a nonzero
subrepresentation, and the normal core of its finite-index stabilizer
would then act trivially on the whole representation. This would make
its image finite. Direct sums and tensoring with a constant vector
space preserve the absence of finite-orbit vectors. Thus (C44.C03.3)
has none.

For completeness, the space with ordered pole positions and all
reduced odd principal-part coefficients is irreducible: it is an open
in a vector bundle over the ordered configuration space, with each
leading coefficient nonzero. The above local normalization works
throughout it. It maps dominantly to the indicated conductor stratum
of covers. Its universal trace-kernel local system restricts to
(C44.C03.3). A finite-orbit vector for the universal monodromy would
remain one on this parameter curve. There are therefore no such
vectors at the geometric generic point of the stratum. Algebraic
classes, defined over finite extensions of its function field, have
finite monodromy orbit. The actual trace projector and admissibility
argument of H9-T22 and H9-T24 now give (C44.C03.1).

The genus and $2$-rank are $g$ and $s-1$. Pries--Zhu,
\emph{The $p$-rank stratification of Artin--Schreier curves},
Theorem 1.1 and Corollary 1.3, printed p.~708, identify the components
in characteristic two with partitions of $g+1$ into $s$ positive
parts $a_i+1$, all of dimension $g-1+(s-1)$. For rank $g-1$ the
partition is $(2,1,\ldots,1)$, with $g-1$ parts equal to one.
It has a simple pole and is a divisor in the hyperelliptic moduli
space of dimension $2g-1$. Fixing the target adds three dimensions
to both spaces and leaves codimension unchanged. This proves the
last assertion.
\end{proof}

\subsection{C44-C04: extremal triple-cover splitting strata}
\begin{theorem}[C44-C04]\label{thm:extremal-maroni}
Let $e\ge2$ and $g=3e-2$. In every characteristic, let $U_{e,2e}$
denote the framed space of smooth connected triple covers of the
fixed target $\mathbf P^1$ with Tschirnhausen bundle
$\mathcal O(e)\oplus\mathcal O(2e)$, unramified over the old
critical set. At its geometric generic point the new middle
trace kernel of an abelian-fourfold pencil contains no algebraic
classes. The norm-height identity and radical in (C44.C02.1) hold.
These assertions are unconditional; height semipositivity is
conditional on the unrefereed 032/Milne inputs through H9-T8.

Together with C44-C02 this settles the geometric generic point
of every admissible triple-cover splitting stratum of genus at
least two. It does not assert the same vanishing at every closed
point over $\overline{\mathbf F}_p$.
\end{theorem}
\begin{proof}
The cubic coefficient space for $E=\mathcal O(e)\oplus\mathcal O(2e)$
is $H^0(\mathcal O\oplus\mathcal O(e)\oplus\mathcal O(2e)
\oplus\mathcal O(3e))$. Its smooth connected unramified locus is
irreducible if nonempty. We construct a family in it for each
characteristic and prove that its new trace-kernel monodromy has
no finite-orbit vectors. The restriction argument of C44-C02
will then prove the assertion for this whole stratum.

\paragraph{Characteristic different from three.}
Choose a coordinate $t$ with zero and infinity outside the old critical
set $\Sigma$. Choose $3e-3$ nonzero points $b_i$ so that their
$\mu_3$-orbits are pairwise disjoint and avoid the $\mu_3$-orbits of
$\Sigma$. Put $f(t)=\prod_{i=1}^{3e-3}(t-b_i)$ and consider
\[
                   u^3=(t^3-s)f(t).
 \tag{C44.C04.1}
\]
Remove zero and the finitely many parameters for which a moving root
hits a fixed branch point or a critical value. The normalized covers
are smooth and have $3e$ simple branch points, all with exponent one;
infinity is unramified. Their genus is $3e-2$, and the relative
eigensheaf decomposition is
\[
        \pi_*\mathcal O_C=
        \mathcal O\oplus\mathcal O(-e)\oplus\mathcal O(-2e).
 \tag{C44.C04.2}
\]
The normalized local equations are the tame equations of H9-T23.
Because the degree of the branch polynomial is $3e$, this is a
relative decomposition with fixed line bundles, not only a
fibrewise numerical calculation. It supplies a framing into $U_{e,2e}$.

Fix a nontrivial cubic character $\chi$. On the $t$ line let
$G_\chi=\mathcal L_\chi(f(t))$ and
$F_\chi=V\otimes G_\chi$, always taking middle extensions.
For $\rho:\mathbf A^1_t\to\mathbf A^1_v$, $v=t^3$, set
\[
                  M_\chi=\rho_*G_\chi,
                  N_\chi=\rho_*F_\chi.
\]
Both are irreducible middle extensions. Indeed, restrict to the
$\mu_3$-stable open on which $\rho$ is a Galois cover. Each input
is irreducible. Its three conjugates are pairwise nonisomorphic:
at a fixed branch point one has nontrivial scalar inertia $\chi$,
whereas either conjugate is lisse, by the choice of the $b_i$.
The elementary induction criterion for a normal subgroup with
pairwise inequivalent conjugates therefore makes its induced
representation irreducible. Finite pushforward is its middle
extension, as is checked by inertia invariants at the omitted points.

Finite pushforward and projection formula identify the two new
character terms with
\[
          H^3(A)\otimes\operatorname{MC}_\chi(M_\chi)
                    \ \oplus\ 
                    \operatorname{MC}_\chi(N_\chi).
 \tag{C44.C04.3}
\]
This follows from the actual H9-T23 character cohomology by replacing
$\mathcal L_\chi(t-s)$ there with
$\mathcal L_\chi(t^3-s)$ and then applying $\rho_*$. The fibre
cohomology is the middle convolution, including its local invariant
conditions. Both input sheaves are tame. At infinity their
representations are respectively
\[
                 1\oplus\chi\oplus\chi^2,
                (1\oplus\chi\oplus\chi^2)^{r}.
 \tag{C44.C04.4}
\]
Here the original $V$ and $G_\chi$ are unramified at infinity,
and the ramification of $t\mapsto t^3$ supplies the regular
representation of $\mu_3$.

Katz, \emph{Rigid Local Systems}, Theorem 3.3.3 and Corollary 3.3.6
(with Corollary 4.3.11 for the characteristic-zero comparison,
as already used in H9-T23),
give irreducibility of the convolutions and their local invariant
quotient formula. In that formula write
$M_\infty/M_\infty^{I_\infty}$ for the input representation.
Each trivial one-dimensional summand in (C44.C04.4) forces a
$J(1,2)$ block in $M_\infty$. Tensoring by the nontrivial $\chi$
turns it into $J(\chi,2)$, and quotienting by inertia invariants
does not remove it. Thus both outputs have infinite local image.
For a rank check, the first input has finite drops $2$ at zero
and $1$ at each of the $3e-3$ fixed branch images; Corollary
3.3.7 gives output rank $3e-2$. The second has finite drops
$q+(3e-1)r$ in total and output rank $q+(3e-2)r$.
Collisions among old critical images do not change their total
drop; the chosen branch images are disjoint from them.
Irreducibility and infinite image exclude finite-orbit vectors,
as in C44-C03.

\paragraph{Characteristic three.}
Choose a coordinate $t$ with zero and infinity outside $\Sigma$.
Choose $e-1$ distinct nonzero points $b_j$, away from $\Sigma$, so
that the reciprocal points $b_j^{-1}$ and their negatives are
distinct and avoid the reciprocal critical set and its negatives.
For nonzero constants $c_j$, put
\[
 f_0(t)=\sum_{j=1}^{e-1}\frac{c_j}{(t-b_j)^2},
 \qquad w^3-w=f_0(t)+s/t^2,\qquad s\in\mathbf G_m.
 \tag{C44.C04.5}
\]
There are $e$ poles, all of order two, hence genus $3e-2$.
The local order-two normalization and the global basis calculation
of H9-T23--T24, (H9.131) and (H9.144), give a smooth family and
the relative splitting (C44.C04.2). Its pole divisor is fixed,
so once again it maps to the framed stratum.

Put $x=1/t$. For either nonzero additive character index
$\alpha\in\mathbf F_3^\times$, set
\[
 G_\alpha=\mathcal L_\psi(\alpha f_0(1/x)),\qquad
 F_\alpha=(\operatorname{inv}^*V)\otimes G_\alpha,
 \qquad \rho:x\longmapsto v=x^2.
\]
The two finite pushforwards $\rho_*G_\alpha$ and
$\rho_*F_\alpha$ are irreducible by the same induction argument,
now for the tame quadratic cover and its involution $x\mapsto-x$.
A finite wild pole at $b_j^{-1}$ has no pole in its negative orbit,
so the two conjugates cannot be isomorphic. Their ranks are $2$
and $2r$, so neither is an exceptional rank-one additive character.
They are Fourier sheaves. At infinity their representations are
$1\oplus\epsilon$ and $(1\oplus\epsilon)^r$, respectively,
where $\epsilon$ is the nontrivial quadratic tame character.

Projection formula converts the variable term
$\mathcal L_\psi(\alpha s x^2)$ into
$\mathcal L_\psi(\alpha s v)$. The actual new character
cohomology is consequently the sum of $H^3(A)$ tensored with
$\operatorname{NFT}_\psi(\rho_*G_\alpha)$ and
$\operatorname{NFT}_\psi(\rho_*F_\alpha)$, pulled back by
$s\mapsto\alpha s$. Katz, \emph{Exponential Sums and Differential
Equations}, Theorem 7.3.8(3), gives irreducibility. Corollary
7.4.3.1 and Theorem 7.4.4(2)--(3) show that the trivial infinity
summands produce $J(1,2)$ blocks at Fourier parameter zero.
The nontrivial quadratic summands remain nontrivial rank-one
blocks and do not cancel these unipotent blocks. Hence these
transforms also have infinite image and no finite-orbit vectors.

In both characteristics every new character has now been handled.
The direct sum has no finite-orbit vectors. The universal lisse
trace-kernel projector on $U_{e,2e}$ restricts to this sum, so
the restriction argument proves generic algebraic vanishing on
the stratum. The norm-height identity and radical follow by
the same actual-cycle push-pull as in C44-C02. All smooth connected
triple-cover splittings have $e\le f\le2e$; the inequality follows
also directly from the cubic equation, since if $f>2e$ its first
coefficient is zero and the cubic has a linear factor, contradicting
smooth connectedness. C44-C02 handles the strict inequality, while
the present result handles equality for all genus at least two.
\end{proof}

\paragraph{Programme location and residual cover question.}
C44-C03--C04 extend the monodromy method of JAC1964 H9-T21--T24 to
specified special strata. Their height meaning is cote 44, Th.~5,
scan 004, with the primitive qualification of scan 008; the
idempotent $1-b^{-1}g^*g_*$ is the new Chow-motive projector.
The argument concerns lisse sheaves and finite monodromy orbits,
using SGA~7~II, XVIII~6.7.2 for the original pencil variation as
in the predecessor proofs. It gives no positivity theorem for an
extra Tate class at a closed point. Such a class belongs to the
remaining H9-T18 primitive incidence gate, after C44-C01 removes
the positive old surface-axis summand.

% END work/covers/boundary_strata.tex

% BEGIN work/covers/as_conductor_strata.tex
\section{Every Artin--Schreier conductor partition in characteristic two}
\label{sec:c44-c05}

\begin{theorem}[C44-C05: generic variable-kernel vanishing on every conductor stratum]
Work over an algebraically closed field \(k\) of characteristic \(2\), in the
abelian-fourfold Lefschetz-pencil setup of C44-C01 and C44-C03.
Write \(f:\widetilde A\to\mathbf P^1\) for the blowup pencil, \(\Sigma\)
for its critical-value set, and \(\mathcal V\) for the geometrically
irreducible variable summand of \(R^3f_*\mathbf Q_\ell\), where
\(\ell\ne2\). At each point of \(\Sigma\) its local monodromy is a
nontrivial tame transvection; at smooth pencil values it is unramified.
If \(\mathcal V=0\), omit the variable monodromy argument below.

Fix odd positive integers \(d_1,\ldots,d_s\), with
\(g=-1+\frac12\sum_i(d_i+1)\ge2\). Let \(\mathcal H_{\mathbf d}\) be the
irreducible Artin--Schreier conductor stratum of connected degree-two
covers \(\pi:C\to\mathbf P^1\) with these pole orders, restricted to
covers branched at smooth pencil values. At its geometric generic point,
put \(X_C=\widetilde A\times_{\mathbf P^1}C\), with incidence embedding
\(k_C:X_C\hookrightarrow A\times C\). A superscript \({\rm new}\)
denotes the anti-invariant summand for the deck involution. Then
\[
 \Alg^2(X_C)\cap
 \bigl(\ker(k_{C*}:H^4(X_C)\to H^6(A\times C))\bigr)^{\rm new}=0.
\]
Consequently the intersection form on the algebraic incidence kernel
is positive definite: only the old surface-axis contribution of
C44-C01 remains. This kernel-sign assertion is unconditional.

If \(s\ge2\), the stronger assertion
\(\Alg^2(X_C)^{\rm new}=0\) holds. The degree-two norm-height and
radical identities of C44-C03 consequently apply on the same domains,
with factor \(1/2\).

If \(s=1\), no vanishing of the ambient new summand is asserted and
no degree-two norm-height identity is inferred. The height
semipositivity conclusion on the Albanese-primitive domain follows
from the preceding kernel assertion and H9-T18 only
\emph{conditionally on the unrefereed preprints 032/Milne2026}.
\end{theorem}

\begin{proof}
Choose a smooth pencil value as infinity. Put \(d=d_1\).
When \(\mathcal V\ne0\), choose a distinguished branch point
\(a\notin\Sigma\) such that
\[
 \Sigma_a=\{(\sigma-a)^{-1}:\sigma\in\Sigma\}\subset\mathbf G_m
\]
has trivial stabilizer under \(\mu_d\). Such \(a\) exist: an equality
\(\zeta\Sigma_a=\Sigma_a\), with \(\zeta\ne1\), induces the homothety
\(t\mapsto a+\zeta^{-1}(t-a)\) permuting \(\Sigma\).
For each \(\zeta\) and permutation, the equation for one point
determines \(a\), so only finitely many values are excluded.
If \(\mathcal V=0\), choose any smooth \(a\) and omit this condition.
Set \(x=(t-a)^{-1}\) and \(\rho(x)=x^d\); this map is tame since \(d\)
is odd. Choose the other branch points so their \(x\)-coordinates
have pairwise disjoint \(\mu_d\)-orbits, disjoint also from
\(\mu_d\Sigma_a\).

Choose \(f_0(t)\) with the prescribed odd pole orders
\(d_2,\ldots,d_s\) at the other branch points, regular at \(a\)
and infinity. For \(s=1\), take \(f_0=0\). Consider
\[
       w^2-w=f_0(t)+\frac{\lambda}{(t-a)^d},
       \qquad \lambda\in\mathbf G_m.                    \tag{C05.1}
\]
Every fibre has exactly the required partition. Its smooth projective
normalizations form a smooth family, with the usual restriction to
the parameter open needed for the cohomological construction.
Indeed at a pole of order \(2b+1\), write the right side as
\(c(t)t^{-(2b+1)}\), \(c(0)\ne0\), and put \(v=t^{b+1}w\).
The equation
\(v^2-t^{b+1}v=t c(t)\) is smooth at the ramified point because its
\(t\)-derivative there is \(-c(0)\ne0\). This is uniform while
the leading coefficient is a unit. Away from the poles the equation
is etale. Branching avoids \(\Sigma\), so the incidence fourfolds
obtained by base change are smooth too.

Put \(\mathcal G=\mathcal L_\psi(f_0(a+1/x))\) and
\(\mathcal F=\mathcal V_a\otimes\mathcal G\), where
\(\mathcal V_a\) is the transformed pencil local system.
Both are unramified at \(x=\infty\).
Take middle extensions and set
\(\mathcal N=\rho_*\mathcal F\) on the affine \(y\)-line.
The restriction of \(\mathcal F\) to its lisse open is irreducible.
Its different deck translates are nonisomorphic: the set where the
local monodromy has the pencil's nontrivial tame unipotent part is
exactly \(\Sigma_a\), and no nontrivial deck transformation preserves
it. The scalar wild ramification at the separate other poles cannot
remove or imitate those tame transvections.
Mackey's criterion for induction along the Galois cover
\(x\mapsto x^d\) therefore gives geometric irreducibility of
\(\mathcal N\); for \(d=1\) this is immediate.

At \(y=\infty\), the local representation of \(\mathcal N\) is the
regular representation of \(\mu_d\), repeated
\(r=\operatorname{rank}\mathcal V\) times. Its \(r\) trivial tame
blocks of length one become \(r\) unipotent blocks of length two at
Fourier parameter \(\lambda=0\). The Fourier middle extension stays
irreducible. Its input has rank \(dr\ge2\), since the nonzero variable
odd cohomology is symplectic, so it is not an exceptional punctual
transform. The displayed unipotent blocks also witness its nonzero
lisse part and infinite geometric monodromy.

The Fourier cohomology is exactly
\[
 H^1\!\left(\mathbf P^1,
   j_*(\mathcal V\otimes
      \mathcal L_\psi(f_0(t)+\lambda/(t-a)^d))\right).     \tag{C05.2}
\]
Finite pushforward by \(\rho\), followed by tensoring with
\(\mathcal L_\psi(\lambda y)\), gives the Fourier integral.
For \(\lambda\ne0\) the sheaf is wild at \(x=\infty\), with no
boundary invariants there, identifying the affine Fourier cohomology
with projective middle-extension cohomology. Finite pushforward
preserves the middle extensions at the other singular points.

An irreducible characteristic-zero representation of infinite image
has no nonzero finite-orbit vector. Otherwise the invariants of the
core of a finite-index vector stabilizer would be a nonzero invariant
subspace, hence the entire representation, forcing finite image.
Thus the variable family (C05.2) has no such vector.

If \(s\ge2\), apply the same argument to
\(\mathcal M=\rho_*\mathcal G\). Its wild pole set is nonempty
and its chosen disjoint \(\mu_d\)-orbits distinguish all deck
translates, so \(\mathcal M\) is irreducible. Its infinity inertia
is the regular representation of \(\mu_d\); Fourier therefore has
a nontrivial length-two unipotent block at zero.
For \(d=1\), the rank-one input has a finite wild pole and is
unramified at infinity, excluding both the constant sheaf and the
linear Artin--Schreier exceptional input. Consequently
\[
 H^1\!\left(\mathbf P^1,
 j_*\mathcal L_\psi(f_0(t)+\lambda/(t-a)^d)\right)
       =H^1(C_\lambda)^{\rm new}                         \tag{C05.3}
\]
has no nonzero finite-orbit vector when \(s\ge2\).
Tensoring with the geometrically constant \(H^3(A)\) preserves
this assertion.

The new degree-four Leray cohomology splits as
\[
 H^4(X_C)^{\rm new}
 =\bigl(H^3(A)\otimes H^1(C)^{\rm new}\bigr)
 \oplus H^1(\mathbf P^1,j_*(\mathcal V\otimes\mathcal L_\pi)).
                                                            \tag{C05.4}
\]
The fibrewise projector
\(\operatorname{res}\circ L_h^{-1}\circ\operatorname{Gysin}\)
canonically splits \(R^3\) into the constant \(H^3(A)\) and
\(\mathcal V\) pieces. The outer Leray terms are deck-invariant,
so only \(E_2^{1,3}\) occurs in the nontrivial deck character.
Thus (C05.4) uses no arbitrary choice of Leray splitting.
Its first summand is \(k_C^*H^4(A\times C)^{\rm new}\).
The incidence divisor has ample class
\(p_A^*h+p_C^*\pi^*c_1\mathcal O_{\mathbf P^1}(1)\).
Hence \(k_{C*}k_C^*:H^4(A\times C)\to H^6(A\times C)\)
is its hard-Lefschetz isomorphism. On the new ambient part it is
\(L_h:H^3(A)\to H^5(A)\) tensored with \(H^1(C)^{\rm new}\),
so is injective. The second summand is killed by Gysin:
the fibrewise Gysin kills variable \(H^3\), and there are no
other new Leray terms to receive its image. It is therefore
exactly \((\ker k_{C*})^{\rm new}\).

The partition stratum is irreducible. Work if necessary on its
normal-form coefficient parameter space with ordered distinct
branch points, which dominates the stratum and carries the family.
The local systems restrict along (C05.1) to those just considered.
A finite-orbit vector for the full parameter monodromy would retain
a finite orbit on this witness curve, so none exists in the variable
system, or, when \(s\ge2\), in the full new degree-four system.
A geometric-generic algebraic class is defined after a finite
extension of the parameter function field. It thus has finite orbit
under its absolute Galois group, which surjects onto the fundamental
group of the smooth coefficient space. Purely inseparable extensions
do not affect this argument. Its projection to these summands must
vanish, proving the two vanishing assertions.

C44-C01 identifies the old incidence kernel with the surface-axis
divisor contribution, whose required sign is the Hodge index theorem.
Explicitly, if \(q:X_C\to\widetilde A\) is the degree-two map,
\(z_{\rm old}=\frac12q^*q_*z\). Proper functoriality keeps \(q_*z\)
in the old incidence kernel, and pullback multiplies the form by
the positive degree \(2\).

For \(s=1\), the equation \(w^2-w=\lambda x^d\) becomes a constant
curve after adjoining a \(d\)-th root of \(\lambda\) and rescaling
\(x\). The scalar system can thus have finite monodromy and its
ambient contribution cannot be discarded. This leaves the kernel
identification unaffected, but prevents full-new vanishing and
the associated norm identity. H9-T18 and the kernel sign yield
the stated Albanese-primitive height conclusion only under the
explicitly stated unrefereed 032/Milne2026 hypotheses.
\end{proof}

\paragraph{Source and scope audit.}
The Fourier inputs are Katz, \emph{Exponential Sums and Differential
Equations}, 7.3.2, 7.3.8(3), 7.4.3.1, and 7.4.4(2)--(3),
equivalently 7.5.4(7); they impose no exclusion of characteristic two.
Conductor-stratum irreducibility, including the one-pole,
\(p\)-rank-zero case, is Pries--Zhu, Proposition 3.6, Remark 3.9
and Theorem 3.10, printed pp.~716--717.
These are geometric-generic statements on every stratum, not
assertions at every closed finite-field point.
For cote 44, scans 004 and 008, this extends the H9-T18 primitive
height criterion and the H9-T22 monodromy calculation. It narrows
QH9-cover-primitivep to further special points inside these strata,
without closing that original pointwise question.

% END work/covers/as_conductor_strata.tex

% BEGIN work/covers/separable_cover_lifts.tex
% C44_POS /root/axis_independent; GPT-6 Astra Pro; effort not exposed.
% Root's separable-cover deformation idea and incidence determinant extension.
\section{Lifting separable covers and a fixed-incidence sign criterion}
\subsection*{C44-C06. Every separable map of a smooth curve to the line lifts}

\textbf{Theorem.}
Let $k$ be a perfect field of positive characteristic, let $C/k$
be a smooth projective geometrically connected curve, and let
$f:C\to\mathbf P^1_k$ be a finite separable morphism of degree $d$.
There is a smooth projective curve $\mathcal C/W(k)$ and a finite
flat morphism
\[
                   \widetilde f:\mathcal C\longrightarrow
                                    \mathbf P^1_{W(k)}
                                                               \tag{C06.1}
\]
of degree $d$ reducing to the given map.
No Galois action or branch configuration is required to lift.

For the incidence consequence assume $k$ is finite.
Let $A/k$ be an abelian fourfold, let $L$ be ample,
and let $s_0,s_1\in H^0(A,L)$ span a pencil.
Write $M=f^*\mathcal O_{\mathbf P^1}(1)$ and let
$u,v\in H^0(C,M)$ be the pullbacks of the two homogeneous
coordinates. Suppose the incidence fourfold
\[
 Y=\{s_0\otimes u+s_1\otimes v=0\}\subset A\times C
                                                               \tag{C06.2}
\]
is smooth. After a finite residue extension,
there is an embedded smooth projective mixed-characteristic
lift of $Y\subset A\times C$, including its defining ample
divisor class. Thus every smooth incidence fourfold of this
form has a geometric lift, including those formed with wild
separable covers.

\textbf{Proof.}
Set $W_n=W(k)/p^n$. We construct compatible smooth proper
curves $C_n/W_n$ and maps $f_n:C_n\to\mathbf P^1_{W_n}$.
The initial pair is $(C,f)$. Suppose it is constructed over
$W_n$, and put $I=p^nW_{n+1}$, a one-dimensional $k$-space.
The small extension $W_{n+1}\to W_n$ has square-zero kernel $I$.

The obstruction to lifting the smooth curve $C_n$ lies in
$H^2(C,T_C)\otimes_k I$, which is zero. Choose one smooth lift
$C'_{n+1}$. It is proper because properness is unchanged under
a nilpotent thickening. The isomorphism classes of such lifts,
with their specified reduction, form a torsor under
$H^1(C,T_C)\otimes_k I$. These are SGA~1, Expos\'e III,
6.9--6.10; the underlying-curve algebraization appears in
III 7.2--7.4. We will instead algebraize the finite maps
directly after constructing them.

On affine opens of $C'_{n+1}$ the map $f_n$ lifts because the
target $\mathbf P^1_{W_{n+1}}$ is smooth. The differences of
local lifts on overlaps form a Cech cocycle, with obstruction
class
\[
                       \beta\in H^1(C,f^*T_{\mathbf P^1})\otimes I.
                                                               \tag{C06.3}
\]
This is the morphism-lifting obstruction of SGA~1 III 5.1--5.6.
Changing the lift of $C_n$ by a class
$\alpha\in H^1(C,T_C)\otimes I$ changes $\beta$ by
$df(\alpha)$, up to the harmless choice of sign for the
Cech cocycles. Indeed a gluing automorphism $1+\alpha_{ij}$
changes the difference of the two local map lifts by
composition with the differential $df$.

Since $f$ is separable, the nonzero map of line bundles
\[
       0\longrightarrow T_C\xrightarrow{df}f^*T_{\mathbf P^1}
                         \longrightarrow N_f\longrightarrow0
                                                               \tag{C06.4}
\]
is injective and its cokernel $N_f$ is supported at finitely
many ramification points. This includes wild ramification;
the torsion sheaf need not be reduced.
Hence $H^1(C,N_f)=0$, and the long exact sequence makes
\[
              H^1(C,T_C)\longrightarrow H^1(C,f^*T_{\mathbf P^1})
\]
surjective. Choose $\alpha$ to cancel (C06.3), modify the
source lift by that class, and glue the local morphisms.
This constructs $(C_{n+1},f_{n+1})$ with the required reduction.
Induction produces the compatible system for every $n$.

Each $f_n$ is finite: it is proper, and its underlying
topological map is the finite special-fibre map, so it is
quasi-finite. It is flat as well. The special map between
smooth curves is finite flat. Inductively, for the first-order
thickenings $C_n\subset C_{n+1}$ and
$\mathbf P^1_{W_n}\subset\mathbf P^1_{W_{n+1}}$, flatness
over $W_{n+1}$ identifies their ideals with their respective
special-fibre structure sheaves tensored with $I$.
The natural conormal map for $f_{n+1}$ is therefore an
isomorphism. Stacks, Lemma 37.10.1, Tag 063Y, proves its
flatness. Thus $(f_n)_*\mathcal O_{C_n}$ is locally free of
rank $d$ on $\mathbf P^1_{W_n}$. No miracle-flatness argument
over the nonregular Artin base is being used.

Stacks, Lemma 30.28.2, Tag 09ZT, algebraizes the compatible
finite maps to a finite morphism
$\widetilde f:\mathcal C\to\mathbf P^1_{W(k)}$.
Explicitly, Grothendieck existence gives the coherent sheaf
from the compatible finite locally free algebras
$(f_n)_*\mathcal O_{C_n}$; full faithfulness algebraizes
their multiplication and unit, and taking the relative
spectrum gives the map. Its completed local modules along
the special fibre are flat of rank $d$, so the algebraized
module is locally free there. The non-locally-free locus
is closed and proper over $W(k)$ and misses the special
fibre; it must therefore be empty. Thus the map is finite
flat of degree $d$ and $\mathcal C$ is flat over $W(k)$.

The special fibre of $\mathcal C$ is smooth. The nonsmooth
locus of $\mathcal C/W(k)$ is closed, and is proper over
the base because $\mathcal C$ is finite over the projective
line. As it misses the special fibre it is empty.
The source is projective: the pullback of
$\mathcal O_{\mathbf P^1}(1)$ under a finite morphism is
ample. This proves (C06.1).

For (C06.2), use the polarized abelian lift C44-A05.
After completing its discrete valuation ring $R$ and,
if necessary, enlarging the perfect residue field, the
Cohen subring $W(k)\subset R$ allows base change of
$(\mathcal C,\widetilde f)$ to $R$.
The ample line bundle $L$ extends to $\mathcal L$ on
$\mathcal A/R$. Ample line bundles on abelian varieties
have no higher cohomology, so cohomology and base change
lift both $s_0,s_1$ to sections $\widetilde s_0,\widetilde s_1$.
The cover lift supplies
$\mathcal M=\widetilde f^*\mathcal O(1)$ and sections
$\widetilde u,\widetilde v$.
Their incidence section
\[
                 \widetilde s_0\otimes\widetilde u+
                 \widetilde s_1\otimes\widetilde v
\]
defines a relative Cartier divisor $\mathcal Y$ in
$\mathcal A\times_R\mathcal C_R$. Its special section is a
nonzerodivisor, so the divisor is flat over $R$.
Its special fibre is the prescribed smooth $Y$.
The same properness argument for the nonsmooth locus
makes $\mathcal Y$ smooth over $R$. It is projective and
its inclusion in the lifted product is the required
embedded lift. \hfill$\square$

\paragraph{Scope of C06.}
This lifts the morphism as a morphism of curves.
It does not preserve a deck group, the branch divisor,
ramification multiplicities, or a chosen Hurwitz stratum.
In particular it does not assert a solution to the
Galois-cover lifting problem. Nor does it lift an
individual middle cycle on the incidence fourfold.

\subsection*{C44-C06a. Removing the separability restriction}

\textbf{Corollary.}
Every finite nonconstant morphism $f:C\to\mathbf P^1_k$
from a smooth projective curve over a perfect field of
characteristic $p$ has a finite flat lift of the same
degree over $W(k)$ with smooth projective source.
The embedded geometric lifting conclusion of C06
therefore holds for every smooth incidence (C06.2)
over a finite field.

\textbf{Proof.}
Let $t$ be the rational coordinate on the line and put
$a=f^*t\in k(C)$. Because $k$ is perfect, the kernel
of $d:k(C)\to\Omega^1_{k(C)/k}$ is $k(C)^p$.
This is Stacks, Lemma 10.158.2, Tag 031W.
Choose the largest $e\ge0$ for which
$a=b^{p^e}$ with $b\in k(C)$. Such a largest exponent
exists: a nonzero pole order of the nonconstant $a$
is divisible by $p^e$ for every such exponent and
therefore bounds it. Maximality gives $db\ne0$.
Thus $b$ defines a finite separable morphism
$h:C\to\mathbf P^1_k$, and
\[
 f=F_e\circ h,\qquad
 F_e:[x:y]\longmapsto[x^{p^e}:y^{p^e}].                      \tag{C06a.1}
\]
Here $F_e$ is the displayed $k$-morphism, acting as
the identity on constants; it is not being confused
with absolute Frobenius on a non-prime field.
The separating-function criterion follows equivalently
from $\Omega^1_{k(C)/k(b)}=0$ for this finite extension.

Lift $h$ by C06 and compose with the same homogeneous
power formula on $\mathbf P^1_{W(k)}$. That formula
is finite flat of degree $p^e$: on each standard
chart its coordinate algebra is free with basis
$1,t,\ldots,t^{p^e-1}$ over the target chart.
The composite is therefore finite flat of degree
$p^e\deg h=\deg f$, and it has the prescribed reduction.
Its source is the smooth projective lift supplied
by C06. The incidence argument of C06 then applies
unchanged. \hfill$\square$

\subsection*{C44-C07. A known positive span and one residual potential Tate direction}

\textbf{Theorem.}
Assume the data of (C06.2) are defined over a finite field,
with $Y$ smooth, and let $k:Y\hookrightarrow T=A\times C$.
All cohomology in this statement is geometric.
For every coefficient prime $\ell\ne p$ put
\[
 V_\ell=\ker\bigl(k_*:H^4(Y,\mathbf Q_\ell(2))
                         \to H^6(T,\mathbf Q_\ell(3))\bigr),
 \qquad q_Y(x,y)=\int_Yxy.                                  \tag{C07.1}
\]
Fix one coefficient prime for
$A_Y=\operatorname{Alg}^2(Y)\cap V_\ell$.
Let $K\subset A_Y$ be a finite-dimensional rational algebraic
subspace on which $q_Y$ is positive definite; $K=0$ is allowed
with the empty determinant equal to one.
After applying the common algebraic projector described
below to cycle representatives, extend the finite field
so that these representatives are defined.
Let $\tau$ be the total algebraic multiplicity of roots
of unity in the normalized Frobenius polynomial of
$K_\ell^\perp\cap V_\ell$.
If $\tau\le1$, then $q_Y$ is positive definite on all of
$A_Y$. More precisely,
\[
 A_Y=K\quad\hbox{or}\quad A_Y=K\perp\mathbf Qz,\qquad q_Y(z,z)>0,
                                                               \tag{C07.2}
\]
and only the first case occurs if $\tau=0$.
The known span $K$ need not lift.
No Tate conjecture, Frobenius semisimplicity, standard
conjecture $D(Y)$, or 032 input is used.

Writing $g=g(C)$, $d=\deg f$ and $D=\int_A h^4$,
the full variable rank is
\[
                  \dim V_\ell=D(5d+2g-2)-14-28g,
                                                               \tag{C07.3}
\]
which is even. Therefore the test with $K=0$ only detects
vanishing of $A_Y$; a residual multiplicity-one test can
be nontrivial only when $\dim K$ is odd.

\textbf{Proof.}
Choose a rational zero-cycle of degree one on $C$, after
a finite extension if desired, and let $\eta$ denote its
cohomology class. The incidence class is
\[
                          a=[Y]=h+d\eta.
                                                               \tag{C07.4}
\]
It is ample, since $L$ and $f^*\mathcal O(1)$ are ample.
Weak and hard Lefschetz give the orthogonal decomposition
\[
                    H^4(Y)=k^*H^4(T)\perp V_\ell,
 \qquad k_*k^*=L_a:H^4(T)\xrightarrow{\sim}H^6(T).
                                                               \tag{C07.5}
\]
We give one inverse correspondence valid for every
coefficient prime, rather than assuming kernel membership
is independent of the realization.

Use the algebraic Chow--Kunneth projectors of the abelian
fourfold and of the curve. On the curve the latter are
$e\times C$, $C\times e$, and their complement in the
diagonal, where $e$ is the rational degree-one zero-cycle.
Kunnemann's motivic hard Lefschetz theorem gives common
algebraic inverse operators
\[
 J=(L_h^2:H^2(A)\to H^6(A))^{-1},\qquad
 J_3=(L_h:H^3(A)\to H^5(A))^{-1}.                            \tag{C07.6}
\]
One may replace $L$ by $L\otimes[-1]^*L$ to obtain a
symmetric ample line bundle and then rescale the operators
by the appropriate rational powers of two.
The exact source is Kunnemann, \emph{Invent.\ Math.}
113 (1993), Theorems 5.1--5.2, printed p.~96.
Only his absolute Chow-motive theorem is used.

In the Kunneth decomposition, an element of $H^6(T)$
is written as $(x,y,z\eta)$ with
\[
 x\in H^6(A),\quad y\in H^5(A)\otimes H^1(C),\quad z\in H^4(A).
 \qquad b=Jx,\quad c=J(hz).
 \]
Define its inverse image in $H^4(T)$ by
\[
 \left(hb+d^{-1}(z-hc),\
                 (J_3\otimes1)y,\
                 (c-db)\eta\right).                        \tag{C07.7}
\]
Since $h(z-hc)=hz-h^2J(hz)=0$, multiplication by
$h+d\eta$ sends the first component to $x$, the odd
component to $y$, and the last total component to
\[
              d\big(hb+d^{-1}(z-hc)\big)+h(c-db)=z.
 \]
This proves the inverse formula. Every operation in it
is an algebraic correspondence, including the product
Kunneth projectors, multiplication by $h$, and the rational
factor $d^{-1}$. Thus one rational correspondence $\Gamma$
on $T$ realizes $L_a^{-1}$ simultaneously for all $\ell$.
The correspondence
\[
                       P=\Delta_Y-k^*\Gamma k_*
                                                               \tag{C07.8}
\]
is the common cohomological orthogonal projector onto
$V_\ell$. Applying it to chosen cycle representatives
preserves their initially defined classes in $A_Y$ and
their rational intersection matrix, and puts them in
$V_\ell$ in every realization. Its nonzero determinant
makes their span independent in every realization.

The embedded lift from C06 and C06a provides one rational Betti
quadratic space for $(V_\ell,q_Y)$ at all $\ell\ne p$.
The Betti kernel is primitive for $a|_Y$: the Gysin map
above the middle degree is injective by dual weak
Lefschetz, so $k_*v=0$ implies $a|_Y\,v=0$.
Hodge--Riemann makes the ordinary Gram determinant
$D_V$ of this rational Betti space positive. Indeed
type $(2,2)$ is positive, and every negative block
comes in a complex-conjugate pair of even real dimension.
This argument only lifts the geometry, not the cycles
spanning $K$.

Here the determinant proof of C44-A15 applies with
the common projector (C07.8). We spell out its sign
step to make the implication precise. Write $G$ for
the positive rational Gram matrix of $K$.
Given any algebraic class, apply $P$ and subtract its
$K$-projection using $G^{-1}$ and rational intersection
numbers. This produces one rational cycle $Z$ in the
orthogonal residual space for every $\ell$.
Algebraic cycles have finite normalized Frobenius orbit.
If $\tau=0$ all residual classes are zero.
If $\tau=1$ and a residual class is nonzero, its potential
Tate line is simple and nondegenerate; its nonzero rational
square $a_Z=q_Y(Z,Z)$ makes the same cycle nonzero in
every realization. It spans that line in each.
Here rationality makes the unique cyclotomic root equal
to $1$ or $-1$; simplicity and orthogonality of the other
generalized eigenspaces prove nondegeneracy without
semisimplicity.

After squaring normalized Frobenius $F$, put $S=F^2$.
On the orthogonal complement $U_\ell$ of that line
there are no eigenvalues $1$ or $-1$.
The operator $S-S^{-1}$ is invertible and skew-adjoint,
so its product with the Gram matrix is alternating.
Its determinant is a square. If
$\psi(T)=\det(T-S|U_\ell)$, reciprocal pairing and purity
give the common rational positive number
\[
              \delta=\det(S-S^{-1})=\psi(1)\psi(-1)>0.
 \]
Both factors are products of positive complex-conjugate
pairs; the empty product is one. Consequently
\[
                  a_Z\det(G)\delta/D_V
                           \in(\mathbf Q_\ell^*)^2
                  \quad\hbox{for every }\ell\ne p.
 \]
A rational number that is a square in all these local
fields is a rational square: otherwise its nontrivial
quadratic extension has an inert prime different from
$p$. Thus $a_Z>0$. Every other algebraic residual class
is a rational multiple of $Z$, its coefficient being
$q_Y(Z',Z)/a_Z$. This proves (C07.2).
The rationality, purity and coefficient independence
of the Frobenius polynomials are Deligne, \emph{Weil I},
Theorem 1.6, printed p.~276; quotienting the middle
polynomial by the ambient factor and $(T-1)^{\dim K}$
gives the residual polynomial. None of these steps
requires semisimplicity or realizes an algebraic class
from a Frobenius root.

For the rank formula, $c(T_{A\times C})=1+(2-2g)\eta$.
The normal sequence for the divisor (C07.4) yields
\[
 \chi_{\rm top}(Y)
     =\int_{A\times C}\bigl(a^5-(2-2g)\eta a^4\bigr)
     =D(5d+2g-2).
 \]
Weak Lefschetz and Kunneth give
\[
 b_1(Y)=8+2g,\qquad b_2(Y)=29+16g,\qquad b_3(Y)=64+56g.
 \]
Poincare duality then gives
$b_4(Y)=\chi_{\rm top}(Y)+84+84g$.
The ambient middle rank is
$b_4(A\times C)=98+112g$, proving (C07.3).
Abelian Riemann--Roch gives $D=24\chi(A,L)$, so that
rank is even. The non-root-of-unity eigenvalues of any
orthogonal transformation occur in reciprocal pairs.
Hence residual cyclotomic multiplicity has the parity
of $\dim V_\ell-\dim K$. This proves the stated parity
qualification. \hfill$\square$

\paragraph{Primary sources and programme scope.}
The deformation inputs were read directly in SGA~1,
Expos\'e III 5.1--5.6, 6.9--6.10 and 7.2--7.4:
\texttt{https://arxiv.org/pdf/math/0206203}.
These are revised pp.~59--62, 67--68 and 70,
respectively (original pp.~71--75, 82--83 and 85--86).
The finite-map algebraization and conormal flatness
criteria are
\texttt{https://stacks.math.columbia.edu/tag/09ZT}
and \texttt{https://stacks.math.columbia.edu/tag/063Y}.
The separating-function input for C06a is
\texttt{https://stacks.math.columbia.edu/tag/031W}.
For (C07.6), the primary scanned source is
\texttt{https://gdz.sub.uni-goettingen.de/download/pdf/PPN356556735\_0113/LOG\_0011.pdf}.
The local-to-global squareclass implication is also
explicit in Ancona--Marmora, \emph{Comment.\ Math.\ Helv.}
100 (2025), proof of Corollary 4.4, printed p.~629;
only that elementary implication is used.

C06 separates geometric lifting from cycle lifting
for the smooth incidences of JAC1964 H9-T18.
C07 supplies a sufficient fixed-cover sign test on
exactly its algebraic kernel, hence a new case of
cote~44, Theorem~5, scan~004 after imposing the
Albanese-primitive domain of H9-T17 and the primitive
qualification on scan~008. It does not settle all
special covers, does not assert that $\tau\le1$
always holds, and makes no claim about singular
incidences requiring a separate resolution.
The kernel sign in C07 is unconditional. Any full-height
conclusion that also uses the separate abelian primitive
terms remains conditional on the unrefereed 032 and
Milne2026 inputs specified in the prompt.

% END work/covers/separable_cover_lifts.tex

% BEGIN work/arithmetic/contribution.tex
% C44_POS, arithmetic continuation, /root/fermat_aj.
% GPT-6 Astra Pro; effort setting not exposed. 2026-11-22 (Europe/Berlin).
% New proofs below are subject to the recorded whole-proof check.

\section{The Fermat regulator: an integral projector and the exact
unramified obstruction}

The notation in this section is that of JAC1964 JAC-P13, JAC-P17,
AR-16a and AR-23.  In particular, $J/\mathbf Q$ is the Jacobian of
$x^4+y^4=z^4$, $N=h^3(J)^G$ with $G=\mu_4^2$, and
$\eta\in V:=\operatorname{CH}^2(N)_{\mathbf Q}$ is the specified
Fourier--Ceresa class.  The nonzero vector
$\alpha_5(\eta)$ spans
$S_5=H^1_f(\mathbf Q,H^3(N_{\overline{\mathbf Q}},\mathbf Q_5)(2))$
by the carried input AR-23, using AR-19 and H9-T15.  We do not reprove
that Selmer calculation, nor its nonzero cycle.  None of the deductions
in this section uses 032 or Milne's report.

The question $\ker\alpha_5=0$ remains unresolved.  The findings here
locate its integral obstruction more precisely.  They do not identify
a Selmer dimension with a Chow dimension.

\subsection{An integral model at the prime 5}

\begin{proposition}[C44-F01]\label{C44-F01}
Let $R=\mathbf Z_{(5)}$ and $S=\operatorname{Spec}\mathbf Z[1/10]$.
Let $\mathcal J/S$ be the Jacobian of the smooth Fermat curve
$\mathcal C:x^4+y^4=z^4$.  There is an exact idempotent, with
coefficients in $\mathbf Z[1/2]\subset R$, given by
\[
 e=\frac{1-[-1]^*}{2}\,\frac1{16}\sum_{g\in G}g^*.
 \tag{F01.1}
\]
The sum is formed after $S'=\operatorname{Spec}\mathbf Z[i,1/10]$
and descended to $S$.  Its rational generic-fibre motive is exactly
$N$.  Put
\[
 A=e\operatorname{CH}^2(J;R).
\]
Then restriction is an isomorphism
\[
 e\operatorname{CH}^2(\mathcal J;R)
       \xrightarrow{\ \sim\ } A,
 \qquad A\otimes_R\mathbf Q=V.                 \tag{F01.2}
\]
Moreover $e$ acts on geometric $5$-adic cohomology only in degree
three, where its image is a free rank-two direct summand.  Thus
\[
 T=eH^3_{\mathrm{et}}(J_{\overline{\mathbf Q}},\mathbf Z_5)(2)
 \tag{F01.3}
\]
is an actual $\mathbf Z_5$ lattice for the representation in the task,
and $eH^b_{\mathrm{et}}(J_{\overline{\mathbf Q}},\mathbf Z/5^n(2))=0$
for $b\ne3$ and every $n\ge1$.
\end{proposition}

\begin{proof}
The Fermat equation is smooth over $\mathbf Z[1/2]$: a singular
projective geometric point would have $x=y=z=0$, since all three
partial derivatives are units times the corresponding cubes.
The section $(1:0:1)$ gives its relative Jacobian, an abelian scheme.
The group action and inversion extend over $S$.  The two averaging
operators in (F01.1) commute and are idempotent.  They use only
powers of two as denominators.  Galois conjugation on $S'/S$ permutes
the sixteen group elements.  Restriction and norm, whose composite
is multiplication by two, descend the idempotent with the asserted
coefficients.  This descent is valid for Chow correspondences and
also for their actions on the cohomology theories used below.

The six characters of $G$ on $H^1(J)$ are
\[
 (1,1),(1,2),(2,1),(3,3),(3,2),(2,3)\quad\text{in }(\mathbf Z/4)^2;
 \tag{F01.4}
\]
these are the carried calculation in JAC-P17.  In their exterior
powers, an odd-degree invariant occurs only in degree three.  The
two invariant triples are the first three and the last three
characters.  Their types are $(3,0)$ and $(0,3)$.  The exact rational
identification with the projector of $N$ is JAC-P13/JAC-P17; in
particular (F01.1) is an integral-at-5 representative of the same
rational summand, not a new motive inferred solely from its
realization.  More explicitly, averaging is zero on $h^1(J)$
because it is induced by the zero norm endomorphism
$\sum_{g\in G}g$.  It is zero on $h^5(J)$ by the duality
identification through the $G$-invariant principal polarization.
Inversion removes the even Deninger--Murre pieces.  Thus the
rational idempotent on the full motive equals $\pi_3p_G$.
The geometric cohomology of an abelian variety is the
exterior algebra of its free degree-one cohomology.  Because the
projector has denominators prime to five, its integral image is a
direct summand and has no torsion.  The character calculation
therefore proves the assertions with both $\mathbf Z_5$ and
$\mathbf Z/5^n$ coefficients.

We next verify that $e$ annihilates the vertical terms in Chow
localization.  On any fibre of $\mathcal J$, the operator
$(1-[-1]^*)/2$ sends divisor classes into
$\operatorname{Pic}^0\otimes R$.  The endomorphism
$\sum_{g\in G}g$ of $\mathcal J_{S'}$ is zero.  Indeed its generic
fibre has zero tangent map by (F01.4), so is zero, and equality of
endomorphisms of an abelian scheme is detected on its dense generic
fibre.  Its dual is also zero, which says that
$\sum g^*$ annihilates $\operatorname{Pic}^0$.  Consequently
\[
 e\operatorname{CH}^1(\mathcal J_s;R)=0
 \quad\text{for every closed }s\in S.            \tag{F01.5}
\]
When a fibre does not contain $i$, check this after the extension
of degree at most two and descend using restriction--norm.

Chow localization gives
\[
 \bigoplus_{s\in S^{(1)}}\operatorname{CH}^1(\mathcal J_s;R)
 \longrightarrow\operatorname{CH}^2(\mathcal J;R)
 \longrightarrow\operatorname{CH}^2(J;R)\longrightarrow0.
 \tag{F01.6}
\]
The action of every automorphism in (F01.1) commutes with these
maps.  Applying the exact direct-summand functor $e$ and (F01.5)
proves (F01.2).  Notice that this argument is integral at five;
it does not invoke rational vanishing of the whole Chow group of
a special fibre.
\end{proof}

\paragraph{Locators and scope.}
The localization sequence is Fulton, \emph{Intersection theory},
Proposition~1.8 and its arithmetic-base version in \S20.2.
The Picard-scheme construction is FGA, Expos\'e~232,
Theorem~3.1 (printed pp.~148--149), applied to the smooth
projective relative curve; the canonical abelian component is
FGA, Expos\'e~236, Theorem~3.5 (printed p.~239), with the smooth
Picard schemes of the curve fibres.  The latter fibre smoothness
also follows from the vanishing of their $H^2(\mathcal O)$.
The relative-motive formalism over schemes smooth over a Dedekind
domain is recalled in the supplied Laga--Shnidman source,
\emph{Ceresa cycles of bielliptic Picard curves}, \S4.1--4.5;
the extension of Deninger--Murre to that base is explained at the
start of \S4.2, citing K\"unnemann, \S1, Remark~1.1.
The generic projector, the character list, and its rational field
of definition are JAC1964 JAC-P13 and JAC-P17.  The new assertion
here is the integral-at-5 localization isomorphism for that
projector.  The finite character check is reproduced by
\texttt{work/arithmetic/check\_fermat\_projector.py}; the script
does not prove any statement about rational equivalence.

\begin{lemma}[C44-F02]\label{C44-F02}
For the lattice $T$ in (F01.3),
\[
 H^0(\mathbf Q,T/5^nT)=0\quad(n\ge1),
 \qquad H^0(\mathbf Q,T\otimes\mathbf Q_5/\mathbf Z_5)=0.
 \tag{F02.1}
\]
Consequently
\[
 eH^3_{\mathrm{et}}(J,\mathbf Z/5^n(2))=0,
 \qquad
 eH^4_{\mathrm{et}}(J,\mathbf Z/5^n(2))
       \simeq H^1(\mathbf Q,T/5^nT).             \tag{F02.2}
\]
\end{lemma}

\begin{proof}
We make AR-23's Frobenius-at-11 check independent of a choice of
integral lattice.  By JAC-P17, the weight-three character is
$\psi_+^3$, where $\psi_+$ is the character of
$E_+:y^2=x^3+x$.  Direct counting gives
$\#E_+(\mathbf F_{11})=12$, so the degree-one Frobenius
polynomial is $X^2+11$.  Cubing its two roots gives
$X^2+11^3$ on $N$.  After the twist $(2)$, the polynomial is
\[
                 X^2+\frac1{11}.
\]
It follows that
\[
       \det(1-\operatorname{Fr}_{11}\mid T)
              =\frac{12}{11}\in\mathbf Z_5^\times.
 \tag{F02.3}
\]
For any stable $\mathbf Z_5$ lattice, this determinant condition
makes $1-\operatorname{Fr}_{11}$ invertible.  Its reductions are
invertible modulo every $5^n$, and it is invertible on the torsion
quotient as well.  Invariance under the entire Galois group implies
invariance under this Frobenius, proving (F02.1).  The conclusion
does not change if arithmetic rather than geometric Frobenius is
used, since both Frobenius and its inverse are integral units.

Apply the Hochschild--Serre spectral sequence with finite
coefficients and the integral projector $e$.  By C44-F01 the
projected geometric cohomology occurs only in row $b=3$.
The spectral sequence therefore degenerates on that row and
identifies its total degrees three and four with the two Galois
groups in (F02.1)--(F02.2).
\end{proof}

\paragraph{Check scope.}
The numerical part of (F02.3), including the point count and its
reduction $12/11\equiv2\pmod5$, is an exact integer/rational
computation in the accompanying script.  This is a lattice-uniform
form of the Frobenius test used in AR-23, not a new proof of
AR-23's nonzero Abel--Jacobi class.

\subsection{The precise obstruction after passing to finite coefficients}

Write
\[
 B=eH^4_{\mathrm{et}}(J,\mathbf Z_{(5)}(2)),
 \qquad
 C=eH^3_{\mathrm{nr}}(J,\mathbf Q_5/\mathbf Z_5(2)).
 \tag{F03.1}
\]
The cohomology defining $B$ is \emph{\'etale motivic cohomology},
not continuous $5$-adic cohomology.  The distinction is essential.
The group $C$ is unramified cohomology over the ground field
$\mathbf Q$, with residues at the codimension-one points of $J$.
The operations defining $e$ act on it by restriction, norm and
the automorphisms of $J_{\mathbf Q(i)}$.

\begin{theorem}[C44-F03]\label{C44-F03}
There is a canonical short exact sequence
\[
                 0\longrightarrow A\longrightarrow B
                   \longrightarrow C\longrightarrow0,
 \tag{F03.2}
\]
and both $A$ and $B$ have no nonzero $5$-primary torsion.
For every $n\ge1$, the kernel of the finite-coefficient cycle map
\[
 c_n:A/5^nA\longrightarrow H^1(\mathbf Q,T/5^nT)
\]
is canonically $C[5^n]$.  In particular, for
$\widehat A=\varprojlim_n A/5^nA$, the completed cycle map satisfies
\[
 \ker\!\left(\widehat c:\widehat A\longrightarrow
                    H^1_{\mathrm{cont}}(\mathbf Q,T)\right)
          = T_5C:=\varprojlim_n C[5^n],            \tag{F03.3}
\]
where the transition on $C[5^n]$ is multiplication by five.

Let
\[
 \rho:A\otimes_R\mathbf Q_5\longrightarrow
                  \widehat A\otimes_{\mathbf Z_5}\mathbf Q_5,
 \qquad K_{\mathrm{comp}}=\ker\rho.
\]
The full scalar-extended Abel--Jacobi kernel in the question has
the exact description
\[
 0\longrightarrow K_{\mathrm{comp}}
 \longrightarrow\ker\alpha_5
 \xrightarrow{\rho}
 \operatorname{im}\rho\cap
             (T_5C\otimes_{\mathbf Z_5}\mathbf Q_5)
 \longrightarrow0.                              \tag{F03.4}
\]
The intersection is taken inside
$\widehat A\otimes\mathbf Q_5$.  Neither group on the two ends
of (F03.4) has been proved to vanish.
\end{theorem}

\begin{proof}
For a smooth variety over a field, the codimension-two
comparison sequence is
\[
 0\longrightarrow\operatorname{CH}^2(J)
 \longrightarrow H^4_{\mathrm{et}}(J,\mathbf Z(2))
 \longrightarrow H^3_{\mathrm{nr}}(J,\mathbf Q/\mathbf Z(2))
 \longrightarrow0.
\]
The exact locator is Kahn, \emph{On the injectivity and
non-injectivity of the $l$-adic cycle class maps}, (1.5),
author p.~3, citing his \emph{Classes de cycles motiviques
\'etales}, Proposition~2.9.  Localizing at five and applying
the idempotent $e$ gives (F03.2).

The coefficient triangle for $\mathbf Z_{(5)}(2)$ has
$\mathbf Z/5^n(2)=\mu_{5^n}^{\otimes2}$ as its finite-coefficient
term.  Its long exact sequence and (F02.2) give
$B[5^n]=0$ and an injection
\[
 B/5^nB\hookrightarrow
       eH^4_{\mathrm{et}}(J,\mathbf Z/5^n(2)).     \tag{F03.5}
\]
The relevant coefficient comparison is the Merkurjev--Suslin
case of the norm-residue theorem; it is recorded in Kahn,
loc.~cit., (3.2) and (4.1), author pp.~7 and~10.
Since $A\hookrightarrow B$, also $A[5^n]=0$.

Apply multiplication by $5^n$ to (F03.2).  The snake lemma gives
\[
 0\longrightarrow C[5^n]\xrightarrow{\delta_n}A/5^nA
 \longrightarrow B/5^nB.
 \tag{F03.6}
\]
Explicitly, for $c\in C[5^n]$, choose $b\in B$ above $c$;
then $5^n b\in A$, and
$\delta_n(c)=5^n b\bmod5^nA$.  This is independent of the
choice of lift.  The absence of $5^n$-torsion in $B$ proves
its injectivity.  Composing with (F03.5) and the
Hochschild--Serre identification gives precisely $c_n$.
Thus (F03.6) proves the finite-coefficient assertion.

If $c\in C[5^{n+1}]$ and $b$ is a lift, then the lift $5b$
of $5c$ shows that reduction of $\delta_{n+1}(c)$ modulo $5^n$
equals $\delta_n(5c)$.  The transition is therefore multiplication
by five.  Taking inverse limits in (F03.6) is exact at the first
two terms, so the kernel of the completed finite-coefficient
cycle map is $T_5C$.  This uses left exactness only; no
unverified surjectivity of transitions is needed.

Continuous cohomology has the comparison exact sequence
\[
 0\longrightarrow\varprojlim{}^1_n
       H^0(\mathbf Q,T/5^nT)
 \longrightarrow H^1_{\mathrm{cont}}(\mathbf Q,T)
 \longrightarrow\varprojlim_n H^1(\mathbf Q,T/5^nT)
 \longrightarrow0.
\]
Its left term is zero by C44-F02.  Equivalently, in total
degree four on $J$ the potential $\varprojlim{}^1$ term is the
limit of the zero groups in (F02.2).  Thus the completed map
just constructed is the continuous cycle map, proving (F03.3).
This is why an inverse limit is legitimate in this particular
summand although it is not interchangeable with continuous
cohomology for an arbitrary variety over a number field.

After tensoring by $\mathbf Q_5$, the cycle map from $A$
factors as $\widehat c\circ\rho$.  Its rational projected
Hochschild--Serre component is the Abel--Jacobi map
$\alpha_5$ from the task.  The inclusion of its image in
$H^1_f$ is the carried AR-16a formalism and does not change
the kernel.  Flatness of $\mathbf Q_5$ over $\mathbf Z_5$
identifies the kernel of $\widehat c\otimes\mathbf Q_5$ with
$T_5C\otimes\mathbf Q_5$.  Taking the inverse image under
$\rho$ and then its image yields (F03.4).
\end{proof}

\paragraph{What this attack removes, and what it leaves.}
C44-F01 removes the ambiguity of a varying arithmetic spread for
the Chow source.  C44-F02 removes both the global degree-three
term and the continuous-cohomology $\varprojlim{}^1$ ambiguity
from the finite-coefficient calculation.  The remaining group
in (F03.3) is the actual Tate module of one specified unramified
group.  It is not the Selmer group.  A Selmer bound does not bound
this Tate module or prove the tensor-to-completion injectivity in
(F03.4).  Specialization of rational odd-motive Chow classes,
already ruled out as a detector by AR-20, supplies neither
missing assertion.

\subsection{A bounded-denominator formulation on rational planes}

For a finite free $R$-submodule $M\subset A$, let
$W_M=M\otimes_R\mathbf Q\subset V$.  Equation (F03.2) and
the absence of $5$-torsion identify $B$ with a subgroup of $V$:
if $b\in B$, choose $5^n b=a\in A$ and send $b$ to $a/5^n$.
This map is well defined and injective, since $C$ is $5$-primary
torsion.  This convention makes $B\cap W_M$ unambiguous.

\begin{proposition}[C44-F04]\label{C44-F04}
The following are equivalent for every finite free $M\subset A$:
\begin{enumerate}
\item $\alpha_5$ is injective on $M\otimes_R\mathbf Q_5$;
\item $T_5(B/M)=0$;
\item $(B\cap W_M)/M$ is finite;
\item there is an integer $b_M\ge0$ such that
      $5^{b_M}(B\cap W_M)\subset M$.
\end{enumerate}
More precisely,
\[
 \ker(\alpha_5|_{M\otimes\mathbf Q_5})
       \simeq T_5(B/M)\otimes_{\mathbf Z_5}\mathbf Q_5.
 \tag{F04.1}
\]
Choose a nonzero multiple $a_\eta\in A$ of the carried class
$\eta$.  It is enough to verify these equivalent conditions for
the modules
\[
                  M_z=R a_\eta+Rz\quad(z\in A).
 \tag{F04.2}
\]
The condition holds when $z$ lies in $\mathbf Q\eta$, by AR-23.
For arbitrary $z$ the bounded-denominator condition in (4) is
unproved.
\end{proposition}

\begin{proof}
Apply the argument of (F03.6) to
$0\to M\to B\to B/M\to0$.  Both $M$ and $B$ have no
$5$-torsion.  It gives
\[
 0\longrightarrow T_5(B/M)\longrightarrow\widehat M
                         \longrightarrow\widehat B.
\]
Since $M$ is finite free over $R$, its completion is
$M\otimes_R\mathbf Z_5$.  The injection (F03.5) implies
that the kernel of the map to $\widehat B$ is exactly the
kernel of the continuous cycle map.  Rationalizing proves
(F04.1).

All $5$-power-torsion elements of $B/M$ lie in
\[
 D_M=(B\cap W_M)/M\ \subset\ W_M/M
          \simeq(\mathbf Q_5/\mathbf Z_5)^{\operatorname{rank}M}.
\]
Conversely every element of $D_M$ is $5$-power torsion, since
$R$ already inverts the other primes.  Thus $T_5(B/M)=T_5D_M$.
A subgroup of a finite direct sum of $\mathbf Q_5/\mathbf Z_5$
is a direct sum of finitely many copies of
$\mathbf Q_5/\mathbf Z_5$ and a finite group.  One proof is to
take its Pontryagin dual: it is a quotient of the finite free
$\mathbf Z_5$ module, hence is finitely generated; the structure
theorem over the discrete valuation ring $\mathbf Z_5$ and
duality give the assertion.  It follows that $T_5D_M$ is finite
free, and vanishes exactly when $D_M$ is finite.  Finiteness of
this subgroup is equivalent to a bound on its exponent, giving
the equivalence with (4) and also with (1), by (F04.1).

Finally AR-23 makes $\alpha_5(a_\eta)\ne0$ and gives a
one-dimensional target.  If all the two-generator modules
(F04.2) have injective scalar-extended map, none can have rank
two.  Hence every $z\in A$ belongs to $\mathbf Q\eta$, so
$V=\mathbf Q\eta$ and the full map is injective.  The converse
is immediate.  On the line $\mathbf Q\eta$, AR-23 already gives
injectivity, so the equivalent bounded-denominator conclusion
is proved on that line without any hypothesis about the rest
of the Chow group.
\end{proof}

\begin{proposition}[C44-F05: an algebraic model of the missing step]
\label{C44-F05}
Even finite generation of the integral Chow-like source,
one-dimensional surjective regulator image, a nonzero distinguished
class, and a positive definite real form do not formally imply
the scalar-extended injectivity in C44-F04.  There is an explicit
model satisfying all those properties in which the completed
kernel is one-dimensional and arises from a nonzero unramified-like
Tate module.
\end{proposition}

\begin{proof}
Choose the root $a\in\mathbf Z_5$ of $a^2=6$ with $a\equiv1\pmod5$.
Hensel's lemma gives this root, and $a\notin\mathbf Q$ because
six is not a rational square.  Put
\[
 A_0=R^2,\qquad
 B_0=\{(r,s)\in\mathbf Q^2:r+as\in\mathbf Z_5\},
 \qquad f(r,s)=r+as.
\]
Then $A_0\subset B_0$, both are $5$-torsion-free, and
\[
  B_0/A_0\simeq\mathbf Q_5/\mathbf Z_5.
 \tag{F05.1}
\]
To see the last statement, take the kernel of the surjection
$(\mathbf Q_5/\mathbf Z_5)^2\to\mathbf Q_5/\mathbf Z_5$,
$(r,s)\mapsto r+as$, whose restriction on the first coordinate
is the identity.  The inverse image of this kernel in
$\mathbf Q^2$ is exactly $B_0$.

The scalar-extended regulator is
$\mathbf Q_5^2\to\mathbf Q_5$, $(r,s)\mapsto r+as$.
It is surjective, sends $\eta_0=(1,0)$ to one, and has kernel
$\mathbf Q_5(-a,1)$.  It is even injective on the underlying
rational vector space $\mathbf Q^2$.  The form
$h((r,s),(r',s'))=rr'+ss'$ is positive definite and
$h(\eta_0,\eta_0)=1$.  Neither fact removes the scalar-extended
kernel.

For a completely explicit division tower, choose integers $a_n$
with $a_n\equiv a\pmod{5^n}$ and
$a_{n+1}\equiv a_n\pmod{5^n}$, and set
\[
 b_n=5^{-n}(-a_n,1)\in B_0.
\]
Then $b_n\bmod A_0$ has order $5^n$ and
$5b_{n+1}-b_n\in A_0$.  The resulting nonzero element of
$T_5(B_0/A_0)$ maps to $(-a,1)\in\widehat A_0$.
Thus the precise obstruction in C44-F03 is present in the
model, despite all the other listed data.  This is a model of
an invalid inference, not a geometric counterexample to
Abel--Jacobi injectivity or to a height conjecture.
\end{proof}

\subsection{What finite coefficients see of the fourfold component}

\begin{proposition}[C44-F06]\label{C44-F06}
Let $F/\mathbf Q$ be finite and let $D/F$ be a smooth projective
geometrically connected curve with a point $d_0\in D(F)$.
On $X=D\times J_F$, put
\[
 q=\pi_1^D\otimes e,\qquad
 A_D=q\operatorname{CH}^2(X;R),\qquad
 B_D=qH^4_{\mathrm{et}}(X,R(2)),
 \quad R=\mathbf Z_{(5)}.
\]
Then $B_D$ is uniquely divisible by five and hence is a
$\mathbf Q$-vector space.  The motivic comparison identifies it
canonically with the rational Chow component:
\[
       B_D\simeq A_D\otimes_R\mathbf Q.
                                                        \tag{F06.1}
\]
In particular $A_D$ has no $5$-primary torsion, and the
finite-coefficient cycle map on $A_D/5^nA_D$ is zero for every $n$.
Writing
\[
 C_D=qH^3_{\mathrm{nr}}(X,\mathbf Q_5/\mathbf Z_5(2)),
\]
one has
\[
 C_D=B_D/A_D,\qquad
 C_D[5^n]\simeq A_D/5^nA_D,
 \qquad T_5C_D\simeq\widehat A_D.             \tag{F06.2}
\]
Thus finite generation of $B_D$ as an $R$-module would force
the entire rational component in JAC-P18 to vanish.  Such finite
generation is not established here; finite generation of $A_D$
alone would not have that consequence.
\end{proposition}

\begin{proof}
The point projector $\pi_1^D$ is integral.  By C44-F01 and
the K\"unneth formula, $q$ selects geometric cohomology only in
degree four, where the integral lattice after twist is
\[
 T_D=H^1_{\mathrm{et}}(D_{\overline F},\mathbf Z_5)
                  \otimes_{\mathbf Z_5}T.
\]
It is free, of rank $4g(D)$.  At any place of $F$ over five,
its rational realization has Hodge--Tate weights
$-2,-1,1,2$ (each with multiplicity $g(D)$, up to the
simultaneous opposite weight convention).  Indeed the two
factors before twist have weights $\{0,1\}$ and $\{0,3\}$.
None is zero.  The Hodge--Tate comparison for smooth proper
varieties therefore implies
\[
                 H^0(F,T_D)=0.                 \tag{F06.3}
\]
Here one may use Scholze, \emph{$p$-adic Hodge theory for
rigid-analytic varieties}, Corollary~1.8, author p.~4, applied
to the analytification, together with functoriality for the
algebraic projectors.  This use needs neither good reduction
of $D$ at five nor a Hodge or Tate conjecture.

Finite-coefficient Hochschild--Serre now gives
\[
 qH^3_{\mathrm{et}}(X,\mathbf Z/5^n(2))=0,
 \qquad
 qH^4_{\mathrm{et}}(X,\mathbf Z/5^n(2))
       =H^0(F,T_D/5^nT_D).
                                                        \tag{F06.4}
\]
As in C44-F03, the coefficient triangle implies that
$B_D[5^n]=0$ and injects $B_D/5^nB_D$ into the second group
in (F06.4).  Taking inverse limits gives
\[
 \widehat B_D\hookrightarrow
        \varprojlim_n H^0(F,T_D/5^nT_D)
        =H^0(F,T_D)=0.
\]
The maps $B_D/5^{n+1}B_D\to B_D/5^nB_D$ are surjective,
so each projection from their inverse limit is surjective.
Consequently $B_D/5^nB_D=0$ for every $n$.  Together with
$B_D[5^n]=0$ this proves unique divisibility by five.
The other primes are already inverted in $R$.

The comparison sequence (F03.2) is valid in codimension two
on $X$ as well, and is compatible with $q$.  Its quotient
is $5$-primary torsion.  Rationalizing its injection
$A_D\hookrightarrow B_D$, and using that $B_D$ is already
a rational vector space, proves (F06.1).  The finite-coefficient
cycle map factors through $B_D/5^nB_D=0$.  Applying the snake
lemma to multiplication by $5^n$ gives (F06.2), including the
transition by multiplication by five; its inverse limit is
the final displayed isomorphism.

A nonzero rational vector space is not finitely generated as
a $\mathbf Z_{(5)}$-module.  Thus the stated finite generation
condition forces $B_D=0$.  JAC-P18 has already proved that
every algebraic class in this projected component is
homologically zero, so this is the vanishing that its
componentwise Murre~(D) would require.  The elementary inclusion
$R\subset\mathbf Q$ shows why finite generation of the smaller
group $A_D$ cannot substitute for finite generation of $B_D$.
\end{proof}

\paragraph{Scope of the additional attack.}
C44-F06 does not prove Murre~(D).  It identifies the relevant
\'etale motivic group as entirely divisible, with all of its
completed Chow source accounted for by unramified cohomology.
This makes precise why torsion-freeness or finite-coefficient
cycle computations alone cannot kill the JAC-P18 obstruction.
It concerns cote~44, Th.~1--3 (scan~002) and the boundary of
Th.~5 (scan~004); it does not change the known polarization
on the incidence higher Picard factor.  In the frameworks of
SGA and Chow motives, it is the tensor summand
$h^1(D)\otimes N$ of weight four with no $(2,2)$ part,
followed by the same codimension-two \`etale-motivic comparison
as C44-F03.  Its vanishing is a rational-equivalence issue,
not the semisimplicity or signed $I$ assertion of cote~16,
scan~049.  No 032/Milne input is used.

\subsection{The remaining finiteness question on one arithmetic model}

\begin{proposition}[C44-F07]\label{C44-F07}
For the abelian scheme $\mathcal J/S$ of C44-F01, restriction
also gives an isomorphism on the projected \emph{\'etale
motivic} groups:
\[
 eH^4_{\mathrm{et}}(\mathcal J,R(2))
                 \xrightarrow{\ \sim\ } B.
                                                        \tag{F07.1}
\]
The completion $\widehat B$ is a free $\mathbf Z_5$-module
of rank exactly one.  Its rational realization is the Selmer
line already realized by $\eta$.

The following two conditions are equivalent:
\begin{enumerate}
\item the scalar-extended Fermat Abel--Jacobi map is injective;
\item $eH^4_{\mathrm{et}}(\mathcal J,R(2))$ is finitely
      generated over $R$.
\end{enumerate}
Condition (2) is not proved here.  In particular the proved
finite generation of the completion is not being substituted
for it.
\end{proposition}

\begin{proof}
Purity and localization for \'etale motivic cohomology
give the exact segment
\[
 \bigoplus_{s\in S^{(1)}}
 H^2_{\mathrm{et}}(\mathcal J_s,R(1))
 \longrightarrow H^4_{\mathrm{et}}(\mathcal J,R(2))
 \longrightarrow H^4_{\mathrm{et}}(J,R(2))
 \longrightarrow\bigoplus_{s\in S^{(1)}}
 H^3_{\mathrm{et}}(\mathcal J_s,R(1)).
                                                        \tag{F07.2}
\]
An exact source for the needed purity is Geisser,
\emph{Motivic cohomology over Dedekind rings},
Theorem~1.2(1), published p.~774: after inverting the residue
characteristic no truncation is needed.  All residue
characteristics of $S$ differ from five and are units in $R$.
The Bloch--Kato premise of that 2004 statement is supplied by
the norm-residue theorem (in this weight, Merkurjev--Suslin).
The localization of the cycle complex and passage to the
generic fibre give the displayed direct sum.  This is also
the localization used in Geisser, \emph{Hasse principles for
\'etale motivic cohomology}, \S4, diagram~(2).

The first outer group in (F07.2) is a sum of Picard groups,
because $\mathbf Z(1)_{\mathrm{et}}=\mathbf G_m[-1]$.
Its $e$-part vanishes by C44-F01.  The last is a sum of
cohomological Brauer groups localized at five.  For a fibre
$\mathcal J_s$, finite-coefficient Hochschild--Serre and the
character calculation of C44-F01 give
\[
           eH^2_{\mathrm{et}}(\mathcal J_s,\mu_{5^n})=0
           \quad(n\ge1).
\]
Indeed its projected geometric cohomology starts in degree
three.  Kummer's sequence implies
$e\operatorname{Br}(\mathcal J_s)[5^n]=0$ for every $n$.
The Brauer group is torsion, so its localization at five has
zero $e$-part.  Apply $e$ to (F07.2) to obtain (F07.1).

The lattice $T$ is a lisse $\mathbf Z_5$-sheaf on $S$.
Proper smooth base change and the same single-row calculation
give
\[
 eH^4_{\mathrm{et}}(\mathcal J,\mathbf Z/5^n(2))
          =H^1_{\mathrm{et}}(S,T/5^nT),
 \qquad eH^3_{\mathrm{et}}(\mathcal J,\mathbf Z/5^n(2))=0.
                                                        \tag{F07.3}
\]
For the second identity, Frobenius at eleven still belongs
to $\pi_1(S)$, so C44-F02 applies.  Finite-coefficient
cohomology of this arithmetic scheme is finite.  Alternatively,
in degree one this follows from the finiteness of number
field extensions of bounded degree unramified outside
$\{2,5,\infty\}$.
It follows from the coefficient comparison that each
$B/5^nB$ is finite, and that
\[
 \widehat B\hookrightarrow
           H^1_{\mathrm{cont}}(S,T).
                                                        \tag{F07.4}
\]
There is no inverse-limit correction by (F07.3).
The group on the right is a finite free $\mathbf Z_5$-module:
it is a compact inverse limit of finite groups, its quotient
modulo five injects into $H^1(S,T/5T)$, and hence compact
Nakayama gives finite generation; its $5$-torsion vanishes
because $H^0(S,T/5T)=0$.  Restriction to
$H^1_{\mathrm{cont}}(\mathbf Q,T)$ is injective by inflation.

The group $B$ is rationally $V$, by (F03.2).  Therefore the
rational image of every element of $B$ in (F07.4) belongs to
the Selmer line of AR-23.  Its intersection with the finite
free lattice $H^1_{\mathrm{cont}}(S,T)$ is a closed rank-one
lattice.  The image of $B$ is dense in $\widehat B$, whose
image in (F07.4) is closed by compactness.  Hence
$\widehat B$ lies in that rank-one lattice.  It is nonzero,
because $a_\eta\in A\subset B$ has nonzero regulator.
This proves its asserted rank, including its rational image.

If $B$ is finitely generated over $R$, torsion-freeness makes
it finite free.  Its rank equals the rank of $\widehat B$,
namely one.  As $A\subset B$ and $A\otimes_R\mathbf Q=V$,
this forces $V=\mathbf Q\eta$, and AR-23 then proves
injectivity.  Conversely, injectivity and AR-23 give
$V=\mathbf Q\eta$.  Apply C44-F04 to $M=Ra_\eta$: since
its rational span is all of $V$, the bounded-denominator
conclusion places $B$ inside $5^{-b_M}M$.  Thus $B$ is a
submodule of a finite free $R$-module and is finitely
generated.  This proves the equivalence.
\end{proof}

\paragraph{Meaning of the arithmetic finiteness gate.}
C44-F07 makes the finiteness question a statement about one
specific summand on one scheme over $\mathbf Z[1/10]$.
Neither finite generation of ordinary Chow groups nor
finite generation of continuous \(5\)-adic cohomology proves
it.  Kahn, \S4, author pp.~10--11, explicitly isolates this
distinction; his Proposition~4.1 gives counterexamples to
unqualified \`etale-motivic finiteness on arithmetic schemes.
Geisser, \emph{Hasse principles}, Conjecture~4.3, states
Lichtenbaum's finiteness expectation for regular schemes
proper over $\mathbf Z$; that conjecture is not a theorem
for the present summand, and $\mathcal J/S$ is proper over
the indicated open base, not over all of $\mathbf Z$.
No claim of direct applicability or proof of that conjecture
is made.  In cote~44 this refines the arithmetic gate of
Th.~5, scan~004; in cote~16 it specifies information about
rational equivalence that conditions $C$ and $I$ on pure
realizations do not provide.  The equivalence does not by
itself count as a strictly smaller main gate under the
progress rule of the prompt.  No 032/Milne input is used.

\subsection{The three proposed arithmetic routes: exact scope}

\paragraph{C44-FA01: Euler systems.}
The supplied Lilienfeldt--Shnidman paper,
\emph{Derivatives of Rankin--Selberg $L$-functions and heights of
generalized Heegner cycles}, arXiv:2408.04375,
Theorem~1.1/13.4, identifies a height on a Hecke-generated cycle
space modulo the real-height radical.  Its Corollary~1.4 expressly
assumes the injectivity of its scalar-extended Abel--Jacobi map.
In \S15, author pp.~71--73, equation~(99) on p.~72 is the Euler-system
conclusion that a nonzero generalized Heegner cohomology class
spans the Selmer space; the next paragraph invokes the assumed
injectivity to pass back to Chow.  The displayed implication chain
on p.~73 marks both uses of Abel--Jacobi injectivity.  That is the exact missing
step, and C44-F03--F04 describe its integral obstruction for the
Fermat summand.  No Euler-system relation at the level of
Galois cohomology bounds the denominators in (F04.2).

There are also independent obstacles to a direct substitution
in the displayed theorem: its standing assumptions in \S1.1
require an imaginary quadratic field of odd discriminant and
an unramified Hecke character, whereas the Fermat character
over $\mathbf Q(i)$ has conductor $(1+i)^4$ by JAC-P17/AR-17.
The stated range is $0<t\le k-1$, so the weight-four central
case with $k+t=2$ is not a direct instance.  These scope failures
are not the reason the existing Selmer result fails---that result
is already available as AR-19/H9-T15---but they prevent the paper
from being treated as a new Chow-generation theorem here.

The Loeffler--Zerbes constructions consulted in this run
interpolate or bound Galois-cohomology classes, not the entire
Chow group.  The exact example is Jetchev--Loeffler--Zerbes,
\emph{Heegner points in Coleman families}, Theorem~A
(published p.~125), equivalently Theorem~5.4.1, whose
specializations are \emph{Abel--Jacobi images} of generalized
Heegner cycles.  An interpolation statement about these images
does not supply the missing denominator bound.  The existing
JAC1964 audit AR-17 and its
\texttt{LATEST\_RANK\_SOURCE\_SCOPE.tex} remain in force; their
statements are not reproved here.

\paragraph{C44-FA02: Gross--Schoen heights.}
AR-24--AR-25 are carried as the exact projected-diagonal and
sixteen-intersection identities.  Zhang,
\emph{Gross--Schoen cycles and dualising sheaves},
Theorem~2.3.5, author p.~27, proves the triple-intersection
identity; Conjecture~2.4.1, author p.~28, asserts semipositivity
of the tensor-correspondence pairing.  The latter is not proved
by the former.  Even adding that conjecture would give the
semipositivity conclusion recorded in AR-25, without an
equality criterion or the injectivity in task~3.  In particular
one must not substitute a real-height calculation for
the absent faithfulness of the $5$-adic regulator.  C44-F05
shows that positivity and the stated Selmer data alone are not
a formal bridge.  The supplied Gu draft, Theorem~1.3,
Corollary~5.4 and Conjecture~1.1, gives formulas and
conditional positivity, not that missing faithfulness.  Its
statements are treated here as supplied draft inputs, not as
independently verified public results, and none is used in the
proofs of C44-F01--F05.

\paragraph{C44-FA03: the algebraically trivial radical.}
JAC-P18 isolates the possible cycle
$Z=(\pi_1^C\otimes\pi_N)_*Z$ on $C\times J$, fixed by the
degree-four product projector.  Xu--Xu,
\emph{Remarks on Murre's conjecture on Chow groups},
Corollary~3.9, author p.~9, gives (A) and (B) for the product;
the earlier definition of (D), author p.~2, is the extra
cycle-class injectivity needed to annihilate this component.
Theorem (B) leaves degree four acting on codimension two and
therefore does not remove this component.  The special
Laga--Shnidman isomorphism in their Theorem~5.2 concerns a
different invariant summand of types $(2,1),(1,2)$, identified
with $h^1(E)(-1)$.  Such a Picard-motive isomorphism cannot hold
for the Fermat $N$, whose types are $(3,0),(0,3)$.
No proof of the needed componentwise (D) was found in these
sources or in the other primary sources checked in this run.

\paragraph{Programme crosswalk and current gate.}
C44-F01--F04 concern the arithmetic extension of cote~44,
scan~004, Th.~5, and the incidence/higher-Picard boundary of
Th.~1--3 on scan~002, as made explicit in the letter on
scans~008--009.  They do not alter JAC-T4.3 or establish a new
polarization of $P^i(X)$.  The actual cycle lives in the
extreme-type motive outside that Picard component.
In cote~16, scans~049--050, the conclusion is a limitation of
the passage from $C$, $I$, and semisimplicity of cohomological
correspondences to rational Chow groups: the kernel measured
by (F03.4) is invisible to those pure-realization assertions.
The corrected $A$, $A'$, $A''$ conditions of SN-C04 are
unchanged.

In Grothendieck's later language, C44-F01 is a statement over
one scheme over $\mathbf Z$; (F03.2) is the comparison from
codimension-two cycles to \emph{\'etale motivic} cohomology;
(F02.2)--(F03.4) concern $\ell$-adic sheaves, weights, and their
arithmetic realization.  For the generic fibre, SGA~6,
Expos\'e~XIV, \S4.1--4.2 (especially (4.1) and (4.4)),
identifies the rational Chow group in codimension two with
degree two of the topological and $\lambda$ filtrations on
$K_0(J)$.  The operators in (F01.1) act there by pullback.
This $\lambda$-ring translation supplies no bound on rational
equivalences.  No new EGA/FGA theorem or standard-conjecture
implication is asserted: the modern comparison theorems are
cited under their actual sources above.

The open rows \texttt{QFermat-AJ} and \texttt{Q35-full} are
\emph{not closed}.  Their exact new test is (F03.4), or,
equivalently, bounded five-power denominators in $B$ on every
rational plane (F04.2).  The original analytic order-one result
AR-11, positive value AR-13, and Selmer-line result AR-23
remain intact.  None of the present findings proves
$\operatorname{rank}\operatorname{CH}^2(N)=1$.

% END work/arithmetic/contribution.tex

% BEGIN work/arithmetic/fermat_excess_picard.tex
% C44_POS: arithmetic algebraic-radical continuation, /root/fermat_aj.
% All statements here are unconditional.  No 032 or Milne input is used.
% This is a reduction of JAC-P18, not a proof of full arithmetic AJ injectivity.

\section{Removing ambient Picard families from the Fermat radical}

We retain JAC-P13, JAC-P17 and JAC-P18 without reproving them.
In particular $J$ is the Fermat-quartic Jacobian and $p_N$ is its
actual rational Chow projector, with realization of type
$(3,0)+(0,3)$.  Its construction from origin-preserving automorphisms
and the canonical abelian Chow--K\"unneth projectors makes it
symmetrically distinguished.  The results below are independent of
032, of the supplied Milne report, and of Milne2022 Proposition~4.2.

For the geometric argument work over an algebraically closed field
$k\subset\mathbf C$ of characteristic zero; in the arithmetic
application take $k=\overline{\mathbf Q}$.  All Chow groups and
homomorphism groups are tensored with $\mathbf Q$.  The use of this
algebraic closure loses no nonzero rational Chow class defined over
a number field: any rational equivalence over the algebraic closure
is defined over a finite extension, and restriction followed by norm
is multiplication by its degree.

\subsection{An exact correspondence identity}

\begin{lemma}[C44-F15, divisor identity]\label{C44-F15-identity}
Let $A/k$ be an abelian threefold.  Let $p$ be a symmetrically
distinguished Chow projector with $p=\pi_3^A p=p\pi_3^A$ and with
Betti realization contained in
$H^{3,0}(A)\oplus H^{0,3}(A)$.  For every
$d\in\operatorname{CH}^1(A)$, write $L_d$ for multiplication by $d$,
viewed as a correspondence of degree one.  Then
\[
                 p L_d\pi_1^A=0
       \quad\hbox{in }\operatorname{CH}^4(A\times A).
       \tag{F15.1}
\]
Also $p_*\operatorname{CH}^1(A)=0$.
\end{lemma}

\begin{proof}
The canonical abelian projectors are polynomials in the pullback
correspondence of $[n]$, for any fixed integer $n\ge2$; thus they
are symmetrically distinguished.  Graphs of homomorphisms, their
transposes, products, and compositions are distinguished by
O'Sullivan's theorem.  Here composition uses only pullback and
pushforward by projections between abelian varieties and their
intersection products.

Decompose
\[
 d=d_++d_-,\qquad d_\pm=\tfrac12(d\pm[-1]^*d).
 \tag{F15.2}
\]
The class $d_+$ is symmetric and $d_-$ lies in
$\operatorname{Pic}^0(A)_{\mathbf Q}$.  Symmetric divisor classes
are distinguished: the distinguished lift of the numerical class
is symmetric, and the difference from $d_+$ lies in
$\operatorname{Pic}^0(A)_{\mathbf Q}$, on which $[-1]^*$ acts by
$-1$, so a symmetric such difference is zero.
Consequently $p L_{d_+}\pi_1^A$ is distinguished.  Its cohomological
realization is zero.  Indeed, multiplication by the divisor class
maps $H^1(A)$ into the $(2,1)+(1,2)$ part of $H^3(A)$, and $p$
kills that part.  A cohomologically trivial cycle is numerically
trivial.  Injectivity of the distinguished subalgebra into the
numerical Chow ring now proves
$p L_{d_+}\pi_1^A=0$ in rational Chow.

For $d_-$ no distinguished-cycle assertion is needed.
The theorem of the square gives $[n]^*d_-=n d_-$.
Multiplication in the Chow motive is compatible with $[n]^*$,
and $[n]^*$ acts as $n$ on $h^1(A)$.  Hence
$L_{d_-}\pi_1^A$ has image in the exact $n^2$ eigensummand
$h^2(A)$, and is killed by $p\pi_2^A=0$.
Equivalently one may multiply on the left by each of the
polynomial projectors and compare the distinct eigenvalues
$n^j$.  This proves (F15.1).

Finally, on divisor classes the symmetric and antisymmetric
parts have $[n]^*$ eigenvalues $n^2$ and $n$ respectively.  Thus
only $\pi_2^A$ and $\pi_1^A$ act on $\operatorname{CH}^1(A)$,
which proves the last assertion.
\end{proof}

\subsection{The exact excess-Picard quotient}

For a pointed smooth projective connected curve $(C,c_0)$ let
\[
 q_C=\pi_1^C\otimes p,
 \qquad W_C(A,p)=(q_C)_*\operatorname{CH}^2(C\times A).
 \tag{F15.3}
\]
Let $i:D\longrightarrow A$ be a proper morphism from a smooth
projective connected surface, generically finite onto an integral
surface.  Put
\[
 B_i=\operatorname{im}\bigl(i^*:\operatorname{Pic}^0(A)
                  \longrightarrow\operatorname{Pic}^0(D)\bigr),
 \qquad Q_i=\operatorname{Pic}^0(D)/B_i.
 \tag{F15.4}
\]
These are an abelian subvariety and an abelian quotient, respectively;
they are not Chow groups or algebraic representatives of codimension-two
cycles.

\begin{proposition}[C44-F15, family reduction]\label{C44-F15}
With the hypotheses of the preceding lemma there is a well-defined
linear map
\[
 \overline\Psi_{C,i}:
 \operatorname{Hom}(\operatorname{Jac}(C),Q_i)_{\mathbf Q}
                  \longrightarrow W_C(A,p)
 \tag{F15.5}
\]
whose image is exactly
\[
 (q_C)_*(1_C\times i)_*
                    \operatorname{CH}^1(C\times D).
 \tag{F15.6}
\]
In particular (F15.6) is zero if $i^*:\operatorname{Pic}^0(A)
\longrightarrow\operatorname{Pic}^0(D)$ is surjective.

For each integral divisor $D_0\subset A$, choose a resolution
$i_{D_0}:\widetilde D_0\longrightarrow D_0\hookrightarrow A$.
The sum of (F15.5) over these divisors is surjective:
\[
 \bigoplus_{D_0\subset A}
 \operatorname{Hom}\bigl(\operatorname{Jac}(C),Q_{i_{D_0}}\bigr)_{\mathbf Q}
       \twoheadrightarrow W_C(A,p).
 \tag{F15.7}
\]
No injectivity of this surjection is claimed.
\end{proposition}

\begin{proof}
Choose a point $d_0\in D(k)$.  The normalized Poincar\'e
description of divisors on a product gives the exact decomposition
\[
 \operatorname{Pic}(C\times D)_{\mathbf Q}
 =\operatorname{pr}_C^*\operatorname{Pic}(C)_{\mathbf Q}
 \oplus\operatorname{pr}_D^*\operatorname{Pic}(D)_{\mathbf Q}
 \oplus\operatorname{Hom}(\operatorname{Jac}(C),
                          \operatorname{Pic}^0(D))_{\mathbf Q}.
 \tag{F15.8}
\]
For clarity, a line bundle in the last summand is trivial on
$\{c_0\}\times D$ and on $C\times\{d_0\}$.  Its family of
restrictions to $D$ defines a pointed morphism
$C\to\operatorname{Pic}^0(D)$, hence a homomorphism from
$\operatorname{Jac}(C)$.  Conversely the universal Poincar\'e
bundle constructs this line bundle.  Subtracting the two restrictions
from any line bundle yields the asserted decomposition, with
uniqueness from the seesaw principle.  This is an assertion about
\emph{divisors on the source product}, where Picard theory is exact.

The second summand of (F15.8), after pushforward, is killed by
$\pi_1^C$, which annihilates the fundamental class $[C]$.
An element of the first summand pushes forward to
$a\boxtimes i_*[D]$, with $a\in\operatorname{CH}^1(C)$.
Its projection by $p$ on $A$ is zero by the last assertion of the
lemma, since $i_*[D]\in\operatorname{CH}^1(A)$.
Notice that $\pi_1^C$ by itself need not kill a degree-zero
$a$; it is the target projector that kills this term.

It remains to treat the normalized third summand.  Suppose its
homomorphism $\phi$ takes values in $B_i$.  A surjection of abelian
varieties splits up to isogeny: choose a polarization and its
orthogonal complementary abelian subvariety, so that the addition
map from the kernel and its complement is an isogeny.  It follows
that $\phi$ lifts, with rational coefficients, to a homomorphism
$\widetilde\phi:\operatorname{Jac}(C)\to\operatorname{Pic}^0(A)$.
Let $\ell\in\operatorname{CH}^1(C\times A)$ be its normalized
Poincar\'e divisor.  Its pullback to $C\times D$ has the required
cross homomorphism.  Adjusting its normalization at $d_0$ changes
it only by the first two summands of (F15.8), already disposed of.
The projection formula gives
\[
 (1_C\times i)_*(1_C\times i)^*\ell
       =\ell\cdot\operatorname{pr}_A^*(i_*[D]).
 \tag{F15.9}
\]
The class $\ell$ is in the exact $\pi_1^A$ summand on its
$A$ factor: pullback by $1_C\times[n]$ multiplies its normalized
Poincar\'e class by $n$.  Applying $q_C$ to (F15.9) therefore
gives zero by (F15.1).  This proves that the map on cross
homomorphisms kills
$\operatorname{Hom}(\operatorname{Jac}(C),B_i)_{\mathbf Q}$.

The isogeny category of abelian varieties is semisimple, by the
same polarization-complement argument.  Hence applying
$\operatorname{Hom}(\operatorname{Jac}(C),-)_{\mathbf Q}$ to
$0\to B_i\to\operatorname{Pic}^0(D)\to Q_i\to0$ is exact,
including surjectivity onto the last term.  This proves the
well-defined factor map (F15.5), independent of a rational
splitting, and the image assertion (F15.6).

To prove (F15.7), take an integral codimension-two subvariety
$Z\subset C\times A$, so $\dim Z=2$.  Its image in $A$ has
dimension one or two.  In dimension one, writing that image as
$\Gamma$, dimension and irreducibility imply
$Z=C\times\Gamma$, so its $\pi_1^C$ projection is zero.
In dimension two its image is an integral divisor $D_0$.
Over the dense open where $\widetilde D_0\to D_0$ is an
isomorphism, take the inverse image of $Z$ and then its closure
$\widetilde Z\subset C\times\widetilde D_0$.
It is an integral divisor on a smooth variety, and proper
pushforward gives $[Z]$ with degree one.  Thus the projected
class of $Z$ belongs to (F15.6).  These classes generate the
Chow group, which proves the surjection.
\end{proof}

\subsection{Only singular ample surfaces with excess irregularity remain}

\begin{proposition}[C44-F16]\label{C44-F16}
Under the hypotheses of C44-F15, the right side of (F15.7) is
generated already by those integral divisors $D_0\subset A$ which
are ample and whose resolution has irregularity
\[
            q(\widetilde D_0)>3.
 \tag{F16.1}
\]
For such an ample divisor,
\[
       \dim Q_{i_{D_0}}=q(\widetilde D_0)-3.
 \tag{F16.2}
\]
All non-ample image surfaces contribute zero.  Every normal ample
divisor with rational singularities contributes zero, including
every smooth ample divisor.  Therefore any nonzero class in
$W_C(A,p)$ requires a singular ample image surface with nonzero
excess Picard quotient and a nonzero map from $\operatorname{Jac}(C)$
to that quotient.

Applied to $(J,p_N)$, this is a strict restriction of the family
component left open by JAC-P18.  It does not establish its vanishing
and does not settle the scalar-extended arithmetic Abel--Jacobi
injection of AR-16a/C44-F07.
\end{proposition}

\begin{proof}
First suppose $D_0$ is not ample.  The connected translation
stabilizer $B$ of $D_0$ has positive dimension, and $D_0$ is the
inverse image of a divisor in $A/B$.  Indeed, the usual
ampleness criterion for an effective divisor on an abelian variety
identifies ampleness with finiteness of its translation stabilizer;
stability under $B$ then makes the divisor saturated for the
smooth quotient map.  Choose a complementary abelian subvariety
$A'$; the sum map
\[
                    u:B\times A'\longrightarrow A
 \tag{F16.3}
\]
is an origin-preserving isogeny.  Flat pullback of $D_0$ is a union
of products $B\times\Gamma_0$, with $\Gamma_0$ an integral
divisor of $A'$.  The dimension of $B$ is one or two.

Pullback by $u$ is an isomorphism of rational Chow motives.
For example a complementary isogeny composed with $u$ is $[n]$,
and $[n]^*$ is invertible on every canonical eigensummand
$h^j$; the inverse is a polynomial in the canonical projectors
and homomorphism graphs.  Consequently the transported projector
$p'=u^*p(u^*)^{-1}$ is distinguished, remains a summand of
$h^3(B\times A')$, and has only the two extreme Hodge types.
It suffices to kill each pulled-back family, since $u^*$ is
injective and intertwines the two projectors.

If $\dim B=2$, then $A'$ is an elliptic curve and each
$\Gamma_0$ is a point.  Its product with $B$ is smooth, and
restriction $\operatorname{Pic}^0(B\times A')\to
\operatorname{Pic}^0(B\times\Gamma_0)$ is surjective.  C44-F15
kills its entire family contribution.

If $\dim B=1$, let $\Gamma\to\Gamma_0\subset A'$ be the
normalization.  The surface $B\times\Gamma$ is smooth, and
\[
 \operatorname{Pic}^0(B\times\Gamma)
       =\operatorname{Pic}^0(B)\times\operatorname{Jac}(\Gamma).
 \tag{F16.4}
\]
After the two pullback summands in (F15.8) are removed, a
divisor on $C\times(B\times\Gamma)$ is the sum of a normalized
cross divisor from $C\times B$ and one from $C\times\Gamma$.
The first homomorphism has values in $\operatorname{Pic}^0(B)$,
which is inherited from $B\times A'$, so C44-F15 kills it.
The pushforward of the second is of the form
\[
                        [B]\boxtimes\gamma,
       \qquad \gamma\in\operatorname{CH}^2(C\times A'),
 \tag{F16.5}
\]
after permuting the factors.  This statement concerns the actual
cycle: the source divisor is pulled back from $C\times\Gamma$,
and the map to $B\times A'$ is the identity on $B$.

The canonical projector $\pi_0^B$ fixes $[B]$.  On the other hand
\[
                  p'(\pi_0^B\otimes\Delta_{A'})=0
       \quad\hbox{as a rational Chow correspondence}.
 \tag{F16.6}
\]
Both factors in (F16.6) are distinguished.  In cohomology only
$H^0(B)\otimes H^3(A')$ can map into $p'$, and the third
cohomology of an abelian surface has types $(2,1)+(1,2)$.
The map is therefore zero in cohomology and hence zero in Chow
by O'Sullivan.  Equations (F16.5)--(F16.6) kill the remaining
term.  Resolving a family over $B\times\Gamma_0$ by the
degree-one strict-transform construction in F15 shows that this
argument covers every family on the original non-ample image
surface, not only a previously chosen resolution.

Now suppose $D_0$ is ample and $i:\widetilde D_0\to A$ is a
resolution.  The image of the homomorphism
$\operatorname{Alb}(\widetilde D_0)\to A$ contains a translate
of $D_0$.  It must be all of $A$: otherwise dimension forces
its image to be an abelian surface and $D_0$ to be a translate
of that surface, which is not ample.  Duality therefore shows
that $i^*:\operatorname{Pic}^0(A)\to
\operatorname{Pic}^0(\widetilde D_0)$ has finite kernel.
The image has dimension three, proving (F16.2).  F15 kills
every term with $q(\widetilde D_0)=3$.

Finally suppose that $D_0$ is normal with rational singularities.
For a resolution $r:\widetilde D_0\to D_0$, these hypotheses
give $r_*\mathcal O_{\widetilde D_0}=\mathcal O_{D_0}$ and
$R^1r_*\mathcal O_{\widetilde D_0}=0$.  Hence
$H^1(\widetilde D_0,\mathcal O)=H^1(D_0,\mathcal O)$.
Since $K_A$ is trivial, Kodaira vanishing and Serre duality give
\[
 H^1(A,\mathcal O_A(-D_0))=
 H^2(A,\mathcal O_A(-D_0))=0.
 \tag{F16.7}
\]
The exact sequence
$0\to\mathcal O_A(-D_0)\to\mathcal O_A\to
\mathcal O_{D_0}\to0$ then identifies
$H^1(A,\mathcal O_A)$ with
$H^1(\widetilde D_0,\mathcal O_{\widetilde D_0})$.
Thus $q(\widetilde D_0)=3$.  Equivalently the pullback of
Picard varieties has isomorphic tangent spaces and is an
isogeny in characteristic zero.  The asserted vanishing follows
from F15.  A smooth divisor is a special case.

Combining these eliminations with (F15.7) proves every assertion.
The last implication is only necessary: the existence of a map
to $Q_i$ does not assert that its image under
$\overline\Psi_{C,i}$ is nonzero.
\end{proof}

\paragraph{Application to the arithmetic radical, with exact scope.}
JAC-P18 says that a nonzero geometrically algebraically trivial
$N$-class would be detected by a nonzero class of
$W_C(J,p_N)$ for a curve over a finite extension.  F15--F16
remove all its families whose image surfaces are non-ample, or
whose resolutions have no Picard quotient beyond the ambient
Jacobian.  It is sufficient to kill the maps (F15.5) for
singular ample surfaces with $q(\widetilde D_0)>3$.
This is a genuinely smaller geometric family gate than demanding
Murre(D) for every $C\times J$: several complete source classes
have been annihilated in rational Chow, not merely in a realization.
It is not a proof that the remaining gate holds.  In particular
normalization of a singular divisor can increase irregularity,
and neither the Picard decomposition of source divisors nor
Poincar\'e reducibility makes their codimension-two pushforward
injective.  The full arithmetic AJ5 finite-generation problem in
C44-F07 remains unchanged.

\paragraph{Primary locators.}
O'Sullivan, \emph{Algebraic cycles on an abelian variety},
arXiv:0908.0626v1, 5 August 2009, 73 pages: introductory
Theorem (i)--(iii), printed p.~2; Theorem~5.1.6, printed
p.~55, for the multiplicative canonical eigendecomposition;
Theorem~6.2.5 and Corollary~6.2.6, printed pp.~70--71,
for distinguished cycles and the numerical section.
Primary URL: \url{https://arxiv.org/pdf/0908.0626}.
The published article is J. reine angew. Math. 654 (2011),
1--81, DOI 10.1515/crelle.2011.025; the locators just given
are explicitly those of the checked author version.

Bost--Charles, \emph{Some remarks concerning the Grothendieck
period conjecture}, author copy of the 2014 online article,
\S3.2.2, formula~(3.1), printed p.~17, and Propositions~3.5--3.6,
printed pp.~17--18, give the exact Picard-product and
Poincar\'e identifications used in (F15.8).
Primary URL:
\url{https://www.imo.universite-paris-saclay.fr/~francois.charles/crelle-periods.pdf}.

Debarre, \emph{Singularities of divisors on abelian varieties},
20 March 2006, introductory paragraph, printed p.~1, explicitly
states the quotient description of a non-ample effective divisor.
Primary author URL:
\url{https://www.math.ens.psl.eu/~debarre/DivVarAb.pdf}.
The proof above explains the subsequent isogeny, normalization,
and exact projector arguments; none of those Chow conclusions is
being attributed to that source.

% END work/arithmetic/fermat_excess_picard.tex

% BEGIN work/arithmetic/fermat_cartier_families.tex
% C44_POS, /root/fermat_aj.  Unconditional family-level P18 elimination.
% No 032 or conditional Milne positivity input.  Full AJ5 remains open.

\section{Cartier families on the singular image surface}

Keep the abelian threefold $A$ and its distinguished
extreme projector $p\le\pi_3^A$ from C44-F15.
All schemes and rational Chow groups in this section
are over $\overline{\mathbf Q}$.  The Picard group
of a singular surface is used directly; it must
not be confused with that of its normalization
or resolution.

\begin{proposition}[C44-F22]\label{C44-F22}
Let $\iota:D_0\hookrightarrow A$ be an integral
ample divisor, with no normality or smoothness
assumption.  Then
\[
 \iota^*:\operatorname{Pic}^0(A)
                       \longrightarrow\operatorname{Pic}^0(D_0)
 \tag{F22.1}
\]
is an isogeny of abelian varieties.  In particular
the target is proper, although this need not hold
for the Picard variety of an arbitrary non-normal
projective variety.

For every smooth projective curve $C$ and line
bundle $\mathcal L$ on $C\times D_0$, one has
\[
 (\pi_1^C\otimes p)_*(1_C\times\iota)_*
       \bigl(c_1(\mathcal L)\cap[C\times D_0]\bigr)=0
                         \quad\hbox{in }\operatorname{CH}^2(C\times A).
 \tag{F22.2}
\]
The same holds for a rational linear combination
of Cartier divisor cycles.  Hence it holds for
an integral cycle $Z\subset C\times D_0$ if
some nonzero integer multiple of $[Z]$ is the
cycle of a Cartier divisor on $C\times D_0$.
In particular it holds if the scheme morphism
$Z\to D_0$ is finite and flat.

Thus a surviving P18 family on the residual
surfaces of F16--F21 must fail this Cartier-cycle
condition.  It cannot be represented by an
integral finite-flat family on its image surface.
No assertion that every Weil divisor family is
Cartier, or that every generically finite family
is flat, is used.
\end{proposition}

\begin{proof}
Since $A$ is smooth, $D_0$ is an effective
Cartier divisor on $A$.  The restriction exact
sequence is
\[
 0\longrightarrow\mathcal O_A(-D_0)
  \longrightarrow\mathcal O_A
  \longrightarrow\mathcal O_{D_0}\longrightarrow0.
 \tag{F22.3}
\]
Serre duality on the threefold $A$, whose
canonical bundle is trivial, and Kodaira
vanishing for the ample bundle
$\mathcal O_A(D_0)$ give
\[
 H^1(A,\mathcal O_A(-D_0))=
 H^2(A,\mathcal O_A(-D_0))=0.
 \tag{F22.4}
\]
Consequently restriction induces an isomorphism
\[
                  H^1(A,\mathcal O_A)
                    \xrightarrow{\sim}H^1(D_0,\mathcal O_{D_0}).
 \tag{F22.5}
\]
The Picard scheme exists for the projective
integral scheme $D_0$; its identity component
is of finite type and separated, and has
tangent space $H^1(D_0,\mathcal O_{D_0})$.
In characteristic zero it is smooth.
The morphism (F22.1) has isomorphism on
tangent spaces by (F22.5).  Its kernel is
therefore zero-dimensional and finite.  Its
image is closed, because the source is proper
and the target is separated, and has the full
dimension three of the smooth connected target.
Thus its image is the whole identity component.
The target is proper and the map is an isogeny,
which proves (F22.1).  No statement about
$\operatorname{Pic}^0(\widetilde D_0)$ follows
from this argument; that group may be larger.

We spell out the divisor-family decomposition
to avoid using a smoothness assumption on
$D_0$ implicitly.  Choose $c_0\in C(k)$ and
put $\mathcal L_0=\mathcal L|_{\{c_0\}\times D_0}$.
The line bundle
$\mathcal L\otimes\operatorname{pr}_{D_0}^*
\mathcal L_0^{-1}$ gives a pointed morphism
\[
                         C\longrightarrow\operatorname{Pic}^0(D_0).
 \tag{F22.6}
\]
Since the target is now known to be abelian,
it factors through a homomorphism
$\operatorname{Jac}(C)\to\operatorname{Pic}^0(D_0)$.
The isogeny (F22.1) has a rational inverse,
so after multiplying by an integer this
homomorphism lifts to
$\operatorname{Pic}^0(A)$.
The normalized Poincar\'e divisor for the
lift on $C\times A$ gives a rational divisor
class $\ell$ satisfying
\[
                         (1_C\otimes\pi_1^A)_*\ell=\ell.
 \tag{F22.7}
\]
The exact projector assertion follows either
from the Poincar\'e construction or from
target multiplication by $n$, whose pullback
on this cross divisor is multiplication by $n$.

By construction, the rational line bundle
$\mathcal L-\ell|_{C\times D_0}-
\operatorname{pr}_{D_0}^*\mathcal L_0$ has
trivial class on every fiber over $C$.
Here subtraction denotes the tensor operation
in the rational Picard group, after clearing
the same integer denominator.
The seesaw principle for the proper integral
variety $D_0$ therefore makes the difference
the pullback of a rational line bundle on $C$.
It follows that the Cartier divisor cycle of
$\mathcal L$ is a rational sum of the following
three types:
\[
  \ell|_{C\times D_0},\qquad
  \operatorname{pr}_{D_0}^*b,\qquad
  \operatorname{pr}_C^*a.
 \tag{F22.8}
\]
This argument uses line bundles and the
Picard functor of $D_0$ itself, not a
representability assertion for codimension-two
Chow groups.

The projection formula sends the first term
to
\[
                             \ell\cdot\operatorname{pr}_A^*[D_0]
                 \quad\hbox{on }C\times A.
 \tag{F22.9}
\]
It is killed by $1_C\otimes p$, using (F22.7)
and the exact identity
$pL_{[D_0]}\pi_1^A=0$ from C44-F15.
The second term in (F22.8) pushes to
$[C]\boxtimes\iota_*(c_1(b)\cap[D_0])$
and is killed by $\pi_1^C[C]=0$.
The third pushes to $a\boxtimes[D_0]$
and is killed by $p\operatorname{CH}^1(A)=0$.
This proves (F22.2), and rational linearity
gives the Cartier-cycle assertion for $Z$.

Finally suppose $Z\to D_0$ is finite flat.
The ambient morphism $C\times D_0\to D_0$
is smooth of relative dimension one, hence
flat and of finite presentation.  Each fiber
$Z_d$ is a finite closed subscheme of the
smooth curve $C_{\kappa(d)}$, and so is an
effective Cartier divisor on that curve.
The fiberwise criterion for relative effective
Cartier divisors applies: the closed immersion
is of finite presentation, and $Z$ is flat
over the base.  Therefore $Z$ is an effective
Cartier divisor on $C\times D_0$, with
$[Z]=c_1(\mathcal O(Z))\cap[C\times D_0]$.
Equation (F22.2) applies.  This criterion
requires no regularity of the base $D_0$.
\end{proof}

\paragraph{Relation with the Albanese-cycle criterion.}
For F20's surface cycle $z_{(0)}$ on $B$, let
$z_{\rm dist}$ be O'Sullivan's unique distinguished
lift of its numerical class.  In characteristic
zero numerical and homological equivalence
coincide on abelian varieties (Kleiman,
\emph{Algebraic cycles and the Weil conjectures},
Remark~3.10, printed p.~380, using Lieberman's
Theorem~2A11 on p.~377), so
\[
                 z_{(0)}-z_{\rm dist}
                    \in\operatorname{CH}^{b-2}_{(0)}(B)_{\rm hom}.
 \tag{F22.10}
\]
The distinguished lift is itself in grade zero:
multiplication by $n$ and uniqueness of the
distinguished lift give the required eigenvalue.
The operator (F20.3) kills $z_{\rm dist}$ by
the same distinguished/Hodge argument as in
F20.  Thus it depends only on the difference
in (F22.10).  Fourier identifies the latter
with a homologically trivial group
$\operatorname{CH}^2_{(0)}(B^\vee)$.
This makes the relationship to AR-15a precise;
it is a reformulation, not an additional
vanishing theorem or a discharge of AJ5.

\paragraph{Primary sources and exact scope.}
The classical characteristic-zero numerical/homological
comparison just used is checked in Kleiman's original
\emph{Algebraic cycles and the Weil conjectures},
\emph{Dix expos\'es sur la cohomologie des sch\'emas}
(1968), 359--386, Remark~3.10 and Corollary~3.9,
printed p.~380, together with Theorem~2A11
(Lieberman), printed p.~377.  The retrieved scan is
\url{https://download.uni-mainz.de/mathematik/Algebraische%20Geometrie/Lehre/WS23.Padische.1966.kleiman.weil.pdf}.
These are unconditional classical Betti statements.

For the singular Picard functor, Kleiman,
\emph{The Picard scheme}, arXiv:math/0504020v1
(1 April 2005), Theorem~5.4, printed p.~37,
gives existence and quasi-projectivity for
projective geometrically integral schemes;
Theorem~5.11, printed pp.~42--44, identifies
the tangent space; and Corollary~5.14,
printed p.~44, gives smoothness in
characteristic zero.  No smoothness or
normality of $D_0$ is imposed in these
assertions.  Primary author preprint:
\url{https://arxiv.org/pdf/math/0504020}.
The universal rigidified Picard family under
the same hypotheses is also stated explicitly
in Conrad, \emph{Math 248B: Picard functors
for curves}, Theorem~2.2, printed p.~1;
Proposition~3.1 on p.~2 computes the tangent
space before the notes specialize to curves:
\url{https://virtualmath1.stanford.edu/~conrad/248BPage/handouts/pic.pdf}.

For the finite-flat implication, Stacks
Project, Lemma~31.19.9(2) and its final
statement, Tag~062Y, gives exactly the
fiberwise criterion used in the proof:
\url{https://stacks.math.columbia.edu/tag/062Y}.
Kimura, \emph{Murre's conjectures for certain
product varieties}, arXiv:math/0509249v2
(20 September 2005), Lemma~2.2, printed p.~5,
records this implication in the same
curve-times-image-surface setting; the
stronger hypotheses of that paper's
Theorem~2.1 are not used here.
For the seesaw principle on a possibly
singular complete integral variety, Milne's
\emph{Abelian Varieties} (1986), corrected
transcription dated 2 January 2022,
Theorem~5.1 and Corollary~5.2, printed p.~7,
have the required scope:
\url{https://www.jmilne.org/math/xnotes/AVs.pdf}.
This is a classical Picard-theoretic input,
not a conditional Milne positivity statement.

\paragraph{Remaining gap.}
Finite-flat and Cartier families on the
singular ample image are now eliminated
without a hypothesis on the irregularity
of its resolution or the distinguishedness
of its Albanese surface cycle.  Families
that do not descend as Cartier cycles
remain possible.  Passing to a resolution
or flattening blowup can enlarge the Picard
variety; it does not repair the argument
for those families automatically.
The other residual P18 classes and the
full actual \'etale-motivic finite-generation
gate C44-F07 remain open.

% END work/arithmetic/fermat_cartier_families.tex

% BEGIN work/arithmetic/fermat_albanese_surfaces.tex
% C44_POS, /root/fermat_aj.  Further unconditional P18 reductions.
% No 032 or conditional Milne positivity input.  Full AJ5 remains open.

\section{Albanese surface cycles and an excess-irregularity example}

We retain the notation and hypotheses of C44-F15: $A$ is
an abelian threefold over $k=\overline{\mathbf Q}\subset
\mathbf C$, and $p\le\pi_3^A$ is a distinguished projector
with only Hodge types $(3,0),(0,3)$.  The Fermat pair
$(J,p_N)$ satisfies these hypotheses.  All Chow groups
have rational coefficients.  These results concern the
algebraically trivial obstruction of JAC-P18; they do not
prove full arithmetic Abel--Jacobi injectivity.

\subsection{A vanishing criterion on the Albanese image}

\begin{proposition}[C44-F20]\label{C44-F20}
Let $i:D\to A$ be a resolution of an integral ample
surface $D_0\subset A$.  Put $B=\operatorname{Alb}(D)$,
$b=\dim B$, and $P=B^\vee$.  The Albanese homomorphism
$\rho:B\to A$ induced by $i$ is surjective.  Choose a
translate $j:D\to B$ of the pointed Albanese map such
that $i=\rho\circ j$, and put
\[
                        z=j_*[D]\in\operatorname{CH}^{b-2}(B).
 \tag{F20.1}
\]
Suppose that its Beauville grade-zero component $z_{(0)}$
is symmetrically distinguished.  Then, for every curve $C$,
\[
 (\pi_1^C\otimes p)_*(1_C\times i)_*
                    \operatorname{CH}^1(C\times D)=0.
 \tag{F20.2}
\]
Consequently the universal excess-Picard cycle of C44-F17
for this surface is zero, and the surface contributes
nothing to the P18 obstruction.

The hypothesis holds if $z$ belongs to the subalgebra
of $\operatorname{CH}^*(B)$ generated by divisor classes.
More generally the same conclusion holds for any
factorization $i=\rho j$ through an abelian variety $B$
such that $j^*:\operatorname{Pic}^0(B)\to
\operatorname{Pic}^0(D)$ is an isogeny and $z_{(0)}$ is
distinguished; the source map $i$ may then be generically
finite onto its image rather than birational.

In particular, if $D$ is a smooth complete-intersection
surface of ample divisors in an abelian variety $B$,
and a homomorphism $\rho:B\to A$ is generically finite
on $D$ with surface image, then all projected families
supported over that image vanish.  This statement does
not assume that $D$ or its image is dominated by a
product of curves.
\end{proposition}

\begin{proof}
The surjectivity of $\rho$ for the ample image is the
argument already checked in F16: if its image were a
proper abelian subvariety, the image surface would be
a translate of an abelian surface and hence non-ample.
For a pointed Albanese map $a$, one has
$i=\rho a+x$ for a constant $x\in A(k)$.  Choose
$b_0\in B(k)$ with $\rho(b_0)=x$, and take
$j=a+b_0$.  Pullback on $\operatorname{Pic}^0$ is unchanged
by this translation and identifies $P$ with
$\operatorname{Pic}^0(D)$.

Let $\ell\in\operatorname{CH}^1(P\times B)$ be the
normalized Poincar\'e divisor, and let $r_B:P\times B\to B$
be the projection.  The all-Picard universal family
after target projection is
\[
 \Gamma=(\pi_1^P\otimes p)_*
     (1_P\times\rho)_*\bigl(\ell\cdot r_B^*z\bigr)
                       \in\operatorname{CH}^2(P\times A).
 \tag{F20.3}
\]
To verify the formula before drawing a conclusion from
it, the normalized Poincar\'e family on $P\times D$
is the pullback of $\ell$ by $1_P\times j$, up to a
divisor pulled back from $P$.  The projection formula
gives
\[
 (1_P\times j)_*(1_P\times j)^*\ell
                            =\ell\cdot r_B^*j_*[D].
 \tag{F20.4}
\]
Composing pushforwards with $1_P\times\rho$ yields
(F20.3).  The normalization term pushes to a divisor
on $P$ externally multiplied by the divisor $i_*[D]$
on $A$ and is killed by $p\operatorname{CH}^1(A)=0$.
If $j^*$ is merely an isogeny, the same construction
parametrizes every cross divisor after rational
linear combinations: homomorphisms into
$\operatorname{Pic}^0(D)$ lift up to isogeny to $P$.

We now check precisely which part of $z$ can contribute.
The Poincar\'e divisor is in Beauville grade zero on
$P\times B$, since simultaneous multiplication by $n$
pulls it back to $n^2\ell$.  Multiplication by this
divisor and proper pushforward by a homomorphism of
abelian varieties preserve the Beauville grade.
The codimension in (F20.3) is correct:
$\ell\cdot r_B^*z$ has codimension $b-1$, and
$1_P\times\rho$ has source minus target dimension
$b-3$ (and is smooth of that relative dimension in
the surjective case).
The output projector $\pi_1^P\otimes p$ has total
Chow--K\"unneth index $1+3=4$ in codimension two,
so it annihilates every Beauville grade except zero.
It follows that (F20.3) is unchanged on replacing
$z$ by $z_{(0)}$.

Under the hypothesis, every cycle and operation now
appearing in (F20.3) is distinguished: the symmetric
Poincar\'e divisor, $z_{(0)}$, the canonical projector
$\pi_1^P$, $p$, the projections, and the abelian
homomorphism $1_P\times\rho$.
O'Sullivan's functorial distinguished subalgebras
therefore contain $\Gamma$.  Its cohomology lies in
\[
                            H^1(P)\otimes\operatorname{im}p.
 \tag{F20.5}
\]
The possible Hodge types are $(4,0),(3,1),(1,3),(0,4)$,
and none is $(2,2)$.  The class of an algebraic
codimension-two cycle in (F20.5) is consequently zero.
A homologically trivial distinguished cycle is zero
in rational Chow, by injectivity of the distinguished
subalgebra into numerical cycles.  Thus
\[
                                  \Gamma=0.
 \tag{F20.6}
\]

For a homomorphism $\phi:\operatorname{Jac}(C)\to P$,
pull back (F20.3) along the pointed map
$C\to\operatorname{Jac}(C)\xrightarrow{\phi}P$ and
then apply $\pi_1^C\otimes p$.  By the Poincar\'e
identification and the projection formula this is
the projected pushforward of the cross divisor
associated with $j^*\phi$.  The observation is
linear for rational homomorphisms.  The two
single-factor terms in $\operatorname{Pic}(C\times D)$
are killed exactly as in F15.  Equation (F20.6)
therefore proves (F20.2).  The universal quotient
cycle of F17 is zero by its curve-detection statement,
or directly by restricting (F20.3) to a rational
section of the excess Picard quotient.

Every divisor on an abelian variety splits into its
grade-zero symmetric part and its grade-one
$\operatorname{Pic}^0$ part.  The grade-zero part of a
polynomial in divisors is thus the same polynomial
in their symmetric parts, in the prescribed
codimension.  Symmetric divisors are distinguished,
and distinguished cycles form a subalgebra.  This
proves the stated sufficient condition on $z$.

For the complete-intersection assertion, write
$D=H_1\cap\cdots\cap H_{b-2}$ scheme-theoretically.
Its pushforward fundamental
cycle is the product of the divisor classes $H_i$.
Moreover restriction on $\operatorname{Pic}^0$ is an
isogeny.  This last assertion can be seen directly
without a conjecture: at each smooth ample divisor
$Y\subset X$ with $\dim X\ge3$, the sequence
$0\to\mathcal O_X(-Y)\to\mathcal O_X\to
\mathcal O_Y\to0$, Serre duality and Kodaira vanishing
give
\[
                            H^1(X,\mathcal O_X)
                              \simeq H^1(Y,\mathcal O_Y).
 \tag{F20.7}
\]
Indeed $H^1(X,\mathcal O_X(-Y))$ and
$H^2(X,\mathcal O_X(-Y))$ vanish, since their duals
are positive-degree cohomology groups of
$\omega_X(Y)$.  Thus restriction on the Picard
varieties has an isomorphism on tangent spaces;
in characteristic zero it is an isogeny.
If the successive intersections are smooth, iteration
proves the assertion for $D$.  The general smooth
complete-intersection case follows from its Koszul
resolution on $B$: every nontrivial Koszul term is
a sum of negative ample line bundles, whose cohomology
vanishes in degrees below $b$ by Serre duality and
Kodaira vanishing on $B$.  The resolution has length
$b-2$, so $H^1(B,\mathcal I_D)$ and
$H^2(B,\mathcal I_D)$ vanish.  Hence restriction
$H^1(B,\mathcal O_B)\to H^1(D,\mathcal O_D)$ is an
isomorphism in this case as well.
The general factorization case of the proof applies.
Finally, every integral family cycle over its image
lifts after multiplication by the generic degree
to a divisor on $C\times D$, by the finite-flat-open
and closure argument (F19.15).  Hence the conclusion
holds for all projected image families, not just
the families already presented on $D$.
\end{proof}

\paragraph{Independence of the translated Albanese choice.}
The grade-zero condition in F20 does not depend on the
translation used in (F20.1).  For completeness this
uses only a known low-dimensional Fourier bound, not
the general Beauville conjecture.  A surface cycle on
an abelian variety of dimension $b$ has codimension
$b-2$, and Fourier sends its grade $s$ to
$\operatorname{CH}^{2+s}_{(s)}(B^\vee)$.
For $s=-1$ this is the zero negative-grade part of
the divisor group; for $s=-2$ it is a nonzero-grade
part of $\operatorname{CH}^0$, also zero; for smaller
$s$ the codimension is negative.  Thus every surface
cycle has only grades $s\ge0$.
Under Fourier, translation is multiplication by
the exponential of a $\operatorname{Pic}^0$ divisor
(up to the irrelevant sign determined by pushforward
or pullback).  Such a divisor has grade one.
Translation therefore leaves the grade-zero part
of every surface cycle unchanged.  This also shows
why no inference about arbitrary higher-dimensional
negative-grade cycles is being made here.

\subsection{An explicit class of singular ample images with excess Picard}

\begin{proposition}[C44-F21]\label{C44-F21}
Let $E$ be any elliptic curve and let
$B=A\times E$, with projection $\rho:B\to A$.
Choose very ample line bundles $H$ on $A$ and $M$ on
$E$, and put $L=H\boxtimes M$.
For a general pair of divisors $H_1,H_2\in|L|$,
the smooth connected surface
\[
                          S=H_1\cap H_2\subset B
 \tag{F21.1}
\]
has the following properties:
\begin{enumerate}
\item The projection $S\to D_0:=\rho(S)\subset A$
is finite birational, so $S$ is the normalization of
the integral image surface $D_0$.
\item The divisor $D_0$ is ample and non-normal,
and $q(S)=4$.  In particular its smooth resolution
has irregularity greater than three.
\item Every P18 family over $D_0$ is killed by
$\pi_1^C\otimes p$, by C44-F20.
\end{enumerate}
Thus the vanishing criterion actually applies inside
F16's residual singular-ample, excess-irregularity
class.  No claim is needed or made that these surfaces
are, or are not, dominated by products of curves.
\end{proposition}

\begin{proof}
The line bundle $L$ is very ample, by the product
embeddings followed by the Segre embedding.  Bertini
and connectedness of ample complete intersections
give a smooth connected surface for a general pair.
Let $N=\dim|L|$, so the pair parameter space has
dimension $2N$.  Consider the off-diagonal relative
pair variety
\[
 R=(B\times_A B)\setminus\Delta_B
                \simeq A\times(E\times E\setminus\Delta_E).
 \tag{F21.2}
\]
It has dimension five.  For any pair $(x,y)\in R$,
the two points of $B$ are distinct, so very ampleness
implies that evaluation of sections of $L$ at $x,y$
is surjective.  Requiring a divisor in $|L|$ to
contain both points imposes two independent linear
conditions.  Requiring both $H_1$ and $H_2$ to do
so imposes four.  The incidence variety of
$(x,y,H_1,H_2)$ thus has dimension
\[
                              5+2N-4=2N+1.
 \tag{F21.3}
\]
By the dimension-of-fibers theorem, for a general
pair the off-diagonal fiber product
$(S\times_A S)\setminus\Delta_S$ has dimension at
most one (or is empty).

The image of $S$ must have dimension two: if its
dimension were one or zero, the generic fibers
would have positive dimension, giving an
off-diagonal pair locus of dimension at least
three or four, respectively.  If the generically
finite map onto its surface image had degree at
least two, characteristic zero makes it generically
separable, and its off-diagonal pair locus would
have dimension two.  Both contradict the bound
just obtained.  Its generic degree is therefore one.
There are also no positive-dimensional fibers
anywhere: a positive-dimensional closed subset of
the elliptic fiber $\{a\}\times E$ must be that
whole fiber.  Such a fiber inside $S$ would yield
an off-diagonal pair locus of dimension two.
Hence the map is quasi-finite and proper, and
therefore finite.  Since $S$ is smooth and integral,
a finite birational map makes it the normalization
of $D_0$.

Write $h=\operatorname{pr}_A^*c_1(H)$ and
$m=\operatorname{pr}_E^*c_1(M)$.  In Chow one has
\[
 [S]=(h+m)^2=h^2+2hm,\qquad
 \rho_*[S]=2\deg(M)c_1(H).
 \tag{F21.4}
\]
Here $m^2=0$, and $\rho_*h^2=0$ by the projection
formula.  The generic degree is one, so
$[D_0]=\rho_*[S]$ as divisor classes.  Thus $D_0$
is linearly equivalent to a positive multiple of
the ample divisor $H$, and is ample.

The restriction argument (F20.7), twice, gives
$H^1(S,\mathcal O_S)\simeq H^1(B,\mathcal O_B)$,
so $q(S)=4$.  Alternatively the Koszul resolution
of the two divisors uses the vanishing
$H^i(B,L^{-r})=0$ for $i<4$ and $r=1,2$, and
gives the same isomorphism.
If $D_0$ were normal, its finite normalization
map would be an isomorphism, making it a smooth
ample divisor in $A$.  The same restriction
argument on $A$ would then give $q(S)=q(A)=3$,
a contradiction.  It is consequently non-normal
and singular, with smooth normalization of
irregularity four.

Finally $[S]$ is a product of divisor classes in
$B$ and restriction on $\operatorname{Pic}^0$
is an isogeny.  C44-F20 therefore annihilates
all projected families over $D_0$.
Every construction and the nonempty open choice
of the pair can be made over $\overline{\mathbf Q}$.
\end{proof}

\paragraph{Source scope and remaining obstruction.}
The distinguished-cycle inputs are exactly the
O'Sullivan theorem and functoriality already
located in C44-F15.  For the grading and the
translation observation we use Beauville (1986),
Propositions~1--2, printed p.~648, and Beauville
(1983), Proposition~4(i),(ii), printed p.~246,
with its proof on p.~247.  The latter states the
Fourier transform of a point and the translation
formula; all identities are in rational Chow.
Primary author PDFs:
\url{https://math.univ-cotedazur.fr/u/beauvill/pubs/Chow.pdf}
and
\url{https://math.univ-cotedazur.fr/u/beauvill/pubs/fourier.pdf}.
The Picard-family identification is the exact
Picard-product statement already used in F15.
The complete-intersection and incidence arguments
above are proved explicitly, using Bertini,
Kodaira vanishing, and standard proper
pushforward and finite-morphism facts.

F20's distinguished Albanese-cycle hypothesis
is a sufficient condition that has just been
verified for a specified class of image surfaces.
It is not asserted for every Albanese image.
The remaining F17 universal cycles for other
singular ample surfaces with excess irregularity
remain open, as does the full AJ5 finite-generation
gate C44-F07.  No step in this file uses 032 or a
conditional Milne positivity statement.

% END work/arithmetic/fermat_albanese_surfaces.tex

% BEGIN work/arithmetic/fermat_product_surfaces.tex
% C44_POS, /root/fermat_aj.  Unconditional further P18 reduction.
% No 032 or conditional Milne positivity input.  Full AJ5 remains open.

\section{Eliminating surfaces dominated by products of curves}

The excess Picard variety in C44-F16 need not be zero for a
surface obtained by adding two curves in an abelian threefold.
The following argument nevertheless kills all its projected
divisor families.  It uses the exact divisor correspondence
identity of C44-F15, rather than an assertion that the Picard
variety of the surface is inherited from the ambient variety.
Work over $k=\overline{\mathbf Q}\subset\mathbf C$ and with
rational Chow groups.

\begin{proposition}[C44-F19]\label{C44-F19}
Let $A$ be an abelian threefold, and let $p$ be a symmetrically
distinguished projector below $\pi_3^A$ whose Betti realization
has only Hodge types $(3,0)$ and $(0,3)$.  Let
\[
                 f:C_1\times C_2\longrightarrow A
 \tag{F19.1}
\]
be a morphism from a product of smooth projective connected
curves whose image is an integral surface $D_0$.
For every smooth projective connected curve $C$ and every
divisor class $W\in\operatorname{CH}^1(C\times C_1\times C_2)$,
one has
\[
 (\pi_1^C\otimes p)_*(1_C\times f)_*W=0
                    \quad\hbox{in }\operatorname{CH}^2(C\times A).
 \tag{F19.2}
\]
Moreover every integral codimension-two cycle
$Z\subset C\times A$ whose projection onto $A$ has image
$D_0$ satisfies
\[
                         (\pi_1^C\otimes p)_*[Z]=0.
 \tag{F19.3}
\]

In particular this applies to the sum of any two curves in
$A$ when their sum has dimension two.  For the Fermat
projector $p_N$ on $J$, all such image surfaces can be
removed from the residual family obstruction of F16--F17,
even if their smooth resolutions have irregularity greater
than three.  More generally it applies whenever there is a
generically finite dominant rational map from a product of
two smooth projective curves to the image surface: its
composition with the inclusion in $J$ extends to a morphism.
It is not asserted that every
singular ample surface with excess irregularity has this
property.
\end{proposition}

\begin{proof}
Choose points on the two curves.  The Albanese variety of
their product is the product of their Jacobians.  The
universal property of the Albanese therefore writes the
morphism (F19.1) as
\[
                       f(x,y)=f_1(x)+f_2(y),
 \tag{F19.4}
\]
where the constant translation has been absorbed in one
of the two maps.  The maps $f_j:C_j\to A$ need not send
the chosen points to zero.  Equivalently (F19.4) follows
from the rigidity lemma for maps to an abelian variety.

The Picard decomposition of a product, applied to
$C\times(C_1\times C_2)$, expresses every divisor class as
a sum of the following types:
\[
 \begin{gathered}
 \operatorname{pr}_C^*a,\qquad
 \operatorname{pr}_{C_1\times C_2}^*b,\qquad
 \operatorname{pr}_{C\times C_1}^*\mathcal L_1,\qquad
 \operatorname{pr}_{C\times C_2}^*\mathcal L_2. 
 \end{gathered}
 \tag{F19.5}
\]
Here the $\mathcal L_j$ are normalized cross divisors,
corresponding to homomorphisms
$\operatorname{Jac}(C)\to\operatorname{Jac}(C_j)$, and are
of degree zero on every fiber $\{c\}\times C_j$.
This decomposition follows from
$\operatorname{Pic}^0(C_1\times C_2)
 =\operatorname{Pic}^0(C_1)\times\operatorname{Pic}^0(C_2)$;
any cross divisor between $C_1$ and $C_2$ is already included
in $b$.  The first term of (F19.5) pushes to
$a\boxtimes f_*[C_1\times C_2]$.  Its second factor is a
divisor on $A$, and is killed by
$p\operatorname{CH}^1(A)=0$, from F15.  The second term
pushes to $[C]\boxtimes f_*b$ and is killed by
$\pi_1^C[C]=0$.  In particular we have not asserted that
$\pi_1^C$ kills all divisor classes on $C$.

It remains to treat the third term; the fourth is identical
with the two curves exchanged.  Define
\[
 \alpha=(1_C\times f_1)_*\mathcal L_1
                     \in\operatorname{CH}^3(C\times A),
 \qquad
 \gamma=(f_2)_*[C_2]\in\operatorname{CH}^2(A).
 \tag{F19.6}
\]
Let $\star_A$ denote the relative Pontryagin operation in
the $A$ factor, with the second input pulled back from $A$.
The projection formula and (F19.4) give the exact identity
\[
 (1_C\times f)_*\operatorname{pr}_{C\times C_1}^*\mathcal L_1
                          =\alpha\star_A\gamma.
 \tag{F19.7}
\]
We next isolate its target Chow--K\"unneth indices.  Put
$\alpha_a=(1_C\otimes\pi_a^A)_*\alpha$ and
$\gamma_b=(\pi_b^A)_*\gamma$.
Because $\mathcal L_1$ has degree zero on the $C_1$ fibers,
the relative degree of $\alpha$ is zero.  The projector
$\pi_6^A=A\times\{0\}$ acts on
$\operatorname{CH}^3(C\times A)$ by this relative degree,
followed by the zero section.  Thus
\[
                             \alpha_6=0.
 \tag{F19.8}
\]
Only indices $a\le5$ can occur.  This uses no injectivity
statement about degree-zero zero-cycles.

For the fixed curve class one has
\[
                         \gamma_b=0\quad(b>4).
 \tag{F19.9}
\]
Indeed Beauville's Fourier isomorphism sends
$\operatorname{CH}^2_{(s)}(A)$ to
$\operatorname{CH}^{1+s}_{(s)}(A^\vee)$.
The index $b$ corresponds to $s=4-b$.  For $b=5$ this
would be $\operatorname{CH}^0_{(-1)}(A^\vee)=0$, since
the fundamental class is in grade zero; for $b=6$ the
target has negative codimension.  The transform is
invertible, which proves (F19.9).  This is the elementary
negative-grade vanishing for curves on an abelian
threefold, not an assumption of the general Beauville
conjecture.

The exact relative Pontryagin grading is
\[
       \alpha_a\star_A\gamma_b
        \in(1_C\otimes\pi_{a+b-6}^A)_*
                         \operatorname{CH}^2(C\times A),
 \tag{F19.10}
\]
with zero understood for an index outside $[0,6]$.
For example this follows directly from multiplication by
any integer $n\ge2$: pushforward $[n]_*$ acts on the
$a$th target component by $n^{6-a}$, and it respects
Pontryagin products by the addition square.  Hence the
product has eigenvalue $n^{12-a-b}$, giving (F19.10).
All these are identities of rational Chow eigensummands;
they remain valid with the independent smooth parameter
curve $C$.

Since $p\le\pi_3^A$, only $a+b=9$ can contribute to
$(1_C\otimes p)_*(\alpha\star_A\gamma)$.
The bounds $a\le5$ and $b\le4$ leave just one pair:
\[
 (1_C\otimes p)_*(\alpha\star_A\gamma)
                =(1_C\otimes p)_*(\alpha_5\star_A\gamma_4).
 \tag{F19.11}
\]

Apply partial Fourier transform in the $A$ factor.  Its
component on target index $a$ changes that index to $6-a$
and changes total codimension by $3-a$.  Consequently
\[
 \lambda=\mathcal F_A(\alpha_5)
      \in(1_C\otimes\pi_1^{A^\vee})_*
                              \operatorname{CH}^1(C\times A^\vee),
 \qquad
 d=\mathcal F_A(\gamma_4)
                              \in\operatorname{CH}^1_{(0)}(A^\vee).
 \tag{F19.12}
\]
In particular $d$ is a symmetric divisor class.  Fourier
transform exchanges the relative Pontryagin operation
with intersection with this fixed divisor:
\[
                 \mathcal F_A(\alpha_5\star_A\gamma_4)
                                      =\lambda\cdot d.
 \tag{F19.13}
\]
These formulas are the usual Poincar\'e-correspondence
identities, with $C$ carried along; their direct proofs
use the same addition square and projection formula.

The middle Fourier component is an isomorphism of Chow
motives $h^3(A)\to h^3(A^\vee)$.  Transport $p$ along it
to the projector $p^\vee$.  It is distinguished, because
the Fourier correspondence and its inverse are built
from the distinguished Poincar\'e divisor.  Its Hodge
realization again has only types $(3,0)$ and $(0,3)$:
Fourier sends $(r,s)$ to $(3-s,3-r)$ in this degree.
Thus C44-F15 on $A^\vee$ gives the actual Chow identity
\[
                   p^\vee L_d\pi_1^{A^\vee}=0.
 \tag{F19.14}
\]
Equations (F19.12)--(F19.14) annihilate the Fourier
transform of (F19.11).  Fourier is invertible, so
$(1_C\otimes p)_*(\alpha\star_A\gamma)=0$.
Together with (F19.5)--(F19.7), and then applying the
independent source projector $\pi_1^C$, this proves
(F19.2).  In fact the two cross terms were already
killed by the target projector alone.

Finally, the morphism $C_1\times C_2\to D_0$ is proper
and generically finite, of some degree $e>0$.  There is
a dense open $U\subset D_0$ over which it is finite
flat of degree $e$.  An integral $Z$ as in the statement
has $Z_U=Z\cap(C\times U)$ dense in $Z$.
Flat pullback of $[Z_U]$ to
$C\times f^{-1}(U)$, followed by closure of its cycle
components in the smooth threefold
$C\times C_1\times C_2$, defines a divisor cycle $W$.
Every resulting component dominates $Z$, and the finite
flat degree calculation on $U$ gives
\[
                            (1_C\times f)_*W=e[Z].
 \tag{F19.15}
\]
There is no extra component over $D_0\setminus U$ in
this proper pushforward: the image of each closure is
the closure of its image on $U$, namely $Z$.
Applying (F19.2) and dividing by $e$ proves (F19.3).
This argument does not require a lift of $f$ to a
preselected resolution of $D_0$, and therefore avoids
an unjustified identification of their Picard varieties.

If only a dominant rational map $C_1\times C_2\dashrightarrow
D_0$ is given, compose it with $D_0\hookrightarrow A$.
Weil's extension theorem extends this rational map to a
morphism on the smooth product.  Its image is still $D_0$:
the image lies in the closed subvariety $D_0$ because this
holds on a dense open subset, and it contains a dense subset
of $D_0$.  Properness makes the image closed.  The extension
is therefore a generically finite dominant morphism onto
$D_0$, to which the argument already proved applies.
\end{proof}

\paragraph{Exact residual gate and limits.}
For the algebraically trivial obstruction of JAC-P18,
C44-F16 leaves only singular ample image surfaces
$D_0\subset J$ whose smooth resolutions have
$q(\widetilde D_0)>3$, through the actual excess quotient
and cycle of C44-F17.  F19 removes, in addition, every
such surface rationally dominated by a product of two
curves.  This is a genuine
annihilation of their rational Chow families, including
sums of two curves with possibly excess irregularity.
No representability-to-injectivity implication is used.
It does not eliminate the other surfaces in that gate,
and does not establish the actual \'etale-motivic
finite generation or full arithmetic AJ5 injectivity
of C44-F07.

\paragraph{Primary source locators.}
Beauville, \emph{Sur l'anneau de Chow d'une vari\'et\'e
ab\'elienne}, Mathematische Annalen~273 (1986), 647--651,
Proposition~1 and Proposition~2(a),(b), printed p.~648,
give the exact multiplication eigenspaces, Pontryagin
grading and Fourier codimension change used above.
Primary author PDF:
\url{https://math.univ-cotedazur.fr/u/beauvill/pubs/Chow.pdf}.
The page has been visually checked because the author's
PDF is an image scan.
Beauville, \emph{Quelques remarques sur la transformation
de Fourier dans l'anneau de Chow d'une vari\'et\'e
ab\'elienne}, LNM~1016 (1983), Proposition~3(i),(ii),
printed p.~243, gives the inverse and product identities;
the explicit Poincar\'e-correspondence proof in
Proposition~3', printed pp.~245--246, also shows their
compatibility with an additional parameter factor.
Primary author PDF:
\url{https://math.univ-cotedazur.fr/u/beauvill/pubs/fourier.pdf}.
The Picard-product and distinguished-divisor inputs are
the ones already checked and cited in C44-F15.
For the product splitting and rational-map extension,
Milne, \emph{Abelian Varieties}, Arithmetic Geometry
(1986), Chapter~V, Corollary~2.5 and Theorem~3.1,
are read in the author's corrected transcription dated
2 January 2022, printed p.~4 (numbering unchanged):
\url{https://www.jmilne.org/math/xnotes/AVs.pdf}.
These are the classical rigidity and Weil extension
theorems, not a use of the conditional Milne positivity
inputs flagged elsewhere in this project.

% END work/arithmetic/fermat_product_surfaces.tex

% BEGIN work/arithmetic/fermat_residual_motive.tex
% C44_POS, /root/fermat_aj.  Unconditional refinements of C44-F15--F16.
% No 032 or Milne input.  The full arithmetic AJ5 gate remains open.

\section{A universal residual cycle and the CM elliptic subquotient}

We use C44-F15--F16 and the Fourier/source-transpose conventions of
AR-15a.  We do not repeat the latter's arbitrary parameter-Jacobian
reduction.  The new point is that the parameter can be the actual
excess Picard quotient of an image surface.  Work over
$k=\overline{\mathbf Q}\subset\mathbf C$, with rational Chow groups
and the contravariant Chow-motive convention
$\mathbb L=\mathbf Q(-1)$.

\subsection{One cycle on the actual excess quotient}

\begin{proposition}[C44-F17]\label{C44-F17}
Let $i:D\to J$ be a resolution of an integral image surface, and
put
\[
 P=\operatorname{Pic}^0(D),\qquad
 B=\operatorname{im}(\operatorname{Pic}^0(J)\to P),\qquad
 Q=P/B,\qquad q=\dim Q.
 \tag{F17.1}
\]
If $q>0$, there is a canonical class
\[
 \beta_i\in
 (\pi_1^Q\otimes p_N)_*\operatorname{CH}^2(Q\times J)
 \subset\operatorname{CH}^2_{(0)}(Q\times J)_{\rm hom}.
 \tag{F17.2}
\]
Here canonical means independent of the rational splitting and
Poincar\'e normalizations used in its construction, for the given
morphism $i$.  For every pointed curve $(C,c_0)$, every residual
family map in (F15.5) is obtained by pullback of $\beta_i$,
followed by $\pi_1^C\otimes p_N$, along the pointed map
$C\to Q$ induced by the homomorphism
$\operatorname{Jac}(C)\to Q$, with rational linearity when a
denominator is present.  Thus $\beta_i=0$ kills all families with
this image surface.

Conversely, if the residual family maps vanish for all curves,
then $\beta_i=0$.  It suffices to test a single smooth curve in
$Q$ whose induced Jacobian homomorphism surjects onto $Q$.

Partial Fourier transform on $Q$ identifies the whole component
in (F17.2) with
\[
 \operatorname{Hom}_{\rm CHM}
       \bigl(h^1(Q^\vee)(-1),N\bigr).
 \tag{F17.3}
\]
A polarization can replace $Q^\vee$ by $Q$ up to isogeny.
For an isogeny decomposition $Q\sim\prod_a Q_a^{m_a}$, with
$Q_a$ simple, the component is the corresponding direct sum
of the components for $Q_a$, repeated $m_a$ times.

This describes the remaining classes exactly; it does not prove
that (F17.2) vanishes.  C44-F16 restricts the relevant surfaces
to singular ample surfaces with $q(D)>3$.
\end{proposition}

\begin{proof}
Choose a rational section $s:Q\to P$ of the quotient map
$P\to Q$, in the isogeny category.  The normalized Poincar\'e
line bundle on $P\times D$, pulled back by $s\times1_D$,
defines a rational divisor class
\[
                         \mathcal P_s\in\operatorname{CH}^1(Q\times D).
 \tag{F17.4}
\]
This is well-defined for a rational homomorphism $s$: if $ms$
is an actual homomorphism, take $m^{-1}c_1(\mathcal P_{ms})$.
The homomorphism--Poincar\'e identification is additive, so the
answer is independent of $m$.  Define
\[
 \beta_i=(\pi_1^Q\otimes p_N)_*(1_Q\times i)_*\mathcal P_s.
 \tag{F17.5}
\]
Pullback of the normalized divisor by $[n]\times1_D$ multiplies
it by $n$.  Thus the source factor in (F17.5) already belongs
to the exact $\pi_1^Q$ eigensummand.

If $s'$ is another rational section, $s'-s$ has values in $B$.
Lift this difference up to isogeny to
$\operatorname{Pic}^0(J)$.  Its Poincar\'e divisor on $Q\times D$
is pulled back from $Q\times J$, up to the two single-factor
terms.  The proof of (F15.9) applies with the abelian parameter
$Q$ in place of a curve: it uses only the projection formula
on the target and $p_N L_d\pi_1^J=0$.  Hence its projected
pushforward is zero.  A change of normalization adds a divisor
pulled back from $Q$ or from $D$.  The first is killed by
$p_N\operatorname{CH}^1(J)=0$, and the second by
$\pi_1^Q[Q]=0$.  This proves independence of all these choices.

Let $\phi:\operatorname{Jac}(C)\to Q$ be an actual homomorphism
and $a_C:C\to\operatorname{Jac}(C)$ the pointed Abel map.
Pullback by $(\phi a_C)\times1_D$ of $\mathcal P_s$ is exactly
the cross divisor for $s\phi$.  Its pushforward commutes with
this pullback on the other factor.  After applying the curve
projector it therefore gives
$\overline\Psi_{C,i}(\phi)$, by F15's independence of the
lift.  This also explains the equality for rational $\phi$:
multiply by an integer and divide, using additivity of the
Poincar\'e class.  No map of points into a Chow group is being
assumed to have an injective algebraic representative.

For the converse choose a smooth complete-intersection curve
$C\subset Q$ when $q\ge2$, and take $C=Q$ when $q=1$.
Its Jacobian maps surjectively to $Q$, since its cohomological
restriction on $H^1$ is injective by weak Lefschetz.
After choosing a basepoint and translating the map, the induced
map of Picard motives
$h^1(Q)\to h^1(C)$ is split injective in rational Chow motives:
the homomorphism of Jacobians splits up to isogeny, and the
degree-one motive functor identifies these homomorphisms with
the corresponding Chow-motive maps.  Applying the splitting on
the first tensor factor recovers any class in (F17.2) from its
projected curve pullback.  If that projected pullback is zero,
$\beta_i=0$.
Translation causes no problem: the component used is $h^1$,
and the pointed Poincar\'e description gives the same cross
homomorphism, with only the already killed normalization terms.

The total Chow--K\"unneth index of $\beta_i$ is $1+3=4$,
so its Beauville grade in codimension two is zero.  In Betti
cohomology the component is
\[
                  H^1(Q)\otimes H^3(N).
 \tag{F17.6}
\]
Its Hodge types are $(4,0),(3,1),(1,3),(0,4)$, with possible
multiplicities, and it has no $(2,2)$ part.  Consequently every
algebraic class in this component, in particular $\beta_i$,
is homologically trivial.

For completeness, the source transpose matters.  As a
correspondence, the cycle in (F17.2) has right projector
$\pi_{2q-1}^Q$, not $\pi_1^Q$.  Partial Fourier transform
interchanges the first-factor cycle indices $1$ and $2q-1$.
It increases codimension by $q-1$, giving
\[
 \widehat\beta_i\in\operatorname{CH}^{q+1}(Q^\vee\times J),
 \qquad
 \widehat\beta_i=p_N\circ\widehat\beta_i\circ\pi_1^{Q^\vee}.
 \tag{F17.7}
\]
Thus it is a map $h^1(Q^\vee)\to N(1)$, equivalently the map
in (F17.3).  These are exact Chow assertions.  The Fourier
correspondence is given by the exponential of the Poincar\'e
divisor and has inverse, up to sign, inversion followed by the
dual Fourier transform.  Its component on $h^1$ is the
$(2q-1)$st power of that divisor divided by $(2q-1)!$.
The eigensummand identities can also be checked on cohomology
and lifted to Chow by O'Sullivan, since the Poincar\'e divisor,
Fourier components, and canonical projectors are distinguished.
Finally the degree-one motive of a product is the sum of the
degree-one motives of its factors, and isogenies are isomorphisms
of rational Chow motives.  This proves the last assertion.
\end{proof}

\subsection{One CM elliptic factor: half of the tensor is killed}

Let $E/k$ have $j(E)=1728$, and let $I=[i]^*$ on $h^1(E)$.
Over the coefficient field $K=\mathbf Q(i)$ define
\[
 e_+=\tfrac12(\pi_1^E-iI),\qquad
 e_-=\tfrac12(\pi_1^E+iI),\qquad
 U=(E,e_+),\quad\overline U=(E,e_-).
 \tag{F18.1}
\]
The two lines have types $(1,0)$ and $(0,1)$ after an ordering
of the embeddings.  For $m\ge1$, the sum of the two projectors
$e_+^{\otimes m}+e_-^{\otimes m}$ has rational coefficients.
Let $N_m$ denote the rational motive it defines on $E^m$.
Thus $N_m\otimes K=U^{\otimes m}\oplus\overline U^{\otimes m}$,
and $N_m$ has rank two and types $(m,0),(0,m)$.

\begin{proposition}[C44-F18]\label{C44-F18}
Over $k=\overline{\mathbf Q}$, the Fermat motive $N$ is
isomorphic to $N_3$ by an actual distinguished Chow
correspondence.  There is a rational Chow-motive isomorphism
\[
               h^1(E)\otimes N
                 \simeq N_4\oplus (N_2\otimes\mathbb L).
 \tag{F18.2}
\]
The second summand has zero codimension-two Chow group:
\[
              \operatorname{CH}^2(N_2\otimes\mathbb L)
                   =\operatorname{CH}^1(N_2)=0.
 \tag{F18.3}
\]
Consequently
\[
 (\pi_1^E\otimes p_N)_*\operatorname{CH}^2(E\times J)
                    \simeq\operatorname{CH}^2(N_4).
 \tag{F18.4}
\]
This kills the whole rational rank-two component with Hodge
types $(3,1),(1,3)$; it leaves exactly the rank-two extreme
component with types $(4,0),(0,4)$.

If $Q$ in F17 has an isogeny factor $E^r$, the same
elimination applies to all $r$ copies of this factor.  In
particular, when $Q\sim E^r$, vanishing of
$\operatorname{CH}^2(N_4)$ is sufficient for all residual
families from this $Q$.  Its vanishing has not been proved
here.  The isogeny $J\sim E^3$, CM splitting, and Kimura
nilpotence alone do not prove it.
\end{proposition}

\begin{proof}
The carried Fermat decomposition JAC-P17 gives $J\sim E^3$
over $k$.  Isogeny graphs transport $p_N$ to a distinguished
projector on $h^3(E^3)$.  In $H^3(E^3)$ its two extreme Hodge
lines are exactly $U^{\otimes3}$ and $\overline U^{\otimes3}$;
all remaining summands have types $(2,1),(1,2)$.
The transported projector and the projector defining $N_3$
are therefore identical in cohomology.  Their difference is
distinguished and numerically trivial, so O'Sullivan makes
them identical in rational Chow.  This proves the asserted
actual isomorphism, rather than merely a realization isomorphism.

The polarized pairing on $h^1(E)$ gives
\[
                         U\otimes\overline U\simeq\mathbb L_K.
 \tag{F18.5}
\]
One can verify (F18.5) without a semisimplicity assumption for
the Chow category: take the projected cup-product pairing and
its projected dual diagonal, and rescale their nonzero
composition on $\mathbb L_K$.  They are distinguished; the
other composite is the identity in cohomology on the rank-one
summand, hence the identity in Chow by O'Sullivan, with
$K$-coefficients.  The extension of the latter theorem from
$\mathbf Q$ to $K$ here just tensors an injective map by the
field $K$.

Expanding the left side of (F18.2) over $K$ gives the four
summands
\[
 U^4,\quad\overline U^4,\quad
 U\otimes\overline U^3,\quad
 \overline U\otimes U^3.
 \tag{F18.6}
\]
Using (F18.5), the last two become
$\mathbb L_K\otimes\overline U^2$ and
$\mathbb L_K\otimes U^2$.  The pairing constructions and
tensor permutations commute with conjugation of the coefficient
field, so the two conjugate pairs descend to the rational
isomorphism (F18.2).  All maps used are actual algebraic
correspondences.  Tensor permutations may change a pairing
normalization by a sign; no positivity claim is used and the
isomorphism and vanishing are unaffected.

The motive $N_2$ is a summand of $h^2(E^2)$ with only
types $(2,0),(0,2)$.  On an abelian surface the degree-two
part of its divisor Chow group is its N\'eron--Severi group
tensored with $\mathbf Q$: the homologically trivial divisors
are in $\operatorname{Pic}^0$, on which the exact projector
is $\pi_1$.  Thus the cycle class map is injective on
$(\pi_2^{E^2})_*\operatorname{CH}^1(E^2)$.
The $N_2$ part has zero divisor cohomology by its Hodge types,
and therefore $\operatorname{CH}^1(N_2)=0$.
The Tate twist identity proves (F18.3); applying Chow groups
to (F18.2) proves (F18.4).

Finally the isogeny splitting on $h^1(Q)$ in F17 is additive.
Apply the preceding argument separately to each elliptic
$E$ summand.  For any other simple factor $Q_a$ no vanishing
follows just from the absence of $(2,2)$ classes in
$H^1(Q_a)\otimes H^3(N)$.  Even for $Q_a=E$, the remaining
$N_4$ component sits in the allowed middle, grade-zero part
of codimension two on the abelian fourfold $E^4$.
Neither the threefold Murre(D) results nor divisor
cycle-class injectivity applies to that component.
\end{proof}

\paragraph{What has and has not changed.}
F17 is the exact universal formulation of F16's residual
family obstruction; it is not counted by itself as a new
main-gate advance.  F18 does annihilate an additional specified
CM subcomponent in rational Chow.  It does not show that
every excess quotient has this CM factor, or that the
remaining extreme $N_4$ cycles vanish.  The full AJ5
finite-generation gate of C44-F07 remains open, and even
vanishing of the entire algebraic radical would leave its
possible non-algebraically-trivial kernel to be addressed.

\paragraph{Primary source locators and scope.}
The Fourier input is Beauville, \emph{Quelques remarques sur
la transformation de Fourier dans l'anneau de Chow d'une
vari\'et\'e ab\'elienne}, Lecture Notes in Mathematics~1016
(1983), 238--260: Proposition~1, printed p.~241, for the
degree interchange, and Proposition~3(i), printed p.~243,
for the inverse identity on rational Chow groups.
Primary author PDF:
\url{https://math.univ-cotedazur.fr/u/beauvill/pubs/fourier.pdf}.
The passage from cohomological projector identities to
distinguished Chow identities uses the exact O'Sullivan
locators recorded in F15.  The quotient Poincar\'e construction
uses Bost--Charles \S3.2, also checked in F15.

For the low-dimensional boundary, the carried source audit
of Vial's threefold theorem and Xu--Xu's Corollary~3.9 is
retained.  Vial, \emph{Remarks on motives of abelian type},
author version 20 July 2015, introductory Theorem~3,
printed p.~2, still makes Murre(D) for products of curves
an additional premise.  It does not prove the grade-zero
codimension-two injection on $E^4$ used in the unproved
last step above.
Primary author PDF:
\url{https://www.math.uni-bielefeld.de/~vial/abelian.pdf}.

% END work/arithmetic/fermat_residual_motive.tex

% BEGIN work/arithmetic/fermat_relative_transfer.tex
% C44_POS, /root/fermat_aj. 2026-11-22 UTC.
% This file proves the relative geometric input for the root's
% polarization-comparison argument. It does not assert that arbitrary
% abstract motivic injections preserve a chosen intersection form.

\section{Relative Fermat correspondences with good reduction}

\subsection{The exact geometric input}

\begin{proposition}[C44-F08]\label{C44-F08}
Fix an integer $d\geq2$, put
\[
 S_d=\operatorname{Spec}\mathbf Z[1/d,\zeta_{2d}],\qquad
 \epsilon=\zeta_{2d},\qquad \epsilon^d=-1,
\]
and write
\[
 F_n=\{x_0^d+\cdots+x_{n+1}^d=0\}\subset\mathbf P^{n+1}_{S_d},
 \quad C=F_1,\quad J=\operatorname{Pic}^0_{C/S_d}.
\]
There is a finite direct sum
\[
 M_n=\bigoplus_{(s,t)}h(J^{s}/S_d)(-t)^{\oplus a_{n,s,t}},
 \qquad a_{n,s,t}\in\mathbf Z_{\geq0},\quad s,t\geq0,
 \tag{F08.1}
\]
and relative rational Chow correspondences
\[
 \alpha_n:h(F_n/S_d)\longrightarrow M_n,
 \qquad
 \beta_n:M_n\longrightarrow h(F_n/S_d),\qquad
 \beta_n\alpha_n=\operatorname{id}.
 \tag{F08.2}
\]
Here $J^0=S_d$ and the convention is
$\mathbb L=\mathbf1(-1)$. All schemes in this assertion are smooth
and projective over $S_d$; $J$ is an abelian scheme. The same conclusion
holds for every finite product of the $F_n$, with the corresponding
products of Jacobians. These identities specialize to every geometric
fibre of characteristic $p\nmid d$, including characteristic $2$
when $d$ is odd.

In particular, for every such fibre and every $r$, an algebraic
correspondence extending over the same characteristic-zero model gives
a split injection
\[
 H^{2r}(F_n)\longrightarrow
 \bigoplus_{s,t}H^{2r-2t}(J^s)(-t)^{\oplus a_{n,s,t}}.
 \tag{F08.3}
\]
This assertion and its proof are unconditional. Conditional input 032
is needed only if the root's separate positive-form comparison is
used to infer the Hodge standard conjecture on the special fibre.
No assertion from Milne's unrefereed report is used here.
\end{proposition}

\begin{proof}
All the Fermat schemes are smooth over $S_d$: since $d$ is a unit,
the partial derivatives cannot vanish simultaneously at a projective
geometric point. The scheme $F_0$ consists of $d$ disjoint sections,
because the roots of $u^d=-1$ lie in the base ring and their pairwise
differences are units after $d$ is inverted. For example, this last
assertion follows from the discriminant of $u^d+1$, which is a unit
times $d^d$. Consequently
\[
 h(F_0/S_d)=\mathbf1^{\oplus d}.
 \tag{F08.4}
\]
The curve has the section $e=[1:\epsilon:0]$. Its relative Jacobian
therefore exists as an abelian scheme, and its normalized universal
Poincar\'e bundle exists on $C\times_{S_d}J$.

\paragraph{The relative Shioda--Katsura map.}
For $n\geq2$, use coordinates $[x_0:\cdots:x_n]$ on $F_{n-1}$
and $[y_0:y_1:y_2]$ on $C$. The rational map is
\[
 \phi_n:F_{n-1}\times C\dashrightarrow F_n,
 \quad (x,y)\longmapsto
 [x_0y_2:\cdots:x_{n-1}y_2:
       \epsilon x_ny_0:\epsilon x_ny_1].
 \tag{F08.5}
\]
Indeed, the sum of the $d$-th powers of the displayed coordinates is
\[
 y_2^d\sum_{i=0}^{n-1}x_i^d
      -x_n^d(y_0^d+y_1^d)
 =-y_2^dx_n^d+x_n^dy_2^d=0.
\]
Its base ideal is precisely the ideal $(x_n,y_2)$, with the equations
understood after local trivialization of the two hyperplane bundles.
To check this near the proposed centre, at least one $x_i$ with
$i<n$ is a unit in a projective chart and at least one $y_j$ with
$j<2$ is a unit. The sections defining the map then generate both
$y_2$ and $x_n$. Outside that centre at least one defining section
is nonzero. The centre is
\[
 Z_n=\{x_n=0\}\times\{y_2=0\}
       =F_{n-2}\times F_0
       \subset P_n:=F_{n-1}\times C.
 \tag{F08.6}
\]
It is smooth of relative codimension two. Hence
\[
 b_n:Y_n=\operatorname{Bl}_{Z_n}(P_n)\longrightarrow P_n
\]
is smooth and projective over $S_d$. The defining sections generate
an invertible ideal on the blowup, so (F08.5) extends to a morphism
$f_n:Y_n\longrightarrow F_n$.

This morphism has generic degree $d$ on every geometric fibre. To
see the degree explicitly, work on the nonempty target open set
\[
 A:=\sum_{i=0}^{n-1}z_i^d\ne0.
\]
On the exceptional divisor the first block of coordinates is a
scalar multiple of a point of $F_{n-2}$, so $A=0$ there. The
proper transforms of $x_n=0$ and $y_2=0$ also map into $A=0$.
Consequently every preimage on this open set has $x_ny_2\ne0$. Normalize
$x_n=y_2=1$. A preimage of a representative
$[z_0:\cdots:z_{n+1}]$ is then determined by
\[
 x_i=\lambda z_i\ (i<n),\qquad
 y_j=\lambda\epsilon^{-1}z_{n+j}\ (j=0,1),
 \qquad \lambda^d=-A^{-1}.
 \tag{F08.7}
\]
Both Fermat equations follow from this one scalar equation and the
target equation. Its $d$ solutions are distinct because $d$ is
invertible. They give distinct points, since both projective
normalizations have been fixed. Thus the map is finite \'etale of
degree $d$ on this open set, as also established in
Shioda--Katsura, Lemmas~1.1--1.3. Projection formula, or directly
composition of graph correspondences, now gives
\[
 (f_n)_*f_n^*=d\,\operatorname{id}_{h(F_n/S_d)}.
 \tag{F08.8}
\]
There is no assertion that $f_n$ is finite everywhere. The identity
is valid for a proper generically finite morphism between these
smooth schemes; contracted exceptional strata introduce no term in
the composition $(f_n)_*f_n^*$.

\paragraph{The codimension-two blowup correspondence.}
Let $j:E_n\hookrightarrow Y_n$ and $p:E_n\to Z_n$ be the exceptional
divisor and its projection. The relative blowup formula is
\[
 \Phi=(b_n^*,j_*p^*):
 h(P_n/S_d)\oplus h(Z_n/S_d)(-1)
       \xrightarrow{\ \sim\ }h(Y_n/S_d).
 \tag{F08.9}
\]
In this codimension-two case its inverse is explicitly
\[
 \Psi=((b_n)_*,-p_*j^*).
 \tag{F08.10}
\]
For completeness the four components of $\Psi\Phi$ can be checked
without a general inverse formula. They use
$(b_n)_*b_n^*=1$, $p_*p^*=0$ for the projective-line bundle $p$,
and
\[
 j^*j_*=-\xi\cdot,
 \qquad \xi=c_1(\mathcal O_{E_n}(1)),\qquad
 p_*(\xi p^*a)=a.
\]
The off-diagonal components are zero and the two diagonal
components are identities. The blowup formula for Chow groups
states that $\Phi$ is surjective as well. It holds after product
with every smooth projective test scheme over $S_d$; thus the
relative Manin identity principle gives $\Phi\Psi=1$ as Chow
correspondences. These are the usual intersection-theoretic
blowup identities over a regular base. Over a field they are
displayed, with an explicit inverse, in Vial, Proposition~5.1 and
Theorem~5.3. The same proof uses the relative intersection theory
of Fulton, \S\S20.1--20.2; no resolution of singularities is used.

Combining (F08.8)--(F08.10) gives a concrete split injection
\[
 h(F_n/S_d)\ 
 \underset{\frac1d(f_n)_*\Phi}{\stackrel{\Psi f_n^*}{\rightleftarrows}}\ 
 h(F_{n-1}/S_d)\otimes h(C/S_d)
       \oplus h(F_{n-2}/S_d)(-1)^{\oplus d}.
 \tag{F08.11}
\]
In particular, the projector onto the image in the right-hand
side is
\[
 \frac1d\,\Psi f_n^*(f_n)_*\Phi.
 \tag{F08.12}
\]
It is idempotent by (F08.8). This supplies the requested explicit
Chow projector rather than only an injection of realizations.

\paragraph{Induction to powers of the curve.}
Starting with (F08.4) and $h(F_1)=h(C)$, (F08.11) shows that
$h(F_n)$ is a retract of
\[
 \bigoplus_{t=0}^{\lfloor n/2\rfloor}
     h(C^{n-2t}/S_d)(-t)^{\oplus m_{n,t}},
 \tag{F08.13}
\]
where the nonnegative integers may be chosen recursively as
\[
 m_{0,0}=d,\quad m_{1,0}=1,\quad
 m_{n,t}=m_{n-1,t}+d\,m_{n-2,t-1},
 \tag{F08.14}
\]
with all unspecified indices zero. For example, the right sides
for $n=2,3$ are respectively
$h(C^2)\oplus\mathbb L^{\oplus d^2}$ and
$h(C^3)\oplus h(C)(-1)^{\oplus(d^2+d)}$.
The retractions are composites and tensor products of the
specified correspondences, so remain defined over the original
base.

\paragraph{A direct relative curve-to-Jacobian splitting.}
We record a construction that avoids needing a separate relative
isomorphism theorem for $h^1(C)$ and $h^1(J)$. The section $e$ has
trivial normal bundle. In the affine chart $x_0=1$, write
$u=x_1,v=x_2$ and $1+u^d+v^d=0$. Along $e$ the derivative
$du^{d-1}=d\epsilon^{d-1}$ is a unit, so the conormal of $e$
is freely generated by $v$. Consequently the standard relative
curve projectors
\[
 \pi_0=[e\times C],\qquad \pi_2=[C\times e],\qquad
 \pi_1=[\Delta_C]-\pi_0-\pi_2
 \tag{F08.15}
\]
are orthogonal idempotents. The only extra term that might occur
over a positive-dimensional base in their composition is the
self-intersection of the section, and it is zero by the trivial
normal bundle. Thus
\[
 h(C/S_d)=\mathbf1\oplus h^1(C/S_d)\oplus\mathbb L.
 \tag{F08.16}
\]

Let $a:C\to J$ send a point $x$ to
$\mathcal O_C(x-e)$, and let $\mathcal P$ be the universal
degree-zero line bundle on $C\times J$, normalized along
$e\times J$. The universal property of the relative Picard
scheme gives
\[
 (1_C\times a)^*\mathcal P
  \simeq\mathcal O_{C\times C}
              (\Delta_C-e\times C-C\times e).
 \tag{F08.17}
\]
Indeed, its restriction to a fibre of the second projection is
$\mathcal O_C(x-e)$, and the right-hand bundle is trivial along
$e\times C$, using the same normal-bundle trivialization. Those
two properties characterize the normalized Poincar\'e family.
Put $\gamma=c_1(\mathcal P)\in
\operatorname{CH}^1(C\times J)_{\mathbf Q}$, a degree-zero
correspondence from $C$ to $J$. Composition with the transpose
graph of $a$ is pullback in the second factor, so (F08.17) says
\[
 a^*\gamma=\pi_1.
 \tag{F08.18}
\]
Therefore
\[
 h^1(C/S_d)
 \underset{\pi_1a^*}{\stackrel{\gamma\pi_1}{\rightleftarrows}}
 h(J/S_d)
\]
is a split injection. The assertion also covers genus zero, when
the displayed primitive curve projector is zero. Combining with
(F08.16), $h(C)$ is a retract of
\[
 U=\mathbf1\oplus h(J)\oplus\mathbb L.
 \tag{F08.19}
\]
Expanding $U^{\otimes k}$ shows that each $h(C^k)(-t)$ is a
retract of a finite direct sum of terms $h(J^s)(-u)$ with
$s,u\geq0$. This turns (F08.13) into (F08.1)--(F08.2).

Finally, tensor products of split injections and retractions are
split injections and retractions. This proves the assertion for
products. Pullback to characteristic-zero fibres and to geometric
special fibres respects every graph, first Chern class, blowup
identity, and correspondence composition used above. Since no
prime outside those dividing $d$ was inverted, the same
construction works at each $p\nmid d$. This proves (F08.3).
\end{proof}

\subsection{Relative primitive projectors and scope of transfer}

\begin{lemma}[C44-F09]\label{C44-F09}
Let $n\geq1$, write $h=c_1(\mathcal O_{F_n}(1))$, and use rational
coefficients.
There is an explicit relative projector
\[
 p_n=[\Delta_{F_n}]
       -\frac1d\sum_{i=0}^{n}h^{n-i}\times h^i.
 \tag{F09.1}
\]
Its realization is the primitive middle cohomology of the Fermat
hypersurface. Algebraic relative correspondences represent the
Lefschetz operator and its lowering operator on $H^*(F_n)$.
For products of Fermat hypersurfaces, their tensor constructions
give the corresponding operators and primitive projectors for
every polarization that is a positive rational combination of the
factor hyperplane classes. These assertions are
unconditional and compatible with all the fibres of F08.
\end{lemma}

\begin{proof}
The $n+1$ coordinate sections $x_0,\ldots,x_n$ of
$\mathcal O_{F_n}(1)$ have no common zero on $F_n$. Thus
$h^{n+1}=0$ in the Chow ring. Moreover the degree of $h^n$ on
each fibre is $d$. The correspondences
\[
 q_i=\frac1d h^{n-i}\times h^i\quad(0\leq i\leq n)
\]
are consequently orthogonal idempotents, with image
$\mathbb L^i$. These assertions follow either by composing the
two factorized cycles or by the degree formula; an exponent
greater than $n$ contributes zero by $h^{n+1}=0$, and an exponent
less than $n$ contributes zero under proper pushforward to the
base. Formula (F09.1) is therefore idempotent.

Weak Lefschetz and Poincar\'e duality for a smooth hypersurface
identify all cohomology outside degree $n$ with these ambient
Tate classes. In degree $n$ they remove the ambient class when
$n$ is even. The remaining subspace is exactly the primitive
middle cohomology. This identification is valid in the complex
fibre and, by smooth proper base change and the relative
hyperplane class, in every geometric good fibre.

On the span of $1,h,\ldots,h^n$, set
\[
 L(h^i)=h^{i+1},\qquad
 \Lambda(h^i)=i(n-i+1)h^{i-1},
 \tag{F09.2}
\]
with the endpoint values zero. Set both operators zero on the
image of $p_n$. The maps between the displayed Tate summands
are algebraic relative correspondences. In realization, this
$L$ equals cup product with $h$: primitive middle classes are
annihilated by $h$. The operator $\Lambda$ is the usual
$\mathfrak{sl}_2$ lowering operator. In particular, inverse
Hard Lefschetz maps are algebraic and extend over $S_d$.

For a product polarization $\sum_j a_jh_j$, with rational
$a_j>0$, take $L=\sum_j a_jL_j$ and
$\Lambda=\sum_j a_j^{-1}\Lambda_j$ on the tensor product.
Their commutator is the sum of the factor degree operators,
so these are the raising and lowering operators for the same
$\mathfrak{sl}_2$ representation. The first realizes cup
product by the specified polarization.
Finite-dimensional $\mathfrak{sl}_2$ representation theory over
$\mathbf Q$ expresses the Lefschetz decomposition projectors
and the inverses of each nonzero Lefschetz map as rational
combinations of these operators and degree projectors. This
supplies algebraic relative correspondences with the stated
realizations.

We do not assert that the diagonal cup correspondence composed
with $p_n$ vanishes in the Chow group. The operators just
constructed have the required cohomological actions; the
possible homologically trivial difference is irrelevant to
the claimed realization statements.
\end{proof}

\begin{corollary}[C44-F10: Fermat positivity under the stated abelian input]
\label{C44-F10}
Let $k$ be an algebraic closure of a finite field of characteristic
$p$, let $p\nmid d$, and let
\[
 X=\prod_{j=1}^{m}F_{n_j,d,k},\qquad n_j\geq1.
\]
Conditional on the unrefereed 032/Milne 2026 inputs specified in
C44-A08, the Hodge standard conjecture holds for $X$ with
respect to every ample class $h=\sum_j a_jh_j$, where
$a_j\in\mathbf Q_{>0}$ and $h_j$ is the pullback of the Fermat
hyperplane class. In particular this proves it for every power
of every positive-dimensional Fermat hypersurface with its product hyperplane
polarization. If every $n_j\geq3$, the conclusion holds for
every ample class of $X$.

This deduction uses C44-A08 and F08--F09. It does not use
Milne2022 Proposition~4.2, a finite-cover positivity transfer,
or a claim that all algebraic cycles of the special fibre lift.
\end{corollary}

\begin{proof}
Choose a closed point of $S_d$ above $p$ and localize at that
point. The result is a mixed-characteristic discrete valuation
ring whose fraction field embeds into $\mathbf C$. Base
extension to its geometric special fibre gives the specified
Fermat schemes over $k$. After multiplying $h$ by a positive
integer, it is the class of a relatively ample line bundle on
this model. Scaling by such an integer changes every relevant
primitive intersection form by a positive scalar and hence
does not affect the asserted sign.

For each even degree $2r\leq\dim X$, take the full primitive
summand for this relative ample class. By F09 its projector
and inverse Lefschetz map are represented by algebraic
correspondences over the same base. Poincar\'e duality and
the inverse Lefschetz map give the inverse of the restricted
primitive pairing, also algebraically. F08 gives a split
injection of the entire cohomology of $X$ into the stated
Tate-twisted cohomology of abelian schemes, and hence an
injection on this primitive summand. All the hypotheses of
C44-A08 hold. Its conclusion is precisely positivity of
\[
 (-1)^r\int_X xyh^{\dim X-2r}
\]
on the rational algebraic primitive classes, with exactly the
conditional input recorded there. This proves the first
assertion for all primitive degrees.

Suppose now that every $n_j\geq3$. Weak Lefschetz gives
$H^1(F_{n_j,d,k})=0$ and
$H^2(F_{n_j,d,k},\mathbf Q_\ell)=\mathbf Q_\ell h_j$.
The N\'eron--Severi cycle map is injective after tensoring
with $\mathbf Q_\ell$, so each factor has rational
N\'eron--Severi group spanned by its hyperplane class.
K\"unneth and $H^1=0$ give the same assertion for the product:
its rational N\'eron--Severi group has basis the $h_j$.
Thus any rational ample class is $\sum_j a_jh_j$ with
rational coefficients. Restricting it to a curve in the
$j$-th factor with all other coordinates fixed shows
$a_j>0$. Such a class is covered by the first assertion.
For factors of dimension two, extra divisor classes can
occur and this last argument is not asserted. For products
of curves, cross terms in the N\'eron--Severi group can also
occur; the first assertion remains as stated.
\end{proof}

\paragraph{What these results do and do not transfer.}
F08 provides a relative split injection into cohomology of
abelian schemes with good reduction. F09 provides the relative
primitive projectors and algebraic inverse Lefschetz maps.
The positive-form comparison in C44-A08 compares the
chosen primitive cup form with the pullback of the abelian
Weil-star form, using positivity of the generic complex Hodge
polarizations and constancy of the comparison endomorphism's
characteristic polynomial. Its use of 032/Milne 2026 remains conditional
on those supplied unrefereed preprints. No implication from an
arbitrary abstract motivic injection to preservation of a
chosen positive form is asserted here.

The distinction between finite covers and this construction is
essential: $f_n$ contracts strata. Pullback of an ample class
can be nef without being ample. The proof of F08 never uses
positivity of that pullback, and F09 uses the original ample
hyperplane classes instead. Thus neither an unproved nef limit
argument nor Milne2022 Proposition~4.2 is a hidden input to the
geometric statements.

\subsection{Primary-source locators}

\begin{itemize}
\item T. Shioda and T. Katsura, \emph{On Fermat varieties},
Tohoku Math. J. \textbf{31} (1979), 97--115:
Lemma~1.1 and formula~(1.3), p.~98, give the degree-$d$ map;
Lemma~1.2, p.~99, resolves it by the specified blowup;
Lemma~1.3, p.~100, gives the finite \'etale open part;
Theorem~1.7 and diagram~(1.25), pp.~101--102, give the full
blowup/quotient/blowdown construction. Most decisively,
Remark~1.8, p.~102, explicitly states that the construction
works over $\mathbf Z[1/d,\zeta_{2d}]$ in the category of
smooth projective schemes. The primary PDF was opened and
Remark~1.8 also checked in the page image:
\url{https://www.jstage.jst.go.jp/article/tmj1949/31/1/31_1_97/_pdf}.
\item C. Vial, \emph{Algebraic cycles and fibrations},
Doc. Math. \textbf{18} (2013), 1521--1553. The author PDF
consulted has 24 pages: Proposition~5.1 on p.~13 is the Chow
blowup formula; Theorem~5.3 on p.~14 verifies the explicit
motivic inverse by Manin's identity principle. It is a
field-base source; the relative use in F08 is via the same
intersection-theoretic identities over the regular base,
not an expansion of Vial's stated hypotheses.
\url{https://www.math.uni-bielefeld.de/~vial/algcycles-fibrations.pdf}.
\item W. Fulton, \emph{Intersection theory}, Proposition~6.7
and \S\S20.1--20.2, for the blowup identities and intersection
theory over a regular base. The applicability of this
relative framework to Chow motives over Dedekind bases is
also stated in J. Laga and A. Shnidman, \emph{Ceresa cycles of
bielliptic Picard curves}, \S4.2, citing K. K\"unnemann,
\emph{Arakelov Chow groups of abelian schemes}, \S1,
Remark~1.1. The supplied Laga--Shnidman source was checked;
its special Picard-type motive isomorphism is not used.
\item V. Guletski\u\i, \emph{Chow motives of abelian type over a
base}, Eur. J. Math. \textbf4 (2018), 1065--1086, \S3
defines relative rational Chow correspondences over a regular
Dedekind base, citing precisely Fulton \S\S20.1--20.2 and
K\"unnemann Remark~1.1, p.~367. The paragraph immediately
before Lemma~3.1 identifies the section self-intersection
correction to the two curve Tate projectors. In F08 this
correction is zero by the explicit trivialization of the
Fermat section's normal bundle. Section~2 gives the normalized
relative Poincar\'e bundle framework.
\url{https://link.springer.com/article/10.1007/s40879-018-0270-9}.
\end{itemize}

% END work/arithmetic/fermat_relative_transfer.tex

% BEGIN work/arithmetic/fermat7_frobenius.tex
\section{C44-F11: a degree-seven Fermat block with extra special classes}
\label{sec:C44-F11}
\paragraph{Statement and dependency flags.}
Let \(X/\mathbb F _2\) be the Fermat fourfold
\[
 X:\quad x_0^7+\cdots+x_5^7=0\subset\mathbb P^5,\qquad T=X\times X,
\]
and put \(K=\mathbb Q(\zeta _7)\). Work over \(\mathbb F _8\) when using
the diagonal automorphisms. There is an explicit relative rational Chow
projector \(e_\alpha\), of cohomological rank six, with
\(\alpha=(1,1,1,3,3,5)\), such that:
\begin{enumerate}
\item In the characteristic-zero Fermat lifting its Hodge types are
\((3,1),(2,2),(1,3)\), each twice, and it has no nonzero rational Hodge
class. On the special fibre Frobenius \(8\) is exactly \(8^2\) on this
whole rank-six block.
\item The fourfold is not supersingular: the character
\(\beta=(1,1,1,1,1,2)\) has normalized Newton slopes \(1,3\).
\item The rank-36 block \(W=e_\alpha\otimes e_\alpha\) in the primitive
middle cohomology of \(T\) splits into six rational rank-six blocks
\(W_r\), \(r=1,\ldots,6\). Only \(W_6\) has rational Hodge classes on the
generic fibre; its entire rank-six space is spanned by explicit
correspondence classes.
\item \emph{Conditional on the supplied 032/Milne-2026 input through
Tate for finite-field abelian varieties}, the special block \(e_\alpha\)
and hence \(W\) are spanned by algebraic cycle classes. Thus \(W\)
contains thirty independent special algebraic directions, in
\(\bigoplus_{r=1}^5W_r\), which cannot be specializations of generic
algebraic cycle classes in the fixed Fermat lifting, even after finite
extension of its base trait.
\end{enumerate}
The first three assertions are unconditional. The last assertion does
not construct equations for the thirty cycles. It proves their
existence and independence under the stated conditional input. It makes
no assertion about lifting those classes in a different deformation.
C44-F13 below strengthens this algebraicity assertion unconditionally
by explicit ruled surfaces.

\begin{proof}
We use the relative Fermat model over
\(\mathbb Z[1/7,\zeta _{14}]\) of C44-F08 and a mixed-characteristic trait
above \(2\). The group \(G=\mu _7^6/\mu _7\) acts diagonally.
Its primitive degree-four characters are the tuples
\(a=(a_0,\ldots,a_5)\), with each \(a_i\in\{1,\ldots,6\}\) and
\(\sum a_i\equiv0\pmod7\). Each occurs once. For a character \(\chi_a\),
the usual group-ring idempotent
\[
 e_a=\frac1{|G|}\sum_{g\in G}\chi_a(g)^{-1}g^*
\]
is defined after adjoining coefficient field \(K\). The orbit sum
\[
 e_{[a]}=\sum_{t=1}^{6}e_{ta}
\]
has rational coefficients and is an actual relative Chow idempotent.
Every such orbit has six elements, since one nonzero coordinate already
forces the stabilizer in \((\mathbb Z/7)^\times\) to be trivial.
The projector is self-adjoint: transposition takes \(e_a\) to \(e_{-a}\).
The group algebra acts on its realization through the field \(K\), and
that realization is a rank-one \(K\)-module.

The classical Fermat character formula is
\[
 q(a)=\frac{\sum_i\langle a_i\rangle}{7}-1,\qquad
 H_a\subset H^{4-q(a),q(a)}.
\]
Its primary locator is Ran, \emph{Cycles on Fermat hypersurfaces},
Compositio Math. 42 (1980), Proposition 1.7, printed p.~125:
\url{https://www.numdam.org/article/CM_1980__42_1_121_0.pdf}.
For \(\alpha=(1,1,1,3,3,5)\) one obtains, in the order \(t=1,\ldots,6\),
\[
 (q(t\alpha))_{t=1}^6=(1,2,1,3,2,3).
\]
A rational Hodge class would generate, under \(\mathbb Q[G]\), a
nonzero rational subrepresentation consisting of Hodge classes.
The orbit representation is irreducible over \(\mathbb Q\), so this
would force its whole realization to have type \((2,2)\), contradicting
the displayed types. Thus its rational Hodge space is zero.

Here is an exact Jacobi-sum calculation. Write
\(\mathbb F _8=\mathbb F _2[\tau]/(\tau^3+\tau+1)\), choose
\(\chi(\tau)=\zeta _7\), and let
\(\psi(x)=(-1)^{\operatorname{Tr}_{\mathbb F _8/\mathbb F _2}(x)}\).
Set \(G_a=\sum_{x\in\mathbb F _8^\times}\chi(x)^a\psi(x)\).
Squaring \(x\) shows that \(G_a\) is constant on
\(H=\{1,2,4\}\) and on \(-H=\{3,5,6\}\). Put
\[
 A=\zeta _7+\zeta _7^2+\zeta _7^4,\quad
 B=\zeta _7^3+\zeta _7^5+\zeta _7^6,\quad \delta=A-B.
\]
The seven nonzero elements have additive-character values \(-1\) at
\(1\), \(+1\) at \(\tau,\tau^2,\tau^4\), and \(-1\) at
\(\tau^3,\tau^5,\tau^6\). Hence
\[
 A+B=-1,\quad AB=2,\quad \delta^2=-7,\qquad
 G_H=-1+\delta,\quad G_{-H}=-1-\delta,\quad G_HG_{-H}=8.
\]
The primitive Frobenius eigenvalue in dimension four is
\[
 j(a)=\frac{\prod_{i=0}^5G_{a_i}}8.
\]
This is the even-dimensional convention in Shioda--Katsura,
\emph{On Fermat varieties}, Tohoku Math. J. 31 (1979),
Section 3, pp.~107--108, formulas (3.1)--(3.10):
\url{https://www.jstage.jst.go.jp/article/tmj1949/31/1/31_1_97/_pdf}.
The five-variable convention and its identification with Frobenius
are also stated in van der Geer--Katsura,
\emph{On the height of Calabi--Yau varieties in positive characteristic},
Documenta Math. 8 (2003), Section 5, printed p.~106:
\url{https://ems.press/content/serial-article-files/25920}.
That paragraph treats arbitrary Fermat dimension and degree; the
Calabi--Yau restriction is imposed only on p.~107 after the Jacobi-sum
and valuation formulas. Equivalently, fix the first nonzero coordinate
to \(1\) and use the five-variable Jacobi sum; the product of the five
remaining characters is nontrivial, and \(G_aG_{-a}=8\) gives the formula.

The tuple \(\alpha\) has three entries in \(H\) and three in \(-H\).
Every multiple \(t\alpha\) retains this balance. Consequently
\[
 j(t\alpha)=8^3/8=64=8^2\qquad(t=1,\ldots,6).
\]
For \(\beta\), all six entries lie in \(H\), so
\[
 j(\beta)=(-1+\delta)^6/8=36-20\delta.
\]
It has polynomial \(z^2-72z+4096\). Its Newton polygon at \(2\), with
vertices \((0,12),(1,3),(2,0)\), gives valuations \(9\) and \(3\).
After division by \(v_2(8)=3\) these are slopes \(3\) and \(1\).
This proves nonsupersingularity.

The characters of \(W\) are all \((t\alpha,s\alpha)\). Their rational
orbits are
\[
 W_r:\quad\{(t\alpha,rt\alpha):t=1,\ldots,6\},
 \qquad r\in(\mathbb Z/7)^\times.
\]
They give six mutually orthogonal relative rational projectors.
Their Hodge second indices are
\[
\begin{array}{c|rrrrrr}
 r\backslash t&1&2&3&4&5&6\\ \hline
 1&2&4&2&6&4&6\\
 2&3&5&4&4&3&5\\
 3&2&5&3&5&3&6\\
 4&4&3&3&5&5&4\\
 5&3&3&2&6&5&5\\
 6&4&4&4&4&4&4
\end{array}
\]
The same irreducibility argument gives no rational generic Hodge
class for \(r\ne6\). The block \(W_6\) is entirely of type \((4,4)\).
Let \(g\) multiply the first coordinate of \(X\) by \(\zeta _7\).
The six projected graph correspondences \(e_\alpha\Gamma_{g^j}\),
\(j=0,\ldots,5\), are supported cohomologically in \(W_6\).
Under Poincare duality they realize \(1,g,\ldots,g^5\) on the rank-one
\(K\)-module. They are independent since \(g\) has minimal polynomial
\(\Phi_7\), of degree six; hence they span \(W_6\).
Each factor \(e_\alpha H^4(X)\) is primitive middle cohomology and is
annihilated by its hyperplane class. Thus \(W\subset P^8(T)\) for
every positive weighted product polarization.

For the last assertion only, impose the supplied conditional
032/Milne-2026 input. C44-F08 gives an actual algebraic injection and
retraction from the Fermat motive into a finite sum of Tate twists
of abelian motives. Tate surjectivity passes to the retract: inject
an invariant class, express its image by algebraic classes, and
apply the algebraic retraction. Since all of \(e_\alpha H^4(X)(2)\)
is Frobenius-\(8\)-invariant, it is spanned by six independent
algebraic classes. Their exterior products span all thirty-six
dimensions of \(W(4)\). The projectors for \(r=1,\ldots,5\) produce
the thirty-dimensional algebraic subspace in the statement.

Specialization commutes with the relative projectors. Every generic
algebraic class has zero projection to those five blocks by the Hodge
calculation. Its specialization therefore has zero projection there.
Every nonzero special algebraic class in their sum consequently
fails to lift as an algebraic cycle class in this Fermat model.
The same argument applies after every finite base extension.
\end{proof}

\paragraph{Status restriction.}
No primary source found here certifies that \(I(X\times X)\) for
this degree-seven example was previously unknown. It is therefore
not labelled an established \emph{B-known/I-unknown} example.
The following result proves a sign statement without claiming
historical priority.

\section{C44-F12: Frobenius symmetry removes the rank-six sign ambiguity}
\label{sec:C44-F12}
\paragraph{Statement and separation of equivalences.}
For \(X/\overline{\mathbb F}_2\) above and \(T=X\times X\):
\begin{enumerate}
\item Unconditionally, \(B(X)\), \(B(T)\), \(I_{\rm num}(X)\), and
\(I_{\rm num}(T)\) hold for every ample class.
\item Unconditionally, numerical and \(13\)-adic homological
equivalence agree on all algebraic cycles of \(X\) and \(T\).
Consequently \(I_{13}(X)\) and \(I_{13}(T)\) hold for every ample class.
\item Conditional on the supplied 032/Milne-2026 input, the same
cohomological statement holds for every \(\ell\ne2\), by C44-F10
(or by the algebraic abelian retraction and conditional \(D_\ell\)
for abelian varieties). C44-F14 below removes this condition for
\(X\), using additional explicit cycles; it does not remove it for \(T\).
\end{enumerate}
The unconditional proof uses the published Ancona--Marmora theorem,
not 032, Milne-2026, or Milne-2022 Proposition 4.2. This proof itself
does not assert Tate surjectivity for the block in C44-F11.

\paragraph{Primary source.}
Ancona--Marmora, \emph{The Hilbert symbol in the Hodge standard conjecture},
Comment. Math. Helv. 100 (2025), 619--662, DOI 10.4171/CMH/591:
\url{https://ems.press/content/serial-article-files/51595?nt=1}.
Section 2.3, pp.~623--624, uses homological motives obtained from
Chow motives. Theorem 3.2 and Remark 3.3, p.~626, give the signature
congruence and zero-or-full-dimension alternative; Lemma 3.4,
pp.~626--627, gives the full-rank comparison with realizations.
The matching author version is arXiv:2210.14001v3, 20 January 2025
(internal date 22 January), printed pp.~5--7:
\url{https://irma.math.unistra.fr/~marmora/2210.14001v3.pdf}.

\paragraph{Rank-six numerical sign lemma.}
Let \(M\) be a relative homological motive over the integers of a finite
extension of \(\mathbb Q_2\), with
\(q:\operatorname{Sym}^2M\to\mathbf1\). Suppose its realization has
dimension six, \(K=\mathbb Q(\zeta _7)\) acts on \(M\), the adjoint
on \(K\) is conjugation, and \(q_B\) polarizes its Hodge structure.
Suppose an algebraic endomorphism \(u\) of the special fibre is a
\(q\)-isometry on numerical classes and is semilinear for
\(\sigma_4:\zeta _7\mapsto\zeta _7^4\). Then \(q\) is positive
definite on numerical classes of the special fibre.

\begin{proof}
Numerical classes have finite dimension and descend to finite residue
extensions. Extend the residue field and trait enough to define a
basis. Put \(V=\operatorname{Hom}_{\rm NUM}(\mathbf1,M_k)\).
If \(V=0\), there is nothing to prove. Otherwise the \(K\)-action and
Ancona--Marmora Remark 3.3 give \(\dim_{\mathbb Q}V=6\), so \(V\)
has rank one over \(K\). Lemma 3.4 identifies \(q|_V\), after realization,
with the nondegenerate form of \(M_k\); it is therefore nondegenerate.

Choose a \(K\)-basis. Trace duality and the adjoint relation give
\[
 q(x,y)=\operatorname{Tr}_{K/\mathbb Q}(c\,x\bar y),
 \qquad c\in K^+\setminus\{0\}.
\]
Indeed represent \(q(x,1)\) by the trace pairing, use the adjoint
relation to recover the form, and use symmetry to obtain \(c=\bar c\).
Nondegeneracy forces \(c\ne0\). Over \(\mathbb R\) this is the orthogonal
sum of three planes indexed by embeddings \(\tau:K^+\hookrightarrow
\mathbb R\), with forms \(2\tau(c)\operatorname{Re}(z\bar w)\).
Each plane is definite, of sign \(\tau(c)\).

The restriction of \(\sigma_4\) to \(K^+\) has order three and acts
transitively on its real embeddings. The semilinear isometry \(u\)
therefore sends these three planes cyclically and their signs agree.
The negative index is \(0\) or \(6\). All four hypotheses of
Ancona--Marmora Theorem 3.2 hold, so this index is divisible by four.
It must be \(0\).
\end{proof}

\paragraph{Separate \(13\)-adic nondegeneracy lemma.}
Let \(U\) be rank one over \(K\otimes\mathbb Q_{13}\), with a
nondegenerate symmetric \(\mathbb Q_{13}\)-form \(q\) whose adjoint
on \(K\) is conjugation. Suppose \(U\) is realized by an algebraic
primitive projector, and let \(A\) be the \(\mathbb Q_{13}\)-span
of its algebraic classes. Then \(q|_A\) is nondegenerate. In particular
every numerically zero projected algebraic class is \(13\)-adically zero.

\begin{proof}
Since \(13\equiv-1\pmod7\), \(K\otimes\mathbb Q_{13}\) is the product
of three unramified quadratic fields, each preserved by conjugation.
For an explicit residue-field check,
\[
 \Phi_7(z)\equiv
 (z^2+6z+1)(z^2+5z+1)(z^2+3z+1)\pmod {13}.
\]
The discriminants \(6,8,5\) are nonsquares modulo \(13\).
The three field idempotents are self-adjoint and split \(U\)
orthogonally into three nondegenerate rank-one field modules.

Algebraic classes are stable under the algebraic \(K\)-action, so
their \(\mathbb Q_{13}\)-span is a \(K\otimes\mathbb Q_{13}\)-submodule.
In a rank-one module over a field a submodule is zero or the whole
module. Hence \(A\) is the sum of a subset of the three orthogonal
nondegenerate pieces and \(q|_A\) is nondegenerate.
A numerically zero algebraic class pairs to zero with all algebraic
classes in \(A\), because \(q\) is intersection with algebraic
polarization powers and projectors. It pairs to zero with their
\(\mathbb Q_{13}\)-span and therefore vanishes. No positivity or
Tate assertion is used in this lemma.
\end{proof}

\paragraph{Application to the fourfold and its square.}
\begin{proof}
C44-F09 supplies relative primitive projectors and Lefschetz operators,
hence \(B(X)\) and \(B(T)\) unconditionally. Write
\[
 H^*(X)=A\oplus P,\qquad
 A=\langle1,h,h^2,h^3,h^4\rangle,\quad P=P^4(X).
\]
The Lefschetz representation \(A\) is the length-five irreducible
string; \(P\) is trivial and concentrated in degree four. Every
primitive character occurs once with a free six-element rational
orbit. Each orbit gives a self-adjoint rank-six relative homological
motive, rank one over \(K\). This uses no Tate surjectivity.

For \(L=a h_1+b h_2\) on \(T\), with \(a,b>0\) rational, the
primitive decomposition is
\[
\begin{array}{c|l}
 \text{source}&\text{primitive summands}\\ \hline
 A\otimes A&\text{one Tate line in each degree }0,2,4,6,8\\
 P\otimes A&P\otimes1\text{ in degree }4\\
 A\otimes P&1\otimes P\text{ in degree }4\\
 P\otimes P&P\otimes P\text{ in degree }8 .
\end{array}
\]
The first row is Clebsch--Gordan for two length-five strings;
equivalently each ambient primitive kernel in degree \(2r\le8\)
has dimension one. In the second row \(h_1P=0\), so its length-five
string starts at \(P\otimes1\); the third row is symmetric. Both
operators vanish in the fourth row. These arguments work for all
\(a,b>0\). C44-F09 makes the projectors algebraic and relative at
the homological-motive level.

Split the non-Tate parts by rational characters of \(G\times G\).
Every nontrivial character in rows two and three occurs once.
Every character pair in row four occurs once with orbit length
six: a nonzero coordinate forces its scalar stabilizer to be trivial.
Thus all non-Tate primitive orbit motives have rank one over \(K\).
Distinct orbits are orthogonal because the cup product pairs a
character only with its negative, in the same rational orbit.
Ambient Tate classes are orthogonal to nontrivial characters.

For a rank-six primitive motive in degree \(2r\), take its Tate
twist \((r)\) and the relative algebraic form
\[
 q_r(x,y)=(-1)^r\int_TxyL^{8-2r}
\]
(and the analogous form on \(X\)). Diagonal automorphisms preserve
\(L\), so the adjoint on \(K\) is conjugation. Hodge--Riemann gives
a Betti polarization, and restriction to an orthogonal rational
sub-Hodge structure preserves this. In degree four on \(T\) the
sign is \(+1\): on \(P\otimes1\) the form is \(7b^4\int_Xuv\);
on \(1\otimes P\) it is \(7a^4\int_Xuv\). Indeed only the term
\(b^4h_2^4\), respectively \(a^4h_1^4\), survives, and
\(\int_Xh^4=7\). These are positive multiples of the generic
primitive degree-four polarization; the two pieces are orthogonal.

Let \(F\) be coordinate-squaring on \(X\) or \(T\), the base change
of Frobenius of its \(\mathbb F _2\)-model, considered over the chosen
residue field. For dimension \(N\),
\[
 F^*L=2L,\qquad \deg F=2^N,\qquad
 q_r(F^*x,F^*y)=2^{2r}q_r(x,y).
\]
The last equality follows by pulling back \(xyL^{N-2r}\) and using
the degree formula. Thus \(u=2^{-r}F^*\) is an algebraic isometry
on numerical classes. For any diagonal order-seven automorphism \(g\),
\(F\circ g=g^2\circ F\), hence
\[
 F^*g^*=(g^4)^*F^*.
\]
The orbit projector is invariant under multiplication by four.
Therefore \(F^*\) preserves each orbit motive and is
\(\sigma_4\)-semilinear on its \(K\)-action. The numerical sign
lemma proves positivity on every non-Tate primitive block.

The ambient primitive Tate lines have hyperplane-generated classes
and intersection numbers unchanged by specialization. Their generic
forms are positive by Hodge--Riemann, hence so are their special forms.
Orthogonality proves \(I_{\rm num}\). The algebraic Lefschetz
decomposition also holds modulo numerical equivalence, since
cohomological identities imply numerical identities. Thus the
primitive numerical space defined by the Lefschetz kernel is exactly
the sum of the spaces just treated, rather than a smaller chosen space.

At \(\ell=13\), apply the separate nondegeneracy lemma to every
rank-six primitive block; the ambient one-dimensional blocks also
have nondegenerate algebraic forms. Project a numerically zero cycle
by the algebraic Lefschetz operators to these primitive pieces,
with the appropriate twists. Each projected class is numerically
zero, hence \(13\)-adically zero. Reassembly proves \(D_{13}\)
in every codimension. Combining this with \(I_{\rm num}\) gives \(I_{13}\).

Finally weak Lefschetz gives \(H^1(X)=0\) and
\(\operatorname{NS}(X)_{\mathbb Q}=\mathbb Qh\). Kunneth gives
\(\operatorname{NS}(T)_{\mathbb Q}=\mathbb Qh_1\oplus\mathbb Qh_2\).
Restriction to a curve in either factor shows that every ample
class on \(T\) has positive coefficients \(a,b\); every ample class
on \(X\) is a positive multiple of \(h\). Thus all ample classes
are covered.

This proof does not transfer \(D_{13}\) to other primes or assert
the existence of classes in a Frobenius-invariant block whose
numerical algebraic space could be zero. The all-\(\ell\) result
for \(T\) remains conditional on 032/Milne-2026 as above;
the additional construction below proves more for \(X\).
\end{proof}

\section{C44-F13: explicit ruled surfaces make the extra classes unconditional}
\label{sec:C44-F13}
\paragraph{Statement.}
The special algebraicity assertion of C44-F11 is unconditional.
The ruled surface
\[
 Z=\overline{\left\{
 [\lambda s:\lambda t:\mu w^2:\mu s^2:\mu t^2:\lambda w]:
 s^7+t^7+w^7=0
 \right\}}\subset X
\]
has nonzero projection to every character line in \(e_\alpha H^4(X)\).
If \(g_0\) scales \(x_0\) by \(\zeta _7\), the six rational cycles
\[
 z_j=e_\alpha(g_0^j)^*[Z],\qquad 0\le j\le5,
\]
form a \(\mathbb Q_\ell\)-basis of this block for every \(\ell\ne2\)
and are defined over \(\mathbb F _8\). Their products and the
projectors \(W_r\) give thirty independent algebraic directions in
\(\bigoplus_{r=1}^5W_r\subset P^8(T)\), none of whose nonzero classes
can specialize from generic algebraic classes in the fixed Fermat
lifting, even after finite base extension. This uses neither
032/Milne-2026 nor a Tate conjecture.

\begin{proof}
Let \(C:s^7+t^7+w^7=0\subset\mathbb P^2\). Partition the coordinates
of \(X\) into \((x_0,x_1,x_5)\) and \((x_2,x_3,x_4)\). The character
\(\alpha\) becomes the concatenation
\[
 u=(1,1,5),\qquad v=(1,3,3).
\]
Both sum to zero modulo seven and occur in \(H^1(C)\).

Over the relative Fermat base of C44-F08 form
\[
 Y=\mathbb P_{C\times C}
 \bigl(\operatorname{pr}_1^*\mathcal O_C(-1)
       \oplus\operatorname{pr}_2^*\mathcal O_C(-1)\bigr),
 \qquad p:Y\longrightarrow C\times C.
\]
The tautological lines lie in complementary coordinate spaces
and define a morphism \(j:Y\to X\), the projective join. The
equation of \(X\) holds because it holds separately on each \(C\).
Thus
\[
 J=j_*p^*:H^1(C)\otimes H^1(C)(-1)\longrightarrow H^4(X)
\]
is an algebraic relative correspondence and sends divisor classes
on \(C\times C\) to codimension-two cycle classes on \(X\).

In characteristic zero this is Ran's operation in Section 4,
printed pp.~138--140. Proposition 4.6 and Corollary 4.7 on p.~140
prove injectivity on primitive cohomology and identify character
pairs with their concatenation:
\url{https://www.numdam.org/article/CM_1980__42_1_121_0.pdf}.
Ran works over \(\mathbb C\); the above correspondence, however,
is defined between smooth projective families over the trait.
It induces a map of lisse \(\mathbb Q_\ell\)-cohomology sheaves.
Injectivity and nonvanishing on each character line therefore
persist to the special fibre by smooth proper base change.
This holds for every \(\ell\ne2\), including \(\ell=7\), since the
coefficient field has characteristic zero. Diagonal equivariance
is compatible with concatenation: the characters under consideration
sum to zero on each coordinate block and so respect the two
blockwise diagonal quotients.

On the special curve define
\[
 \sigma(s,t,w)=(w,s,t),\qquad F(s,t,w)=(s^2,t^2,w^2),
 \qquad f=\sigma\circ F.
\]
Here \(F\) is the base change of the Frobenius of the
\(\mathbb F _2\)-model, viewed as a morphism over \(\mathbb F _8\).
Pullback by \(\sigma\) sends \((a,b,c)\) to \((b,c,a)\),
while \(F^*\) sends \(a\) to \(2a\). Therefore
\[
 f^*H_v=H_{2(3,3,1)}=H_{(6,6,2)}=H_{-u}.
\]
For every \(t=1,\ldots,6\), \(f^*:H_{tv}\to H_{-tu}\) is an
isomorphism: \(\sigma\) is an isomorphism and Frobenius pullback
is invertible on \(\ell\)-adic cohomology.

Let \(\Gamma_f\subset C\times C\) have first coordinate the source.
For \(\eta,\theta\in H^1(C)\),
\[
 \int_{C\times C}[\Gamma_f](\eta\otimes\theta)
   =\int_C\eta f^*\theta.
\]
Choose nonzero \(\eta\in H_{-tu}\), \(\theta\in H_{-tv}\).
Then \(f^*\theta\) is nonzero in \(H_{tu}\) and pairs nontrivially
with \(\eta\), since opposite one-dimensional character spaces
have a nondegenerate Poincare pairing. Hence the
\(H_{tu}\otimes H_{tv}\) component of \([\Gamma_f]\) is nonzero
for every \(t\). The join is nonzero on each such line, so
\(J[\Gamma_f]\) has nonzero component in every \(H_{t\alpha}\).
The two other degree-two Kunneth summands of \(C\times C\)
have a trivial character on one factor; equivariance prevents
them from contributing to or cancelling these components.

The inverse image of \(\Gamma_f\) is
\(\mathbb P_C(\mathcal O_C(-1)\oplus\mathcal O_C(-2))\).
Its map to \(X\) is the displayed parameterization of \(Z\),
homogeneous with these two line-bundle weights. It lies in \(X\)
since, in characteristic two,
\[
 \lambda^7(s^7+t^7+w^7)
 +\mu^7(w^{14}+s^{14}+t^{14})=0.
\]
On \(\lambda\ne0\), the first block recovers \([s:t:w]\) and
the second recovers the fibre coordinate. The map is birational
onto \(Z\), so \(J[\Gamma_f]=[Z]\) with multiplicity one.

The projector is explicitly rational:
\[
 e_\alpha=
 \frac1{|\ker\chi_\alpha|}\sum_{g\in\ker\chi_\alpha}g^*
 -\frac1{|G|}\sum_{g\in G}g^*,
 \qquad |G|=7^5,\quad|\ker\chi_\alpha|=7^4.
\]
It is the sum of the six nontrivial characters of the quotient
\(G/\ker\chi_\alpha\), hence precisely the C44-F11 orbit projector.
Every \(z_j\) is an explicit rational combination of diagonal
translates of \(Z\), defined over \(\mathbb F _8\).
Over a coefficient extension containing \(\zeta _7\), write
\[
 e_\alpha[Z]=\sum_{t=1}^6c_t,\qquad
 0\ne c_t\in H_{t\alpha}.
\]
Since \(\alpha_0=1\), \(z_j=\sum_t\zeta _7^{jt}c_t\).
The six-by-six matrix \((\zeta _7^{jt})_{0\le j\le5,\ 1\le t\le6}\)
is Vandermonde with distinct nodes and is invertible.
Thus the \(z_j\) form a basis after coefficient extension,
and hence over \(\mathbb Q_\ell\), for every \(\ell\ne2\).

Their thirty-six exterior products form a basis of
\(e_\alpha H^4(X)\otimes e_\alpha H^4(X)\).
Applying the relative rational projectors \(W_r\) spans each
of the six blocks and in particular the thirty-dimensional
sum with \(r\ne6\). The generic Hodge obstruction and the
specialization argument of C44-F11 apply unchanged to these
five blocks, proving the claimed nonliftability in this
Fermat model after every finite base extension.
\end{proof}

\section{C44-F14: all Tate classes of the degree-seven fourfold}
\label{sec:C44-F14}
\paragraph{Statement.}
For \(X/\mathbb F _2\), the Tate conjecture holds unconditionally
in every codimension and over every finite extension of
\(\mathbb F _2\). The codimension-two geometric Tate space has
dimension \(2221\): the ambient line and \(2220\) primitive
character lines. Of those primitive lines, \(1860\) come from
linear planes and \(360\) from diagonal translates and coordinate
permutations of the C44-F13 ruled surface.
Consequently \(D_\ell(X)\) and \(I_\ell(X)\) hold unconditionally
for every \(\ell\ne2\) and every ample class.
The \(360\) exceptional directions have zero generic rational
Hodge space in the fixed Fermat lifting and cannot specialize
from generic algebraic classes. This does not establish all-\(\ell\)
\(D\) or \(I\) for \(X\times X\).

\begin{proof}
Use \(G_H=-1+\delta\), \(G_{-H}=-1-\delta\),
\(H=\{1,2,4\}\), and \(\delta^2=-7\) from C44-F11.
For a primitive character \(a\), let \(h(a)\) count its entries in
\(H\). The Jacobi-sum formula gives
\[
 j(a)=\frac{G_H^{h(a)}G_{-H}^{6-h(a)}}8
     =64\,\rho^{h(a)-3},\qquad \rho=G_H/G_{-H}.
\]
The polynomial of \(G_H\) is \(z^2+2z+8\), with \(2\)-adic
Newton valuations \(1,2\). Therefore \(\rho\) has valuation
\(1\) or \(-1\) at a prime above \(2\) and is not a root of
unity. A primitive character is geometric Tate exactly when
\(h(a)=3\), and then its Frobenius-\(8\) eigenvalue is \(64\).
All geometric Tate classes are thus fixed over \(\mathbb F _8\)
after the Tate twist.

Here is a hand classification. Let \(a_1,a_2,a_4\) count the
entries \(1,2,4\), and let \(b_1,b_2,b_4\) count the opposite
entries \(6,5,3\). Both sums equal three. Put \(d_i=a_i-b_i\).
The character-sum condition is
\[
 d_1+d_2+d_4=0,\qquad d_2+3d_4\equiv0\pmod7,\qquad
 -3\le d_i\le3.
\]
For \(d_4=0\), these force all \(d_i=0\).
For \(d_4=1,2,3\), the range of \(d_2\) forces, respectively,
\[
 (d_1,d_2,d_4)=(2,-3,1),\quad(-3,1,2),\quad(-1,-2,3).
\]
Negative \(d_4\) gives their negatives. Each nonzero vector
has positive and negative parts of total three, so it uniquely
determines all the counts \(a_i,b_i\). The zero vector gives
the ten common count vectors \(a_i=b_i\) of sum three.

The ten zero-difference cases are precisely three opposite pairs.
The six exceptional sorted tuples are
\[
\begin{gathered}
 (1,1,1,3,3,5),\quad(1,1,4,5,5,5),\quad(1,2,2,3,3,3),\\
 (2,2,2,3,6,6),\quad(2,4,4,6,6,6),\quad(4,4,4,5,5,6).
\end{gathered}
\]
They are exactly the sorted scalar multiples of \(\alpha\).
Every exceptional character is therefore a coordinate permutation
of a scalar multiple of \(\alpha\). C44-F13, coordinate permutations,
rational orbit projectors, and diagonal translates algebraically
span the entire exceptional subspace.

The opposite-pair characters are spanned by linear-plane classes.
Indeed iterate the projective join of three copies of the
zero-dimensional degree-seven Fermat variety in \(\mathbb P^1\).
Its primitive degree-zero characters are \((b,-b)\), and point
classes span its degree-zero cohomology. Joining three chosen
points gives a plane in \(X\). Ran Proposition 4.6 and Corollary 4.7,
spread out as in C44-F13, show that these primitive join maps
are nonzero on every concatenated character line. Thus linear
planes span all opposite-pair characters in every good fibre.
This uses an actual join correspondence, not a conjectural
specialization-surjectivity assertion.

For the ordered count, three paired patterns have multiplicities
\((3,3)\), six have \((2,2,1,1)\), and one has six distinct entries.
Their contribution is
\[
 3\frac{6!}{3!3!}+6\frac{6!}{2!2!}+6!
 =60+1080+720=1860.
\]
Each exceptional pattern has multiplicities \((3,2,1)\);
the six patterns contribute \(6\cdot6!/(3!2!)=360\).
Thus the primitive Tate space has dimension \(2220\),
and adding \(h^2\) gives \(2221\).

All constructed cycles are defined over \(\mathbb F _8\) and
span the geometric Tate space after every \(\ell\)-adic realization.
Over another finite field \(\mathbb F _{2^m}\), pass to the common
extension containing \(\mathbb F _8\) and average the cycles over
its Galois group. With rational coefficients the normalized norm
projects onto the relevant Frobenius invariants. This proves
codimension-two Tate over every finite extension. The other
codimensions have only ambient one-dimensional cohomology,
spanned by powers of the hyperplane class.

The primitive geometric Tate space is nondegenerate for cup
product: it consists exactly of \(h(a)=3\) character lines,
is closed under \(a\mapsto-a\), and each line pairs perfectly
with its opposite. Since algebraic classes span it, every
numerically zero cycle has zero class for every \(\ell\ne2\).
The ambient pieces are nondegenerate as well. Thus \(D_\ell(X)\)
holds in every codimension, and the unconditional numerical
positivity of C44-F12 gives \(I_\ell(X)\) for every \(\ell\ne2\).
All ample classes are positive multiples of \(h\), as in C44-F12.

Finally a rational generic Hodge character must satisfy
\(\sum_i\langle ta_i\rangle=21\) for every unit \(t\) modulo seven.
Summing over \(t\in H\) gives
\[
 63=7h(a)+14(6-h(a)),
\]
so \(h(a)=3\). Our classification therefore includes every
possible generic rational Hodge character. The opposite-pair
characters have type \((2,2)\) throughout their rational orbit.
Every exceptional character is a coordinate permutation and
scalar multiple of \(\alpha\), whose orbit has mixed Hodge
types by C44-F11 and no nonzero rational Hodge class.
Thus the entire \(360\)-dimensional exceptional special subspace
has zero generic algebraic-class specialization image in
the fixed Fermat lifting.
\end{proof}

% END work/arithmetic/fermat7_frobenius.tex

% BEGIN work/scan_w/contribution.tex
% C44_POS /root/scan_w.  New deductions, 2026-11-22 Europe/Berlin.
% The model is GPT-6 Astra Pro.  Runtime reasoning-effort setting is not exposed.
% This file is a paper fragment.  The root paper owns theorem environments.

\section{An algebraic replacement for Wirtinger in the Weil-plane detector argument}

The statements in this section distinguish an unconditional intersection
estimate from the positivity and algebraicity inputs.  Every invocation of
032, Corollaries 8.3 or 8.4, is conditional on the unrefereed preprints
specified by the user.  The characteristic-zero Theorem W in
\texttt{HODGE\_DIVISION\_ATTACK\_20261121.md}, R4, is not being reproved.
We give its missing positive-characteristic degree estimate and then record
the resulting, slightly differently normalized, detector theorem.

\subsection*{C44-W01. A characteristic-free intersection bound on an abelian variety}

\begin{proposition}[C44-W01]
Let $A$ be an abelian variety of dimension $2n$ over an algebraically
closed field, and let $h=c_1(L)$ be an integral ample class.  Put
\[
 D(Z)=\deg(h^n\cdot Z),\qquad
 V=\frac{\deg h^{2n}}{(2n)!}=\chi(L)\in\mathbf Z_{>0},\qquad
 b_h=\frac{2^n}{n!\sqrt V}.
\]
For effective integral cycles $Z,Z'$ of dimension $n$, including cycles
with multiplicities,
\begin{equation}\label{c44:algebraic-bezout}
 0\leq\deg(Z\cdot Z')\leq b_h^2D(Z)D(Z').
\end{equation}
No standard conjecture and no lifting hypothesis is used.
\end{proposition}

\begin{proof}
It suffices by bilinearity to treat integral subvarieties.  On
$Z\times Z'$ consider $f(z,z')=z-z'$ and $s(z,z')=z+z'$.  Translation
moving, or equivalently refined pullback of a general point of $A$ by
$f$, gives
\[
 f_*[Z\times Z']=e[A],\qquad e=\deg(Z\cdot Z')\geq0.
\]
If $f$ is not dominant, both sides have coefficient $e=0$.  If it is
dominant, it is generically finite and $e$ is its degree, including any
inseparable multiplicity.  To see explicitly that the intersection
coefficient agrees with this degree, a general translate of $Z'$ meets
$Z$ properly and avoids the singular loci of both subvarieties: the
bad incidence loci have dimension at most $2n-1$.  At the resulting
zero-dimensional intersections the two smooth subvarieties are local
complete intersections in the smooth variety $A$, and properness gives
the usual length intersection multiplicity.  The same description
applies to the generic fibre of $f$.  Translation preserves cycle
classes.  This also proves the assertion without separability or
transversality assumptions.

Let $h_1,h_2$ be the pullbacks of $h$ to $Z\times Z'$, and set
$\alpha=f^*h$, $\beta=s^*h$.  Both are nef.  The theorem of the cube,
in its parallelogram consequence, gives the numerical identity
\[
 \alpha+\beta=2(h_1+h_2).
\]
Here $[-1]^*h=h$ numerically; no symmetric representative of the
original line bundle is required.  A convenient exact source is
Milne, \emph{Abelian Varieties}, Corollary 6.5, applied to the maps
$z,z',-z'$; the numerical consequence is the displayed identity.
All mixed intersections of nef divisors with an effective cycle are
nonnegative, by approximation by ample rational classes.  Therefore
\begin{align*}
 e\deg h^{2n}
 &=\deg(\alpha^{2n}\cap[Z\times Z'])\\
 &\leq\deg((\alpha+\beta)^{2n}\cap[Z\times Z'])\\
 &=2^{2n}\binom{2n}{n}D(Z)D(Z').
\end{align*}
The last equality uses $\dim Z=\dim Z'=n$.  Division by
$\deg h^{2n}=(2n)!V$ proves the bound.  The asserted formula and
integrality for $V$ are the abelian Riemann--Roch theorem, Milne,
Theorem 13.3(b); ampleness gives $V>0$.
\end{proof}

\paragraph{Realizations and historical scope.}
Equation \eqref{c44:algebraic-bezout} is an inequality between rational
intersection numbers.  It is valid for classes represented in either
$\ell$-adic or crystalline cohomology because the cycle map preserves
these products and trace.  It does not put an order on $\mathbf Q_\ell$.
This is precisely the kind of degree bound needed below.  In
Grothendieck's framework its operations are proper pushforward, refined
intersection, and the biextension identity; SGA 7 I, Exp. VIII,
4.12.1 and 2.3.11 provide the biextension convention already used by
JAC-T4.3.  It supplies a new quantitative input to the positive
incidence calculation of cote 44, scans 008--009, and to Theorem W,
step 3.  It makes no assertion that a new axis class lifts.

\subsection*{C44-W02. The signed Weil-plane Gram matrix in any characteristic}

Fix an imaginary quadratic field $K$ and an order
$\mathcal O\subseteq K\cap\operatorname{End}(A)$ such that the Rosati
involution of $h$ induces conjugation on $K$.  Choose
$a\in\mathcal O$ with $a/\bar a$ not a root of unity.  Write
\[
 \nu=a\bar a\in\mathbf Z_{>0},\quad q=\nu^n,\quad
 s=a^{2n}+\bar a^{2n}\in\mathbf Z,\quad H=4q^2-s^2>0.
\]
For an invertible prime $\ell$, define
\[
 W_\ell=\bigwedge_{K\otimes\mathbf Q_\ell}^{2n}H^1(A,\mathbf Q_\ell)
 \subset H^{2n}(A,\mathbf Q_\ell).
\]
It has dimension two over $\mathbf Q_\ell$.  Let
\[
 f(t)=(t-a^{2n})(t-\bar a^{2n}),\quad
 g(t)=\prod_{j=1}^{2n-1}(t-a^j\bar a^{2n-j}).
\]
Both polynomials have integer coefficients.  Choose $R>0$ and
$u,v\in\mathbf Z[t]$ with $uf+vg=R$, for example the absolute
resultant with the sign of the Bezout relation adjusted.  Put
\[
 p(t)=v(t)g(t)=\sum_jc_jt^j,\qquad
 \kappa=\sum_j|c_j|\nu^{jn},\qquad e_W=R^{-1}p([a]^*).
\]
These are exactly the projector data of the supplied Theorem W.
We use them, with the same normalization, on the new realization.

\begin{proposition}[C44-W02; conditional on unrefereed 032, Corollary 8.4]
Let $Y\subset A$ be an integral $n$-dimensional subvariety with
$e_W[Y]\ne0$, and let $D=D(Y)$.  Then
\[
 y_1=p([a]^*)[Y],\qquad y_2=[a]^*y_1
\]
are independent integral algebraic cycle classes in $W_\ell$.
The form
\[
 Q(x,y)=(-1)^n\deg(xy)
\]
is positive definite on their rational span.  With $d=Q(y_1,y_1)$,
its Gram matrix is exactly
\begin{equation}\label{c44:weil-gram}
 G=d\begin{pmatrix}1&s/2\\s/2&q^2\end{pmatrix},\qquad
 \det G=d^2H/4\in\mathbf Z_{>0},\qquad
 0<d\leq b_h^2R\kappa D^2\leq b_h^2\kappa^2D^2.
\end{equation}
Moreover
\[
 W_{\rm alg}:=e_W\operatorname{Alg}^n(A)
    =\mathbf Qy_1\oplus\mathbf Qy_2,
 \qquad W_{\rm alg}\otimes\mathbf Q_\ell=W_\ell.
\]
The two classes have integral signed-cycle representatives whose
positive and negative parts have degrees at most $\kappa D$ and
$q\kappa D$, respectively.
\end{proposition}

\begin{proof}
After extending the coefficient field, write $H^1=V_a\oplus V_{\bar a}$.
Both eigenspaces have dimension $2n$, and the Rosati condition says
that $h$ lies in $V_a\otimes V_{\bar a}$.  Each of the two extreme
exterior powers is therefore annihilated by $h$.  Thus $W_\ell$ is
middle primitive.  The eigenspaces of $[a]^*$ in middle cohomology
have eigenvalues $a^j\bar a^{2n-j}$; all are distinct by the choice
of $a$.  Hence the integral polynomial $p$ is $R$ on the extreme
two lines and zero on the other weight spaces.  This verifies the
projector on this realization, rather than importing a Betti lattice.
The exterior-algebra description is valid integrally for etale
cohomology by Milne, Theorem 15.1.

Every $[a^j]$ is a finite flat isogeny, so the pullbacks in
$p([a]^*)[Y]$ are effective integral cycles with their flat
multiplicities.  Rosati gives $[a^j]^*h=\nu^jh$ numerically and
$\deg[a^j]=\nu^{2nj}$.  Projection formula consequently gives
\[
 D([a^j]^*Y)=\nu^{jn}D.
\]
This proves both asserted degree budgets.  The minimal polynomial
of $[a]^*$ on $W_\ell$ is the irreducible quadratic $f$: indeed
$a^{2n}$ is nonreal.  Hence $y_1,y_2$ are rationally independent.

The signed Hodge standard conjecture, used with exactly the sign
$(-1)^n$, is positive definite on every rational subspace of
primitive algebraic classes.  This is the sole input for positivity
here.  It is the even-degree version of the primitive sign in
cote 44, scan 008; JAC-T4.3 establishes the corresponding incidence
polarization by a curve test.  The unsigned form printed in
Milne's report, Section 3.6, is not used.

Let $T=[a]^*|W_\ell$.  Projection formula gives
$Q(Tx,Ty)=q^2Q(x,y)$.  Thus $T^\dagger=q^2T^{-1}$ and
$T+T^\dagger=s$.  It follows that
$Q(y_1,y_2)=sd/2$ and $Q(y_2,y_2)=q^2d$, proving the matrix
formula.  Its entries are integers, since the $y_i$ are integral
cycle classes, and its determinant is a positive integer.
There is a useful improvement over expanding every pair of effective
pullbacks.  On all of $H^{2n}(A)$, the intersection adjoint of
$[a]^*$ is $q^2[a]^{*-1}$.  The pairing therefore matches the
weight indexed by $j$ with the weight indexed by $2n-j$.
The projector $e_W$ selects both extreme weights and annihilates
their orthogonal complement, so it is self-adjoint.  Consequently
$p([a]^*)^\dagger p([a]^*)=R p([a]^*)$ on middle cohomology.
Apply C44-W01 to each term of $[Y]\,p([a]^*)[Y]$ to obtain
\[
 d=|\deg(y_1^2)|=R\left|\deg\bigl([Y]p([a]^*)[Y]\bigr)\right|
 \leq b_h^2R\kappa D^2.
\]
Finally $p(a^{2n})=R$ and $|a^{2n}|=q$ give $R\leq\kappa$
by the ordinary triangle inequality on this explicit algebraic number.
This numerical inequality on projector coefficients is not an
archimedean realization of a cycle class.
This proves the quantitative replacement for Wirtinger.

Finally any finite rationally independent collection of classes in
$W_{\rm alg}$ has a positive definite rational Gram matrix by the
same Hodge-standard input.  Its determinant is nonzero, so those
classes are also $\mathbf Q_\ell$-independent.  There can be at most
two.  The two classes already constructed therefore span
$W_{\rm alg}$ and its realization.
\end{proof}

\paragraph{What ``integral'' means in positive characteristic.}
There is no canonical integral Betti lattice over $\overline{\mathbf F}_p$.
Let $\Lambda$ denote the group of integral algebraic middle classes
modulo numerical equivalence, viewed in cohomology using 032's
numerical--homological equality.  It is a lattice: choose finitely
many integral test cycles giving a nondegenerate rational pairing;
the intersection map embeds $\Lambda$ in a finite-rank free abelian
group.  Write
\[
 \Lambda^\sharp=\{x\in\operatorname{Alg}^n(A):\deg(x\lambda)\in\mathbf Z
                       \text{ for all }\lambda\in\Lambda\}.
\]
The denominator condition needed below is $\delta t\in\Lambda^\sharp$.
The stronger condition $\delta t\in\Lambda$ suffices.  If an integral
compatible realization is specified, its rational classes with integral
pairings also satisfy this condition.  Integrality only at one prime
$\ell$ is insufficient to make the Cramer numerators integers.

\subsection*{C44-W03. A quadratic detector bound in characteristic $p$}

\begin{theorem}[C44-W03; conditional on unrefereed 032, Corollary 8.4]
Retain the preceding notation and suppose $W_{\rm alg}\ne0$.
For $0\ne t\in W_{\rm alg}$ let $\delta\geq1$ be an integer with
$\delta t\in\Lambda^\sharp$, and put $T_0=Q(t,t)>0$.
Explicitly, the detector degree and division complexity are
\begin{align*}
 \delta_W&=\min\{D(Y):Y\subset A\text{ integral of dimension }n,
                                      \ e_W[Y]\ne0\},\\
 c_t&=\min\{\max(N,D(Z^+),D(Z^-)):
       N\in\mathbf Z_{\geq1},\ Z^\pm\text{ effective integral cycles},
                              \ Nt=[Z^+]-[Z^-]\}.
\end{align*}
These are the definitions in the supplied Theorem W; an empty minimum
is interpreted as infinity.  Cycle equalities here are equalities of
classes, and an integral cycle may have several components with
nonnegative integer multiplicities.
Set $B=b_h\sqrt{R\kappa}$.  Then
\begin{equation}\label{c44:weil-quadratic}
 \delta_W\leq c_t\leq
 \delta\max\{HB^2,\,4q\kappa B\sqrt{HT_0}\}\,\delta_W^2.
\end{equation}
In particular one can replace $B$ everywhere by the uniform constant
$B_0=2^n\sqrt{R\kappa}/n!$, independent of the abelian variety, the
characteristic, and the polarization type.  The quartic version
\begin{equation}\label{c44:weil-quartic}
 c_t\leq\delta q^2\max\{B^4,\,4\kappa B^3\sqrt{T_0}\}\,\delta_W^4
\end{equation}
also holds.  Thus the characteristic-$p$ division theorem has both
the requested detector formulation and a quadratic sharpening,
with fully algebraic, explicit constants.
\end{theorem}

\begin{proof}
The hypothesis $W_{\rm alg}\ne0$ supplies a detector: express a
nonzero class as a rational combination of integral effective
cycles and apply $e_W$.  Choose a detector $Y$ of minimal degree
$D=\delta_W$, and form $y_1,y_2$ as in C44-W02.  Put
$b_i=Q(\delta t,y_i)\in\mathbf Z$.  Inverting the exact Gram matrix
\eqref{c44:weil-gram} gives the identity
\begin{equation}\label{c44:exact-clearing}
 \delta Hd\,t=(4q^2b_1-2sb_2)y_1+(4b_2-2sb_1)y_2.
\end{equation}
All coefficients and the positive multiplier $N=\delta Hd$ are
integers.  If $t=t_1y_1+t_2y_2$, Cauchy--Schwarz in this
positive definite rational plane, or the diagonal entries of
$G^{-1}$, gives
\[
 |t_1|\leq\frac{2q\sqrt{T_0}}{\sqrt{dH}},\qquad
 |t_2|\leq\frac{2\sqrt{T_0}}{\sqrt{dH}}.
\]
Therefore the two integral coefficients on the right of
\eqref{c44:exact-clearing} have absolute values at most
$2\delta q\sqrt{HT_0d}$ and $2\delta\sqrt{HT_0d}$.
Using the degree budgets of C44-W02, the positive and negative
parts of the resulting integral cycle each have degree at most
\[
 \kappa D\bigl(2\delta q\sqrt{HT_0d}
                     +2\delta q\sqrt{HT_0d}\bigr)
 \leq4\delta q\kappa B\sqrt{HT_0}\,D^2.
\]
Also $N\leq\delta HB^2D^2$.  These two estimates prove the
upper bound in \eqref{c44:weil-quadratic}.

For completeness the ordinary two-by-two adjugate clearing with
multiplier $\delta\det G$ gives
$N\leq\delta q^2B^4D^4$.  Its two coefficients are bounded by
$2\delta\sqrt{T_0}\|y_1\|\|y_2\|^2$ and
$2\delta\sqrt{T_0}\|y_1\|^2\|y_2\|$, where
$\|y_1\|=\sqrt d\leq BD$ and $\|y_2\|=q\sqrt d$.
Their degree budget is at most
$4\delta q^2\kappa B^3\sqrt{T_0}D^4$, proving
\eqref{c44:weil-quartic}.

Conversely any class equation
$Nt=[Z^+]-[Z^-]$, with $N\geq1$ and effective integral cycles,
has a nonzero Weil projection on at least one integral component
of $Z^+$ or $Z^-$.  That component has degree at most the
corresponding cycle degree.  Hence $\delta_W\leq c_t$.
\end{proof}

\paragraph{Exact change to the supplied Theorem W.}
The original complex constant $|R|/n!$ in the norm estimate is
replaced by $2^n\sqrt{R\kappa}/(n!\sqrt V)$ for the integral numerator
$y_1$.  We do not claim that Wirtinger's sharper constant survives
unchanged.  All other projector, degree, independence, and integer
clearing steps are now valid without an archimedean realization.
The bound is on cohomological or numerical cycle classes, exactly
as in the supplied definition of $c_t$; it is not an equality in
the Chow group modulo rational equivalence.

\subsection*{C44-W04. The precise balanced-signature hypothesis over a finite field}

\begin{proposition}[C44-W04; Tate-algebraicity conclusion conditional on unrefereed 032]
Let $(A,h,K)$ descend to a finite field $\mathbf F_Q$, with the
$K$-action and polarization defined there, and $\dim A=2n$.
If $p$ splits in $K$, assume the two rational $K\otimes\mathbf Q_p$
components of $A[p^\infty]$ both have dimension $n$ (each has
height $2n$).  If $p$ is inert or ramified, no additional determinant
balance is necessary.  Then Frobenius acts with finite order on
$W_\ell(n)$.  Conditional on 032, Corollary 8.3, the whole
$W_\ell$ is algebraic, $W_{\rm alg}$ is its two-dimensional
rational cycle structure, $\delta_W<\infty$, and C44-W03 applies
to every nonzero class in that rational Weil plane.

When $p$ is split and $p$ does not divide the conductor of
$\mathcal O$, the stated $p$-divisible balance is exactly the usual
$(n,n)$ multiplicity condition on $\operatorname{Lie}A$.
It also holds for the special fibre of a balanced polarized
$\mathcal O_K\otimes\mathbf Z_p$-linear abelian scheme in mixed
characteristic.  These hypotheses give the usual meaningful
all-characteristic version of the signature condition.
\end{proposition}

\begin{proof}
Write $\pi$ for the $Q$-power Frobenius endomorphism.  Its
$K$-linear determinant on $H^1$ is an element $\alpha\in K$,
independent of the realization.  One can see its rationality
without choosing a complex lift: the traces of $b\pi^j$ for
$b\in K$ are rational, by the characteristic polynomial theorem
for endomorphisms, and the nondegenerate trace pairing of $K$
recovers the $K$-linear traces.  Newton identities recover the
$K$-linear characteristic polynomial and its determinant.
The same argument in crystalline cohomology gives the same
determinant: in any of these realizations, the exterior-algebra
description of abelian cohomology and the top-degree action give
$\det(m-[b])=\deg([m]-[b])$ for every integer $m$ for which the
map is an isogeny.  Polynomial identity then makes the
characteristic polynomial independent of the realization.
The trace/characteristic-polynomial input is
Milne, \emph{Abelian Varieties}, Proposition 12.9; it is the
degree-defined characteristic polynomial of an endomorphism.

The determinant $\alpha$ is an algebraic integer, since its
eigenvalues are products of Frobenius eigenvalues.  The
polarization pairing gives
\[
 \alpha\bar\alpha=Q^{2n},\qquad |\alpha|=Q^n
\]
at the two complex embeddings.  The absolute-value assertion is
Weil's theorem for abelian varieties, Milne, Theorem 19.1(c).
In particular $\alpha$ is a unit away from primes over $p$.
If $p$ is nonsplit, its unique prime is fixed by conjugation, so
the norm identity implies
$v(\alpha)=n\,v(Q)$ there.  If $p$ splits, the determinant of
Frobenius on each rational Dieudonne summand has normalized
determinant valuation equal to that summand's dimension:
the Newton polygon of a $p$-divisible group of height $h_0$
and dimension $d_0$ ends at $(h_0,d_0)$, so its slopes sum
to $d_0$.  For exact source statements see Oort,
\emph{Abelian Varieties over Finite Fields}, Sections
21.16--21.19 and Corollary 9.3.  Equivalently each
Dieudonne--Manin factor of dimension $d_0$ and height
$d_0+c_0$ has $d_0+c_0$ slopes $d_0/(d_0+c_0)$.
The balance hypothesis therefore gives
$v(\alpha)=n\,v(Q)$ at both primes.

Consequently $\alpha/Q^n$ is a unit in the imaginary quadratic
field $K$ of absolute value one, hence a root of unity.  On the
two Weil lines the Frobenius eigenvalues are $\alpha$ and
$\bar\alpha$.  A finite extension of the finite field kills
their normalized root-of-unity factors.  The Tate conjecture
in the stated 032 input then gives algebraic classes spanning
$W_\ell$: apply the algebraic projector $e_W$ to cycles spanning
the Tate classes.  C44-W02 identifies their rational cycle
span as a plane.  This proves the assertions.

For split $p$ and a $p$-maximal order, the two integral
idempotents act on the $p$-divisible group and on its Lie
algebra, so the dimensions agree with the two Lie
multiplicities.  In an integral balanced family these ranks
are constant.  Isogenies preserve the dimensions and heights
of the rational $p$-divisible summands.
\end{proof}

\begin{proposition}[C44-W05: a bad-order ambiguity in the word signature]
A determinant identity on $\operatorname{Lie}A$ for an order whose
conductor is divisible by a split prime $p$ does not by itself
imply the balance hypothesis of C44-W04.  In particular the
literal naive such identity cannot replace that hypothesis.
\end{proposition}

\begin{proof}
For a completely concrete example take
$E:y^2=x^3-x$ over $\mathbf F_5$, $K=\mathbf Q(i)$.
It has eight rational points, hence Frobenius trace $-2$;
in particular it is ordinary.  Its order-four automorphism
gives the $\mathbf Z[i]$-action, and its Frobenius is one
of $-1\pm2i$.  Thus $p=5$ splits in $K$, as required.
The following construction also works with any ordinary
CM elliptic curve at a split prime.  After a finite
extension all data are defined over $\mathbf F_Q$.
Let $A=E^{2n}$ with the scalar $K$-action and product polarization,
and restrict the action to $\mathcal O=\mathbf Z+p\mathcal O_K$
(or to $\mathbf Z+p\operatorname{End}(E)$ if necessary).
For $a=m+pb$ the action on the Lie algebra is multiplication by
$m\bmod p$, and the reductions of $a$ and $\bar a$ agree.
Thus
\[
 \det(T-a\mid\operatorname{Lie}A)
 =(T-a)^n(T-\bar a)^n
\]
holds after reduction for every $a\in\mathcal O$, although the
two rational $p$-divisible components have dimensions $2n$ and
$0$.  If $\pi$ is the ordinary Frobenius of $E$, the normalized
Weil determinant is $\pi^{2n}/Q^n$; its valuations at the two
primes over $p$ are nonzero and opposite.  It is not a root of
unity, and the two Weil lines are not Tate.  The example
therefore proves the claimed insufficiency.  It is not a
counterexample to C44-W03 or to the usual balanced PEL meaning
of Weil type.
\end{proof}

\section{Geometric positive transfers and special complete intersections}

The abstract abelian-type transfer remains distinct from a transfer that
preserves the actual incidence form. The following statements give
such geometric transfers and identify the extra special-fibre directions
that they do not reach. Their positivity conclusions are conditional
on the unrefereed 032 input whenever that qualification is displayed.

\subsection*{C44-W06. A finite abelian cover gives an explicit inherited polarization}

\begin{theorem}[C44-W06; conditional on unrefereed 032, Corollary 8.4]
Let $X$ be smooth projective of odd dimension $d=2m-1$ over an
algebraically closed field. Suppose that there is a finite surjective
morphism $f:A\to X$ from an abelian variety. For every rational ample
class $h$ on $X$, put $h_A=f^*h$ and $e=\deg f$. The primitive
incidence quotient $B_X=P^m(X)_{\rm prim}$ has the signed rational
polarization of cote 44, Th.~4. More precisely the induced map
\[
 i_f:B_X\longrightarrow B_A=P^m(A)_{\rm prim}
\]
has finite kernel, and its actual incidence classes satisfy
\begin{equation}\label{c44:finite-cover-polarization}
 i_f^*\ell_{B_A}=e\ell_{B_X},\qquad
 \ell_{B_X}=\frac{(-1)^m}{2}\Delta^*d_{B_X}.
\end{equation}
In particular this is a positive transfer on a concrete class of
nonabelian varieties of geometric abelian type. It does not use
Milne's additional transfer assertion from JAC-P20.
\end{theorem}

\begin{proof}
A finite morphism between smooth varieties of the same dimension is
flat; its degree includes inseparable multiplicity. Push--pull gives
$f_*f^*=e$, and $h_A$ is ample. Lefschetz inverses for $X$ are
algebraic: in degree $i\leq d$, an inverse to $L_h^{d-i}$ is
\[
 \frac1e f_* (L_{h_A}^{d-i})^{-1}f^*.
\]
Indeed $f^*$ commutes with the indicated Lefschetz operators and
the formula composes with $L_h^{d-i}$ to the identity. The inverse
on the abelian variety is algebraic. The same pull--push construction
with its Kunneth projectors supplies those of $X$. Thus the middle
primitive projector is algebraic and self-transpose. Saito's
incidence construction and transpose compatibility, used in
JAC-T4.1, give its nondegenerate factor $B_X$.

Pullback takes primitive cycle families on $X$ to primitive cycle
families on $A$, since $h_Af^*=f^*h$. It therefore induces $i_f$.
Pushforward induces a rational left inverse up to $e$, so $i_f$ has
finite kernel. For two normalized families $z,z'$ on parameter
varieties, projection formula gives the equality of divisorial
incidence correspondences
\[
 (p_{12})_*\bigl((1\times f)^*z\cdot(1\times f)^*z'\bigr)
   =e\,(p_{12})_*(z\cdot z').
\]
By the incidence functoriality of JAC-T4.1, this is the equality
$(i_f\times i_f)^*d_{B_A}=e d_{B_X}$. Diagonal pullback gives
\eqref{c44:finite-cover-polarization}, with the same sign and
normalization as JAC-T4.3. That theorem makes $\ell_{B_A}$ ample
conditional on 032. Pullback by a homomorphism with finite kernel
is ample on its source, and division by $e>0$ preserves rational
ampleness. This proves the conclusion. At the motive level
$e^{-1}f_*f^*=1$ exhibits $h(X)$ as an actual Chow-motive summand
of $h(A)$, so the stated abelian-type scope is also explicit.
\end{proof}

\paragraph{Concrete nonabelian examples.}
In characteristic different from two, take
$A=E_0\times E_1\times\cdots\times E_{d-1}$ with $d\geq3$ odd,
choose a nonzero two-torsion point $a\in E_0$, and quotient by
\[
 (x_0,x_1,x_2,\ldots)\longmapsto(x_0+a,-x_1,x_2,\ldots).
\]
The action is free because of its first coordinate. Its quotient
is smooth projective, the cover is finite etale, and the invariant
$H^1(\mathcal O)$ has dimension $d-1$, so the quotient is not an
abelian variety. In characteristic two use a nonzero three-torsion
translation on $E_0$ and an order-three automorphism of
$E_1:y^2+y=x^3$, acting by $x\mapsto\zeta_3x$. The same argument
works. These examples prove that C44-W06 has nonabelian content in
every characteristic.

\subsection*{C44-W07. The fixed part for every smooth complete intersection}

\begin{proposition}[C44-W07; fixed-part positivity conditional on unrefereed 032]
Let $A$ have dimension $N=2m+c-1$, let $h$ be ample, and let
$k:Y\hookrightarrow A$ be a smooth successive intersection of
$c$ divisors of class $h$, with $\dim Y=n=2m-1$. There is no
generic-last-section hypothesis in this statement. Define
\[
 r=(L_h^c)^{-1}k_*,\qquad
 F=k^*\ker L_h^{c+1},\qquad E=\ker k_*\subset H^n(Y),
\]
and let $e_A$ be the primitive projector on $H^n(A)$. Then
\[
 \Pi_F=k^*e_A r,\quad \Pi_E=1-k^*r,\qquad
 H^n(Y)_{\rm prim}=F\mathbin\perp E
\]
are algebraic self-adjoint projectors and an orthogonal decomposition.
They induce a rational isogeny decomposition of primitive incidence
quotients
\[
 P^m(Y)_{\rm prim}\sim B_F\times B_E.
\]
Its signed incidence class is nondegenerate on both factors and is
a rational polarization on $B_F$, conditional on 032. Consequently
the full Th.~4 assertion for this particular $Y$ is equivalent to
ampleness of its signed incidence class on $B_E$. In particular it
holds whenever $B_E=0$.
\end{proposition}

\begin{proof}
The fixed-and-variable projector proof of H9-T13 uses smooth weak
Lefschetz and the identity $[Y]=h^c$, before any pencil monodromy is
introduced. These hypotheses hold here. Explicitly,
$rk^*=1$, and $k^*r$ is the orthogonal projector onto the ambient
middle cohomology. The primitive condition for an ambient class
$k^*u$ is $L_h^{c+1}u=0$: the Gysin map on $H^{n+2}(Y)$ is
injective by dual weak Lefschetz. The same injectivity makes every
element of $\ker k_*$ primitive. The primitive projector of $A$
is algebraic; the projectors of $Y$ outside its middle degree come
from the weak-Lefschetz identifications with $A$ and their duals.
Thus the displayed middle projectors are actual algebraic
correspondences. Self-adjointness and nondegeneracy give the
incidence decomposition exactly as in JAC-T4.1.

For clarity, the positive scalar in the fixed-part test also remains
unchanged. If a normalized fixed family pulled to a pointed curve
$C$ of genus $b\geq1$ has class $\xi$, put $\eta=(1\otimes r)\xi$
and lift its parameter degree-one component to an algebraic class
$w$ on $J(C)\times A$. Then $h^{c+1}w=0$ and, for
$H=\theta+h$,
\[
 H^{b+c}w=0,\qquad
 \int_{J(C)\times A}w^2H^{b+c-1}
   =\frac{(b+c-1)!}{c!}\int_{C\times Y}\xi^2.
\]
This is H9-T13, equation (H9.47); its proof needs no genericity.
The coefficient is positive. Conditional on 032, the signed curve
degree is nonnegative, with equality only for a zero mixed class.
The nondegenerate signed divisor is therefore ample by the
nef-plus-nondegenerate criterion of JAC-T4.3. Orthogonality makes
ampleness of the full signed divisor equivalent to ampleness of
both factors, proving the final assertions.
\end{proof}

\subsection*{C44-W08. Specialization of the inherited incidence polarization}

\begin{theorem}[C44-W08; conditional on unrefereed 032]
Let $\mathscr Y\to U$ be the smooth locus of a genuine Lefschetz
pencil in the setting of H9-T13, and let $s\in U$ be any geometric
point. Write $B_{\bar\eta}=P^m(Y_{\bar\eta})_{\rm prim}$ for the
geometric generic primitive incidence quotient. After a finite
extension of the local generic field, $B_{\bar\eta}$ has good
reduction at $s$, and its special fibre maps with finite kernel to
$P^m(Y_s)_{\rm prim}$. The actual signed incidence form is positive
on this inherited image. Hence Th.~4 holds for $Y_s$ whenever
\begin{equation}\label{c44:no-incidence-jump}
 \dim P^m(Y_s)_{\rm prim}=\dim P^m(Y_{\bar\eta})_{\rm prim}.
\end{equation}
More particularly, if the generic variable incidence factor of the
pencil is nonzero, then Th.~4 holds for every smooth geometric fibre
of the pencil. If that generic factor is zero, only the new
variable incidence directions that occur at a special fibre remain
outside this theorem.
\end{theorem}

\begin{proof}
Work over the strict henselian local trait at $s$, or its completion,
and pass to a finite extension of its fraction field over which an
abelian representative $B$ and a normalized rational Poincare cycle
$z$ are defined. These are finite data; rational descent is the one
fixed in JAC-P00. As in H9-T5, equations (H9.10)--(H9.11), the
mixed class of $z$ injects the Tate module of $B$ into the appropriate
Tate twist of $H^n(Y_{\bar\eta})$. Indeed its self-incidence is
nondegenerate, so a vector killed by the family map would be in
the radical of that incidence pairing. Smooth proper base change
makes the target unramified at this smooth fibre. Thus $B$ has
good reduction by the Neron--Ogg--Shafarevich criterion, just as in
H9-T5. One exact primary author source for the criterion is Milne,
\emph{Complex Multiplication}, Theorem 6.12, p.~51, which cites
Serre--Tate (1968), Theorem 1.

Let $\mathscr B$ denote the resulting abelian scheme over the trait.
Take the closure of a rational-cycle representative of $z$ in
$\mathscr B\times\mathscr Y$. These schemes are smooth over a
regular trait, so the refined fibre restrictions and intersection
products are defined. Its special-fibre mixed class is the
specialization of the generic mixed class. In particular it is
still in parameter degree one and is primitive on $Y_s$; these
conditions are kernels of maps of the smooth-proper cohomology
systems. Rigidification or degree-one projection, if desired,
alters no incidence class.

Form the actual relative self-incidence divisor by the same proper
pushforward of the product of the two cycle families. Clear its
rational denominator. Its signed half-diagonal pullback is a line
bundle $\mathscr L$ on $\mathscr B$, and refined base change says
that its two fibre classes are exactly the corresponding signed
incidence classes, up to the one positive clearing factor. The
generic class is ample by H9-T13.

Here is the specialization step without an assertion that arbitrary
special cycles lift. Choose a relatively ample line bundle
$\mathscr H$ on $\mathscr B$. Flatness and the Hilbert polynomial
make
\[
 \chi(L_t+uH_t)=\frac{(c_1(L_t)+u c_1(H_t))^{\dim B}}{(\dim B)!}
\]
the same real polynomial in $u$ for the generic and special fibres
$t$. On the generic fibre this polynomial is positive for every
$u\geq0$, since $L_{\eta}+uH_{\eta}$ is ample. Thus on the
special fibre $\chi(L_s)\ne0$ and the polynomial has no positive
real root. Mumford's index theorem gives index zero for $L_s$;
the nondegenerate effective criterion makes it ample. These are
exactly the two abelian-variety criteria used in the checked proof
of JAC-T4.3 (Mumford, \emph{Abelian Varieties}, Chapter III,
Section 16, and Chapter II, Section 6, Application 1).

The specialized family induces a homomorphism
$u:B_s\to P^m(Y_s)_{\rm prim}$. Saito's transpose and incidence
compatibility identifies $u^*\ell_s$ with the ample class just
constructed. Therefore $u$ has finite kernel: restriction to any
positive-dimensional connected kernel would make an ample class
zero. This proves positivity on the inherited image. Under
\eqref{c44:no-incidence-jump} the map is an isogeny, and ampleness
descends through that isogeny.

For the final assertion apply the same construction to the generic
variable factor alone. If that factor is nonzero, the irreducibility
after finite-index restriction in H9-T13, or H9-T5 (H9.11), forces
its family map to fill the entire variable local system $E$.
Specialization preserves this rank. The special variable incidence
factor injects into $E_s$ by its own nondegenerate self-incidence,
so its dimension is at most $\operatorname{rank}(E)/2$. The
inherited image already has that dimension. It therefore maps
isogenously onto the special variable incidence factor and gives
its positive signed form. C44-W07 supplies the fixed factor for
every smooth fibre. This proves the stated all-smooth-fibre scope.
\end{proof}

\paragraph{The remaining special-fibre gate.}
When the generic variable incidence quotient vanishes, the
construction specializes only the fixed incidence quotient. A
special fibre may have a new variable incidence factor. Neither
smooth proper base change nor constancy of the rank of $E$ makes
that factor extend as algebraic cycle families through the pencil.
Its signed positivity is the smaller unresolved statement. Thus
C44-W08 does not turn H9-T13 into an unconditional assertion about
every smooth complete intersection; it proves precisely the
rank-preserving specializations and the nonzero-generic-variable
case, with 032 retained as the positivity input.

\subsection*{C44-W09. The full geometric-abelian-type extension with its extra input exposed}

\begin{proposition}[C44-W09; conditional on unrefereed 032 and the separate JAC-P20 transfer]
Assume, in addition to 032, Milne's Proposition 4.2 in
\emph{Grothendieck's standard conjecture of Lefschetz type over
finite fields}, version 14 July 2022, in the total-cohomology form
specified and audited in JAC-P20. Then JAC-T4.3 holds for every
smooth projective odd-dimensional variety of geometric abelian
type, with its actual signed primitive incidence divisor.
\end{proposition}

\begin{proof}
JAC-P16 gives algebraic Lefschetz and primitive projectors. For
every pointed test curve $C$, the product $J(C)\times X$ is still
of geometric abelian type. JAC-P20's actual abelian cohomology
injection and the additional transfer assertion give the Hodge
standard conjecture for this product. The construction of
JAC-T4.2, using the algebraic Abel inverse on the first factor,
now gives an algebraic primitive class $w$ on $J(C)\times X$ and
the exact identity
\[
 \int_{J(C)\times X}w^2(\theta+h)^{b-1}
       =(b-1)!\int_{C\times X}\xi^2.
\]
The Hodge-standard input on this particular product makes its
signed value nonnegative. Saito's nondegenerate incidence factor
and JAC-T4.3's curve criterion then make the signed divisor ample.
This argument proves the stated conditional corollary; it does
not establish the additional transfer assertion itself.
\end{proof}

\paragraph{Framework and precise dependency.}
C44-W06 proves a positive scalar comparison of the actual forms,
not merely an embedding of motives. C44-W07--W08 use smooth
proper base change, SGA 7's fixed and variable projectors, and
the extension of actual cycle correspondences across a regular
trait. Their target is cote 44, scans 008--009, Th.~4 and its
primitive incidence higher Picard variety. The new residual gate
is the positive sign on a special new variable incidence factor,
not the already checked generic pencil sign. For all varieties
of geometric abelian type, JAC-P20 still leaves the comparison
between an abstract transported positive form and the particular
geometric form as an additional dependency. C44-W09 records
the exact implication if that dependency is supplied.


\section{The supersingular Fermat test for cote 16}

\subsection*{C44-S01. The proposed Fermat candidate already has the required sign}

\begin{proposition}[C44-S01; unconditional]
Let $X=X_N^{2r}\subset\mathbf P^{2r+1}$ be the smooth Fermat
hypersurface of degree $N\geq3$ in characteristic $p\nmid N$,
with $r\geq2$.  Suppose $p^\nu\equiv-1\pmod N$ for some
$\nu\geq1$.  Then all even cohomology is algebraic and all odd
cohomology vanishes.  Numerical and homological equivalence
coincide.  The signed form
\[
 Q_V(x,y)=(-1)^r\deg(xy),\qquad
 V=\ker(h:H^{2r}(X)\longrightarrow H^{2r+2}(X)),
\]
is positive definite on its rational algebraic cycle structure.
In particular the Hodge standard conjecture for $X$ is known
unconditionally; it is not a consequence that needs 032.
\end{proposition}

\begin{proof}
Put $q=p^\nu$ and $b=(q+1)/N$.  The Hermitian Fermat
hypersurface $H=X_{q+1}^{2r}$ maps finitely and surjectively to
$X$ by
\[
 [x_0:\cdots:x_{2r+1}]
 \longmapsto[x_0^b:\cdots:x_{2r+1}^b].
\]
Its degree is $b^{2r+1}$: $f^*h_X=bh_H$ and
$\int_Hh_H^{2r}=q+1=bN$ determine the degree by projection
formula.  In particular $f_*f^*=b^{2r+1}$ and $f^*$ is injective.
Shimada, \emph{Lattices of Algebraic Cycles on Fermat Varieties
in Positive Characteristics}, Theorem 1.1(1)--(2) and
Corollary 1.2, proves that linear cycles span the middle
cohomology of $H$ and that its signed primitive intersection
form is positive definite.  Weak Lefschetz gives the other
cohomology groups.  Pushing forward these algebraic classes
shows algebraicity of all cohomology of $X$.

If $x\in V$, then
$h_Hf^*x=b^{-1}f^*(h_Xx)=0$, and
\[
 (-1)^r\int_H(f^*x)^2
  =b^{2r+1}(-1)^r\int_Xx^2.
\]
The injectivity of $f^*$ and Shimada's positive form prove the
claim.  In all other degrees $X$ has the ambient projective
space Lefschetz strings, whose sign is immediate.  The
intersection pairing is nondegenerate on the rational
algebraic classes because these span cohomology; thus
numerical and homological equivalence agree.  Since
$\operatorname{NS}(X)_\mathbf Q=\mathbf Qh$, this also covers
every ample polarization.
\end{proof}

\paragraph{Exact status of the proposed test.}
Shioda--Katsura, \emph{On Fermat Varieties}, Theorem II,
printed p.~97, says for even dimension and $N\geq4$ that
algebraic generation of cohomology is equivalent to the
displayed congruence.  Thus, in that range, the entire stated
supersingular Fermat candidate class is covered by C44-S01.
The new argument above is only the finite-cover descent
of Shimada's already known sign.  It neither attributes a
new theorem to Grothendieck nor treats algebraicity alone as
a positivity theorem.  The square needed in scan 049 is
handled explicitly below.

\subsection*{C44-S02. An explicit smooth Segre hyperplane}

\begin{proposition}[C44-S02; unconditional]
Let $X=X_N^d$ be a smooth Fermat hypersurface with $N\geq3$.
Choose a finite field $k_0$ containing every $(N-2)$-nd root
of unity in $\overline{\mathbf F}_p$.  Let $M>N(d+1)$ and
choose $\beta$ of degree $M$ over $k_0$ in its extension of
degree $M$.  Define $c_i=\beta^{(N-2)i}$ for
$0\leq i\leq d+1$.  Then
\[
 W=\left\{\sum_{i=0}^{d+1}c_ix_iy_i=0\right\}
       \subset X\times X
\]
is a smooth hyperplane section for the Segre polarization
$L=h_1+h_2$.  Also $Y=X\cap\{x_{d+1}=0\}$ is a smooth
hyperplane section of $X$.
\end{proposition}

\begin{proof}
The assertion for $Y$ follows from $p\nmid N$.  If $W$ had
a singular point, the differential of its equation would
annihilate both projective tangent spaces.  Since the
Fermat gradients are $N(x_i^{N-1})$ and $N(y_i^{N-1})$,
there would be scalars $\lambda,\mu$ with
\[
 c_iy_i=\lambda x_i^{N-1},\qquad
 c_ix_i=\mu y_i^{N-1}\quad\text{for every }i.
\]
Both scalars are nonzero, since all $c_i$ are nonzero and
neither projective coordinate vector is zero.  The nonzero
coordinates of $x$ and $y$ have the same nonempty support
$S$.  Substitution gives, for $i\in S$,
\[
 (x_i^N)^{N-2}=(\mu\lambda^{N-1})^{-1}c_i^N.
\]
Choose $\rho$ with $\rho^{N-2}=(\mu\lambda^{N-1})^{-1}$.
It follows that
$x_i^N=\rho\zeta_i\beta^{Ni}$ for roots of unity
$\zeta_i\in k_0$ with $\zeta_i^{N-2}=1$.
This remains valid if $p\mid(N-2)$, since only the distinct
roots are used.  The Fermat equation for $x$ now gives
\[
 \sum_{i\in S}\zeta_i\beta^{Ni}=0.
\]
This is a nonzero polynomial over $k_0$ of degree at most
$N(d+1)<M$ vanishing at $\beta$, a contradiction.  The
Jacobian criterion therefore proves smoothness.
\end{proof}

\subsection*{C44-S03. The exact square sign and all three corrected scan-050 conditions}

\begin{theorem}[C44-S03; unconditional]
For $X=X_N^{2r}$ of C44-S01, put $d=2r$, $T=X\times X$,
$L=h_1+h_2$, and take $Y,W$ as in C44-S02.  Then
\[
 C_{\rm G}(X),\quad C_{\rm G}(Y),\quad R_0(T,W),\quad
 G_{L,0}(T,W),\quad A_0(T),\quad G_0(T,W),\quad A_0(W)
\]
all hold with the meanings fixed by SIGN\_NEXT SN-C03--SN-C04.
The Hodge standard conjecture $I(T)$ holds.  In its middle
degree the primitive space and form are explicitly
\begin{gather}
 P^{2d}(T)=\mathbf Qw\ \perp\ (V\otimes V),\qquad
 w=\sum_{i=0}^{d}(-1)^ih_1^ih_2^{d-i},\label{c44:fermat-square-primitive}\\
 \int_T(cw+z)^2=(d+1)N^2c^2+(Q_V\otimes Q_V)(z,z)>0
 \quad\text{if }(c,z)\ne(0,0).\label{c44:fermat-square-sign}
\end{gather}
The ring of degree-zero rational correspondence classes
modulo homological equivalence is
\[
 \prod_{i=0}^{d}\operatorname{End}_{\mathbf Q}
       \bigl(\operatorname{Alg}^i(X)\bigr)
 \simeq \mathbf Q^d\times M_{b_d(X)}(\mathbf Q),
\]
and is therefore semisimple.  These are the actual sign and
semisimplicity conclusions of cote 16, scan 049, for this
nonabelian family.
\end{theorem}

\begin{proof}
Weak Lefschetz gives the orthogonal decomposition
\[
 H^{\rm even}(X)=U\perp V,\qquad
 U=\mathbf Q\langle1,h,\ldots,h^d\rangle,
\]
where $V$ is in cohomological degree $d$ and $hV=0$.
The rational structures here are given by cycle classes;
they become the full cohomology after scalar extension.
The middle cohomology of $T$ consequently consists of the
ambient terms, the two terms $V\otimes h^r$ and
$h^r\otimes V$, and $V\otimes V$.  On the ambient
terms the kernel of multiplication by $L$ is the single
line displayed in \eqref{c44:fermat-square-primitive}:
successive coefficients satisfy $c_i+c_{i-1}=0$.
Multiplication by $L$ is injective on each of the two
cross terms and their images have different Kunneth
degrees.  It is zero on $V\otimes V$.  Orthogonality
follows from $\int_Xvh^r=0$ for $v\in V$.

In $w^2$ the $d+1$ contributing pairs have indices
$i,d-i$ and sign $(-1)^d=1$, giving
$\int_Tw^2=(d+1)N^2$.  The pairing on $V\otimes V$
is the tensor product of the two raw pairings; their
two factors $(-1)^r$ cancel.  It is therefore
$Q_V\otimes Q_V$, a positive definite form.  This proves
the exact critical sign.

For the full Hodge standard statement, decompose
$H^*(T)$ into the four $L$-stable blocks
$U\otimes U$, $V\otimes U$, $U\otimes V$, and
$V\otimes V$.  They are orthogonal.  The first has
the Lefschetz module and intersection forms of
$\mathbf P^d\times\mathbf P^d$, multiplied by $N^2$;
its signed forms are positive by characteristic-zero
Hodge--Riemann, with identical rational matrices in
every characteristic.  In $V\otimes U$ the primitive
classes are exactly the base $V\otimes1$, in
codimension $r$.  Their signed form is $NQ_V$;
the same is true for $U\otimes V$.  The last block
is primitive in middle degree and was just treated.
This proves $I(T)$ in all degrees.  Every ample
class on $T$ is a positive linear combination of
$h_1,h_2$; the same argument, with positive rescalings,
also handles these polarizations if required.

For a smooth hypersurface $Z$ of degree $N$ and
dimension $e$, the Lefschetz inverses below its middle
degree are already explicit.  For even $i<e$ the
correspondence
\[
 \frac1N h_Z^{i/2}\times h_Z^{i/2}
\]
induces the inverse of
$L_Z^{e-i}:H^i(Z)\longrightarrow H^{2e-i}(Z)$;
for odd $i<e$ both spaces vanish.  The middle inverse
is the identity.  Thus $C_{\rm G}(X)=B(X)$ and
$C_{\rm G}(Y)=B(Y)$.  The product stability of $B$
(Kleiman 1968, Proposition 2.12), or the cycle-basis
matrix units below, gives $B(T)$ and in particular
the algebraicity of $(L^2)^{-1}$ in the critical
degrees of SN-C03.

Write $j:W\hookrightarrow T$.  If $j^*x$ is algebraic
for $x\in H^{2d-2}(T)$, then
$L j_*j^*x=L^2x$ is algebraic; applying the algebraic
inverse proves $R_0(T,W)$.  For
$\gamma\in\operatorname{Alg}^{d+1}(T)$, set
$x=(L^2)^{-1}\gamma$ and $y=j^*x$.  Both are
algebraic and $Lj_*y=\gamma$, proving $G_{L,0}$.
The critical $A_0(T)$ was used here; the equivalence
in SN-C03 supplies $G_0$ and $A_0(W)$ as well.
This verifies the maps with the required extra $L$;
no unweighted lower-Gysin replacement is made.

Finally choose rational cycle bases for each
$\operatorname{Alg}^i(X)$ and the intersection-dual
bases in degree $d-i$.  Exterior products of these
classes give all rational matrix units on
$H^{2i}(X)$.  Every degree-zero correspondence
acts degreewise; Kunneth and Poincare duality make
this action faithful modulo homological equivalence.
It is therefore exactly the product of the stated
full matrix algebras.  Since all nonmiddle Betti
numbers in even degree are one, the displayed
description follows.
\end{proof}

\paragraph{What task 4 does and does not close.}
C44-S01--C44-S03 rule out the suggested supersingular Fermat
family as an example whose square-sign status is unknown.
They verify the corrected scan-050 proposition and the
scan-049 sign and correspondence ring explicitly and
unconditionally on that family.  A different smooth
projective $X$ with a genuinely unestablished $I(X\times X)$
is still required for the adversarial part of task 4; absence
of a source for one named equation would not certify that
status.  None of these findings closes the special-axis
sign for a complete intersection in an abelian variety.

\paragraph{Framework and programme crosswalk.}
For C44-S01--S03 the relevant archive locators are
\texttt{ARCHIVE\_16\_scan049.tex}, the semisimple
correspondence paragraph, and
\texttt{ARCHIVE\_16\_scan050.tex}, the proposition and
final weighted diagram.  They instantiate SIGN\_NEXT
SN-C03--SN-C04 and its successor question 4.  Their
Lefschetz inverses also give the $C_{\rm G}$ input to
cote 44, Th.~6, scan 004, through JAC-P02.  Weak
Lefschetz and Gysin are the Weil-cohomology operations
in SGA 5, Exp. VII, Section 7 (complete intersections).
The finite power maps and explicit hyperplanes are
scheme morphisms over finite extensions of $\mathbf F_p$;
the identities survive every Weil cohomology realization.
At the level of homological motives the Fermat object
is a direct sum of Tate objects with the positive
Lefschetz forms just computed.  No equality with its
Chow motive modulo rational equivalence is silently
deduced.  C44-W01--W05 use the exterior-power, hence
$\lambda$-operation, of $h^1(A)$, the algebraic
endomorphism projector, and rational cycle intersections.
No additional $K$-theory assertion, Selmer injection,
or motivic $L$-function nonvanishing is inferred.

% END work/scan_w/contribution.tex

% BEGIN work/scan_w/relative_th4.tex
% C44_POS /root/scan_w; GPT-6 Astra Pro; effort not exposed.
\section{Theorem 4 from an extending odd-degree abelian correspondence}

\begin{theorem}[C44-W10: a relative criterion for the actual incidence polarization]
Let $R$ be a complete discrete valuation ring of mixed characteristic
$(0,p)$ with algebraically closed residue field $k$ and fraction field
admitting an embedding into $\mathbf C$. Fix $\ell\ne p$.
Let $\mathscr X/R$ be smooth projective of relative dimension
$n=2m-1$, with relatively ample class $h$, and write $X=\mathscr X_k$.
Suppose a relative rational algebraic correspondence $e$ realizes the
orthogonal primitive projector
\[
 P^n(\mathscr X)=\ker(L_h:H^n(\mathscr X)\longrightarrow
                                      H^{n+2}(\mathscr X)),
 \tag{C44.W10.1}
\]
extended by zero in other degrees. Suppose also that relative
algebraic correspondences give an injection on every geometric fibre
\[
 \alpha:P^n(\mathscr X)\longrightarrow
       \bigoplus_\nu H^{2s_\nu-1}(\mathscr A_\nu)(s_\nu-m),
 \tag{C44.W10.2}
\]
where the $\mathscr A_\nu/R$ are abelian schemes and $s_\nu\geq1$.

Then, conditional on the unrefereed 032/Milne 2026 input in C44-A08,
cote~44, Theorem~4 holds for $X$ with this primitive polarization.
Precisely, on Saito's primitive incidence quotient
$B=P^m(X)_{\rm prim}$ the actual rational incidence divisor is ample:
\[
       \ell_B=\frac{(-1)^m}{2}\,\Delta_B^*d_B
                  \in\operatorname{NS}(B)_{\mathbf Q}.
 \tag{C44.W10.3}
\]
After multiplication by a positive integer it defines a polarization,
including when $p=2$. No extension to characteristic zero of a test
map $C\to B$ or of its cycle family is assumed. Milne 2022,
Proposition~4.2, is not an additional input.
\end{theorem}

\begin{proof}
The correspondence $e$ induces an idempotent on $P^m(X)$ by the
incidence-detection result JAC-P01: $e^2-e$ has zero relevant odd
cohomological action. In middle degree, primitive cohomology is the
orthogonal complement of $hH^{n-2}(X)$, and hard Lefschetz makes both
summands nondegenerate. Thus $e$ is self-transpose on the relevant
cohomology. Saito's transpose compatibility for the canonical
incidence duality makes its image $B$ a nondegenerate incidence
factor. This is exactly the argument of JAC-T4.1; that argument uses
the algebraicity and self-transpose property just supplied, not that
$X$ itself is abelian. It gives a normalized rational cycle family
$z$ on $B\times X$, with class entirely in
\[
                H^1(B)\otimes P^n(X),
\]
whose self-incidence is $d_B$. The associated symmetric rational
homomorphism is an isogeny.

Let $f:C\to B$ be the normalization of an arbitrary integral curve
in $B$, translated so that a chosen point maps to the origin. Its
genus $b$ is positive, since an abelian variety has no rational
curves. The pullback $\xi=(f\times1)^*z$ is a rational algebraic
class in
\[
                  H^1(C)\otimes P^n(X).
 \tag{C44.W10.4}
\]
We will prove the actual geometric sign
\[
                         (-1)^m\int_{C\times X}\xi^2\geq0,
 \tag{C44.W10.5}
\]
with strict inequality if its cohomology class is nonzero.

\paragraph{Lifting the test curve alone.}
Every smooth proper curve over $k$ lifts to a smooth proper curve
$\mathscr C/R$. Here is a precise way to use the curve deformation
theorem over this particular trait. The stack of smooth proper
curves is smooth over $\operatorname{Spec}\mathbf Z$ (Stacks
Project, Tag~0E1N, Lemma~109.16.2). Choose a scheme atlas near the
point represented by $C$. Since $k$ is algebraically closed, after
passing to an atlas point there is a $k$-point representing this
curve. The atlas is smooth over $R$. A smooth scheme over a
henselian local ring has an $R$-point through any given residue-field
point: use an etale coordinate map to affine space, lift its
coordinates, and apply the henselian lifting property. Pulling back
the universal curve gives $\mathscr C$. The chosen point of $C$
lifts by the same argument to a section. Its divisor has degree one
on each fibre and is relatively ample, so this curve model is
projective. Write $u$ for that relative ample divisor class.

The relative Jacobian $\mathscr J=\operatorname{Pic}^0_{\mathscr C/R}$
is an abelian scheme. Its Abel and Poincare correspondences identify
$H^1(\mathscr C)$ with $H^1(\mathscr J)$ on every geometric fibre,
by algebraic correspondences over $R$. Concretely, if $a$ is the
relative Abel map and $\theta$ the canonical principal polarization,
an inverse to $a^*$ on degree one is
\[
        j_C=(b-1)!\,(L_\theta^{b-1})^{-1}a_*.
 \tag{C44.W10.6}
\]
The inverse Lefschetz correspondence is relative algebraic by the
explicit abelian operators in C44-A08. The formula follows fibrewise
from $a_*[C]=\theta^{b-1}/(b-1)!$ and the projection formula, as in
JAC-T4.2. A relatively ample representative of the polarization can
be chosen over this complete trait after a finite extension if
necessary; the argument is unchanged by that extension.

\paragraph{A primitive middle summand with its actual duality.}
On $\mathscr C\times_R\mathscr X$ put
\[
                  V=H^1(\mathscr C)\otimes P^n(\mathscr X).
\]
Its degree is $1+n=2m$, the middle degree of this $2m$-fold.
The product of the relative degree-one curve projector and $e$
is an algebraic projector with image $V$. Both factors are
self-transpose for their middle Poincare dualities, so the product
is self-transpose for middle Poincare duality on the product.
Its restricted pairing is nondegenerate. The inverse duality map
is represented algebraically by this projector together with the
middle Poincare duality correspondence. Thus the algebraic-inverse
hypothesis of C44-A08 holds here, rather than being inferred from
an abstract choice of positive form.

For the product ample class $u+h$ we have
\[
                       (u+h)V=0.
 \tag{C44.W10.7}
\]
Indeed $uH^1(\mathscr C)=0$ by the degree on a curve and
$hP^n(\mathscr X)=0$ by (C44.W10.1). Hence $V$ is primitive.
Tensoring (C44.W10.2) with (C44.W10.6), and then using the Kunneth
summand of a product of abelian schemes, gives an extending
algebraic injection
\[
 V\longrightarrow
 \bigoplus_\nu H^{2s_\nu}(\mathscr J\times_R\mathscr A_\nu)
                                  (s_\nu-m).
 \tag{C44.W10.8}
\]
It is injective on every geometric fibre. All hypotheses of
C44-A08 are now verified with $r=m$ and dimension $2m$.
Its signed primitive form is exactly
$(-1)^m\int_{C\times X}\xi^2$, with no additional Lefschetz
factor or comparison scalar. It proves (C44.W10.5).
The cycle $\xi$ itself, and the map $f$, have been used only in the
special fibre; neither was lifted.

\paragraph{The incidence divisor.}
The normalized incidence identity and refined pullback give
\[
                  \deg_C f^*\ell_B
                         =\frac{(-1)^m}{2}
                                       \int_{C\times X}\xi^2\geq0.
 \tag{C44.W10.9}
\]
The degree-one projection of the parameter family removes the
degree-zero and degree-two parameter components in this formula.
Since every integral curve in $B$ has been tested, $\ell_B$ is nef.
Its associated symmetric homomorphism is an isogeny by the
nondegeneracy of the primitive incidence factor proved above.
A nef nondegenerate rational divisor on an abelian variety is ample,
by the criterion used in JAC-T4.3 (Mumford's index theorem and the
effective nondegenerate criterion). This proves (C44.W10.3).
The zero-dimensional case for $B$ is immediate. All normalizations
are rational; division by $2$ therefore creates no exception at
characteristic $2$.
\end{proof}

\begin{corollary}[C44-W11: odd Fermat varieties and their product families]
Let $k=\overline{\mathbf F}_p$, let $d\geq2$ with $p\nmid d$, and let
\[
             X=\prod_{j=1}^{a}F_{n_j,d,k},\qquad
             n_j\geq1,\qquad \sum_jn_j=2m-1.
\]
For every polarization $h=\sum_j b_jh_j$ with $b_j\in\mathbf Q_{>0}$
and factor hyperplane classes $h_j$, the actual primitive incidence
divisor (C44.W10.3) is ample, conditional on the unrefereed
032/Milne 2026 input. In particular cote~44, Theorem~4 holds for
every odd-dimensional smooth Fermat hypersurface in every good
characteristic, with every rational ample class. For the displayed
products it holds for every rational ample class if all $n_j\geq3$;
when curve or surface factors occur, only the displayed product
polarizations are asserted.
\end{corollary}

\begin{proof}
C44-F08 gives an actual relative split Chow-motive injection of
the Fermat product into a sum of Tate twists of powers of the
relative Fermat Jacobian, over
$\operatorname{Spec}\mathbf Z[1/d,\zeta_{2d}]$.
Restriction to odd middle degree gives (C44.W10.2). C44-F09 gives
relative Lefschetz operators and the primitive projector for every
displayed $h$. In middle degree its primitive image is orthogonal
to the Lefschetz image, so this projector is self-transpose as
required in C44-W10. For a single odd Fermat hypersurface this can
also be read directly from
\[
 [\Delta]-\frac1d\sum_{i=0}^{n}h^{n-i}\times h^i;
\]
transposition permutes the terms and its image is the primitive
middle cohomology.

Choose a prime of the cyclotomic base above $p$. Its local ring is
a discrete valuation ring. Complete its maximal unramified
extension and choose the algebraic closure of the residue field
to identify that residue field with $k$. This gives a complete
mixed-characteristic discrete valuation ring to which all the
relative correspondences and abelian schemes pull back. Its
fraction field has cardinality at most that of $\mathbf C$ and
characteristic zero, so it admits an abstract embedding into
$\mathbf C$: choose a transcendence basis and extend across the
remaining algebraic extension. A common such embedding is used
for the whole construction. An embedding is only needed to apply
complex Hodge--Riemann to this common generic fibre; no continuity
with a $p$-adic topology is required. Consequently C44-W10 applies.

For one odd-dimensional hypersurface of dimension at least three,
weak Lefschetz gives rational Neron--Severi rank one, generated by
the hyperplane class. In dimension one, numerical divisor classes
are likewise determined by their degree, so every rational ample
class has the same Lefschetz action as a positive rational multiple
of the hyperplane class. These observations prove the single-factor
assertion for all ample classes. For products with all $n_j\geq3$,
the Neron--Severi computation in C44-F10 shows that every ample
class is a positive rational combination of the factor hyperplane
classes. No such computation is invoked for products containing
curves or surfaces.
\end{proof}

\paragraph{Scope and the remaining geometric-type gate.}
C44-W10 is a sufficient relative condition, and C44-W11 supplies a
nonabelian family meeting it. It strengthens C44-W09 on this family
by replacing its additional Milne2022 transfer input with actual
extending correspondences and an explicit geometric-form proof.
It does not show that an arbitrary variety of geometric abelian
type has such a lift or extending injection. The input from 032
is still conditional; the constructions F08--F09 and the test-curve
lifting are unconditional. The result concerns the incidence
quotient of cote~44, scans~008--009, with the primitive qualification
in scan~008, and supplies a positive sign on its actual divisor,
not a separately transported polarization.


% END work/scan_w/relative_th4.tex

% BEGIN work/scan_w/sharp_cycle_bound.tex
% C44_POS /root/scan_w; GPT-6 Astra Pro; effort not exposed.
\section{Sharp complementary-cycle degrees and the refined detector constant}

\begin{proposition}[C44-W12: a sharp algebraic complementary-cycle estimate]
Let $A$ be an abelian variety of dimension $g\geq1$ over an
algebraically closed field of arbitrary characteristic, and let
$h=c_1(L)$ be an integral ample class. Put
\[
                   V=\frac{\int_Ah^g}{g!}=\chi(L)>0.
\]
If $Z,Z'$ are effective cycles of dimensions $r$ and $g-r$ with
nonnegative integer coefficients, then
\[
 0\leq\deg(Z\cdot Z')
   \leq\frac{\deg(h^rZ)\deg(h^{g-r}Z')}{r!(g-r)!V}.
 \tag{C44.W12.1}
\]
The coefficient is sharp. Neither the standard conjectures nor a
characteristic-zero lift is an input.
\end{proposition}

\begin{proof}
By bilinearity it suffices to treat integral subvarieties. Their
product $T=Z\times Z'$ is an integral projective variety, which may
be singular. Set $e=\deg(Z\cdot Z')$. The translation-moving
argument in C44-W01, with complementary dimensions $r,g-r$, gives
\[
 f_*[T]=e[A],\qquad f(z,z')=z-z',\qquad e\geq0.
\]
Inseparable generic degrees are included. Let $s(z,z')=z+z'$.
The degree of this map, interpreted as zero when it is not dominant,
is $\deg(Z\cdot[-1]_*Z')$. In a Weil cohomology of characteristic
different from the coefficient prime, $[-1]^*$ acts by $(-1)^i$
on $H^i(A)$, and therefore acts identically on every cycle class.
It follows that
\[
                         s_*[T]=e[A].
 \tag{C44.W12.2}
\]
This is an equality of top-dimensional pushforward cycles, whose
coefficients have just been computed as intersection numbers; no
algebraic equivalence between $Z'$ and $[-1]_*Z'$ is asserted.

Write $h_1,h_2$ for the pullbacks to $T$ and put
$\alpha=f^*h$, $\beta=s^*h$. The two classes are nef and the
numerical parallelogram identity gives
\[
 \alpha+\beta=2(h_1+h_2),\qquad
 \int_T\alpha^g=\int_T\beta^g=e\int_Ah^g.
\]
The Khovanskii--Teissier inequality for two nef divisors gives
\[
 \left(\int_T(\alpha+\beta)^g\right)^{1/g}
       \geq\left(\int_T\alpha^g\right)^{1/g}
            +\left(\int_T\beta^g\right)^{1/g}.
 \tag{C44.W12.3}
\]
It applies to integral possibly singular projective varieties in
arbitrary characteristic: see Lazarsfeld, \emph{Positivity in
Algebraic Geometry I}, Corollary~1.6.3(iii), pp.~90--91, and
Remark~1.6.5, p.~91. Equivalently, Cutkosky,
\emph{Teissier's problem on inequalities of nef divisors over an
arbitrary field}, Corollary~6.3, formula~(20), gives precisely this
statement for a complete integral variety over any field.
In particular no smoothness of $Z$ or $Z'$ is being assumed here.

Substitute the preceding equalities into (C44.W12.3) and cancel
the factor $2$. Raising to the $g$-th power gives
\[
 e\int_Ah^g\leq\int_T(h_1+h_2)^g
       =\binom gr\deg(h^rZ)\deg(h^{g-r}Z').
\]
Using $\int_Ah^g=g!V$ proves (C44.W12.1).

For sharpness, take a product of principally polarized abelian
varieties of dimensions $r$ and $g-r$, with its product principal
polarization, and take the two complementary factors as $Z,Z'$.
Then $e=1$, $V=1$, and their degrees are $r!$ and $(g-r)!$,
respectively. Equality holds. The cases $r=0$ or $r=g$ are given
by a point and the whole variety.
\end{proof}

\begin{theorem}[C44-W13: refined Gram and division constants]
Retain all the data and hypotheses of C44-W02--C44-W03, including
the explicit integral polynomial $p$, the dual numerical-lattice
condition $\delta t\in\Lambda^\sharp$, and the conditional
unrefereed 032 positivity input. Define
\[
 \widehat B=\frac{\sqrt{R\kappa/2}}{n!\sqrt V},\qquad
 \varepsilon=
 \begin{cases}1,&s\text{ even},\\2,&s\text{ odd}.
 \end{cases}
 \tag{C44.W13.1}
\]
For an integral detector $Y$ of degree $D$, its Gram entry obeys
\[
                  0<d=Q(y_1,y_1)\leq\widehat B^2D^2.
 \tag{C44.W13.2}
\]
The exact Gram matrix remains
$d\left(\begin{smallmatrix}1&s/2\\s/2&q^2\end{smallmatrix}\right)$.
For every nonzero rational Weil class $t$ under the stated lattice
condition, with $T_0=Q(t,t)$, the detector complexity satisfies
\[
 \delta_W\leq c_t\leq
 \delta\varepsilon\max\left\{
          \frac H4\widehat B^2,
          \frac{q\kappa}{2}\widehat B\sqrt{HT_0}\right\}
                                        \delta_W^2.
 \tag{C44.W13.3}
\]
The corresponding quartic bound is
\[
 c_t\leq\delta\max\left\{
          \frac H4\widehat B^4,
          \frac{q\kappa}{2}\widehat B^3\sqrt{HT_0}\right\}
                                        \delta_W^4.
 \tag{C44.W13.4}
\]
Both bounds remain valid with the variety-independent constant
$\widehat B_0=\sqrt{R\kappa/2}/n!$ in place of $\widehat B$.
\end{theorem}

\begin{proof}
The middle $K$-weight is $a^n\bar a^n=q$, so $g(q)=0$ and hence
$p(q)=0$. Use the actual effective cycle representatives
\[
 Y^+=\sum_{c_j>0}c_j[a^j]^*Y,\qquad
 Y^-=\sum_{c_j<0}(-c_j)[a^j]^*Y.
\]
They have difference representing $y_1$. The degree formula from
C44-W02 and $p(q)=0$ imply the exact identities
\[
       D(Y^+)=D(Y^-)=\frac{\kappa D}{2}.
 \tag{C44.W13.5}
\]
Thus $\kappa$ is even, and $[a]^*Y^\pm$ represent $y_2$ with
each part of degree $q\kappa D/2$. These are explicit signed
representatives; cancellation of common components is unnecessary.

The self-adjoint projector identity from C44-W02 is
$p(T)^\dagger p(T)=Rp(T)$. Consequently
\[
 d=R\left|\deg\bigl(Y\cdot(Y^+-Y^-)\bigr)\right|.
\]
By C44-W12 with $g=2n,r=n$, the two numbers
$I_\pm=\deg(Y\cdot Y^\pm)$ are nonnegative and satisfy
\[
              I_\pm\leq\frac{\kappa D^2}{2(n!)^2V}.
\]
Use $|I_+-I_-|\leq\max(I_+,I_-)$, rather than adding the two
upper bounds. This proves (C44.W13.2). All intersection estimates
in this paragraph are unconditional; 032 is used for the positive
Gram form and the subsequent Cauchy--Schwarz inequality.

Choose $D=\delta_W$ and put $b_i=Q(\delta t,y_i)\in\mathbf Z$.
Exact inversion of the Gram matrix gives
\[
 \frac{\delta\varepsilon Hd}{4}\,t
   =\varepsilon\left(q^2b_1-\frac{s}{2}b_2\right)y_1
        +\varepsilon\left(b_2-\frac{s}{2}b_1\right)y_2.
 \tag{C44.W13.6}
\]
Every displayed coefficient, including the multiplier on the left,
is an integer. If $s$ is even then $4\mid H$. If $s$ is odd,
integrality of the Gram entry $sd/2$ forces $2\mid d$, and
$\varepsilon=2$ makes both the multiplier and coefficients integral.
The multiplier is positive.

Writing $t=t_1y_1+t_2y_2$, positive definiteness gives
\[
 |t_1|\leq\frac{2q\sqrt{T_0}}{\sqrt{dH}},\qquad
 |t_2|\leq\frac{2\sqrt{T_0}}{\sqrt{dH}}.
\]
The two integer coefficients in (C44.W13.6) therefore have
absolute values at most
$\delta\varepsilon q\sqrt{HT_0d}/2$ and
$\delta\varepsilon\sqrt{HT_0d}/2$, respectively. With the exact
part-degrees (C44.W13.5), each positive or negative part of their
cycle combination has degree at most
\[
 \frac{\delta\varepsilon q\kappa}{2}
            D\sqrt{HT_0d}
    \leq\frac{\delta\varepsilon q\kappa}{2}
            \widehat B\sqrt{HT_0}\,D^2.
\]
The multiplier is at most
$\delta\varepsilon H\widehat B^2D^2/4$. These give the upper
bound in (C44.W13.3). The lower bound is exactly the integral
component argument of C44-W03.

For (C44.W13.4), use the integer multiplier
$N=\delta\det G=\delta d^2H/4$. The two adjugate coefficients
are $Nt_1,Nt_2$ and are integers by the integral Gram matrix.
Their absolute values are at most
$\delta qd^{3/2}\sqrt{HT_0}/2$ and
$\delta d^{3/2}\sqrt{HT_0}/2$. The degree of either signed part
is consequently at most
\[
        \frac{\delta q\kappa}{2}D d^{3/2}\sqrt{HT_0}
       \leq\frac{\delta q\kappa}{2}
                         \widehat B^3\sqrt{HT_0}D^4,
\]
whereas $N\leq\delta H\widehat B^4D^4/4$.
Finally $V\geq1$ proves the uniform-constant assertion.
\end{proof}

\paragraph{Exact relation to the earlier constants.}
C44-W12 sharpens the valid elementary estimate C44-W01, and
C44-W13 uses additionally the vanishing of the middle weight
$p(q)$ and the parity of the integral Gram matrix. No alteration
of the projector, detector degree, denominator convention, or
class-level meaning of $c_t$ is involved. In the characteristic-$p$
Theorem W deduction, the strongest constant established here for
$\sqrt{Q(y_1,y_1)}$ is $\widehat B D$. An identification with the
original analytic Wirtinger constant is not asserted.

The primary nef-inequality sources inspected were
\url{https://www.math.stonybrook.edu/~roblaz/Reprints/Laz.PAG.Chapt1.pdf},
Theorem~1.6.1 and Corollary~1.6.3 with their proofs and
Remark~1.6.5, and
\url{https://arxiv.org/pdf/1304.1218}, Theorem~6.1 and
Corollary~6.3, especially formula~(20) and its proof.


% END work/scan_w/sharp_cycle_bound.tex

% BEGIN work/consolidated_th5.tex
\section{The corrected fifth theorem, with all domains visible}

\begin{theorem}[C44-T05: synthesis of H9-T1, H9-T2, H9-T7,
H9-T17 and H9-T18]\label{thm:final-th5}
Use rational Chow groups and the signed geometric intersection height
of JAC1964 H9-T1. The following is the consolidated form of cote 44,
scan 004, Th.~5. The formulas below are unconditional under their
stated admissibility and geometric hypotheses. Each assertion that
uses the abelian Hodge standard conjecture is explicitly conditional
on the unrefereed 032/Milne inputs.

\begin{enumerate}
\item Let $f:\mathcal X\to C$ be projective with $\mathcal X$ smooth
of dimension $2m$, $C$ a smooth projective curve, and smooth generic
fibre $Y$. Let $x_0,y_0$ extend numerically trivial cycles on $Y$.
Choose a finite-dimensional vertical correction space, supported in
finitely many relevant fibres. For each such fibre let $W_s$ be its vertical numerical cycle
space, $R_s$ its radical, and $M_s$ the nonsingular matrix on a chosen
basis of $W_s/R_s$. Suppose the extensions annihilate $R_s$, and
that these spaces account for all vertical pairings. Write $a_s,b_s$
for their pairing vectors. Then admissible corrections exist and
\[
 h_m(x_0|_Y,y_0|_Y)=(-1)^m
 \left(\deg(x_0y_0)-\sum_s a_s^{t}M_s^{-1}b_s\right).
 \tag{C44.T05.1}
\]
If only a smaller vertical space has been used, this expression is
the corrected intersection, and is a height only after the remaining
vertical orthogonality is verified. Failure to annihilate a vertical
radical is an obstruction to correction, not a term one may omit.

\item For a Lefschetz pencil on a smooth $2m$-fold $X$, if
$H^{2m}(Y_s)\to H^{2m}(Y_{\bar\eta_s})$ is injective at every
singular fibre, all corrections vanish for ambient cycles whose
generic restrictions are homologically trivial:
\[
                  h_m(j^*z,j^*z')=(-1)^m\int_Xzz'.
 \tag{C44.T05.2}
\]
The same injectivity makes any global extension of a homologically
trivial generic cycle admissible. In the ordinary Picard--Lefschetz
setting the condition is equivalent to nonzero vanishing cycles;
SGA 7 condition $(A)$ is the characteristic-independent formulation
used in the predecessor. For genuine abelian pencils it holds by
JAC-P07. If $X$ is abelian, conditional on the unrefereed 032/Milne
preprints, (C44.T05.2) is positive definite
on primitive ambient algebraic cohomology.

\item If $X$ is abelian, let $B$ be its smooth axis and use
$\mathcal H,\mathcal H_0,A_X,A_B$ of H9-T7, (H9.17). Then
\[
 \mathcal H/\mathcal H_0\simeq A_X\mathbin\perp A_B,
 \qquad h_m\simeq
 \left((-1)^m\int_Xzz'\right)\mathbin\perp
 \left((-1)^{m-1}\int_Bww'\right).
 \tag{C44.T05.3}
\]
In particular $h_Bw=0$ for $w\in A_B$. Conditional on 032/Milne,
semipositivity on the whole $\mathcal H$ is equivalent to the weak
axis sign on $A_B$. If the axis form is positive definite, the
radical is exactly $\mathcal H_0$. If it is only semidefinite,
the quotient of the radical by $\mathcal H_0$ is the radical of
that axis form; equality with $\mathcal H_0$ is not inferred.
The axis sign is known for $m=2$, for characteristic zero, and in
the generic-axis vanishing cases H9-T8, H9-T8a and JAC-P09.
C44-A04 and C44-A06 give additional explicit positive subspaces,
not a sign theorem for every special axis.

\item For a smooth threefold generic fibre, on the admissible
homologically trivial domain let
\[
 \mathcal H_h^{\rm Alb}=\{z:\operatorname{alb}(hz)=0\},\qquad
 \psi_h(D)=\operatorname{alb}(h^2D),\qquad
 D_z=\psi_h^{-1}\operatorname{alb}(hz).
\]
The rational isogeny $\psi_h:\Pic^0(Y)\to\Alb(Y)$ gives the
orthogonal decomposition
\[
                    z=(z-hD_z)+hD_z.
 \tag{C44.T05.4}
\]
The first term lies in $\mathcal H_h^{\rm Alb}$; the second has
negative definite height modulo the constant trace. In
characteristic zero the first height is semidefinite by Zhang,
Theorem 1, with the precise admissible-model equality criterion
of H9-T17, (H9.77)--(H9.80). This sign statement does not extend
to the whole homologically trivial group after an arbitrary
finite cover: H9-T17 gives the exact negative value $-12d^3$ on
the displayed $E^4$ example, in every characteristic.

\item For an abelian-fourfold pencil and a finite separable cover
$C\to\mathbf P^1$ unramified over the critical values, let
$k:\mathcal X\hookrightarrow A\times C$ be the smooth incidence
fourfold. All homologically trivial generic cycles are admissible.
On $\mathcal H_h^{\rm Alb}$ the primitive correction of H9-T18
gives
\[
              h_2(z,z)=\int_{A\times C}a^2\lambda
                              +\int_{\mathcal X}v^2,
 \quad \lambda^2a=0,\quad k_*v=0.
 \tag{C44.T05.5}
\]
Conditional on 032/Milne the first term is positive definite on
its algebraic domain. The remaining assertion is precisely the
nonnegativity of the algebraic primitive form on $\ker k_*$.
C44-C01 splits off its positive old surface-axis part; the new
trace-kernel part is the remaining special-cover gate. The
generic vanishing results H9-T19--T26 and C44-C02--C04 give
the norm identities in their explicitly stated parameter scopes.
They do not establish this sign for every closed special cover.
\end{enumerate}
\end{theorem}

\begin{proof}
This is an assembly of the specified predecessor results, not a new
proof of them. Part (1) is H9-T1, (H9.3)--(H9.5), equal to
SIGN\_NEXT SN-44-01 after its normalization, with vertical blocks
separated by disjoint fibres. Part (2) is H9-T2, (H9.6)--(H9.7),
and its full-extension argument at the beginning of H9-T7;
the Picard--Lefschetz equivalence is SN-44-03. Part (3) is the
isometry H9-T7, (H9.18), with the elementary radical statement
of an orthogonal sum and the conditional ambient sign. Part (4)
is H9-T17, (H9.79)--(H9.84), including both of Zhang's hypotheses
and Kahn's trace quotient. Part (5) is H9-T12 admissibility and
H9-T18, (H9.86)--(H9.91), followed by the independently checked
old/new decomposition C44-C01. No step asserts algebraic
equivalence from homological triviality.
\end{proof}

\paragraph{The incidence higher Picard variety and the arithmetic reading.}
Grothendieck's $P^m$ in cote 44, scans 002 and 004, is attached
to algebraically trivial cycles via incidence. A homologically
trivial class in the displayed larger Chow groups is not thereby
a point of $P^m$. The Neron--Tate interpretation on a primitive
incidence quotient requires the actual regular homomorphism and
the height comparison; these are the separate inputs of
JAC-T4.1--T4.5 and SN-44-01. The full-middle-factor polarization
is refuted by JAC-P03, and the unrestricted covered height is
refuted by H9-T17. Neither counterexample refutes the primitive
axis sign asked here. The normalization and the Albanese kernel
implement the primitive-decomposition margin of scan 008.

In cote 16, scans 049--050, the algebraicity of Lefschetz operators
is condition $C=C_G=B$, while the height signs use $I$ on the
specified primitive algebraic spaces. Restoring the Lefschetz
factor in $A''$ (SN-C03--C04) is a separate issue from these
signs. Theorem C44-T05 does not deduce $I$ from $C$ or identify
cohomological semisimplicity with injectivity on rational Chow
groups. In the SGA framework, the local input is SGA 7 II,
XVIII 5.3.2, 5.3.4--5.3.5 and 6.2; the variable irreducibility
is XVIII 6.7.2. In the Chow-motive framework the global terms
are the blowup summand $h(B)(1)$ and the old/new finite-cover
projectors. These locators and their roles are retained from
the checked predecessor arguments.

% END work/consolidated_th5.tex

\part{Unchecked candidates and unresolved work}

\section{Unresolved goals}
The fixed-axis sign is not proved for every special axis, and no
wrong-sign algebraic class is constructed. The universal lifting
question for classes remains open. The special-cover sign remains
open at arbitrary closed points with new algebraic trace-kernel
classes. Fermat Abel--Jacobi injectivity remains open, with the
exact gates specified above. Full Th.~4 for arbitrary new special
incidence directions is not proved. The proposed supersingular
Fermat family does not meet the requested initially-unknown-$I$
status. The ledger and question register give the current precise
state of every treated row.
\end{document}
